% Cumulative English reader for Exposes I--VIII.
% The ordered component list is recorded in reader-components-1-to-8.txt.
\documentclass[oneside]{book}
\usepackage{sga-style}
\geometry{text={6.3in,8in},centering}
\setlength{\emergencystretch}{1.5em}
\title{Séminaire de Géométrie Algébrique du Bois-Marie\\[20pt]
Étale Cohomology\\[20pt]
(SGA $4\frac 1 2$)}
\author{by P. Deligne\\[10pt]
with the collaboration of J.~F. Boutot,\\
A. Grothendieck, L. Illusie, and J.~L. Verdier}
\date{1977}

\hypersetup{
  pdftitle={SGA 4 1/2: Étale Cohomology — English Translation},
  pdfauthor={Pierre Deligne; J. F. Boutot; Alexander Grothendieck; Luc Illusie; Jean-Louis Verdier},
  pdfsubject={Algebraic geometry; étale cohomology},
  pdfkeywords={SGA 4 1/2, étale cohomology, algebraic geometry, English translation},
  bookmarksopen=true,
  bookmarksnumbered=true
}

\begin{document}

\pagenumbering{alph}
\maketitle

\pagenumbering{roman}
\tableofcontents
\clearpage

\pagenumbering{arabic}
% BEGIN INCLUDED components/sga4half_1_en.tex
% !TEX root = ../SGA4half_English_Reader_1_to_8.tex

\chapter{Étale cohomology: the starting points}\label{I}
\blfootnote{by P.\ Deligne, written up by J.F.\ Boutot}

\section{Grothendieck topologies}\label{I:1}

Grothendieck topologies first appeared as the foundation of Grothendieck's
theory of descent (cf. SGA 1 VI, VIII); their use in the corresponding
cohomology theories came later. We follow the same path here: by formalizing
the classical notions of localization, local property, and patching
(\ref{I:1-1}, \ref{I:1-2}, \ref{I:1-3}), we obtain the general concept of a
Grothendieck topology (\ref{I:1-6}); to justify its introduction into
algebraic geometry, we prove a faithfully flat descent theorem
(\ref{I:1-4}), generalizing the classical theorem 90 of Hilbert
(\ref{I:1-5}).

The reader will find a more complete, but concise, exposition of the
formalism in Giraud \cite{gi64}. M. Artin's notes, “Grothendieck
topologies” \cite{ar62} (Chapters I to III), also remain useful. The
866 pages of Exposés I to V of SGA 4 are valuable when one considers exotic
topologies, such as the one giving rise to crystalline cohomology; in order
to use the étale topology, which is so close to classical intuition, it is
not essential to read them.

\subsection{Sieves}\label{I:1-1}

Let $X$ be a topological space and let $f:X\to \dR$ be a real-valued
function on $X$. The continuity of $f$ is a local property; in other words,
if $f$ is continuous on every sufficiently small open subset of $X$, then
$f$ is continuous on all of $X$. To formalize the notion of a “local
property,” we introduce some definitions.

A set $\cU$ of open subsets of $X$ is called a \emph{sieve} if, for every
$U\in\cU$ and $V\subset U$, one has $V\in\cU$. A sieve is called
\emph{covering} if the union of all the open subsets belonging to it is
equal to $X$.

Given a family $\{U_i\}$ of open subsets of $X$, the sieve generated by
$\{U_i\}$ is, by definition, the set of open subsets $U$ of $X$ such that
$U$ is contained in one of the $U_i$.

A property $P(U)$, defined for every open subset $U$ of $X$, is called
\emph{local} if, for every covering sieve $\cU$ of every open subset $U$
of $X$, $P(U)$ is true if and only if $P(V)$ is true for every
$V\in \cU$. For example, given $f:X\to \dR$, the property “$f$ is
continuous on $U$” is local.

\subsection{Sheaves}\label{I:1-2}

Let us make precise the notion of a function given locally on $X$.

\subsubsection{The sieve point of view}\label{I:1-2-1}

Let $\cU$ be a sieve of open subsets of $X$. A function given
$\cU$-locally on $X$ consists of a function $f_U$ on $U$ for every
$U\in \cU$, such that if $V\subset U$ then $f_V=f_U|V$.

\subsubsection{The Čech point of view}\label{I:1-2-2}

If the sieve $\cU$ is generated by a family of open subsets $U_i$ of $X$,
then giving a function $\cU$-locally amounts to giving a function $f_i$
on each $U_i$ such that
$f_i|{U_i\cap U_j} = f_j|{U_i\cap U_j}$.

In other words, if $Z=\coprod U_i$, then giving a function
$\cU$-locally amounts to giving a function on $Z$ that is constant on the
fibres of the natural projection $Z\to X$.

\subsubsection{}\label{I:1-2-3}

The continuous functions form a sheaf. This means that for every covering
sieve $\cU$ of an open subset $V$ of $X$, and every function
$\{f_U\}$ given $\cU$-locally such that each $f_U$ is continuous on
$U$, there exists a unique continuous function $f$ on $V$ such that
$f|U=f_U$ for every $U\in \cU$.

\subsection{Stacks}\label{I:1-3}

Let us now make precise the notion of a vector bundle given locally on $X$.

\subsubsection{The sieve point of view}\label{I:1-3-1}

Let $\cU$ be a sieve of open subsets of $X$. A vector bundle given
$\cU$-locally on $X$ consists of
\begin{enumerate}[\indent a)]
  \item a vector bundle $E_U$ on each $U\in \cU$,
  \item if $V\subset U$, an isomorphism
    $\rho_{U,V}:E_V\iso E_U|V$, satisfying
  \item if $W\subset V\subset U$, the diagram
    \[\xymatrix{
      E_W \ar[r]^-{\rho_{U,W}} \ar[dr]_-{\rho_{V,W}}
        & E_U|W \\
      & E_V|W \ar[u]_-{\rho_{U,V}|W}
    }\]
    commutes, that is,
    $\rho_{U,W} = (\mbox{$\rho_{U,V}$ restricted to $W$}) \circ \rho_{V,W}$.
\end{enumerate}

\subsubsection{The Čech point of view}\label{I:1-3-2}

If the sieve $\cU$ is generated by a family of open subsets $U_i$ of
$X$, then giving a vector bundle $\cU$-locally amounts to giving:
\begin{enumerate}[\indent a)]
  \item a vector bundle $E_i$ on each $U_i$,
  \item if $U_{ij}=U_i\cap U_j=U_i\times_X U_j$, an isomorphism
    $\rho_{ji}:E_i|U_{ij}\iso E_j|U_{ij}$, such that
  \item if $U_{ijk}=U_i\times_X U_j\times_X U_k$, the diagram
    \[\xymatrix{
      E_i |U_{ijk} \ar[r]^-{\rho_{ki}|U_{ijk}} \ar[dr]_-{\rho_{ji}|U_{ijk}}
        & E_k|U_{ijk} \\
      & E_j|U_{ijk} \ar[u]_-{\rho_{kj}|U_{ijk}}
    }\]
    commutes, that is, $\rho_{ki}=\rho_{kj}\circ\rho_{ji}$ on
    $U_{ijk}$.
\end{enumerate}

In other words, if $Z=\coprod U_i$ and $\pi:Z\to X$ is the natural
projection, then giving a vector bundle $\cU$-locally amounts to giving:
\begin{enumerate}[\indent a)]
  \item a vector bundle $E$ on $Z$,
  \item if $x$ and $y$ are two points of $Z$ such that
    $\pi(x)=\pi(y)$, an isomorphism $\rho_{yx}:E_x\iso E_y$ between the
    fibres of $E$ at $x$ and $y$, depending continuously on $(x,y)$,
    and such that
  \item if $x$, $y$, and $z$ are three points of $Z$ such that
    $\pi(x)=\pi(y)=\pi(z)$, then
    $\rho_{zx}=\rho_{zy}\circ\rho_{yx}$.
\end{enumerate}

\subsubsection{}\label{I:1-3-3}

A vector bundle $E$ on $X$ defines a vector bundle $E_\cU$ given
$\cU$-locally, namely the system of restrictions $E_U$ of $E$ to the
objects of $\cU$. The fact that the notion of a vector bundle is local can
be expressed as follows: for every covering sieve $\cU$ of $X$, the
functor $E\mapsto E_\cU$ from vector bundles on $X$ to vector bundles
given $\cU$-locally is an equivalence of categories.

\subsubsection{}\label{I:1-3-4}

If in \ref{I:1} one replaces “open subset of $X$” by “subset of $X$,”
one obtains the notion of a sieve of subspaces of $X$. In this setting
there are also patching theorems. For example, let $X$ be a space and
let $\cC$ be a sieve of subspaces of $X$ generated by a locally finite
closed covering of $X$. Then the functor $E\mapsto E_\cC$ from vector
bundles on $X$ to vector bundles given $\cC$-locally is an equivalence
of categories.

In algebraic geometry, it is also useful to consider “sieves of spaces
over $X$”; this is what we shall see in the next paragraph.

\subsection{Faithfully flat descent}\label{I:1-4}

\subsubsection{}\label{I:1-4-1}

In the setting of schemes, the Zariski topology is not fine enough for the
study of nonlinear problems, and one is led to replace the open immersions
in the preceding definitions by more general morphisms. From this point of
view, descent techniques appear as localization techniques. Thus the
following descent statement can be expressed by saying that the properties
under consideration are local for the faithfully flat topology (a morphism
of schemes is called faithfully flat if it is flat and surjective).

\begin{proposition}\label{I:1-4-2}
Let $A$ be a ring and $B$ a faithfully flat $A$-algebra. Then:
\begin{enumerate}[(i)]
  \item A sequence $\Sigma=(M'\to M\to M'')$ of $A$-modules is exact if
    the sequence $\Sigma_{(B)}$ deduced from it by extension of scalars to
    $B$ is exact.
  \item An $A$-module $M$ is of finite type (resp. of finite presentation,
    flat, locally free of finite rank, invertible (i.e. locally free of rank
    one)) if the $B$-module $M_{(B)}$ has the corresponding property.
\end{enumerate}
\end{proposition}
\begin{proof}
(i) Since the functor $M\mapsto M_{(B)}$ is exact (by flatness of $B$),
it suffices to show that if an $A$-module $N$ is nonzero, then
$N_{(B)}$ is nonzero. If $N$ is nonzero, $N$ contains a nonzero cyclic
submodule $A/\fa$; hence $N_{(B)}$ contains the cyclic submodule
$(A/\fa)_{(B)}=B/\fa B$, which is nonzero by surjectivity of the
structural morphism $\varphi:\spec(B)\to\spec(A)$ (if $V(\fa)$ is
nonempty, then $\varphi^{-1}(V(\fa))=V(\fa B)$ is nonempty).

(ii) For every family $(x_i)$ of elements of $M_{(B)}$, there exists a
finitely generated submodule $M'$ of $M$ such that $M'_{(B)}$ contains
the $x_i$. If $M_{(B)}$ is of finite type and the $x_i$ generate
$M_{(B)}$, then $M'_{(B)}=M_{(B)}$, and hence $M'=M$ and $M$ is of
finite type.
\end{proof}

If $M_{(B)}$ is of finite presentation, the preceding argument yields a
surjection $A^n\to M$. If $N$ is the kernel of this surjection, then the
$B$-module $N_{(B)}$ is of finite type, so $N$ is of finite type and
$M$ is of finite presentation. The assertion for “flat” follows at once
from (i); “locally free of finite rank” means “flat and of finite
presentation,” and the rank can be tested after extension of scalars to
fields.

\subsubsection{}\label{I:1-4-3}

Let $X$ be a scheme and let $\cS$ be a class of $X$-schemes stable under
fibre products over $X$. A class $\cU\subset\cS$ is a \emph{sieve} on
$X$ (relative to $\cS$) if, for every morphism $\varphi:V\to U$ of
$X$-schemes with $U,V\in\cS$ and $U\in\cU$, one has $V\in\cU$.
The sieve \emph{generated} by a family $\{U_i\}$ of $X$-schemes in
$\cS$ is the class of $V\in\cS$ for which there exists a morphism of
$X$-schemes from $V$ to one of the $U_i$.

\subsubsection{}\label{I:1-4-4}

Let $\cU$ be a sieve on $X$. A quasi-coherent module given
$\cU$-locally on $X$ consists of
\begin{enumerate}[\indent a)]
  \item a quasi-coherent module $E_U$ on every $U\in\cU$,
  \item for every $U\in\cU$ and every morphism $\varphi:V\to U$ of
    $X$-schemes in $\cS$, an isomorphism
    $\rho_\varphi:E_V\iso\varphi^*E_U$, such that
  \item if $\psi:W\to V$ is a morphism of $X$-schemes in $\cS$, the
    diagram
    \[\xymatrix{
      E_W \ar[r]^-{\rho_{\varphi\circ\psi}} \ar[dr]_-{\rho_\psi}
        & \psi^*\varphi^* E_U \\
      & \psi^* E_V \ar[u]_-{\psi^*\rho_\varphi}
    }\]
    commutes, that is,
    $\rho_{\varphi\circ\psi}=(\psi^*\rho_\varphi)\circ\rho_\psi$.
\end{enumerate}

If $E$ is a quasi-coherent module on $X$, we write $E_\cU$ for the
module given $\cU$-locally whose value is $\varphi_U^*E$ on
$\varphi:U\to X$ and such that, for every morphism $\psi:V\to U$, the
restriction isomorphism $\rho_\psi$ is the canonical isomorphism
$E_V=(\varphi_U\circ\psi)^*E\iso\psi^*\varphi_U^*E=\psi^*E_U$.

\begin{theorem}\label{I:1-4-5}
Let $\{U_i\}\in\cS$ be a finite family of $X$-schemes flat over $X$
such that $X$ is the union of the images of the $U_i$, and let $\cU$
be the sieve generated by $\{U_i\}$. Then the functor
$E\mapsto E_\cU$ is an equivalence from the category of quasi-coherent
modules on $X$ to the category of quasi-coherent modules given
$\cU$-locally.
\end{theorem}
\begin{proof}
We treat only the case in which $X$ is affine and $\cU$ is generated by
an affine $X$-scheme $U$ faithfully flat over $X$. The reduction to this
case is formal. Put $X=\spec(A)$ and $U=\spec(B)$.

If $U\to X$ admits a section, then $X$ belongs to the sieve $\cU$ and
the assertion is clear. We reduce to this case.

A quasi-coherent module given $\cU$-locally defines modules $M'$,
$M''$, and $M'''$ on $U$, $U\times_XU$, and
$U\times_XU\times_XU$, and isomorphisms
$\rho:p^*M^\bullet\simeq M^\bullet$ for every projection morphism $p$
between these spaces; this is a \emph{cartesian diagram}
\[\xymatrix{
  M^* : M' \ar@<2pt>[r] \ar@<-2pt>[r]
    & M'' \ar[r] \ar@<4pt>[r] \ar@<-4pt>[r]
    & M'''
}\]
over
\[\xymatrix{
  U_* : U
    & U\times_X U \ar@<2pt>[l] \ar@<-2pt>[l]
    & U\times_X U\times_X U \ar[l] \ar@<4pt>[l] \ar@<-4pt>[l] \mbox{.}
}\]
Conversely, $M^*$ determines the module given $\cU$-locally: for
$V\in\cU$ there exists $\varphi:V\to U$, and we put
$M_V=\varphi^*M'$; for $\varphi_1,\varphi_2:V\to U$ there is a natural
identification
$\varphi_1^*M'\simeq(\varphi_1\times\varphi_2)^*M''\simeq\varphi_2^*M'$,
and by using $M'''$ one sees that these identifications are compatible,
so the definition is legitimate. In short, giving a module
$\cU$-locally is equivalent to giving a cartesian diagram $M^*$ over
$U_*$.

In algebraic terms, giving $M^*$ amounts to giving a cartesian diagram of
modules
\[\xymatrix{
  M' \ar@<3pt>[r]|-{\partial_0} \ar@<-3pt>[r]|-{\partial_1}
    & M'' \ar@<6pt>[r]|-{\partial_0} \ar[r]|-{\partial_1} \ar@<-6pt>[r]|-{\partial_2}
    & M'''
}\]
over the diagram of rings
\[\xymatrix{
  B \ar@<3pt>[r]|-{\partial_0} \ar@<-3pt>[r]|-{\partial_1}
    & B\otimes_A B \ar@<6pt>[r]|-{\partial_0} \ar[r]|-{\partial_1} \ar@<-6pt>[r]|-{\partial_2}
    & B\otimes_A B\otimes_A B
}\]
(more precisely, $\partial_i(bm)=\partial_i(b)\cdot\partial_i(m)$, the
usual identities such as $\partial_0\partial_1=\partial_0\partial_0$
hold, and “cartesian” means that the morphisms
$\partial_i:M'\otimes_{B,\partial_i}(B\otimes_A B)\to M''$ and
$M''\otimes_{B\otimes_A B,\partial_i}(B\otimes_A B\otimes_A B)\to M'''$
are isomorphisms).

The functor $E\mapsto E_\cU$ becomes the functor that associates to an
$A$-module $M$ the diagram
\[\xymatrix{
  M^* = (M\otimes_A B \ar@<2pt>[r] \ar@<-2pt>[r]
    & M\otimes_A B\otimes_A B \ar@<-4pt>[r] \ar[r] \ar@<4pt>[r]
    & M\otimes_A B\otimes_A B\otimes_A B)\text{.}
}\]
Its right adjoint is the functor
\[\xymatrix{
  (M' \ar@<2pt>[r] \ar@<-2pt>[r]
    & M'' \ar@<-4pt>[r] \ar[r] \ar@<4pt>[r]
    & M''') \ar@{|->}[r]
    & \ker(M' \ar@<2pt>[r] \ar@<-2pt>[r]
    & M''')\text{.}
}\]
We must prove that the adjunction maps
\[
  M \to\ker(M\otimes_A B\rightrightarrows M\otimes_A B\otimes_A B)
\]
and
\[
  \ker(M'\rightrightarrows M'')\otimes_A B \to M'
\]
are isomorphisms. By (\ref{I:1-4-2}.i), it suffices to prove this after
a faithfully flat base change $A\to A'$ (with $B$ becoming
$B'=B\otimes_AA'$). Taking $A'=B$ reduces us to the case in which
$U\to X$ admits a section.
\end{proof}

\subsection{A special case: Hilbert's theorem 90}\label{I:1-5}

\subsubsection{}\label{I:1-5-1}

Let $k$ be a field, let $k'$ be a Galois extension of $k$, and let
$G=\gal(k'/k)$. Then the homomorphism
\begin{align*}
  k'\otimes_k k' &\to \oplus_{\sigma\in G} k' \\
  x\otimes y     &\mapsto \{x\cdot \sigma(y)\}_{\sigma\in G}
\end{align*}
is bijective.

It follows that giving a module locally for the sieve on $\spec(k)$
generated by $\spec(k')$ is equivalent to giving a $k'$-vector space
equipped with a semilinear action of $G$, that is:
\begin{enumerate}[\indent a)]
  \item a $k'$-vector space $V'$,
  \item for every $\sigma\in G$, an endomorphism $\varphi_\sigma$ of the
    underlying group of $V'$ such that
    $\varphi_\sigma(\lambda v)=\sigma(\lambda)\varphi_\sigma(v)$ for all
    $\lambda\in k'$ and $v\in V'$, satisfying the condition
  \item for all $\sigma,\tau\in G$,
    $\varphi_{\tau\sigma}=\varphi_\tau\circ\varphi_\sigma$.
\end{enumerate}

Let $V={V'}^G$ be the group of invariants under this action of $G$. It
is a $k$-vector space, and by Theorem \ref{I:1-4-5} one has:

\begin{proposition}\label{I:1-5-2}
The inclusion of $V$ in $V'$ defines an isomorphism
$V\otimes_k k'\iso V'$.
\end{proposition}

In particular, if $V'$ has dimension $1$ and $v'\in V'$ is nonzero, then
$\varphi_\sigma$ is determined by the constant
$c(\sigma)\in {k'}^\times$ such that
$\varphi_\sigma(v')=c(\sigma)v'$, and condition c) becomes
\[
  c(\tau\sigma) = c(\tau)\cdot \tau(c(\sigma))\text{.}
\]
By the proposition, there exists a nonzero invariant vector
$v=\mu v'$, $\mu\in {k'}^\times$. Hence, for every $\sigma\in G$,
\[
  c(\sigma) = \mu\cdot \sigma(\mu^{-1})\text{.}
\]
In other words, every $1$-cocycle of $G$ with values in
${k'}^\times$ is a coboundary:

\begin{corollary}\label{I:1-5-3}
One has $\h^1(G,{k'}^\times)=0$.
\end{corollary}

\subsection{Grothendieck topologies}\label{I:1-6}

We now transpose the definitions of the preceding paragraphs to an
abstract setting that encompasses both topological spaces and schemes.

\subsubsection{}\label{I:1-6-1}

Let $\cS$ be a category and let $U$ be an object of $\cS$. A
\emph{sieve} on $U$ is a subset $\cU$ of $\ob(\cS/U)$ such that if
$\varphi:V\to U$ belongs to $\cU$ and $\psi:W\to V$ is a morphism in
$\cS$, then $\varphi\circ\psi:W\to U$ belongs to $\cU$.

If $\{\varphi_i:U_i\to U\}$ is a family of morphisms, the sieve generated
by the $U_i$ is, by definition, the set of morphisms $\varphi:V\to U$
that factor through one of the $\varphi_i$.

If $\cU$ is a sieve on $U$ and $\varphi:V\to U$ is a morphism, the
restriction $\cU_V$ of $\cU$ to $V$ is, by definition, the sieve on
$V$ consisting of the morphisms $\psi:W\to V$ such that
$\varphi\circ\psi:W\to U$ belongs to $\cU$.

\subsubsection{}\label{I:1-6-2}

A \emph{Grothendieck topology} on $\cS$ consists of the choice, for every
object $U$ of $\cS$, of a set $C(U)$ of sieves on $U$, called
covering sieves, such that the following axioms hold:
\begin{enumerate}[\indent a)]
  \item The sieve generated by the identity of $U$ is covering.
  \item If $\cU$ is a covering sieve on $U$ and $V\to U$ is a
    morphism, then the sieve $\cU_V$ is covering.
  \item A locally covering sieve is covering. In other words, if $\cU$
    is a covering sieve on $U$ and $\cU'$ is a sieve on $U$ such that,
    for every $V\to U$ belonging to $\cU$, the sieve $\cU_V'$ is
    covering, then $\cU'$ is covering.
\end{enumerate}

A category equipped with a Grothendieck topology is called a \emph{site}.

\subsubsection{}\label{I:1-6-3}

Given a site $\cS$, a presheaf on $\cS$ is a contravariant functor
$\sF$ from $\cS$ to the category of sets. For every object $U$ of
$\cS$, the elements of $\sF(U)$ are called sections of $\sF$ over
$U$. For every morphism $V\to U$ and every $s\in\sF(U)$, we denote by
$s|V$ ($s$ restricted to $V$) the image of $s$ in $\sF(V)$.

If $\cU$ is a sieve on $U$, a section given $\cU$-locally consists,
for every $V\to U$ belonging to $\cU$, of a section
$s_V\in\sF(V)$ such that, for every morphism $W\to V$,
$s_V|W=s_W$. The presheaf $\sF$ is called a \emph{sheaf} if, for every
object $U$ of $\cS$, every covering sieve $\cU$ on $U$, and every
section $\{s_V\}$ given $\cU$-locally, there exists a unique section
$s\in\sF(U)$ such that $s|V=s_V$ for every $V\to U$ belonging to
$\cU$.

\emph{Abelian sheaves} are defined similarly by replacing the category of
sets with the category of abelian groups. One shows that the category of
abelian sheaves on $\cS$ is an abelian category with enough injectives. A
sequence $\sF\xrightarrow f\sG\xrightarrow g\sH$ of sheaves is exact if,
for every object $U$ of $\cS$ and every $s\in\sG(U)$ such that
$g(s)=0$, there exists locally a $t$ such that $f(t)=s$; i.e. there
exists a covering sieve $\cU$ on $U$ and, for every $V\in\cU$, a
section $t_V$ of $\sF$ on $V$ such that $f(t_V)=s|V$.

\subsubsection{Examples}\label{I:1-6-4}

We have seen two examples above.

\begin{enumerate}[\indent a)]
  \item Let $X$ be a topological space and let $\cS$ be the category
    whose objects are the open subsets of $X$ and whose morphisms are the
    natural inclusions. The Grothendieck topology on $\cS$ corresponding
    to the usual topology of $X$ is the one for which a sieve $\cU$ on
    an open subset $U$ of $X$ is covering if the union of the open
    subsets belonging to the sieve equals $U$. It is clear that the
    category of sheaves on $\cS$ is equivalent to the category of sheaves
    on $X$ in the usual sense.
  \item Let $X$ be a scheme and let $\cS$ be the category of schemes over
    $X$. The fpqc (faithfully flat quasi-compact) topology on $\cS$ is
    the Grothendieck topology for which a sieve on an $X$-scheme $U$ is
    covering if it is generated by a finite family of flat morphisms whose
    images cover $U$.
\end{enumerate}

\subsubsection{Cohomology}\label{I:1-6-5}

We shall always assume that the category $\cS$ has a final object $X$.
The global sections of an abelian sheaf $\sF$, denoted by $\Gamma\sF$ or
$\h^0(X,\sF)$, form the group $\sF(X)$. The functor
$\sF\mapsto\Gamma\sF$ is a left exact functor from the category of
abelian sheaves on $\cS$ to the category of abelian groups; its derived
functors (or satellites) are denoted by $\h^i(X,-)$. These cohomology
groups represent the obstructions to passing from the local to the global.
By definition, if $0\to\sF\to\sG\to\sH\to0$ is an exact sequence of
abelian sheaves, then there is a long exact cohomology sequence:
\[\xymatrix{
  0 \ar[r]
    & \h^0(X,\sF) \ar[r]
    & \h^0(X,\sG) \ar[r]
    & \h^0(X,\sH) \ar[r]
    & \h^1(X,\sF) \ar[r]
    & \cdots \\
  \ldots \ar[r]
    & \h^n(X,\sF) \ar[r]
    & \h^n(X,\sG) \ar[r]
    & \h^n(X,\sH) \ar[r]
    & \h^{n+1}(X,\sF) \ar[r]
    & \cdots
}\]

\subsubsection{}\label{I:1-6-6}

Given an abelian sheaf $\sF$ on $\cS$, an \emph{$\sF$-torsor} is a
sheaf $\sG$ equipped with an action $\sF\times\sG\to\sG$ of $\sF$
such that locally (after restriction to all the objects of a covering
sieve of the final object $X$), $\sG$ with its $\sF$-action is
isomorphic to $\sF$ with the canonical translation action
$\sF\times\sF\to\sF$.

One can show that $\h^1(X,\sF)$ may be interpreted as the set of
isomorphism classes of $\sF$-torsors.

\section{The étale topology}\label{I:2}

We specialize the definitions of the preceding chapter to the étale
topology of a scheme $X$ (\ref{I:2-1}, \ref{I:2-2}, \ref{I:2-3}). When
$X$ is the spectrum of a field $K$, the corresponding cohomology agrees
with the Galois cohomology of $K$ (\ref{I:2-4}).

\subsection{The étale topology}\label{I:2-1}

We begin by recalling some facts about étale morphisms.

\begin{definition}\label{I:2-1-1}
Let $A$ be a (commutative) ring. An $A$-algebra $B$ is called étale if
$B$ is an $A$-algebra of finite presentation and the following equivalent
conditions hold:
\begin{enumerate}[\indent a)]
  \item For every $A$-algebra $C$ and every square-zero ideal $J$ of $C$,
    the canonical map
    \[
      \hom_{A\textnormal{-alg}}(B,C) \to \hom_{A\textnormal{-alg}}(B,C/J)
    \]
    is a bijection.
  \item $B$ is a flat $A$-module and $\Omega_{B/A}=0$ (where
    $\Omega_{B/A}$ denotes the module of relative differentials).
  \item Let $B=A[X_1,\dotsc,X_n]/I$ be a presentation of $B$. Then, for
    every prime ideal $\fp$ of $A[X_1,\dotsc,X_n]$ containing $I$, there
    exist polynomials $P_1,\dotsc,P_n\in I$ such that $I_\fp$ is generated
    by the images of $P_1,\dotsc,P_n$ and
    $\det(\partial P_i/\partial X_j)\notin\fp$.
\end{enumerate}
\end{definition}
(cf. \cite[I]{sga1} or \cite[V]{ra70})
% [cf. SGA I, Exposé I, or M. {\sc Raynoud}, \emph{Henselian Local Rings}, Chapter V].

A morphism of schemes $f:X\to S$ is called \emph{étale} if, for every
$x\in X$, there exist an affine open neighbourhood $U=\spec(A)$ of
$f(x)$ and an affine open neighbourhood $V=\spec(B)$ of $x$ in
$X\times_SU$ such that $B$ is an étale $A$-algebra.

\subsubsection{Examples}\label{I:2-1-2}
\begin{enumerate}[\indent a)]
  \item If $A$ is a field, an $A$-algebra $B$ is étale if and only if it
    is a finite product of separable extensions of $A$.
  \item If $X$ and $S$ are schemes of finite type over $\dC$, a morphism
    $f:X\to S$ is étale if and only if its analytification
    $f^{\text{an}}:X^{\text{an}}\to S^{\text{an}}$ is a local
    isomorphism.
\end{enumerate}

\subsubsection{Consequences}\label{I:2-1-3}
\begin{enumerate}[\indent a)]
  \item (base change) If $f:X\to S$ is an étale morphism, then so is
    $f_{S'}:X\times_SS'\to S'$ for every morphism $S'\to S$.
  \item If $f:X\to S$ and $g:Y\to S$ are two étale morphisms, then every
    $S$-morphism from $X$ to $Y$ is étale.
  \item (descent) Let $f:X\to S$ be a morphism. If there exists a
    faithfully flat morphism $S'\to S$ such that
    $f_{S'}:X\times_SS'\to S'$ is étale, then $f$ is étale.
\end{enumerate}

\subsubsection{}\label{I:2-1-4}

Let $X$ be a scheme, and let $\cS$ be the category of étale
$X$-schemes. By (\ref{I:2-1-3}.c), every morphism in $\cS$ is étale.
The \emph{étale topology} on $\cS$ is the topology for which a sieve on
$U$ is covering if it is generated by a finite family of morphisms
$\varphi_i:U_i\to U$ such that the union of the images of the
$\varphi_i$ covers $U$. The \emph{étale site} of $X$, denoted by
$\et X$, is the site defined by $\cS$ with the étale topology.

\subsection{Examples of sheaves}\label{I:2-2}

\subsubsection{Constant sheaf}\label{I:2-2-1}

Let $C$ be an abelian group and, for simplicity, suppose that $X$ is
noetherian. We denote by $\const C_X$ (or simply $C$ if there is no
ambiguity) the sheaf defined by $U\mapsto C^{\pi_0(U)}$, where
$\pi_0(U)$ is the (finite) set of connected components of $U$. The most
important case will be $C=\dZ/n$. Thus, by definition,
\[
  \h^0(X,\dZ/n)=\left(\dZ/n\right)^{\pi_0(X)}\text{.}
\]
Moreover, $\h^1(X,\dZ/n)$ is the set of isomorphism classes of
$\dZ/n$-torsors (\ref{I:1-6-6}), in other words, of finite Galois étale
coverings of $X$ with group $\dZ/n$. In particular, if $X$ is connected
and $\pi_1(X)$ is its fundamental group for a chosen base point, then
\[
  \h^1(X,\dZ/n) = \hom(\pi_1(X),\dZ/n)\text{.}
\]

\subsubsection{Multiplicative group}\label{I:2-2-2}

We denote by $\dG_{m,X}$ (or $\dG_m$ if there is no ambiguity) the sheaf
defined by $U\mapsto\Gamma(U,\sO_U^\times)$; it is indeed a sheaf by the
faithfully flat descent theorem (\ref{I:1-4-5}). By definition,
\[
  \h^0(X,\dG_m) = \h^0(X,\sO_X)^\times\text{;}
\]
in particular, if $X$ is reduced, connected, and proper over an
algebraically closed field $k$, then
\[
  \h^0(X,\dG_m) = k^\times\text{.}
\]

\begin{proposition}\label{I:2-2-3}
There is an isomorphism
\[
  \h^1(X,\dG_m) = \pic(X)\text{,}
\]
where $\pic(X)$ is the group of isomorphism classes of invertible sheaves
on $X$.
\end{proposition}
\begin{proof}
Let $*$ be the functor that associates to an invertible sheaf $\sL$ on
$X$ the following presheaf $\sL^*$ on $\et X$: for
$\varphi:U\to X$ étale,
\[
  \sL^*(U)=\operatorname{Isom}_U(\sO_U,\varphi^*\sL)\text{.}
\]
% The original references here are just to (4.2) and (4.5), but the context
% makes it seem that they refer to the results in Chapter 1. Likewise a little
% later.
By (\ref{I:1-4-2}.i) and (\ref{I:1-4-5}) (full faithfulness), this
presheaf is a sheaf; indeed, it is a $\dG_m$-torsor. One checks
immediately that
\begin{enumerate}[\indent a)]
  \item the functor $*$ is compatible with (étale) localization;
  \item it induces an equivalence from the category of trivial invertible
    sheaves (i.e. isomorphic to $\sO_X$) to the category of trivial
    $\dG_m$-torsors: $\sL$ is trivial if and only if $\sL^*$ is trivial.
\end{enumerate}

Moreover, by (\ref{I:1-4-2}.ii) and (\ref{I:1-4-5}),
\begin{enumerate}[\indent a)]
\setcounter{enumi}{2}
  \item the notion of an invertible sheaf is local for the étale topology.
\end{enumerate}

It follows formally from a), b), and c) that $*$ is an equivalence between
the category of invertible sheaves on $X$ and that of $\dG_m$-torsors on
$\et X$; it induces the desired isomorphism. The inverse equivalence is
constructed as follows. If $\sT$ is a $\dG_m$-torsor, there exists a
finite étale covering $\{U_i\}$ of $X$ such that the torsors
$\sT/U_i$ are trivial; $\sT$ is then trivial on every $V$ étale over
$X$ belonging to the sieve $\cU\subset\et X$ generated by $\{U_i\}$.
On every $V\in\cU$, $\sT|V$ corresponds to an invertible sheaf $\sL_V$
(by b)), and the $\sL_V$ form an invertible sheaf given $\cU$-locally,
$\sL_\cU$ (by a)). By c), the latter comes from an invertible sheaf
$\sL(\sT)$ on $X$, and $\sT\mapsto\sL(\sT)$ is the desired inverse
of $*$.
\end{proof}

\subsubsection{Roots of unity}\label{I:2-2-4}

For every integer $n>0$, the sheaf of $n$th roots of unity, denoted by
$\dmu_n$, is the kernel of the $n$th-power map on $\dG_m$.
If $X$ is a scheme over a separably closed field $k$ and $n$ is
invertible in $k$, the choice of a primitive $n$th root of unity
$\zeta\in k$ defines an isomorphism $i\mapsto\zeta^i$ from $\dZ/n$ to
$\dmu_n$.

The relation between cohomology with coefficients in $\dmu_m$ and
cohomology with coefficients in $\dG_m$ is given by the cohomology exact
sequence deduced from the following theorem.

\begin{theorem}[Kummer theory]\label{I:2-2-5}
If $n$ is invertible on $X$, the $n$th-power map on $\dG_m$ is an
epimorphism of sheaves. Hence there is an exact sequence
\[
  0 \to \dmu_n \to \dG_m \to \dG_m \to 0\text{.}
\]
\end{theorem}
\begin{proof}
Let $U\to X$ be an étale morphism and let
$a\in\dG_m(U)=\Gamma(U,\sO_U^\times)$. Since $n$ is invertible on $U$,
the equation $T^n-a=0$ is separable; in other words,
$U'=\spec\left(\sO_U[T]/(T^n-a)\right)$ is étale over $U$. Moreover,
$U'\to U$ is surjective and this element has an $n$th root on $U'$, which proves
the result.
\end{proof}

\subsection{Fibres, direct images}\label{I:2-3}

\subsubsection{}\label{I:2-3-1}

A \emph{geometric point} of $X$ is a morphism $\bar x\to X$, where
$\bar x$ is the spectrum of a separably closed field $k(\bar x)$. By
abuse of notation, we denote it by $\bar x$, leaving the morphism
$\bar x\to X$ implicit. If $x$ is the image of $\bar x$ in $X$, we say
that $\bar x$ is centred at $x$. If $k(\bar x)$ is an algebraic extension
of the residue field $k(x)$, then $\bar x$ is called an
\emph{algebraic} geometric point of $X$.

An \emph{étale neighbourhood} of $\bar x$ is a commutative diagram
\[\xymatrix{
  & U \ar[d] \\
  \bar x \ar[ur] \ar[r]
  & X\text{,}
}\]
where $U\to X$ is étale.

The \emph{strict localization} of $X$ at $\bar x$ is the ring
$\sO_{X,\bar x}=\varinjlim\Gamma(U,\sO_U)$, where the direct limit is
taken over the étale neighbourhoods of $\bar x$. It is a strictly henselian
local ring whose residue field is the separable closure in $k(\bar x)$
of the residue field $k(x)$ of $X$ at $x$. It plays the role of the local
ring for the étale topology.

\subsubsection{}\label{I:2-3-2}

Given a sheaf $\sF$ on $\et X$, the \emph{fibre} of $\sF$ at $\bar x$
is the set (resp. group, \dots)
$\sF_{\bar x}=\varinjlim\sF(U)$, where the direct limit is again taken
over the étale neighbourhoods of $\bar x$.

A homomorphism of sheaves $\sF\to\sG$ is a mono-/epi-/isomorphism if and
only if the induced morphisms $\sF_{\bar x}\to\sG_{\bar x}$ on fibres
have the corresponding property at every geometric point of $X$. If $X$
is of finite type over an algebraically closed field, it suffices that
this hold at the rational points of $X$.

\subsubsection{}\label{I:2-3-3}

If $f:X\to Y$ is a morphism of schemes and $\sF$ is a sheaf on
$\et X$, the \emph{direct image} $f_*\sF$ of $\sF$ under $f$ is the
sheaf on $\et Y$ defined by
$f_*\sF(V)=\sF(X\times_YV)$ for every $V$ étale over $Y$. The functor
$f_*:\sh(\et X)\to\sh(\et Y)$ is left exact. Its right derived functors
$\eR^qf_*$ are called the higher direct images. If $\bar y$ is a
geometric point of $Y$, then
\[
  (\eR^q f_* \sF)_{\bar y} = \varinjlim \h^q(V\times_Y X,\sF)\text{,}
\]
where the direct limit is taken over the étale neighbourhoods $V$ of
$\bar y$.

Let $\sO_{Y,\bar y}$ be the strict localization of $Y$ at $\bar y$,
let $\widetilde Y=\spec(\sO_{Y,\bar y})$, and let
$\widetilde X=X\times_Y\widetilde Y$. The sheaf $\sF$ may be extended
to $\et{\widetilde X}$ (this is a special case of the general notion of
inverse image) as follows. Let $\widetilde U$ be a scheme étale over
$\widetilde X$. Then there exist an étale neighbourhood $V$ of $\bar y$
and a scheme $U$ étale over $X\times_YV$ such that
$\widetilde U=U\times_V\widetilde Y$; set
\[
  \sF(\widetilde U)=\varinjlim \sF\left(U\times_V V'\right)\text{,}
\]
where the direct limit is taken over the étale neighbourhoods $V'$ of
$\bar y$ that dominate $V$. With this definition,
\[
  \left(\eR^q f_* \sF\right)_{\bar y} = \h^q(\widetilde X,\sF)\text{.}
\]

The functor $f_*$ has a left adjoint $f^*$, the “inverse image” functor.
If $\bar x$ is a geometric point of $X$ and $f(\bar x)$ its image in
$Y$, then $(f^*\sF)_{\bar x}=\sF_{f(\bar x)}$. This formula shows that
$f^*$ is exact. Consequently, $f_*$ sends injective sheaves to injective
sheaves, and the spectral sequence for the composite functor
$\Gamma\circ f_*$ (resp. $g_*f_*$) gives the following theorem.

\begin{theorem}[Leray spectral sequence]\label{I:2-3-4}
Let $\sF$ be an abelian sheaf on $\et X$ and let $f:X\to Y$ be a
morphism of schemes (resp. let
$X\xrightarrow fY\xrightarrow gZ$ be morphisms of schemes). There is a
spectral sequence
\begin{align*}
  E_2^{pq} &= \h^p(Y,\eR^q f_* \sF) \Rightarrow \h^{p+q}(X,\sF) \\
  \text{(resp.}\qquad E_2^{pq} &= \eR^p g_* \eR^q f_* \sF \Rightarrow \eR^{p+q}(gf)_* \sF\text{).}
\end{align*}
\end{theorem}

\begin{corollary}\label{I:2-3-5}
If $\eR^qf_*\sF=0$ for every $q>0$, then
$\h^p(Y,f_*\sF)=\h^p(X,\sF)$ (resp.
$\eR^pg_*(f_*\sF)=\eR^p(gf)_*\sF$) for every $p\geqslant0$.
\end{corollary}

This applies in particular in the following case:

\begin{proposition}\label{I:2-3-6}
Let $f:X\to Y$ be a finite morphism (or even, by passage to the limit,
an integral morphism), and let $\sF$ be an abelian sheaf on $X$. Then
$\eR^qf_*\sF=0$ for every $q>0$.
\end{proposition}

Indeed, let $\bar y$ be a geometric point of $Y$, let $\widetilde Y$ be
the spectrum of the strict localization of $Y$ at $\bar y$, and let
$\widetilde X=X\times_Y\widetilde Y$. By the preceding discussion, it
suffices to prove that $\h^q(\widetilde X,\sF)=0$ for every $q>0$. But
$\widetilde X$ is the spectrum of a product of strictly henselian local
rings (cf. \cite[I]{ra70}), and the functor
$\Gamma(\widetilde X,-)$ is exact because every surjective étale
$\widetilde X$-scheme admits a section. This proves the assertion.

\subsection{Galois cohomology}\label{I:2-4}

For $X=\spec(K)$, the spectrum of a field, we shall see that étale
cohomology can be identified with Galois cohomology.

\subsubsection{}\label{I:2-4-1}

We begin with a topological analogy. If $K$ is the function field of an
integral affine algebraic variety $Y=\spec(A)$ over $\dC$, then
$K=\varinjlim_{f\in A}A[1/f]$.

In other words, $X=\varprojlim U$, where $U$ runs over the open subsets
of $Y$. It is known that there exist arbitrarily small Zariski open
subsets that are $K(\pi,1)$ spaces for the classical topology. It is
therefore natural to regard $\spec(K)$ itself as a $K(\pi,1)$, where
$\pi$ is the fundamental group (in the algebraic sense) of $X$, namely
the Galois group of $\bar K/K$, where $\bar K$ is a separable closure of
$K$.

\subsubsection{}\label{I:2-4-2}

More precisely, let $K$ be a field, let $\bar K$ be a separable closure
of $K$, and let $G=\gal(\bar K/K)$ be the topological Galois group. To
every finite étale $K$-algebra $A$ (a finite product of separable
extensions of $K$), associate the finite set $\hom_K(A,\bar K)$. The
Galois group $G$ acts on this set through a discrete (hence finite)
quotient. If $A=K[T]/(F)$, this set is identified with the set of roots
of $F$ in $\bar K$. Galois theory, in the form given to it by
Grothendieck, says:

\begin{proposition}\label{I:2-4-3}
The functor
\[
  \left\{\begin{array}{c}
          \textnormal{$K$-algebras} \\
          \textnormal{finite étale}
        \end{array}\right\}
  \to
  \left\{\begin{array}{c}
          \textnormal{finite sets on which} \\
          \textnormal{$G$ acts continuously}
        \end{array}\right\}
\]
that associates $\hom_K(A,\bar K)$ to an étale algebra $A$ is an
anti-equivalence of categories.
\end{proposition}

This yields an analogous description of sheaves for the étale topology on
$\spec(K)$:

\begin{proposition}\label{I:2-4-4}
The functor
\[
  \left\{\begin{array}{c}
          \textnormal{étale sheaves} \\
          \textnormal{on $\spec(K)$}
        \end{array}\right\}
  \to
  \left\{\begin{array}{c}
          \textnormal{sets on which} \\
          \textnormal{$G$ acts continuously}
        \end{array}\right\}
\]
that associates to a sheaf $\sF$ its fibre $\sF_{\bar K}$ at the
geometric point $\spec(\bar K)$ is an equivalence of categories.
\end{proposition}

The group $G$ is said to act continuously on a set $E$ if the stabilizer
of every element of $E$ is an open subgroup of $G$. The inverse functor
is described in the evident way: let $A$ be a finite étale
$K$-algebra, let $U=\spec(A)$, and let
$U(\bar K)=\hom_K(A,\bar K)$ be the corresponding $G$-set; then
$\sF(U)=\hom_{G\text{-set}}(U(\bar K),\sF_{\bar K})$.

In particular, if $X=\spec(K)$, then $\sF(X)=\sF_{\bar K}^G$. Restricting
to abelian sheaves and passing to derived functors gives canonical
isomorphisms
\[
  \h^q(\et X,\sF) = \h^q(G,\sF_{\bar K})
\]

\subsubsection{Examples}\label{I:2-4-5}

\begin{enumerate}[\indent a)]
  \item The constant sheaf $\dZ/n$ corresponds to $\dZ/n$ with the
    trivial action of $G$.
  \item The sheaf $\dmu_n$ of $n$th roots of unity corresponds to the
    group $\dmu_n(\bar K)$ of $n$th roots of unity in $\bar K$, with
    the natural action of $G$.
  \item The sheaf $\dG_m$ corresponds to the group $\bar K^\times$ with
    the natural action of $G$.
\end{enumerate}

\section{The cohomology of curves}\label{I:3}

For topological spaces, unwindings using the K\"unneth formula and
simplicial decompositions reduce the calculation of cohomology to the
interval $I=[0,1]$, for which $\h^0(I,\dZ)=\dZ$ and
$\h^q(I,\dZ)=0$ for $q>0$.

In our case, the unwindings lead to more complicated objects, namely curves
over an algebraically closed field; in this chapter we calculate their
cohomology. The situation is more complicated than in the topological case,
because the cohomology groups vanish only for $q>2$. The essential
ingredient in the calculations is the vanishing of the Brauer group of the
function field of such a curve (Tsen's theorem, \ref{I:3-2}).

\subsection{The Brauer group}\label{I:3-1}

We first recall its classical definition.

\begin{definition}\label{I:3-1-1}
Let $K$ be a field and let $A$ be a finite-dimensional $K$-algebra. The
algebra $A$ is called a central simple algebra over $K$ if the following
equivalent conditions hold:
\begin{enumerate}[\indent a)]
  \item $A$ has no nontrivial two-sided ideal, and its centre is $K$.
  \item There exists a finite Galois extension $K'/K$ such that
    $A_{K'}=A\otimes_KK'$ is isomorphic to an algebra of square matrices
    over $K'$.
  \item $A$ is $K$-isomorphic to an algebra of square matrices over a
    division algebra with centre $K$.
\end{enumerate}
\end{definition}

Two such algebras are called equivalent if the division algebras associated
with them by c) are $K$-isomorphic. If the algebras have the same
dimension, this is equivalent to their being $K$-isomorphic. Tensor product
induces, on passing to the quotient, an abelian group structure on the set
of equivalence classes. This group is classically called the
\emph{Brauer group} of $K$ and is denoted by $\br(K)$.

\subsubsection{}\label{I:3-1-2}

We denote by $\br(n,K)$ the set of $K$-isomorphism classes of
$K$-algebras $A$ such that there exists a finite Galois extension $K'$
of $K$ for which $A_{K'}$ is isomorphic to the algebra $M_n(K')$ of
$n\times n$ square matrices over $K'$. By definition, $\br(K)$ is the
union of the subsets $\br(n,K)$ for $n\in\dN$. Let $\bar K$ be an
algebraic closure of $K$ and let $G=\gal(\bar K/K)$. The set $\br(n,K)$
is the set of “forms” of $M_n(\bar K)$, and hence is canonically
isomorphic to $\h^1\left(G,\aut(M_n(\bar K))\right)$.

Every automorphism of $M_n(\bar K)$ is inner. Consequently,
$\aut(M_n(\bar K))$ is identified with the projective linear group
$\pgl(n,\bar K)$, and there is a canonical bijection
\[
  \theta_n : \br(n,K) \iso \h^1\left(G,\pgl(n,\bar K)\right)\text{.}
\]

On the other hand, the exact sequence
\begin{equation*}\tag{3.1.2.1}\label{I:eq:3-1-2-1}
  1 \to \bar K^\times \to \gl(n,\bar K) \to \pgl(n,\bar K) \to 1\text{,}
\end{equation*}
defines a coboundary operator
\[
  \Delta_n : \h^1\left(G,\pgl(n,\bar K)\right) \to \h^2(G,\bar K^\times)\text{.}
\]

Composing $\theta_n$ and $\Delta_n$ gives a map
\[
  \delta_n : \br(n,K) \to \h^2(G,\bar K^\times)\text{.}
\]
It is readily checked that the maps $\delta_n$ are compatible and define
a group homomorphism
\[
  \delta : \br(K) \to \h^2(G,\bar K^\times)\text{.}
\]

\begin{proposition}\label{I:3-1-3}
The homomorphism $\delta:\br(K)\to\h^2(G,\bar K^\times)$ is bijective.
\end{proposition}

This follows from the next two lemmas.

\begin{lemma}\label{I:3-1-4}
The map
$\Delta_n:\h^1\left(G,\pgl(n,\bar K)\right)\to
\h^2(G,\bar K^\times)$ is injective.
\end{lemma}

By \cite{se94}, Corollary to Proposition I-44, it suffices to check that
whenever the exact sequence (\ref{I:eq:3-1-2-1}) is twisted by an element
of $\h^1\left(G,\pgl(n,\bar K)\right)$, the $\h^1$ of the middle group
is trivial. This middle group is the group of $\bar K$-points of the
multiplicative group of a central simple algebra $A$ of rank $n^2$ over
$K$. To prove that $\h^1(G,A_{\bar K}^\times)=0$, interpret $A^\times$
as the automorphism group of the free $A$-module $L$ of rank $1$, and
$\h^1$ as the set of “forms” of $L$---$A$-modules of rank $n^2$ over
$K$, which are automatically free.

\begin{lemma}\label{I:3-1-5}
Let $\alpha\in\h^2(G,\bar K^\times)$, let $K'$ be a finite extension of
$K$ contained in $\bar K$, let $n=[K':K]$, and let
$G'=\gal(\bar K/K')$. If the image of $\alpha$ in
$\h^2(G',\bar K^\times)$ is zero, then $\alpha$ belongs to the image of
$\Delta_n$.
\end{lemma}

First note that
\[
  \h^2(G',\bar K^\times) \simeq \h^2\left(G,(\bar K\otimes_K K')^\times\right)\text{.}
\]
(Geometrically, if $x=\spec(K)$, $x'=\spec(K')$, and
$\pi:x'\to x$ is the canonical morphism, then
$\eR^q\pi_*(\dG_{m,x'})=0$ for $q>0$ and consequently
$\h^q(x',\dG_{m,x'})\simeq\h^q(x,\pi_*\dG_{m,x'})$ for
$q\geqslant0$.)

Moreover, the choice of a basis of $K'$ as a vector space over $K$
defines a homomorphism
\[
  (\bar K\otimes_K K')^\times \to \gl(n,\bar K)
\]
that sends an element $x$ to the endomorphism of $\bar K\otimes_KK'$
given by multiplication by $x$. There is then a commutative diagram with
exact rows:
\[\xymatrix{
  1 \ar[r]
    & \bar K^\times \ar[r] \ar@{=}[d]
    & (\bar K\otimes_K K')^\times \ar[r] \ar[d]
    & (\bar K\otimes_K K')^\times / \bar K^\times \ar[r] \ar[d]
    & 1 \\
  1 \ar[r]
    & \bar K^\times \ar[r]
    & \gl(n,\bar K) \ar[r]
    & \pgl(n,\bar K) \ar[r]
    & 1
}\]
The lemma follows from the commutative diagram obtained by passing to
cohomology:
\[\xymatrix{
  \h^1\left(G,(\bar K\otimes_K K')^\times/\bar K^\times\right) \ar[r] \ar[d]
    & \h^2(G,\bar K^\times) \ar[r] \ar@{=}[d]
    & \h^2\left(G,(\bar K\otimes_K K')^\times\right) \\
  \h^1\left(G,\pgl(n,\bar K)\right) \ar[r]^-{\Delta_n}
    & \h^2(G,\bar K^\times) \text{.}
}\]

\begin{proposition}\label{I:3-1-6}
Let $K$ be a field, let $\bar K$ be an algebraic closure of $K$, and let
$G=\gal(\bar K/K)$. Suppose that $\br(K')=0$ for every finite extension
$K'$ of $K$. Then:
\begin{enumerate}[\indent i)]
  \item $\h^q(G,\bar K^\times)=0$ for every $q>0$.
  \item $\h^q(G,\sF)=0$ for every torsion $G$-module $\sF$ and every
    $q\geqslant2$.
\end{enumerate}
\end{proposition}

(For the proof, cf. \cite{se94}.)

\subsection{Tsen's theorem}\label{I:3-2}

\begin{definition}\label{I:3-2-1}
A field $K$ is called $C_1$ if every nonconstant homogeneous polynomial
$f(x_1,\dotsc,x_n)$ of degree $d<n$ has a nontrivial zero.
\end{definition}

\begin{proposition}\label{I:3-2-2}
If a field $K$ is $C_1$, then $\br(K)=0$.
\end{proposition}

We must show that every division algebra $D$ with centre $K$ and finite
over $K$ is equal to $K$. Let $r^2$ be the degree of $D$ over $K$ and
let $\operatorname{Nrd}:D\to K$ be the reduced norm.

(Locally for the étale topology on $K$, $D$ is isomorphic---not
canonically---to a matrix algebra $M_r$, and the reduced norm agrees with
the determinant map. This is well defined independently of the chosen
isomorphism between $D$ and $M_r$, because every automorphism of $M_r$
is inner and similar matrices have the same determinant. This map, defined
locally for the étale topology, descends, by its local uniqueness, to a map
$\operatorname{Nrd}:D\to K$.)

The only zero of $\operatorname{Nrd}$ is the zero element of $D$, since
for $x\ne0$ one has
$\operatorname{Nrd}(x)\cdot\operatorname{Nrd}(x^{-1})=1$. On the other
hand, if $\{e_1,\dotsc,e_{r^2}\}$ is a basis of $D$ over $K$ and
$x=\sum x_ie_i$, then $\operatorname{Nrd}(x)$ is a homogeneous
polynomial $\operatorname{Nrd}(x_1,\dotsc,x_{r^2})$ of degree $r$
(this is clear locally for the étale topology). Since $K$ is $C_1$, one
has $r^2\leqslant r$, that is, $r=1$ and $D=K$.

\begin{theorem}[Tsen]\label{I:3-2-3}
Let $k$ be an algebraically closed field and let $K$ be an extension of
$k$ of transcendence degree $1$. Then $K$ is $C_1$.
\end{theorem}

First suppose that $K=k(X)$. Let
\[
  f(T) = \sum a_{i_1,\dotsc i_n} T_1^{i_1} \dotsm T_n^{i_n}
\]
be a homogeneous polynomial of degree $d<n$ with coefficients in
$k(X)$. After multiplying the coefficients by a common denominator, we
may suppose that they lie in $k[X]$. Let
$\delta=\sup\deg(a_{i_1\dotsc i_n})$. We seek a nontrivial zero in
$k[X]$ by the method of undetermined coefficients, writing each $T_i$
($i=1,\dotsc,n$) as a polynomial of degree $N$ in $X$. The equation
$f(T)=0$ then becomes a system of homogeneous equations in the
$n\times(N+1)$ coefficients of the polynomials $T_i(X)$, expressing the
vanishing of the coefficients of the polynomial in $X$ obtained by
replacing $T_i$ with $T_i(X)$. This polynomial has degree at most
$\delta+Nd$, so there are $\delta+Nd+1$ equations in
$n\times(N+1)$ variables. Since $k$ is algebraically closed, this system
has a nontrivial solution if $n(N+1)>Nd+\delta+1$, which holds for
sufficiently large $N$ when $d<n$.

To prove the theorem in general, it plainly suffices to prove it when
$K$ is a finite extension of the purely transcendental extension
$k(X)$ of $k$. Let $f(T)=f(T_1,\dotsc,T_n)$ be a homogeneous
polynomial of degree $d<n$ with coefficients in $K$. Let
$s=[K:k(X)]$ and let $e_1,\dotsc,e_s$ be a basis of $K$ over $k(X)$.
Introduce $sn$ new variables $U_{ij}$ such that
$T_i=\sum U_{ij}e_j$. For $f(T)$ to have a nontrivial zero in $K$, it
suffices that the polynomial $g(U_{ij})=N_{K/k(X)}(f(T))$ have a nontrivial
zero in $k(X)$. But $g$ is a polynomial of degree $sd$ in $sn$
variables, which proves the result.

\begin{corollary}\label{I:3-2-4}
Let $k$ be an algebraically closed field and let $K$ be an extension of
$k$ of transcendence degree $1$. Then the étale cohomology groups
$\h^q(\spec(K),\dG_m)$ vanish for every $q>0$.
\end{corollary}

\subsection{Cohomology of smooth curves}

From now on, unless expressly stated otherwise, all cohomology groups are
étale cohomology groups.

\begin{proposition}\label{I:3-3-1}
Let $k$ be an algebraically closed field and let $X$ be a connected
nonsingular projective curve over $k$. Then
\begin{align*}
  \h^0(X,\dG_m) &= k^\times \text{,} \\
  \h^1(X,\dG_m) &= \pic(X) \text{,} \\
  \h^q(X,\dG_m) &= 0 \text{ for $q\geqslant 2$.}
\end{align*}
\end{proposition}

Let $\eta$ be the generic point of $X$, let $j:\eta\to X$ be the
canonical morphism, and let $\dG_{m,\eta}$ be the multiplicative group of
the function field $K(X)$. For every closed point $x$ of $X$, let
$i_x:x\to X$ be the canonical immersion and let $\dZ_x$ be the constant
sheaf with value $\dZ$ on $x$. Thus $j_*\dG_{m,\eta}$ is the sheaf of
nonzero meromorphic functions on $X$ and
$\oplus_{x\in X}i_{x*}\dZ_x$ is the sheaf of divisors; hence there is an
exact sequence of sheaves:
\begin{equation*}\tag{3.3.1.1}\label{I:eq:3-3-1-1}
\xymatrix{
  0 \ar[r]
    & \dG_m \ar[r]
    & j_* \dG_{m,\eta} \ar[r]^-{\dv}
    & \displaystyle\bigoplus_{x\in X} i_{x*} \dZ_x \ar[r]
    & 0\text{.}
}
\end{equation*}

\begin{lemma}\label{I:3-3-2}
One has $\eR^qj_*\dG_{m,\eta}=0$ for every $q>0$.
\end{lemma}

It suffices to show that the fibre of this sheaf at every closed point
$x$ of $X$ is zero. If $\tilde\sO_{X,x}$ is the henselization of $X$ at
$x$ and $K$ is the fraction field of $\tilde\sO_{X,x}$, then
\[
  \spec(K) = \eta\times_X \spec(\tilde\sO_{X,x})\text{,}
\]
and hence $(\eR^qj_*\dG_{m,\eta})_x=\h^q(\spec(K),\dG_m)$.

Now $K$ is an algebraic extension of $k(X)$, hence an extension of $k$
of transcendence degree $1$; the lemma follows from (\ref{I:3-2-4}).

\begin{lemma}\label{I:3-3-3}
One has $\h^q(X,j_*\dG_{m,\eta})=0$ for every $q>0$.
\end{lemma}

Indeed, (\ref{I:3-3-2}) and the Leray spectral sequence for $j$ give
\[
  \h^q(X, j_* \dG_{m,\eta}) = \h^q(\eta,\dG_{m,\eta})
\]
for every $q\geqslant0$, and the right-hand side vanishes for $q>0$ by
(\ref{I:3-2-4}).

\begin{lemma}\label{I:3-3-4}
One has
$\h^q\left(X,\bigoplus_{x\in X}i_{x*}\dZ_x\right)=0$ for every $q>0$.
\end{lemma}

Indeed, for every closed point $x$ of $X$ one has
$\eR^qi_{x*}\dZ_x=0$ for $q>0$, because $i_x$ is a finite morphism
(\ref{I:2-3-6}), and
\[
  \h^q(X,i_{x*}\dZ_x) = \h^q(x,\dZ_x)\text{.}
\]
The right-hand side vanishes for every $q>0$ because $x$ is the spectrum
of an algebraically closed field. (This also shows that the lemma holds
more generally for every skyscraper sheaf on $X$.)

The preceding lemmas and the exact sequence (\ref{I:eq:3-3-1-1}) give
\[
  \h^q(X,\dG_m) = 0 \text{ for $q\geqslant 2$,}
\]
and the low-degree exact cohomology sequence
\[
  1 \to \h^0(X,\dG_m) \to \h^0(X, j_*\dG_{m,\eta}) \to \h^0\left(X,\bigoplus_{x\in X} i_{x*}\dZ_x\right) \to \h^1(X,\dG_m) \to 1
\]
which is precisely the exact sequence
\[
  1 \to k^\times \to k(X)^\times\to \Div(X) \to \pic(X) \to 1\text{.}
\]
Proposition \ref{I:3-3-1} implies that the cohomology groups of $X$ with
values in $\dZ/n$, for $n$ prime to the characteristic of $k$, have
reasonable values:

\begin{corollary}\label{I:3-3-5}
If $X$ has genus $g$ and $n$ is invertible in $k$, then
$\h^q(X,\dZ/n)$ vanishes for $q>2$ and is free over $\dZ/n$ of ranks
$1,2g,1$ for $q=0,1,2$. Replacing $\dZ/n$ by the isomorphic group
$\dmu_n$, there are canonical isomorphisms
\begin{align*}
  \h^0(X,\dmu_n) &= \dmu_n \\
  \h^1(X,\dmu_n) &= \pic^0(X)_n \\
  \h^2(X,\dmu_n) &= \dZ/n \text{.}
\end{align*}
\end{corollary}

Since $k$ is algebraically closed, $\dZ/n$ is isomorphic
(noncanonically) to $\dmu_n$. The Kummer exact sequence
\[
  0 \to \dmu_n \to \dG_m \to \dG_m \to 0\text{,}
\]
and Proposition \ref{I:3-3-1} yield the equalities
\[
  \h^q(X,\dZ/n) = 0 \text{ for $q>2$,}
\]
and, in low degrees, the exact sequences
\[\xymatrix{
  0 \ar[r]
    & \h^0(X,\dmu_n) \ar[r]
    & k^\times \ar[r]^-n
    & k^\times \ar[r]
    & 0
}\]
\[\xymatrix{
  0 \ar[r]
    & \h^1(X,\dmu_n) \ar[r]
    & \pic(X) \ar[r]^-n
    & \h^2(X,\dmu_n) \ar[r]
    & 0\text{.}
}\]
Moreover, there is an exact sequence
\[\xymatrix{
  0 \ar[r]
    & \pic^0(X) \ar[r]
    & \pic(X) \ar[r]^-{\deg}
    & \dZ \ar[r]
    & 0 \text{,}
}\]
and $\pic^0(X)$ is identified with the group of $k$-rational points of
an abelian variety of dimension $g$, the Jacobian of $X$. In such a
group, multiplication by $n$ is surjective and its kernel is a free
$\dZ/n\dZ$-module of rank $2g$ (because $n$ is invertible in $k$);
this proves the corollary.

A clever unwinding using the “trace method” yields the following
proposition as a corollary.

\begin{proposition}[{\cite[IX 5.7]{sga4}}]\label{I:3-3-6}
Let $k$ be an algebraically closed field, let $X$ be an algebraic curve
over $k$, and let $\sF$ be a torsion sheaf on $X$. Then:
\begin{enumerate}[\indent i)]
  \item $\h^q(X,\sF)=0$ for $q>2$.
  \item If $X$ is affine, then even $\h^q(X,\sF)=0$ for $q>1$.
\end{enumerate}
\end{proposition}

For the proof, as well as for an exposition of the “trace method,” we
refer to \cite[IX 5]{sga4}.

\subsection{Unwindings}\label{I:3-4}

To calculate the cohomology of varieties of dimension $>1$, one uses
fibrations by curves, thereby reducing to the study of morphisms whose
fibres have dimension $\leqslant1$. This principle has several variants;
we mention some of them.

\subsubsection{}\label{I:3-4-1}

Let $A$ be a $k$-algebra of finite type, and let $a_1,\dotsc,a_n$ be
generators of $A$. If we put $X_0=\spec(k)$,
$X_i=\spec(k[a_1,\dotsc,a_i])$, and $X_n=\spec(A)$, the canonical
inclusions
$k[a_1,\dotsc,a_i]\to k[a_1,\dotsc,a_i,a_{i+1}]$ define morphisms
$X_n\to X_{n-1}\to\cdots\to X_1\to X_0$ whose fibres have dimension
$\leqslant1$.

\subsubsection{}\label{I:3-4-2}

For a smooth morphism, one can be more precise. An
\emph{elementary fibration} is a morphism of schemes $f:X\to S$ that
can be embedded in a commutative diagram
\[\xymatrix{
  X \ar@{^{(}->}[r]^-j \ar[dr]_-f
    & \bar X \ar[d]^-{\bar f}
    & \ar@{_{(}->}[l]_-i \ar[dl]^-g Y \\
  & S
}\]
satisfying the following conditions:
\begin{enumerate}[\indent i)]
  \item $j$ is an open immersion, dense in every fibre, and
    $X=\bar X\setminus Y$.
  \item $\bar f$ is smooth and projective, with irreducible geometric
    fibres of dimension $1$.
  \item $g$ is an étale covering and no fibre of $g$ is empty.
\end{enumerate}

A \emph{good neighbourhood} relative to $S$ is an $S$-scheme $X$ for
which there exist $S$-schemes $X=X_n,\dotsc,X_0=S$ and elementary
fibrations $f_i:X_i\to X_{i-1}$, $i=1,\dotsc,n$. One can show
\cite[XI 3.3]{sga4} that if $X$ is a smooth scheme over an algebraically
closed field $k$, then every rational point of $X$ has an open
neighbourhood that is a good neighbourhood (relative to $\spec(k)$).

\subsubsection{}\label{I:3-4-3}

A proper morphism $f:X\to S$ can be unwound as follows. By Chow's lemma,
there is a commutative diagram
\[\xymatrix{
  X \ar[dr]_-f
    & \ar[l]_-\pi \ar[d]^-{\bar f} \bar X \\
  & S
}\]
in which $\pi$ and $\bar f$ are projective morphisms, with $\pi$ in
addition an isomorphism over a dense open subset of $X$. Locally on
$S$, $\bar X$ is a closed subscheme of a standard projective space
$\dP_S^n$.

The latter is unwound by considering the projection
$\varphi:\dP_S^n\to\dP_S^1$ that sends the point with homogeneous
coordinates $(x_0,x_1,\dotsc,x_n)$ to $(x_0,x_1)$. This is a rational
map defined outside the closed subset $Y\simeq\dP_S^{n-2}$ of
$\dP_S^n$ given by the homogeneous equations $x_0=x_1=0$. Let
$u:P\to\dP_S^n$ be the blow-up with centre $Y$; the fibres of $u$ have
dimension $\leqslant1$. Moreover, there is a natural morphism
$v:P\to\dP_S^1$ extending the rational map $\varphi$, and $v$ makes
$P$ a $\dP_S^1$-scheme locally isomorphic to the standard projective
space $\dP^{n-1}$, which may in turn be projected onto a $\dP^1$, and
so on.

\subsubsection{}\label{I:3-4-4}

A smooth projective variety $X$ can be swept out by a
\emph{Lefschetz pencil}. The blow-up $\widetilde X$ of the intersection
of the axis of the pencil with $X$ maps to $\dP^1$, and the fibres of
this map are the hyperplane sections of $X$ by the hyperplanes of the
pencil.

\section{The proper base change theorem}\label{I:4}

\subsection{Introduction}\label{I:4-1}

This chapter is devoted to the proof and applications of the following
theorem.

\begin{theorem}\label{I:4-1-1}
Let $f:X\to S$ be a proper morphism of schemes and let $\sF$ be a
torsion abelian sheaf on $X$. Then, for every $q\geqslant0$, the fibre
of $\eR^qf_*\sF$ at a geometric point $s$ of $S$ is isomorphic to the
cohomology $\h^q(X_s,\sF)$ of the fibre
$X_s=X\otimes_S\spec k(s)$ of $f$ at $s$.
\end{theorem}

For a proper separated continuous map $f:X\to S$ between topological
spaces (separated means that the diagonal of $X\times_SX$ is closed) and
an abelian sheaf $\sF$ on $X$, the analogous result is well known and
elementary. Since $f$ is closed, the $f^{-1}(V)$, with $V$ ranging over
the neighbourhoods of $s$, form a fundamental system of neighbourhoods
of $X_s$, and one checks that
$\h^\bullet(X_s,\sF)=\varinjlim_U\h^\bullet(U,\sF)$, where $U$ ranges
over the neighbourhoods of $X_s$. In practice, $X_s$ even has a
fundamental system $\cU$ of neighbourhoods $U$ of which it is a
deformation retract, and for $\sF$ constant one therefore has
$\h^\bullet(X_s,\sF)=\h^\bullet(U,\sF)$. Figuratively speaking, the
special fibre swallows the general fibre.

For schemes the proof is more delicate, and it is essential to suppose
that $\sF$ is torsion (\cite[XII.2]{sga4}). In view of the description of
the fibres of $\eR^qf_*\sF$ (\ref{I:2-3-3}), Theorem
(\ref{I:4-1-1}) is essentially equivalent to the following theorem.

\begin{theorem}\label{I:4-1-2}
Let $A$ be a strictly henselian local ring and let $S=\spec(A)$. Let
$f:X\to S$ be proper, and let $X_0$ be the closed fibre of $f$. Then, for every
torsion abelian sheaf $\sF$ on $X$ and every $q\geqslant0$, one has
$\h^q(X,\sF)\iso\h^q(X_0,\sF)$.
\end{theorem}

By passage to the limit, it suffices to prove the theorem when $A$ is the
strict henselization of a finitely generated $\dZ$-algebra at a prime
ideal. We first treat the case $q=0$ or $1$ and $\sF=\dZ/n$
(\ref{I:4-2}). An argument based on constructible sheaves (\ref{I:4-3})
shows moreover that it suffices to consider constant $\sF$. On the other
hand, the unwinding in (\ref{I:3-4-3}) allows us to suppose that $X_0$
is a curve; in that case it remains only to prove the theorem for $q=2$
(\ref{I:4-4}).

Among other applications (\ref{I:4-6}), the theorem allows one to define
cohomology with proper support (\ref{I:4-5}).

\subsection{Proof for \texorpdfstring{$q=0$}{q=0} or \texorpdfstring{$1$}{1} and \texorpdfstring{$\sF=\dZ/n$}{\sF=Z/n}}\label{I:4-2}

The result for $q=0$ and $\sF$ constant is equivalent to the following
proposition (Zariski's connectedness theorem).

\begin{proposition}\label{I:4-2-1}
Let $A$ be a noetherian henselian local ring and let $S=\spec(A)$. Let
$f:X\to S$ be a proper morphism, and let $X_0$ be the closed fibre of
$f$. Then the sets of connected components $\pi_0(X)$ and
$\pi_0(X_0)$ are in bijection.
\end{proposition}

Equivalently, we must show that the sets $\operatorname{Of}(X)$ and
$\operatorname{Of}(X_0)$ of open-and-closed subsets are in bijection.
The set $\operatorname{Of}(X)$ is known to correspond bijectively to the
set of idempotents of $\Gamma(X,\sO_X)$, and similarly
$\operatorname{Of}(X_0)$ corresponds bijectively to the set of
idempotents of $\Gamma(X_0,\sO_{X_0})$. Thus we must show that the
canonical map
\[
  \operatorname{Idem} \Gamma(X,\sO_X) \to \operatorname{Idem}\Gamma(X_0,\sO_{X_0})
\]
is bijective.

Let $\fm$ denote the maximal ideal of $A$, let
$\Gamma(X,\sO_X)^\wedge$ denote the completion of $\Gamma(X,\sO_X)$
for the $\fm$-adic topology, and, for every integer $n\geqslant0$, let
$X_n=X\otimes_AA/\fm^{n+1}$. By the finiteness theorem for proper
morphisms \cite[III.3.2]{ega}, $\Gamma(X,\sO_X)$ is a finite
$A$-algebra; since $A$ is henselian, the canonical map
\[
  \operatorname{Idem} \Gamma(X,\sO_X) \to \operatorname{Idem}\Gamma(X,\sO_X)^\wedge
\]
is bijective. In particular, the canonical map
\[
  \operatorname{Idem} \Gamma(X,\sO_X)^\wedge \to \varprojlim \operatorname{Idem}\Gamma(X_n,\sO_{X_n})
\]
is bijective. But since $X_n$ and $X_0$ have the same underlying
topological space, the canonical map
\[
  \operatorname{Idem}\Gamma(X_n,\sO_{X_n}) \to \operatorname{Idem} \Gamma(X_0,\sO_{X_0})
\]
is bijective for every $n$, which completes the proof.

Since $\h^1(X,\dZ/n)$ is in bijection with the set of isomorphism classes
of Galois étale coverings of $X$ with group $\dZ/n$, the theorem for
$q=1$ and $\sF=\dZ/n$ follows from the next proposition.

\begin{proposition}\label{I:4-2-2}
Let $A$ be a noetherian henselian local ring and let $S=\spec(A)$. Let
$f:X\to S$ be a proper morphism, and let $X_0$ be the closed fibre of $f$. Then
the restriction functor
\[
  \operatorname{Rev.et}(X) \to \operatorname{Rev.et}(X_0)
\]
is an equivalence of categories.
\end{proposition}

(If $X_0$ is connected and a geometric point of $X_0$ has been chosen
as base point, this says that the canonical map
$\pi_1(X_0)\to\pi_1(X)$ on (profinite) fundamental groups is bijective.)

Proposition (\ref{I:4-2-1}) shows that this functor is fully faithful.
Indeed, if $X'$ and $X''$ are two étale coverings of $X$, an
$X$-morphism from $X'$ to $X''$ is determined by its graph, which is an
open-and-closed subset of $X'\times_XX''$.

It remains to show that every étale covering $X_0'$ of $X_0$ extends to
an étale covering of $X$. Étale coverings are insensitive to nilpotent
elements \cite[ch.1]{sga1}; consequently, $X_0'$ lifts uniquely to an
étale covering $X_n'$ of $X_n$ for every $n\geqslant0$, in other words,
to an étale covering $\fX'$ of the formal scheme $\fX$ obtained by
completing $X$ along $X_0$. By Grothendieck's algebraization theorem for
formal coherent sheaves (the existence theorem, \cite[III.5]{ega}),
$\fX'$ is the formal completion of an étale covering $\bar X'$ of
$\bar X=X\otimes_A\hat A$.

By passage to the limit, it suffices to prove the proposition when $A$ is
the henselization of a finitely generated $\dZ$-algebra. We can then
apply Artin's approximation theorem to the functor
$\sF:(\textnormal{$A$-algebras})\to(\textnormal{sets})$ that associates
to an $A$-algebra $B$ the set of isomorphism classes of étale coverings
of $X\otimes_AB$. Indeed, this functor is locally of finite presentation:
if $B_i$ is a filtered direct system of $A$-algebras and
$B=\varinjlim B_i$, then $\sF(B)=\varinjlim\sF(B_i)$. By Artin's
theorem, given an element $\bar\xi\in\sF(\hat A)$, here the isomorphism class
of $\bar X'$, there exists $\xi\in\sF(A)$ having the same image as
$\bar\xi$ in $\sF(A/\fm)$. In other words, there exists an étale
covering $X'$ of $X$ whose restriction to $X_0$ is isomorphic to
$X_0'$.

\subsection{Constructible sheaves}\label{I:4-3}

In this paragraph, $X$ is a \emph{noetherian} scheme, and by a sheaf on
$X$ we mean an \emph{abelian} sheaf on $\et X$.

\begin{definition}\label{I:4-3-1}
A sheaf $\sF$ on $X$ is called \emph{locally constant constructible}
(abbreviated l.c.c.) if it is represented by an étale covering of $X$.
\end{definition}

\begin{definition}\label{I:4-3-2}
A sheaf $\sF$ on $X$ is called \emph{constructible} if it satisfies the
following equivalent conditions:
\begin{enumerate}[\indent (i)]
  \item There exists a finite surjective family of subschemes $X_i$ of
    $X$ such that the restriction of $\sF$ to $X_i$ is l.c.c.
  \item There exist a finite family of finite morphisms $p_i:X_i'\to X$,
    for each $i$ a constant constructible sheaf $C_i$ (i.e. one defined
    by a finite abelian group) on $X_i'$, and a monomorphism
    $\sF\to\prod p_{i*}C_i$.
\end{enumerate}
\end{definition}

It is readily checked that the category of constructible sheaves on $X$
is an abelian category. Moreover, if $u:\sF\to\sG$ is a homomorphism of
sheaves and $\sF$ is constructible, then $\im(u)$ is constructible.

\begin{lemma}\label{I:4-3-3}
Every torsion sheaf $\sF$ is a filtered direct limit of constructible
sheaves.
\end{lemma}

Indeed, if $j:U\to X$ is an étale scheme of finite type over $X$, an
element $\xi\in\sF(U)$ such that $n\xi=0$ defines a homomorphism of
sheaves $j_!\const{\dZ/n}\to\sF$ whose image (the smallest subsheaf of
$\sF$ for which $\xi$ is a local section) is a constructible subsheaf
of $\sF$. Clearly $\sF$ is the direct limit of such subsheaves.

\begin{definition}\label{I:4-3-4}
Let $\cC$ be an abelian category and let $T$ be a functor from $\cC$ to
the category of abelian groups. The functor $T$ is called
\emph{effaceable} in $\cC$ if, for every object $A$ of $\cC$ and every
$\alpha\in T(A)$, there exists a monomorphism $u:A\to M$ in $\cC$ such
that $T(u)\alpha=0$.
\end{definition}

\begin{lemma}\label{I:4-3-5}
The functors $\h^q(X,-)$ for $q>0$ are effaceable in the category of
constructible sheaves on $X$.
\end{lemma}

It suffices to observe that, if $\sF$ is a constructible sheaf, there is
necessarily an integer $n>0$ such that $\sF$ is a sheaf of
$\dZ/n$-modules. There is then a monomorphism $\sF\hookrightarrow\sG$,
where $\sG$ is a sheaf of $\dZ/n$-modules and
$\h^q(X,\sG)=0$ for every $q>0$. For example, one can take for $\sG$
the Godement resolution $\prod_{x\in X}i_{x*}\sF_{\bar x}$, where $x$
runs over the points of $X$ and $i_x:\bar x\to X$ is a geometric point
centred at $x$. By (\ref{I:4-3-3}), $\sG$ is a direct limit of
constructible sheaves, whence the lemma, because the functors
$\h^q(X,-)$ commute with direct limits.

\begin{lemma}\label{I:4-3-6}
Let $\varphi^\bullet:T^\bullet\to{T'}^\bullet$ be a morphism of
cohomological functors defined on an abelian category $\cC$ with values
in the category of abelian groups. Suppose that $T^q$ is effaceable for
$q>0$, and let $\cE$ be a subset of the objects of $\cC$ such that
every object of $\cC$ is contained in an object belonging to $\cE$.
Then the following conditions are equivalent:
\begin{enumerate}[\indent (i)]
  \item $\varphi^q(A)$ is bijective for every $q\geqslant0$ and every
    $A\in\ob\cC$.
  \item $\varphi^0(M)$ is bijective and $\varphi^q(M)$ is surjective for
    every $q>0$ and every $M\in\cE$.
  \item $\varphi^0(A)$ is bijective for every $A\in\ob\cC$, and
    ${T'}^q$ is effaceable for every $q>0$.
\end{enumerate}
\end{lemma}

The proof is by induction on $q$ and presents no difficulty.

\begin{proposition}\label{I:4-3-7}
Let $X_0$ be a subscheme of $X$. Suppose that, for every
$n>0$ and every scheme $X'$ finite over $X$, the canonical map
\[
  \h^q(X',\dZ/n) \to \h^q(X_0',\dZ/n)\text{,}
\]
where $X_0'=X'\times_XX_0$, is bijective for $q=0$ and surjective for
$q>0$. Then, for every torsion sheaf $\sF$ on $X$ and every
$q\geqslant0$, the canonical map
\[
  \h^q(X,\sF) \to \h^q(X_0,\sF)
\]
is bijective.
\end{proposition}

By passage to the limit, it suffices to prove the assertion for
constructible $\sF$. Apply Lemma (\ref{I:4-3-6}) with $\cC$ the category
of constructible sheaves on $X$, $T^q=\h^q(X,-)$,
${T'}^q=\h^q(X_0,-)$, and $\cE$ the set of constructible sheaves of the
form $\prod p_{i*}C_i$, where $p_i:X_i'\to X$ is finite and $C_i$ is a
finite constant sheaf on $X_i'$.

\subsection{End of the proof}\label{I:4-4}

By the method of fibration by curves (\ref{I:3-4-3}), we reduce to proving
the theorem in relative dimension $\leqslant1$. By the preceding
paragraph, it suffices to show the following. If $S$ is the spectrum of a
strictly henselian noetherian local ring, $f:X\to S$ is a proper
morphism whose closed fibre $X_0$ has dimension $\leqslant1$, and $n$
is an integer $>0$, then the canonical homomorphism
\[
  \h^q(X,\dZ/n) \to \h^q(X_0,\dZ/n)
\]
is bijective for $q=0$ and surjective for $q>0$.

The cases $q=0$ and $1$ were treated above, and
$\h^q(X_0,\dZ/n)=0$ for $q\geqslant3$; it therefore suffices to treat
$q=2$. We may clearly suppose that $n$ is a power of a prime. If
$n=p^r$, where $p$ is the characteristic of the residue field of $S$,
Artin--Schreier theory shows that $\h^2(X_0,\dZ/p^r)=0$. If
$n=\ell^r$, $\ell\ne p$, Kummer theory gives a commutative diagram
\[\xymatrix{
  \pic(X) \ar[d] \ar[r]^-\alpha
    & \h^2(X,\dZ/\ell^r) \ar[d] \\
  \pic(X_0) \ar[r]^-\beta
    & \h^2(X_0,\dZ/\ell^r)
}\]
in which $\beta$ is surjective (this was seen in Chapter \ref{I:3} for a
smooth curve over an algebraically closed field, but similar arguments
apply to any curve over a separably closed field).

Thus it suffices to prove:

\begin{proposition}\label{I:4-4-1}
Let $S$ be the spectrum of a henselian noetherian local ring, and let
$f:X\to S$ be a proper morphism whose closed fibre $X_0$ has dimension
$\leqslant1$. Then the canonical restriction map
\[
  \pic(X)\to \pic(X_0)
\]
is surjective. (In fact, it is enough that $f$ be a \emph{separated}
morphism of finite type.)
\end{proposition}

To simplify the proof, we suppose that $X$ is integral, although this is
not necessary. Every invertible sheaf on $X_0$ is associated with a
Cartier divisor (because $X_0$ is a curve, hence quasi-projective), so it
suffices to show that the canonical map $\Div(X)\to\Div(X_0)$ is
surjective.

Every divisor on $X_0$ is a linear combination of divisors whose support
is concentrated at a single nonisolated closed point of $X_0$. Let $x$
be such a point, let $t_0\in\sO_{X_0,x}$ be a regular nonunit in
$\sO_{X_0,x}$, and let $D_0$ be the divisor concentrated at $x$ with
local equation $t_0$. Let $U$ be an open neighbourhood of $x$ in $X$
on which there is a section $t\in\Gamma(U,\sO_U)$ lifting $t_0$. Let
$Y$ be the closed subset of $U$ defined by $t=0$; after making $U$
sufficiently small, we may suppose that $x$ is the only point of
$Y\cap X_0$. Then $Y$ is quasi-finite over $S$ at $x$. Since $S$ is
the spectrum of a henselian local ring, it follows that
$Y=Y_1\amalg Y_2$, where $Y_1$ is finite over $S$ and $Y_2$ does not
meet $X_0$. Moreover, since $X$ is separated over $S$, $Y_1$ is closed
in $X$.

After replacing $U$ by a smaller open neighbourhood of $x$, we may
suppose that $Y=Y_1$, in other words, that $Y$ is closed in $X$. We then
define a divisor $D$ on $X$ lifting $D_0$ by setting
$D|X\setminus Y=0$ and $D|U=\dv(t)$; this makes sense because $t$ is
invertible on $U\setminus Y$.

\subsubsection{Remark}\label{I:4-4-2}

When $f$ is proper, one could also give a proof similar to that of
Proposition (\ref{I:4-2-2}). Indeed, since $X_0$ is a curve, there is no
obstruction to lifting an invertible sheaf on $X_0$ to the infinitesimal
neighbourhoods $X_n$ of $X_0$, and hence to the formal completion $\fX$
of $X$ along $X_0$. One then concludes by applying successively
Grothendieck's existence theorem and Artin's approximation theorem.

\subsection{Cohomology with proper support}\label{I:4-5}

\begin{definition}\label{I:4-5-1}
Let $X$ be a separated scheme of finite type over a field $k$. By a
theorem of Nagata, there exist a scheme $\bar X$ proper over $k$ and an
open immersion $j:X\to\bar X$. For every torsion sheaf $\sF$ on $X$,
let $j_!\sF$ denote the extension by $0$ of $\sF$ to $\bar X$, and
define the groups of \emph{cohomology with proper support} by
\[
  \h_c^q(X,\sF)=\h^q(\bar X, j_! \sF) \text{.}
\]
\end{definition}

Let us show that this definition is independent of the chosen
compactification $j:X\to\bar X$. Let $j_1:X\to\bar X_1$ and
$j_2:X\to\bar X_2$ be two compactifications. Then $X$ maps into
$\bar X_1\times\bar X_2$ by $x\mapsto(j_1(x),j_2(x))$, and
the closed image $\bar X_3$ of $X$ under this map is a compactification
of $X$. We thus have a commutative diagram
\[\xymatrix{
  & \bar X_1 \\
  X \ar[ur]^-{j_1} \ar[rr]^-{(j_1,j_2)} \ar[dr]_-{j_2}
    & & \bar X_3 \ar[ul]_-{p_1} \ar[dl]^-{p_2} \\
  & \bar X_2
}\]
where $p_1$ and $p_2$, the restrictions to $\bar X_3$ of the natural
projections, are proper morphisms. It therefore suffices to treat the case
of a commutative diagram
\[\xymatrix{
  X \ar[r]^-{j_2} \ar[dr]_-{j_1}
    & \bar X_2 \ar[d]^-p \\
  & \bar X_1
}\]
with $p$ proper.

\begin{lemma}\label{I:4-5-2}
One has $p_*(j_{2!}\sF)=j_{1!}\sF$ and
$\eR^qp_*(j_{2!}\sF)=0$ for $q>0$.
\end{lemma}

Let us note at once that the lemma suffices. Using the Leray spectral
sequence for $p$, it gives, for every $q\geqslant0$,
\[
  \h^q(\bar X_2, j_{2!} \sF) = \h^q(\bar X_1, j_{1!} \sF) \text{.}
\]

To prove the lemma, argue fibre by fibre using the base change theorem
(\ref{I:4-1-1}) for $p$. The result is immediate because above a point
of $X$, $p$ is an isomorphism, whereas above a point of
$\bar X_1\setminus X$, $j_{2!}\sF$ vanishes on the fibre of $p$.

\subsubsection{}\label{I:4-5-3}

Similarly, if $f:X\to S$ is a separated morphism of finite type between
noetherian schemes, there exist a proper morphism $\bar f:\bar X\to S$
and an open immersion $j:X\to\bar X$. The higher direct images with
proper support $\eR^qf_!$ are defined, for every torsion sheaf $\sF$ on
$X$, by
\[
  \eR^q f_! \sF = \eR^q \bar f_*(j_! \sF) \text{.}
\]
As before, one checks that this definition is independent of the chosen
compactification.

\begin{theorem}\label{I:4-5-4}
Let $f:X\to S$ be a separated morphism of finite type between noetherian
schemes, and let $\sF$ be a torsion sheaf on $X$. Then the fibre of
$\eR^qf_!\sF$ at a geometric point $s$ of $S$ is isomorphic to the
cohomology with proper support $\h_c^q(X_s,\sF)$ of the fibre $X_s$ of
$f$ at $s$.
\end{theorem}

This is a simple variant of the proper base change theorem
(\ref{I:4-1-1}). More generally, if
\[\xymatrix{
  X \ar[d]^-f
    & X' \ar[l]_-{g'} \ar[d]^-{f'} \\
  S
    & S' \ar[l]^-g
}\]
is a cartesian diagram, there is a canonical isomorphism
\begin{equation*}\tag{4.5.4.1}\label{I:eq:4-5-4-1}
  g^* (\eR^q f_! \sF) \simeq \eR^q f'_!({g'}^* \sF) \text{.}
\end{equation*}

\subsection{Applications}\label{I:4-6}

\begin{theorem}[Vanishing theorem]\label{I:4-6-1}
Let $f:X\to S$ be a separated morphism of finite type whose fibres have
dimension $\leqslant n$, and let $\sF$ be a torsion sheaf on $X$. Then
$\eR^qf_!\sF=0$ for $q>2n$.
\end{theorem}

By the base change theorem, we may suppose that $S$ is the spectrum of a
separably closed field. If $\dim X=n$, there exists an affine open
$U$ of $X$ such that $\dim(X\setminus U)<n$. We then have an exact
sequence $0\to\sF_U\to\sF\to\sF_{X\setminus U}\to0$, and induction on
$n$ reduces us to proving the theorem for $X=U$ affine. The method of
fibration by curves (\ref{I:3-4-1}) and the base change theorem then
reduce us to a curve over a separably closed field, for which the desired
result follows from Tsen's theorem (\ref{I:3-3-6}).

\begin{theorem}[Finiteness theorem]\label{I:4-6-2}
Let $f:X\to S$ be a separated morphism of finite type and let $\sF$ be
a constructible sheaf on $X$. Then the sheaves $\eR^qf_!\sF$ are
constructible.
\end{theorem}

We consider only the case in which $\sF$ is annihilated by an integer
invertible on $X$.

To prove the theorem, one reduces to the case in which $\sF$ is the
constant sheaf $\const{\dZ/n}$ and $f:X\to S$ is proper and smooth with
geometrically connected curve fibres of genus $g$. If $n$ is invertible
on $X$, the sheaves $\eR^qf_*\sF$ are then locally free of finite rank
and vanish for $q>2$ (\ref{I:4-6-1}). Replacing $\dZ/n$ by the sheaf
$\dmu_n$, which is locally isomorphic to it (over $S$), there are
canonical isomorphisms
\begin{equation*}\tag{4.6.2.1}\label{I:eq:4-6-2-1}
\begin{aligned}
  \eR^0 f_* \dmu_n &= \dmu_n \\
  \eR^1 f_* \dmu_n &= \const\pic(X/S)_n \\
  \eR^2 f_* \dmu_n &= \dZ/n \text{.}
\end{aligned}
\end{equation*}

\begin{theorem}[Comparison with classical cohomology]\label{I:4-6-3}
Let $f:X\to S$ be a separated morphism of schemes of finite type over
$\dC$, and let $\sF$ be a torsion sheaf on $X$. Denote by a superscript
$\an{(-)}$ the functor passing to the usual topological spaces, and by
$\eR^q\an f_!$ the derived functors of direct image with proper support
under $\an f$. Then
\[
  \an{(\eR^q f_! \sF)} \simeq \eR^q \an f_! \an \sF \text{.}
\]
In particular, when $S=$ a point and $\sF$ is the constant sheaf
$\dZ/n$,
\[
  \h_c^q(X,\dZ/n) \simeq \h_c^q(\an X,\dZ/n)\text{.}
\]
\end{theorem}

Unwindings using the base change theorem reduce us to the case in which
$X$ is a smooth proper curve, $S=$ a point, and $\sF=\dZ/n$. The
cohomology groups in question then vanish for $q\ne0,1,2$, and we invoke
\cite{se55}: indeed, if $X$ is proper over $\dC$, then
$\pi_0(X)=\pi_0(\an X)$ and $\pi_1(X)=$ the profinite completion of
$\pi_1(\an X)$, which proves the assertion for $q=0,1$. For $q=2$,
one uses the Kummer exact sequence and the fact that, again by
\cite{se55}, $\pic(X)=\pic(\an X)$.

\begin{theorem}[Cohomological dimension of affine schemes]\label{I:4-6-4}
Let $X$ be an affine scheme of finite type over a separably closed field
and let $\sF$ be a torsion sheaf on $X$. Then
$\h^q(X,\sF)=0$ for $q>\dim(X)$.
\end{theorem}

For the very elegant proof, we refer to \cite{sga4}, XIV \S 2 and 3.

\subsubsection{Remark}\label{I:4-6-5}

This theorem is in a sense a substitute for Morse theory. Indeed, consider
the classical case in which $X$ is smooth and affine over $\dC$, embedded
in a standard affine space $\dC^N$. Then, for almost every point
$p\in\dC^N$, the function “distance from $p$” on $X$ is a Morse
function, and the indices of its critical points are smaller than
$\dim(X)$. Thus $X$ is obtained by attaching handles of index less than
$\dim(X)$, which gives the classical analogue of (\ref{I:4-6-4}).

\section{Local acyclicity of smooth morphisms}\label{I:5}

Let $X$ be a complex analytic variety and let $f:X\to D$ be a morphism
from $X$ to the disk. Write $[0,t]$ for the closed line segment in $D$
with endpoints $0$ and $t$, and $(0,t]$ for the half-open segment. If
$f$ is \emph{smooth}, the inclusion
\[
  j : f^{-1}\left((0,t]\right) \hookrightarrow f^{-1}\left([0,t]\right)
\]
is a homotopy equivalence; the special fibre
$X_0=f^{-1}(0)$ can be pushed into $f^{-1}([0,t])$.

In practice, for sufficiently small $t$, $f^{-1}((0,t])$ will be a
bundle over $(0,t]$, so that the inclusion
\[
  X_t = f^{-1}(t) \hookrightarrow f^{-1}\left((0,t]\right)
\]
will also be a homotopy equivalence. The \emph{cospecialization morphism}
is then the homotopy class of maps
\[
  \cosp : X_0 \hookrightarrow f^{-1}\left([0,t]\right) \xleftarrow\sim f^{-1}\left((0,t]\right) \xleftarrow\sim X_t \text{.}
\]
Figuratively speaking,
for a smooth morphism the general fibre swallows the special fibre.

No longer suppose that $f$ is necessarily smooth (but suppose that
$f^{-1}((0,t])$ is a bundle over $(0,t]$). A morphism
$\cosp^\bullet$ can still be defined in cohomology whenever
$j_*\dZ=\dZ$ and $\eR^qj_*\dZ=0$ for $q>0$. Under these hypotheses,
the Leray spectral sequence for $j$ shows that
\[
  \h^\bullet\left(f^{-1}\left([0,t]\right),\dZ\right) \iso \h^\bullet\left(f^{-1}\left((0,t]\right),\dZ\right)
\]
and $\cosp^\bullet$ is the composite morphism
\[
  \cosp^\bullet : \h^\bullet(X_t,\dZ)
    \xleftarrow\sim \h^\bullet\left(f^{-1}\left((0,t]\right),\dZ\right)
    \xleftarrow\sim \h^\bullet\left(f^{-1}\left([0,t]\right),\dZ\right)
    \to \h^\bullet(X_0,\dZ)
\]
The fibre of $\eR^qj_*\dZ$ at a point $x\in X_0$ is calculated as
follows. In an ambient space, take a ball $B_\varepsilon$ centred at $x$
with sufficiently small radius $\varepsilon$, and for sufficiently small
$\eta$ put $E=X\cap B_\varepsilon\cap f^{-1}(\eta t)$; this is the
\emph{vanishing-cycle variety} at $x$. One has
\[
  (\eR^q j_* \dZ)_x \xleftarrow\sim \h^q\left(X\cap B_\varepsilon\cap f^{-1}\left((0,\eta t]\right),\dZ\right) \xrightarrow\sim \h^q(E,\dZ)
\]
and the cospecialization morphism is defined in cohomology whenever the
vanishing-cycle varieties are acyclic
($\h^0(E,\dZ)=\dZ$ and $\h^q(E,\dZ)=0$ for $q>0$), which is expressed
by saying that $f$ is \emph{locally acyclic}.

This chapter is devoted to the analogue of this situation for a smooth
morphism of schemes in étale cohomology. In this setting, however, it is
essential to restrict to torsion coefficients of order \emph{prime to the
residue characteristics}. Paragraph \ref{I:5-1} gives generalities on
locally acyclic morphisms and cospecialization maps. In Paragraph
\ref{I:5-2}, we prove that a smooth morphism is locally acyclic. In
Paragraph \ref{I:5-3}, this result is combined with those of the preceding
chapter to obtain two applications: a specialization theorem for
cohomology groups (the cohomology of the geometric fibres of a smooth
proper morphism is locally constant), and a base change theorem for a
smooth morphism.

Throughout what follows, we fix an integer $n$, and “scheme” means
“scheme on which $n$ is invertible.” A “geometric point” will always mean
an “algebraic geometric point” (\ref{I:2-3-1})
$x:\spec(k)\to X$, with $k$ algebraically closed.

\subsection{Locally acyclic morphisms}\label{I:5-1}

\subsubsection{Notation}\label{I:5-1-1}

Given a scheme $S$ and a geometric point $s$ of $S$, we denote by
$\widetilde S^s$ the spectrum of the strict localization of $S$ at $s$.

\begin{definition}\label{I:5-1-2}
A geometric point $t$ of $S$ is called a \emph{generization} of $s$ if
it is defined by an algebraic closure of the residue field of a point of
$\widetilde S^s$. We also say that $s$ is a \emph{specialization} of
$t$, and call the \emph{specialization map} the $S$-morphism
$t\to\widetilde S^s$.
\end{definition}

\begin{definition}\label{I:5-1-3}
Let $f:X\to S$ be a morphism of schemes. Let $s$ be a geometric point
of $S$, $t$ a generization of $s$, $x$ a geometric point of $X$ over
$s$, and
$\widetilde X_t^x=\widetilde X^x\times_{\widetilde S^s}t$. Then
$\widetilde X_t^x$ is called a \emph{vanishing-cycle variety} of $f$ at
$x$.

The morphism $f$ is called \emph{locally acyclic} if the reduced
cohomology of every vanishing-cycle variety $\widetilde X_t^x$ vanishes:
\begin{equation*}\tag{5.1.3.1}\label{I:eq:5-1-3-1}
  \hat\h^\bullet(\widetilde X_t^x,\dZ/n) = 0\text{,}
\end{equation*}
i.e. $\h^0\left(\widetilde X_t^x,\dZ/n\right)=\dZ/n$ and
$\h^q(\widetilde X_t^x,\dZ/n)=0$ for $q>0$.
\end{definition}

\begin{lemma}\label{I:5-1-4}
Let $f:X\to S$ be a locally acyclic morphism and let $g:S'\to S$ be a
quasi-finite morphism (or a projective limit of quasi-finite morphisms).
Then the morphism $f':X'\to S'$ obtained from $f$ by base change is
locally acyclic.
\end{lemma}

Indeed, every vanishing-cycle variety of $f'$ is a vanishing-cycle
variety of $f$.

\begin{lemma}\label{I:5-1-5}
Let $f:X\to S$ be a locally acyclic morphism. For every geometric point
$t$ of $S$ giving rise to a cartesian diagram
\[\xymatrix{
  X_t \ar[r]^-{\varepsilon'} \ar[d]
    & X \ar[d]^-f \\
  t \ar[r]^-\varepsilon
    & S
}\]
one has $\varepsilon_*'\dZ/n=f^*\varepsilon_*\dZ/n$ and
$\eR^q\varepsilon_*'\dZ/n=0$ for $q>0$.
\end{lemma}

Let $\bar S$ be the closure of $\varepsilon(t)$, let $S'$ be the
normalization of $\bar S$ in $k(t)$, and consider the cartesian diagram
\[\xymatrix{
  X_t \ar[r]^-{i'} \ar[d]
    & X' \ar[r]^-{\alpha'} \ar[d]^-{f'}
    & X \ar[d]^-f \\
  t \ar[r]^-i
    & S' \ar[r]^-\alpha
    & S
}\]
The local rings of $S'$ are normal with separably closed fraction fields.
They are therefore strictly henselian, and the local acyclicity of $f'$
(\ref{I:5-1-4}) gives
$i'_*\dZ/n={f'}^*i_*\dZ/n$ and $\eR^q i'_*\dZ/n=0$ for $q>0$. Since $\alpha$
is integral, it follows that
\[
  \eR^q\varepsilon_*'\dZ/n
    = \alpha_*' \eR^q i_*' \dZ/n
    = \alpha_*' {f'}^* \eR^q i_* \dZ/n
    = f^* \alpha_* \eR^q i_* \dZ/n
    = f^* \eR^q \varepsilon_* \dZ/n \text{,}
\]
which proves the lemma.

\subsubsection{}\label{I:5-1-6}

Given a locally acyclic morphism $f:X\to S$ and a specialization map
$t\to\widetilde S^s$, we shall define canonical homomorphisms, called
cospecialization maps,
\[
  \cosp^\bullet : \h^\bullet(X_t,\dZ/n) \to \h^\bullet(X_s,\dZ/n)\text{,}
\]
relating the cohomology of the general fibre $X_t=X\times_St$ to that of
the special fibre $X_s=X\times_Ss$.

Consider the cartesian diagram
\[\xymatrix{
  X_t \ar[d] \ar[r]^-{\varepsilon'}
    & \widetilde X \ar[d]^-{f'}
    & \ar[l] X_s \ar[d] \\
  t \ar[r]^-\varepsilon
    & \widetilde S^s
    & \ar[l] s
}\]
obtained from $f$ by base change. By (\ref{I:5-1-4}), $f'$ is again
locally acyclic. The definition of local acyclicity immediately shows that
the restrictions to $X_s$ of the sheaves $\eR^q\varepsilon_*'\dZ/n$ are
$\dZ/n$ for $q=0$ and $0$ for $q>0$. By (\ref{I:5-1-5}), we even know
that $\eR^q\varepsilon_*'\dZ/n=0$ for $q>0$. We define
$\cosp^\bullet$ as the composite map
\begin{equation*}\tag{5.1.6.1}\label{I:eq:5-1-6-1}
  \h^\bullet(X_t,\dZ/n) \simeq \h^\bullet(\widetilde X,\varepsilon_*' \dZ/n) \to \h^\bullet(X_s,\dZ/n) \text{.}
\end{equation*}

\paragraph{Variant}
Let $\bar S$ be the closure of $\varepsilon(t)$ in
$\widetilde S^s$, let $S'$ be the normalization of $\bar S$ in $k(t)$,
and let $X'/S'$ be obtained from $X/S$ by base change. Diagram
\eqref{I:eq:5-1-6-1} may also be written
\[
  \h^\bullet(X_t,\dZ/n) \simeq \h^\bullet(X',\dZ/n) \to \h^\bullet(X_s,\dZ/n) \text{.}
\]

\begin{theorem}\label{I:5-1-7}
Let $S$ be a locally noetherian scheme, $s$ a geometric point of $S$,
and $f:X\to S$ a morphism. Suppose that
\begin{enumerate}[\indent a)]
  \item $f$ is locally acyclic,
  \item for every specialization map $t\to\widetilde S^s$ and every
    $q\geqslant0$, the cospecialization maps
    $\h^q(X_t,\dZ/n)\to\h^q(X_s,\dZ/n)$ are bijective.
\end{enumerate}
Then the canonical homomorphism
$(\eR^qf_*\dZ/n)_s\to\h^q(X_s,\dZ/n)$ is bijective for every
$q\geqslant0$.
\end{theorem}

To prove the theorem, we may clearly suppose that
$S=\widetilde S^s$. We shall in fact show that, for every sheaf of
$\dZ/n\dZ$-modules $\sF$ on $S$, the canonical homomorphism
$\varphi^q(\sF):(\eR^qf_*f^*\sF)_s\to\h^q(X_s,f^*\sF)$ is bijective.

Every sheaf of $\dZ/n\dZ$-modules is a filtered direct limit of
constructible sheaves of $\dZ/n\dZ$-modules (\ref{I:4-3-3}). Moreover,
every constructible sheaf of $\dZ/n\dZ$-modules embeds in a sheaf of the
form $\prod i_{\lambda *}C_\lambda$, where
$i_\lambda:t_\lambda\to S$ is a finite family of generizations of $s$
and $C_\lambda$ is a free $\dZ/n\dZ$-module of finite rank on
$t_\lambda$. By the definition of the cospecialization maps, condition
b) says that the homomorphisms $\varphi^q(\sF)$ are bijective when
$\sF$ has this form.

We conclude by a variant of Lemma (\ref{I:4-3-6}):

\begin{lemma}\label{I:5-1-8}
Let $\cC$ be an abelian category in which filtered direct limits exist.
Let $\varphi^\bullet:T^\bullet\to{T'}^\bullet$ be a morphism of
cohomological functors commuting with filtered direct limits, defined on
$\cC$ with values in the category of abelian groups. Suppose that there
are two subsets $\cD$ and $\cE$ of the objects of $\cC$ such that:
\begin{enumerate}[\indent a)]
  \item every object of $\cC$ is a filtered direct limit of objects
    belonging to $\cD$,
  \item every object belonging to $\cD$ is contained in an object
    belonging to $\cE$.
\end{enumerate}
Then the following conditions are equivalent:
\begin{enumerate}[\indent (i)]
  \item $\varphi^q(A)$ is bijective for every $q\geqslant0$ and every
    $A\in\ob\cC$.
  \item $\varphi^q(M)$ is bijective for every $q\geqslant0$ and every
    $M\in\cE$.
\end{enumerate}
\end{lemma}

The proof of the lemma proceeds by passage to the direct limit, induction
on $q$, and repeated application of the five lemma to the diagram of
exact cohomology sequences deduced from an exact sequence
$0\to A\to M\to A'\to0$, with $A\in\cD$, $M\in\cE$, and
$A'\in\ob\cC$.

\begin{corollary}\label{I:5-1-9}
Let $S$ be the spectrum of a strictly henselian noetherian local ring and
let $f:X\to S$ be a locally acyclic morphism. Suppose that for every
geometric point $t$ of $S$ one has $\h^0(X_t,\dZ/n)=\dZ/n$ and
$\h^q(X_t,\dZ/n)=0$ for $q>0$ (in other words, the geometric fibres of
$f$ are acyclic). Then $f_*\dZ/n=\dZ/n$ and
$\eR^qf_*\dZ/n=0$ for $q>0$.
\end{corollary}

\begin{corollary}\label{I:5-1-10}
Let $f:X\to Y$ and $g:Y\to Z$ be morphisms of locally noetherian
schemes. If $f$ and $g$ are locally acyclic, then so is $g\circ f$.
\end{corollary}

We may suppose that $X$, $Y$, and $Z$ are strictly local and that $f$
and $g$ are local morphisms. It is then enough to show that, if $z$ is
an algebraic geometric point of $Z$, then
$\h^0(X_z,\dZ/n)=\dZ/n$ and $\h^q(X_z,\dZ/n)=0$ for $q>0$.

Since $g$ is locally acyclic, one has $\h^0(Y_z,\dZ/n)=\dZ/n$ and
$\h^q(Y_z,\dZ/n)=0$ for $q>0$. Moreover, the morphism
$f_z:X_z\to Y_z$ is locally acyclic (\ref{I:5-1-4}), and its geometric
fibres are acyclic because they are vanishing-cycle varieties of $f$.
By (\ref{I:5-1-9}), $\eR^q{f_z}_*\dZ/n=0$ for $q>0$. In addition,
${f_z}_*\dZ/n$ is constant with fibre $\dZ/n$ on $Y_z$. The result
follows from the Leray spectral sequence for $f_z$.

\subsection{Local acyclicity of a smooth morphism}\label{I:5-2}

\begin{theorem}\label{I:5-2-1}
A smooth morphism is locally acyclic.
\end{theorem}

Let $f:X\to S$ be a smooth morphism. The assertion is local for the étale
topology on $X$ and $S$, so we may suppose that $X$ is standard affine
space of dimension $d$ over $S$. By passage to the limit, we may suppose
that $S$ is noetherian, and the transitivity of local acyclicity
(\ref{I:5-1-10}) shows that it suffices to treat $d=1$.

Let $s$ be a geometric point of $S$ and $x$ a geometric point of $X$
centred at a closed point of $X_s$. We must show that the geometric fibres
of $\widetilde X^x\to\widetilde S^s$ are acyclic. Henceforth put
$S=\widetilde S^s=\spec(A)$ and $X=\widetilde X^x$. Then
$X\simeq\spec A\{T\}$, where $A\{T\}$ is the henselization of $A[T]$
at the point $T=0$ above $s$.

If $t$ is a geometric point of $S$, the fibre $X_t$ is a projective
limit of smooth affine curves over $t$. Thus $\h^q(X_t,\dZ/n)=0$ for
$q\geqslant2$, and it suffices to show that
$\h^0(X_t,\dZ/n)=\dZ/n$ and $\h^1(X_t,\dZ/n)=0$ for $n$ prime to the
residue characteristic of $S$. This follows from the next two
propositions.

\begin{proposition}\label{I:5-2-2}
Let $A$ be a strictly henselian local ring, $S=\spec(A)$, and
$X=\spec A\{T\}$. Then the geometric fibres of $X\to S$ are connected.
\end{proposition}

By passage to the limit, we may reduce to the case in which $A$ is a
strict henselization of a finitely generated $\dZ$-algebra.

Let $\bar t$ be a geometric point of $S$ lying over $t$, and let $k'$
be a finite separable extension of $k(t)$ in $k(\bar t)$. Put
$t'=\spec(k')$ and $X_{t'}=X\times_S\spec(k')$. We must check that, for
all $\bar t$ and $t'$, $X_{t'}$ is connected (by which we mean connected
and nonempty). Let $A'$ be the normalization of $A$ in $k'$, i.e. the
ring of elements of $k'$ integral over the image of $A$ in $k(t)$. One
has $A\{T\}\otimes_AA'\iso A'\{T\}$: indeed, the left-hand side is
henselian local (because $A'$ is finite over $A$ and local) and is a
limit of local étale algebras over
$A'\{T\}=A\{T\}\otimes_AA'$. Thus $X_{t'}$ is still the fibre at $t'$
of $X=\spec\left(A'\{T\}\right)$ over $S'=\spec(A')$. The local scheme
$X'$ is normal, hence integral; its localization $X_{t'}$ is still
integral, and a fortiori connected.

\begin{proposition}\label{I:5-2-3}
Let $A$ be a strictly henselian local ring, $S=\spec(A)$, and
$X=\spec\left(A\{T\}\right)$. Let $\bar t$ be a geometric point of $S$
and $X_{\bar t}$ the corresponding geometric fibre. Then every Galois
étale covering of $X_{\bar t}$ of order prime to the characteristic of
the residue field of $A$ is trivial.
\end{proposition}

\begin{lemma}[Zariski--Nagata purity theorem in dimension $2$] \label{I:5-2-4} % originally 5.2.3.1
Let $C$ be a regular local ring of dimension $2$ and let $C'$ be a
finite normal $C$-algebra étale over the open complement of the closed
point of $\spec(C)$. Then $C'$ is étale over $C$.
\end{lemma}

Indeed, $C'$ is normal of dimension $2$, and hence
$\operatorname{prof}(C')=2$. Since
$\operatorname{prof}(C')+\dim\operatorname{proj}(C')=\dim(C)=2$, it
follows that $C'$ is free over $C$. The set of points of $C$ where $C'$
is ramified is then defined by an equation, the discriminant; since it
contains no point of height $1$, it is empty.

\begin{lemma}[Special case of Abhyankar's lemma]\label{I:5-2-5} % originally 5.2.3.2
Let $S=\spec(V)$ be a trait, let $\pi$ be a uniformizer, let $\eta$ be
the generic point of $S$, let $X$ be smooth over $S$, irreducible, and
of relative dimension $1$, let $\widetilde X_\eta$ be a Galois étale
covering of $X_\eta$ of degree $n$ invertible on $S$, and let
$S_1=\spec\left(V[\pi^{1/n}]\right)$. Use a subscript $(-)_1$ to denote
base change from $S$ to $S_1$. Then $\widetilde X_{1\eta}$ extends to an
étale covering of $X_1$.
\end{lemma}

Let $\widetilde X_1$ be the normalization of $X_1$ in
$\widetilde X_{1\eta}$. By the structure of the tame inertia groups of
the discrete valuation rings obtained by localizing $X$ at the generic
points of the special fibre $X_s$, $\widetilde X_1$ is étale over $X_1$
on the general fibre and at the generic points of the special fibre. By
(\ref{I:5-2-4}), it is étale everywhere.

\subsubsection{}\label{I:5-2-6} % originally 5.2.3.3

Let $t$ denote the point over which $\bar t$ lies. We are free to replace
$A$ by the normalization of $A$ in a finite separable extension $k(t')$ of
$k(t)$ in $k(\bar t)$ (cf. \ref{I:5-2-2}). This and a preliminary
passage to the limit allow us to suppose that
\begin{enumerate}[\indent a)]
  \item $A$ is normal and noetherian, and $t$ is the generic point of $S$.
  \item The étale covering under consideration of $X_{\bar t}$ comes from
    an étale covering of $X_t$.
  \item It comes from an étale covering $\widetilde X_U$ of the inverse
    image $X_U$ of a nonempty open subset $U$ of $S$ (this follows from
    b): $t$ is the limit of the $U$).
  \item The complement of $U$ has codimension $\geqslant2$ (at the price
    of enlarging $k(t)$, by applying (\ref{I:5-2-5}) to the discrete
    valuation rings obtained by localizing $S$ at the points of
    $S\setminus U$ of codimension $1$ in $S$; there are only finitely many
    such points).
  \item The covering $\widetilde X_U$ is trivial over the subscheme $T=0$.
    (This may require enlarging $k(t)$ once more.)
\end{enumerate}

Under these conditions, we shall show that the covering
$\widetilde X_U$ is trivial.

\begin{lemma}\label{I:5-2-7} % originally 5.2.3.4
Let $A$ be a normal noetherian strictly henselian local ring, let $U$
be an open subset of $\spec(A)$ whose complement has codimension
$\geqslant2$, let $V$ be its inverse image in
$X=\spec\left(A\{T\}\right)$, and let $V'$ be an étale covering of
$V$. If $V'$ is trivial over $T=0$, then $V'$ is trivial.
\end{lemma}

Let $B=\Gamma(V',\sO)$. Since $V'$ is the inverse image of $V$ in
$\spec(B)$, it suffices to show that $B$ is finite étale over $A\{T\}$
(and hence split, since $A\{T\}$ is strictly henselian). Let
$\widehat X=\spec A\llbracket T\rrbracket$, and denote by $(-)^\wedge$ base
change from $X$ to $\widehat X$. The scheme $\widehat X$ is faithfully
flat over $X$. Hence
$\Gamma(\widehat V',\sO)=B\otimes_{A\{T\}}A\llbracket T\rrbracket$,
and it suffices to show that this ring $\widehat B$ is finite étale over
$A\llbracket T\rrbracket$.

Let $V_m$ (resp. $V_m'$) be the subscheme of $\widehat V$ (resp.
$\widehat V'$) defined by $T^{m+1}=0$. By hypothesis, $V_0'$ is a
trivial covering of $V_0$: a sum of $n$ copies of $V_0$. The same is true
of $V_m'/V_m$, because étale coverings are insensitive to nilpotents.
We obtain
\[
  \varphi : \Gamma(\widehat V',\sO) \to \varprojlim_m \Gamma(V_m',\sO) = \left(\varprojlim_m \Gamma(V_m,\sO)\right)^n \text{.}
\]
By hypothesis, the complement of $U$ has depth $\geqslant2$: one has
$\Gamma(V_m,\sO)=A[T]/(T^{m+1})$, and $\varphi$ is a homomorphism from
$\widehat B$ to $A\llbracket T\rrbracket^n$. Over $U$, it provides $n$
distinct sections of $\widehat V'/\widehat V$; hence $\widehat V'$ is
trivial, a sum of $n$ copies of $\widehat V$. Since the complement of
$\widehat V$ in $\widehat X$ still has codimension $\geqslant2$ (and
therefore depth $\geqslant2$), it follows that
$\widehat B=A\llbracket T\rrbracket^n$, proving the lemma.

\subsection{Applications}\label{I:5-3}

\begin{theorem}[Specialization of cohomology groups]\label{I:5-3-1}
Let $f:X\to S$ be a proper locally acyclic morphism, for example a smooth
proper morphism. Then the sheaves $\eR^qf_*\dZ/n$ are locally constant
constructible, and for every specialization map $t\to\widetilde S^s$,
the cospecialization maps
$\h^q(X_t,\dZ/n)\to\h^q(X_s,\dZ/n)$ are bijective.
\end{theorem}

This follows immediately from the definition of the cospecialization maps
and the finiteness and proper base change theorems.

\begin{theorem}[Base change by a smooth morphism]\label{I:5-3-2}
Let
\[\xymatrix{
  X' \ar[r]^-{g'} \ar[d]^-{f'}
    & X \ar[d]^-f \\
  S' \ar[r]^-{g}
    & S
}\]
be a cartesian diagram with $g$ smooth. For every sheaf $\sF$ on $X$
that is torsion of order prime to the residue characteristics of $S$,
one has
\[
  g^* \eR^q f_* \sF \iso \eR^q f'_*({g'}^* \sF) \text{.}
\]
\end{theorem}

After passing to an open covering of $X$, we reduce to the case where
$X$ is affine, and then, by passage to the limit, to the case where $X$
is of finite type over $S$. The morphism $f$ then factors as an open
immersion $j:X\to\bar X$ followed by a proper morphism
$\bar f:\bar X\to S$. The Leray spectral sequence for $\bar f\circ j$
and the proper base change theorem show that it suffices to prove the
theorem when $X\to S$ is an open immersion.

In this case, if $\sF$ has the form $\varepsilon_*C$, where
$\varepsilon:t\to X$ is a geometric point of $X$, the theorem is a
corollary of (\ref{I:5-1-5}). The general case follows from Lemma
(\ref{I:5-1-8}).

\begin{corollary}\label{I:5-3-3}
Let $K/k$ be an extension of separably closed fields, let $X$ be a
$k$-scheme, and let $n$ be an integer prime to the characteristic of
$k$. Then the canonical map
$\h^q(X,\dZ/n)\to\h^q(X_K,\dZ/n)$ is bijective for every
$q\geqslant0$.
\end{corollary}

It suffices to observe that $\bar K$ is a direct limit of smooth
$\bar k$-algebras.

\begin{theorem}[Relative purity]\label{I:5-3-4}
Let
\begin{equation*}\tag{5.3.4.1}\label{I:eq:5-3-4-1}\xymatrix{
  U \ar@{^{(}->}[r]^-j \ar[dr]
    & X \ar[d]^-f
    & \ar@{_{(}->}[l]_-i Y \ar[dl]^-h \\
  & S
}\end{equation*}
be a commutative diagram with $f$ smooth of pure relative dimension $N$,
$h$ smooth of pure relative dimension $N-1$, $i$ a closed immersion,
and $U=X\setminus Y$. If $n$ is prime to the residue characteristics of
$S$, then
\begin{equation*}
\begin{aligned}
  j_* \dZ/n     &= \dZ/n \\
  \eR^1 j_*\dZ/n &= \dZ/n(-1)_Y \\
  \eR^q j_*\dZ/n &= 0 \qquad \text{for $q\geqslant 2$}
\end{aligned}
\end{equation*}
\end{theorem}

In these formulas, $\dZ/n(-1)$ denotes the $\dZ/n$-dual of $\dmu_n$.
If $t$ is a local equation for $Y$, the isomorphism
$\eR^1j_*\dZ/n\simeq\dZ/n(-1)_Y$ is defined by the map
$a:\dZ/n\to\eR^1j_*\dmu_n$ sending $1$ to the class of the
$\dmu_n$-torsor of $n$th roots of $t$.

The question is local. This allows us to replace $(X,Y)$ by a locally
isomorphic pair, for example
\[\xymatrix{
  \dA_T^1 \ar@{^{(}->}[r]^-j \ar[dr]_-g
    & \dP_T^1 \ar[d]^-f
    & \ar@{_{(}->}[l]_-i T \ar[dl] \\
  & T
}\]
with $T=\dA_S^{N-1}$ and $i=$ the section at infinity. Corollary
(\ref{I:5-1-9}) applies to $g$ and gives
$\eR^qg_*\dZ/n=\dZ/n$ for $q=0$ and $0$ for $q>0$. For $f$, on the
other hand, \eqref{I:eq:4-6-2-1} gives
\[
  \eR^q f_* \dZ/n =
  \begin{cases}
    \dZ/n & q=0,\\
    0 & q=1,\\
    \dZ/n(-1) & q=2,\\
    0 & q>2.
  \end{cases}
\]
It is readily checked that $j_*\dZ/n=\dZ/n$, and that the
$\eR^qj_*\dZ/n=0$ are concentrated on $i(T)$ for $q>0$. The Leray
spectral sequence
$E_2^{pq}=\eR^pf_*\eR^qj_*\dZ/n\Rightarrow\eR^{p+q}g_*\dZ/n$ therefore
reduces to
\[\xymatrix{
  i^* \eR^q j_* \dZ/n & 0 & \cdots \\
  i^* \eR^1 j_* \dZ/n \ar[drr]^-{d_2} & 0 & \cdots \\
  \dZ/n              & 0 & \dZ/n(-1) & 0 & \cdots
}\]
$\eR^qj_*\dZ/n=0$ for $q\geqslant2$, and $\eR^1j_*\dZ/n$ is the
extension by zero of a locally free sheaf of rank one on $T$ (isomorphic,
via $d_2$, to $\dZ/n(-1)$). Since the map $a$ defined above is injective
(as may be checked fibre by fibre), it is an isomorphism, which proves
(\ref{I:5-3-4}).

\subsubsection{}\label{I:5-3-5}

We refer to \cite[XVI.4,5]{sga4} for proofs of the following applications
of the acyclicity theorem (\ref{I:5-2-1}).

\begin{theorem}\label{I:5-3-6} % originally 5.3.5.1
Let $f:X\to S$ be a morphism of schemes of finite type over $\dC$ and
let $\sF$ be a constructible sheaf on $X$. Then
\[
  \an{(\eR^q f_*\sF)} \sim \eR^q \an f_* (\an\sF)
\]
(cf. \ref{I:4-6-3}: for ordinary cohomology it is necessary to assume
that $\sF$ is constructible, not merely torsion).
\end{theorem}

\begin{theorem}\label{I:5-3-7} % originally 5.3.5.2
Let $f:X\to S$ be a morphism of schemes of finite type over a field $k$
of characteristic $0$, and let $\sF$ be a constructible sheaf on $X$.
Then the $\eR^qf_*\sF$ are also constructible.
\end{theorem}

The proof uses resolution of singularities and (\ref{I:5-3-4}). It is
generalized to a morphism of finite type between excellent schemes of
characteristic $0$ in \cite[XIX.5]{sga4}. Another proof, independent of
resolution, is given in this volume (\nameref{VII}, \ref{VII:1-1})
% NOTE: originally (Th. finitude, 1.1),
and applies to a morphism of schemes of finite type over a field or over
a Dedekind ring.

\section{Poincaré duality}

\subsection{Introduction}\label{I:6-1}

Let $X$ be an oriented topological manifold of pure dimension $N$, and
suppose that $X$ admits a finite open covering
$\cU=(U_i)_{1\leqslant i\leqslant K}$ such that every nonempty
intersection of the open subsets $U_i$ is homeomorphic to a ball. For
such a manifold, the Poincaré duality theorem can be presented as follows:
\begin{enumerate}[\indent A]
  \item The cohomology of $X$ is the \v Cech cohomology associated with
    the covering $\cU$. It is the cohomology of the complex
    \begin{equation}\label{I:eq:6-1}
      0 \to \dZ^{A_0} \to \dZ^{A_1} \to \cdots
    \end{equation}
    where
    $A_k=\{(i_0,\dotsc,i_k):i_0<\cdots<i_k\text{ and }U_{i_0}\cap\cdots\cap U_{i_k}\ne\varnothing\}$.
  \item For $a\in A_k$, $a=(i_0,\dotsc,i_k)$, let
    $U_a=U_{i_0}\cap\cdots\cap U_{i_k}$ and let $j_a$ be the inclusion
    of $U_a$ into $X$. The constant sheaf $\dZ$ on $X$ admits the
    (left) resolution
    \begin{equation}\label{I:eq:6-2}
      \xymatrix{
        \cdots \ar[r]
          & \displaystyle\bigoplus_{a\in A_1} j_{a!}\dZ \ar[r]
          & \displaystyle\bigoplus_{a\in A_0} j_{a!} \dZ \ar[r] \ar[d]
          & 0 \\
        & & \dZ \text{.}
      }
    \end{equation}
    The compactly supported cohomology $\h_c^\bullet(X,j_{a!}\dZ)$ is
    precisely the cohomology with proper support of the (oriented) ball
    $U_a$:
      \[
        \h_c^i(X,j_{a!}\dZ) = \begin{cases}
                                0   & \text{if $i\ne N$} \\
                                \dZ & \text{if $i = N$}
                              \end{cases}
      \]
\end{enumerate}

The hypercohomology spectral sequence for the complex
\eqref{I:eq:6-2} and cohomology with proper support therefore says that
$\h_c^i(X)$ is the $(i-N)$th cohomology group of the complex
\begin{equation}\label{I:eq:6-3}
  \cdots \to \dZ^{A_1} \to \dZ^{A_0} \to 0 \text{.}
\end{equation}
This complex is dual to the complex \eqref{I:eq:6-1}, whence Poincaré
duality.

The essential points of this construction are
\begin{enumerate}[\indent a)]
  \item the existence of a theory of cohomology with proper support;
  \item the fact that every point $x$ of a manifold $X$ of pure dimension
    $N$ has a fundamental system of open neighbourhoods $U$ for which
    \begin{equation}\label{I:eq:6-4}
      \h_c^i(U) = \begin{cases}
                    0   & \text{for $i\ne N$,} \\
                    \dZ & \text{for $i=N$.}
                  \end{cases}
    \end{equation}
\end{enumerate}

Poincaré duality in étale cohomology can be constructed on this model.
For $X$ smooth of pure dimension $N$ over an algebraically closed field
$k$, $n$ invertible on $X$, and $x$ a closed point of $X$, the key
point is to calculate the projective limit over the étale neighbourhoods
$U$ of $x$:
\begin{equation}\label{I:eq:6-5}
  \varprojlim \h_c^i(U,\dZ/n)
    = \begin{cases}
        0     & \text{if $i\ne 2N$} \\
        \dZ/n & \text{if $i=2N$.}
      \end{cases}
\end{equation}

Just as for topological manifolds one must first treat directly the case
of an open ball (or even simply the interval $(0,1)$), here one must
first treat directly the case of curves (\ref{I:6-2}). The local
acyclicity theorem for smooth morphisms then reduces
\eqref{I:eq:6-5} to this special case (\ref{I:6-3}).

The isomorphisms \eqref{I:eq:6-4} and \eqref{I:eq:6-5} are not
canonical: they depend on the choice of an \emph{orientation} of $X$.
For $n$ invertible on a scheme $X$, $\dmu_n$ is a sheaf of free
$\dZ/n$-modules of rank one. We denote its $N$th tensor power by
$\dZ/n(N)$ ($N\in\dZ$). The intrinsic form of the second line of
\eqref{I:eq:6-5} is
\begin{equation}\label{I:eq:6-5'}
  \varprojlim \h_c^{2N} \left(U,\dZ/n(N)\right) = \dZ/n \text{,}
\end{equation}
and $\dZ/n(N)$ is called the \emph{orientation sheaf} of $X$. Since the
sheaf $\dZ/n(N)$ is constant and isomorphic to $\dZ/n$, it may be taken
outside the limit sign (with $N$ fixed), and one may instead write
\begin{equation}\label{I:eq:6-5''}
  \varprojlim \h_c^{2N} (U,\dZ/n) = \dZ/n(-N) \text{,}
\end{equation}
and Poincaré duality will take the form of a perfect duality, with values
in $\dZ/n(-N)$, between $\h^i(X,\dZ/n)$ and
$\h_c^{2N-i}(X,\dZ/n)$.

\subsection{The case of curves}\label{I:6-2}

\subsubsection{}\label{I:6-2-1}

Let $\bar X$ be a smooth projective curve over an algebraically closed
field $k$, and let $n$ be invertible on $\bar X$. The proof of
(\ref{I:3-3-5}) gives, for connected $\bar X$, a canonical isomorphism
\[\xymatrix{
  \h^2(\bar X,\dmu_n) = \pic(\bar X)/n\pic(\bar X) \ar[r]^-{\deg}_-{\sim}
    & \dZ/n
}\]
Let $D$ be a reduced divisor of $\bar X$ and let
$X=\bar X\setminus D$:
\[\xymatrix{
  X \ar@{^{(}->}[r]^-j
    & \bar X
    & \ar@{_{(}->}[l]_-i D \text{.}
}\]
The exact sequence
$0\to j_!\dmu_n\to\dmu_n\to i_*\dmu_n\to0$ and the fact that
$\h^i(\bar X,i_*\dmu_n)=\h^i(D,\dmu_n)=0$ for $i>0$ give an
isomorphism
\[\xymatrix{
  \h_c^2(X,\dmu_n) = \h^2(\bar X,j_! \dmu_n) \ar[r]^-\sim
    & \h^2(\bar X,\dmu_n) \ar[r]^-\sim
    & \dZ/n \text{.}
}\]
For disconnected $\bar X$, similarly
\[
  \h_c^2(X,\dmu_n) = (\dZ/n)^{\pi_0(X)}
\]
and the \emph{trace morphism} is defined to be the sum
\[\xymatrix{
  \tr : \h_c^2(X,\dmu_n) \simeq (\dZ/n)^{\pi_0(X)} \ar[r]^-\Sigma
    & \dZ/n\text{.}
}\]

\begin{theorem}\label{I:6-2-2}
The pairing $\tr(a\smallsmile b)$ identifies each of the two groups
$\h^i(X,\dZ/n)$ and $\h_c^{2-i}(X,\dmu_n)$ with the dual (with values
in $\dZ/n$) of the other.
\end{theorem}

\emph{Transcendental proof.} Suppose that $\bar X$ is a smooth projective
curve over a trait $S$ and that $j:X\hookrightarrow\bar X$ is the
inclusion of the complement of a divisor $D$ étale over $S$. The
cohomology (resp. cohomology with proper support) of the special and
generic geometric fibres of $X/S$ are “the same,” i.e. the fibres of
locally constant sheaves on $S$. This follows from the analogous facts
for $\bar X$ and $D$, by means of the exact sequence
$0\to j_!\dZ/n\to\dZ/n\to{\dZ/n}_D\to0$ (for cohomology with proper
support) and the formulas $j_*\dZ/n=\dZ/n$,
$\eR^1j_*\dZ/n\simeq{\dZ/n}_D(-1)$, and
$\eR^ij_*\dZ/n=0$ ($i\geqslant2$) (for ordinary cohomology)
(\ref{I:5-3-4}).

This specialization principle reduces (\ref{I:6-2-2}) to the case in
which $k$ has characteristic $0$. By (\ref{I:5-3-3}), this case reduces
to $k=\dC$. Finally, when $k=\dC$, the groups
$\h^\bullet(X,\dZ/n)$ and $\h_c^\bullet(X,\dmu_n)$ agree with the
groups of the same name calculated for the classical topological space
$X_{\textnormal{cl}}$, and, through the isomorphism
$\dZ/n\to\dmu_n:x\mapsto\exp\left(\frac{2\pi ix}{n}\right)$, the trace
morphism is identified with “integration over the fundamental class.”
Thus (\ref{I:6-2-2}) follows from Poincaré duality for
$X_{\textnormal{cl}}$.

\subsubsection{Algebraic proof}\label{I:6-2-3}

For a very economical proof, see (\nameref{V}, \S\ref{V:2}).
% originally (Duality \S 2)
% NOTE: should be hyperlinked V:2
Here is another, related to the autoduality of the Jacobian.

Return to the notation of (\ref{I:6-2-1}). We may suppose---and do
suppose---that $X$ is connected. The cases $i=0$ and $i=2$ being
trivial, we also suppose that $i=1$. Define $_D\dG_m$ by the exact
sequence $0\to_D\dG_m\to\dG_m\to i_*\dG_m\to0$ (sections of $\dG_m$
congruent to $1\mod D$). The group $\h^1(\bar X,_D\dG_m)$ classifies
invertible sheaves on $\bar X$ trivialized on $D$. It is the group of
points of $\pic_D(\bar X)$, an extension of $\dZ$ (the degree) by the
group of points of Rosenlicht's generalized Jacobian
$\pic_D^0(\bar X)$ (corresponding to conductor $1$ at each point of
$D$), itself an extension of the abelian variety $\pic^0(\bar X)$ by
the torus $\dG_m^D/(\text{diagonal $\dG_m$})$.

\begin{enumerate}[\indent a)]
  \item The exact sequence
    $0\to j_!\dmu_n\to_D\dG_m\xrightarrow{n}_D\dG_m\to0$ gives an
    isomorphism
    \begin{equation*}\tag{6.2.3.1}\label{I:eq:6-2-3-1}
      \h_c^1(X,\dmu_n) = \pic_D^0(\bar X)_n \text{.}
    \end{equation*}
  \item The map sending $x\in X(k)$ to the class of the invertible sheaf
    $\sO(x)$ on $\bar X$, trivialized by $1$ on $D$, comes from a
    morphism
    \[
      f:X \to \pic_D(\bar X) \text{.}
    \]
\end{enumerate}
For the remainder, fix a base point $0$ and put
$f_0(x)=f(x)-f(0)$. For every homomorphism
$v:\pic_D^0(\bar X)_n\to\dZ/n$, let
$\bar v\in\h^1\left(\pic_D^0(\bar X),\dZ/n\right)$ be the image under
$v$ of the class in
$\h^1\left(\pic_D^0(\bar X),\pic_D^0(\bar X)_n\right)$ of the torsor
defined by the extension
\[\xymatrix{
  0 \ar[r]
    & \pic_D^0(\bar X)_n \ar[r]
    & \pic_D^0(\bar X) \ar[r]^-n
    & \pic_D^0(\bar X) \ar[r]
    & 0
}\]
Geometric class field theory (as presented in Serre \cite{se59}) shows
that the map $v\mapsto f_0^*(\bar v)$
\begin{equation*}\tag{6.2.3.2}\label{I:eq:6-2-3-2}
  \hom\left(\pic_D^0(\bar X)_n,\dZ/n\right) \to \h^1(X,\dZ/n)
\end{equation*}
is an isomorphism. To deduce (\ref{I:6-2-2}) from
\eqref{I:eq:6-2-3-1} and \eqref{I:eq:6-2-3-2}, it remains to know that
\begin{equation*}\tag{2.3.3}\label{I:eq:6-2-3-3}
  \tr\left(u \smallsmile f_0^*(\bar v)\right) = -v(u) \text{.}
\end{equation*}

% originally (Duality, 3.2.4).

\subsection{The general case}\label{I:6-3}

Let $X$ be a smooth algebraic variety of pure dimension $N$ over an
algebraically closed field $k$. To state the Poincaré duality theorem,
one must first define the trace morphism
\[
  \tr : \h_c^{2N}\left(X,\dZ/n(N)\right) \to \dZ/n\text{.}
\]
The definition is a painful unwinding starting from the case of curves
\cite[XVIII \S 2]{sga4}. One then has:

\begin{theorem}\label{I:6-3-1}
The pairing $\tr(a\smallsmile b)$ identifies each of the groups
$\h_c^i\left(X,\dZ/n(N)\right)$ and $\h^{2N-i}(X,\dZ/n)$ with the dual
of the other.
\end{theorem}

Let $x\in X$ be a closed point and let $X_x$ be the strict localization
of $X$ at $x$. For $U$ ranging over the étale neighbourhoods of $x$,
put
\begin{equation*}\tag{6.3.1.1}\label{I:eq:6-3-1-1}
  \h_c^\bullet(X_x,\dZ/n) = \varprojlim \h_c^\bullet(U,\dZ/n) \text{.}
\end{equation*}
It would be preferable to consider instead the pro-object
$\text{“}\varprojlim\text{”}\h_c^\bullet(U,\dZ/n)$, but since the
groups involved are finite, the difference is inessential. As we tried
to explain in the introduction, (\ref{I:6-3-1}) follows from the facts
that
\begin{equation*}\tag{6.3.1.2}\label{I:eq:6-3-1-2}
\begin{aligned}
  \h_c^i(X_x,\dZ/n) &= 0 \text{ for $i\ne 2N$ and that } \\
  \tr : \h_c^{2 N}\left(X_x,  \dZ/n(N)\right) &\iso \dZ/n \text{ is an isomorphism.}
\end{aligned}
\end{equation*}

The case $N=0$ is trivial. If $N>0$, let $Y_y$ be the strict
localization at a closed point of a smooth scheme $Y$ of pure dimension
$N-1$, and let $f:X_x\to Y_y$ be an essentially smooth morphism (of
relative dimension one). The proof uses the Leray spectral sequence for
cohomology with proper support for $f$ to reduce to the case of curves.
Because the “cohomology with proper support” under consideration is
defined as limits \eqref{I:eq:6-3-1-1}, the existence of such a spectral
sequence raises various passage-to-the-limit problems, treated in too much
detail in \cite[XVIII]{sga4}. For every geometric point $z$ of $Y_y$,
one has
\[
  (\eR^if_! \dZ/n)_z = \h_c^i(f^{-1}(z),\dZ/n) \text{.}
\]
The geometric fibre $f^{-1}(z)$ is a projective limit of smooth curves
over an algebraically closed field. It satisfies Poincaré duality. Its
ordinary cohomology is given by the local acyclicity theorem for smooth
morphisms:
\[
  \h^i(f^{-1}(z),\dZ/n)
    = \begin{cases}
        \dZ/n & \text{for $i=0$} \\
        0     & \text{for $i>0$.}
      \end{cases}
\]
By duality,
\[
  \h_c^i(f^{-1}(z),\dZ/n)
    = \begin{cases}
        \dZ/n(-1) & \text{for $i=2$} \\
        0         & \text{for $i\ne 2$.}
      \end{cases}
\]
and the Leray spectral sequence becomes
\[
  \h_c^i\left(X_x,\dZ/n(N)\right) = \h_c^{i-2}\left(Y_y,\dZ/n(N-1)\right) \text{.}
\]
The result follows by induction on $N$.

\subsection{Variants and applications}\label{I:6-4}

One can construct in étale cohomology a “duality formalism”
($=$ functors $\eR f_*$, $\eR f_!$, $f^*$, and $\eR f^!$ satisfying various
compatibilities and adjunction formulas) parallel to the one that exists
in coherent cohomology. In this language, the results of the preceding
paragraph can be rewritten as follows: if $f:X\to S$ is smooth of pure
relative dimension $N$ and $S=\spec(k)$ with $k$ algebraically closed,
then
\[
  \eR f^! \dZ/n = \dZ/n [2N](N) \text{.}
\]
This result holds without any hypothesis on $S$. It has the following
corollary.

\begin{corollary}\label{I:6-4-1}
If $f$ is smooth of pure relative dimension $N$ and the sheaves
$\eR^if_!\dZ/n$ are locally constant, then the sheaves
$\eR^if_*\dZ/n$ are also locally constant, and
\[
  \eR^i f_* \dZ/n = \const\hom\left(\eR^{2 N-i} f_! \dZ/n(N),\dZ/n\right)\text{.}
\]
\end{corollary}

In particular, under the hypotheses of the corollary, the sheaves
$\eR^if_*\dZ/n$ are constructible. From this one can show that, if $S$
is of finite type over the spectrum of a field or of a Dedekind ring,
then, for every morphism of finite type $f:X\to S$ and every constructible
sheaf $\sF$ on $X$, the sheaves $\eR^if_*\sF$ are constructible
(\nameref{VII}, \ref{VII:1-1}).
% originally (Th. finitude, 1.1).

% END INCLUDED components/sga4half_1_en.tex

% BEGIN INCLUDED components/sga4half_2_en.tex
% !TEX root = ../SGA4half_English_Reader_1_to_8.tex

\chapter{Report on the trace formula}\label{II}

In this text, I have attempted to present Grothendieck's cohomological theory
of $L$-functions as directly as possible. I follow very closely some of the
lectures given by Grothendieck at the IHES in the spring of 1966. In the
spirit of this volume, no appeal will be made to SGA 5---apart from two
references to passages in Exposé XI, which is otherwise independent of that
seminar.

In \S 10 (pp. 81--90) of \cite{de73}, the reader will find a summary of the
theory in the crucial case of curves.










\section{Some reminders}\label{II:1}





\subsection{Notation}\label{II:1-1}

$p,q,\dF,\dF_q$: $p$ is a prime number, $q=p^f$ a power of $p$, and $\dF$
an algebraic closure of the prime field $\dF_p$. For every power $p^n$ of
$p$, denote by $\dF_{p^n}$ the subfield with $p^n$ elements of $\dF$.

$X_0,X$: An object over $\dF_q$ will often be denoted by a symbol carrying a
subscript $0$. The same symbol without the $0$ then denotes the object
obtained from it by extension of scalars to $\dF$. For example, if $\sF_0$ is
a sheaf on a scheme $X_0$ over $\dF_q$, denote by $\sF$ its inverse image on
$X=X_0\otimes_{\dF_q} \dF$.

$F$: If $X_0$ is a scheme over $\dF_q$, the endomorphism of $X_0$ that
``sends the point with coordinates $x$ to the point with coordinates $x^q$''
is called the \emph{Frobenius morphism} and denoted by $F$: it is the identity
on the underlying topological space and, for $x$ a local section of the
structure sheaf, $F^* x = x^q$.

The endomorphism of $X$ obtained from it by extension of scalars is also
denoted by $F$. On the set $X(\dF) = X_0(\dF)$ of rational points of $X$,
$F$ acts as the Frobenius substitution
$\varphi\in\gal(\dF/\dF_q)$ defined by $\varphi(x)=x^q$. In particular,
$X^F=X_0(\dF_q)$.

In case of ambiguity, we shall write $F_{(q)}$ rather than $F$.










\subsection{Frobenius correspondence}\label{II:1-2}

Let $X_0$ be a scheme over $\dF_q$, and let $\sF_0$ be a sheaf on $X_0$. I
like to view the \emph{Frobenius correspondence} (an $F$-endomorphism of
$\sF_0$) as follows:
\begin{equation*}\tag{1.2.1}\label{II:eq:1-2-1}
  F^* : F^*\sF_0 \iso \sF_0
\end{equation*}

Regard $\sF_0$ as the sheaf of local sections of a space $[\sF_0]$ étale over
$X_0$ ($[\sF_0]$ is an algebraic space in the sense of M. Artin). The
functoriality of $F$ gives a commutative diagram
\[\xymatrix{
  [\sF_0] \ar[r]^-F \ar[d]
    & [\sF_0] \ar[d] \\
  X_0 \ar[r]^F
    & X_0
}\]
Because $[\sF_0]$ is étale over $X_0$, this diagram is Cartesian; it gives an
isomorphism $[\sF_0]\iso F^*[\sF_0]$ whose inverse is
\eqref{II:eq:1-2-1}. We also denote by $F^*$ the correspondences or morphisms
deduced from $F$ by extension of scalars or by functoriality, such as
$F^*:\h_c^\bullet(X,\sF)\to \h_c^\bullet(X,\sF)$.

For a formal definition, I refer to \cite[XI.1,2]{sga5}. Only the case
$q=p$ is considered there, but the general case is treated in the same way.






\subsection{}\label{II:1-3}

The Frobenius correspondence is functorial in $(X_0,\sF_0)$. If
$f_0:X_0\to Y_0$ is a morphism of $\dF_q$-schemes, $\sF_0$ a sheaf on $X_0$,
$\sG_0$ a sheaf on $Y_0$, and
$u\in\hom(f_0^*\sG_0,\sF_0)$ an $f_0$-morphism from $\sG_0$ to $\sF_0$, there
is a commutative diagram of spaces $f_0 F = F f_0$, and above it a
commutative diagram of sheaves
\[\xymatrix{
  X_0 \ar[r]^-{f_0} \ar[d]^-F
    & Y_0 \ar[d]^-F \\
  X_0 \ar[r]^-{f_0}
    & Y_0
}\qquad\qquad
\xymatrix{
  \sF_0
    & \ar[l]_-u \sG_0 \\
  \sF_0 \ar[u]_-{F^*}
    & \ar[l] \sG_0 \ar[u]^-{F^*}
}\]
In particular, the correspondence on $f_{0*}\sF_0$ deduced from the
Frobenius correspondence of $(X_0,\sF_0)$ is the Frobenius correspondence of
$(Y_0,f_{0*}\sF_0)$.






\subsection{}\label{II:1-4}

The natural morphism $Y\to Y_0$ is a limit of étale morphisms. It follows
that for every abelian sheaf $\sF_0$ on $X_0$, the inverse images on $Y$ of
the $\eR^if_0\sF_0$ are the $\eR^if_*\sF$.

On $\eR^if_*\sF$, one has
\begin{enumerate}[\indent a)]
  \item the correspondence
    $F^*\eR^i f_*\sF \iso \eR^i f_*\sF \xrightarrow{F^*} \eR^if_*\sF$
    deduced by functoriality from the Frobenius correspondence;
  \item the correspondence $F^*\eR^if_*\sF\to\eR^if_*\sF$ that is the
    inverse image on $Y$ of the Frobenius correspondence on
    $\eR^i f_{0*}\sF_0$.
\end{enumerate}

It follows from (\ref{II:1-3}) that these correspondences coincide. They
induce the initial term of a Frobenius action on the entire spectral sequence
\[
  \h^p(Y,\eR^q f_*\sF) \rightarrow \h^{p+q}(X,\sF) \text{.}
\]
Likewise for cohomology with proper support, and for the spectral sequence
\[
  \h_c^p(Y,\eR^q f_* \sF) \rightarrow \h_c^{p+q}(X,\sF)
\]





\subsection{Change of finite field}\label{II:1-5}

Let $(X_0,\sF_0)$ be over $\dF_q$, let $n$ be an integer, and let
$(X_1,\sF_1)$ over $\dF_{q^n}$ be obtained from $(X_0,\sF_0)$ by extension of
scalars. On the one hand there is the Frobenius correspondence
$F_{(q)}:X_0\to X_0)$, $F_{(q)}^*:F_{(q)}^*\sF_0\to \sF_0$, and on the other
there is $F_{(q^n)}:X_1\to X_1$,
$F_{(q^n)}^*:F_{(q^n)}^*\sF_1\to\sF_1$. It is easily verified that
$F_{(q^n)}$ is obtained from the $n$-th power of $F_{(q)}$ by extension of
scalars from $\dF_q$ to $\dF_{q^n}$; after extension of scalars to $\dF$, one
still has the same identity $F_{(q^n)}=F_{(q)}^n$ between endomorphisms of
$X$ and correspondences on $X$.

If $x$ is a closed point of $X_0$, with $[k(x),\dF_q]=n$, and
$\bar x\in X_0(\dF) = X(\dF)$ lies over $x$, then $\bar x$ is a fixed point
of $F_{(q^n)}=F_{(q)}^n$. Denote by $F_x^*$ the endomorphism of
$\sF_{\bar x}$ induced by $F_{(q^n)}^*$. Up to isomorphism,
$(F_{\bar x}^*,\sF_{\bar x})$ does not depend on the choice of $\bar x$.
The trace, etc. of $F_x^*$ is denoted by notation such as
$\tr(F_x^*,\sF)$.






\subsection{\texorpdfstring{$L$}{L}-function}\label{II:1-6}

Let $X_0$ be a scheme of finite type over $\dF_q$. Denote by $|X_0|$ the set
of its closed points, and for $x\in |X_0|$ set
$\deg(x)=[k(x),\dF_p]$. If $\sF_0$ is a $\dQ_\ell$-sheaf on $X_0$ (for this
notion, see below)---or a constructible flat sheaf of modules over a
commutative Noetherian ring $A$---denote by $L(X_0,\sF_0)$ the formal series
with constant term $1$, in $\dQ_\ell\llbracket t\rrbracket$ or in
$A\llbracket t\rrbracket$,
\[
  L(X_0,\sF_0) = \prod_{x\in|X_0|} \det\left(1-F_x^* t^{\deg(x)},\sF\right)^{-1} \text{.}
\]






\subsection{Remark}\label{II:1-7}

In definition (\ref{II:1-6}), one has $\deg(x)=[k(x),\dF_p]$ (absolute
degree), and not $\deg(x)=[k(x),\dF_q]$ (degree relative to $\dF_q$). The
effect is that $L(X_0,\sF_0)$ depends only on the scheme $X_0$ and the sheaf
$\sF_0$, not on $q$ and the structure morphism
$X_0\to\spec(\dF_q)$. The indeterminate $t$ is an avatar of $p^{-s}$.






\subsection{The Galois point of view}\label{II:1-8}

This subsection (n$^\circ$) will not be used in the remainder of the report. For $\sF_0$
an abelian sheaf on $X_0$, the Galois group $\gal(\dF/\dF_q)$ acts on
$\dF$, and by transport of structure on $\h_c^i(X,\sF)$. In particular, the
Frobenius substitution $\varphi$ (\ref{II:1-1}) induces an automorphism, also
denoted by $\varphi$, of $\h_c^i(X,\sF)$. The Frobenius correspondence, on the
other hand, defines an endomorphism $F^*$ of $\h_c^i(X,\sF)$. One has
\begin{equation*}\tag{1.8.1}\label{II:eq:1-8-1}
  {F^*}^{-1} = \varphi \qquad \text{(on $\h_c^i(X,\sF)$).}
\end{equation*}

To see this, apply (\ref{II:1-4}) to the case where $Y_0$ is a point:
$Y_0=\spec(\dF_q)$. One finds that $\h_c^i(X,\sF)$ is the fiber, at the
geometric point $\spec(\dF)$ of $Y_0$, of $\eR^i f_! \sF_0$, and that $F^*$
is the fiber of the Frobenius correspondence of $\eR^i f_!\sF_0$. Let
$[\eR^i f_! \sF_0]$ be as in (\ref{II:1-2}). One has
$\h_c^i(X,\sF) = [\eR^i f_!\sF_0](\dF)$ and
\begin{enumerate}[\indent a)]
  \item $\varphi$ acts by transport of structure, via its action on $\dF$;
  \item $F^*$ is the inverse of
    $F:[\eR^i f_!\sF_0](\dF) \to [\eR^i f_!\sF_0](\dF)$.
\end{enumerate}

The conclusion follows from the identity $F=\varphi$ of (\ref{II:1-1}).

One advantage of the Galois point of view is that it permits reasoning by
transport of structure.










\section{\texorpdfstring{$\dQ_\ell$}{Ql}-sheaves}\label{II:2}

The notion of $\dQ_\ell$-sheaves is developed in detail in
\cite[V,VI]{sga5}. We shall use only the definitions and results below.






\begin{definition_}\label{II:2-1}
Let $X$ be a Noetherian scheme. A $\dZ_\ell$-sheaf $\sF$ on $X$ is an
inverse system of sheaves $\sF_n$ ($n\geqslant 0$), with $\sF_n$ a
constructible sheaf of $\dZ/\ell^{n+1}$-modules \cite[IX.2]{sga4}, such that
the transition morphism $\sF_n\to\sF_{n-1}$ factors through an isomorphism
\[
  \sF_n\otimes_{\dZ/\ell^{n+1}} \dZ/\ell^n \iso \sF_{n-1}\text{.}
\]
The sheaf $\sF$ is said to be \emph{lisse} if the $\sF_n$ are locally
constant.
\end{definition_}

The terminology of \cite{sga5} is different: ``constructible
$\dZ_\ell$-sheaf'' is used there for ``$\dZ_\ell$-sheaf,'' and
``constructibly twisted constant'' for ``lisse.''

Likewise below, for $\dQ_\ell$-sheaves.






\subsection{}\label{II:2-2}

If $k$ is a separably closed field, a sheaf of $\dZ/\ell^n$-modules on
$\spec(k)$ is identified with a $\dZ/\ell^n$-module, and the functor
$\sF\mapsto \varprojlim \sF_n$ identifies the $\dZ_\ell$-sheaves on
$\spec(k)$ with the $\dZ_\ell$-modules of finite type. This justifies the
following definition: the \emph{fiber} of the $\dZ_\ell$-sheaf $\sF$ at the
geometric point $x$ of $X$ is the $\dZ_\ell$-module
$\varprojlim (\sF_n)_x$.






\subsection{}\label{II:2-3}

Suppose $X$ is connected, and let $x$ be a geometric point of $X$. It is
known that the ``fiber at $x$'' functor is an equivalence between the
category of constructible locally constant sheaves of
$\dZ/\ell^n$-modules and the category of finite-type
$\dZ/\ell^n$-modules equipped with a continuous action of $\pi_1(X,x)$.
Passing to the limit gives the following proposition.






\begin{proposition_}\label{II:2-4}
Under the hypotheses above, the ``fiber at $x$'' functor is an equivalence
between the category of lisse $\dZ_\ell$-sheaves on $X$ and that of
finite-type $\dZ_\ell$-modules equipped with a continuous action of
$\pi_1(X,x)$.
\end{proposition_}






\begin{proposition_}\label{II:2-5}
Let $\sF$ be a $\dZ_\ell$-sheaf on a Noetherian scheme $X$. There is a finite
partition of $X$ into locally closed parts $X_i$ such that the $\sF|X_i$ are
lisse.
\end{proposition_}

Set $\gr_\ell^n\sF=\ell^n\sF_m/\ell^{m+1}\sF_m$ for $m\geqslant n$, and
$\gr_\ell\sF=\bigoplus \gr_\ell^n\sF$. It is a sheaf of modules over the
Noetherian ring $\gr_\ell\dZ=\dZ/\ell[t]$, the graded ring of $\dZ$ for the
$\ell$-adic filtration. It is a quotient of
$\gr_\ell^0\sF \otimes_{\dZ/\ell} \dZ/\ell[t]$, hence a constructible sheaf
of $\gr_\ell\dZ$-modules \cite[IX.2]{sga4}; and there is a finite partition
of $X$ into locally closed parts $X_i$ such that
$\gr_\ell \sF|X_i$ is locally constant. Each of the sheaves
$\gr_\ell^n\sF|X_i$ is then locally constant, as are the $\sF_m|X_i$, which
are successive extensions of them.






\subsection{}\label{II:2-6}

If $\sF$ and $\sG$ are two $\dZ_\ell$-sheaves on $X$, set
$\hom(\sF,\sG) = \varprojlim(\sF_n,\sG_n)$; this is a $\dZ_\ell$-module.
Since $\hom(\sF_n,\sG_n) = \hom(\sF_m,\sG_n)$ for $m\geqslant n$, one also
has
\[
  \hom(\sF,\sG) = \varprojlim_n \varinjlim_m \hom(\sF_m,\sG_n) \text{,}
\]
so that the functor $\sF\mapsto \text{``}\varprojlim\text{''}\sF_n$ is fully
faithful from the category of $\dZ_\ell$-sheaves to that of pro-sheaves on
$X$ ($=$ pro-objects of the category of sheaves on $X$). We shall use it to
identify $\dZ_\ell$-sheaves with certain pro-sheaves (those such that every
$\sF_n=\sF/\ell^{n+1}\sF$ is a sheaf, and
$\sF\iso \text{``}\varprojlim\text{''}\sF_n$). It is easily verified that if
$\sF+\text{``}\varprojlim\text{''}\sF_n$, with $\sF_n$ a constructible sheaf
of $\dZ/\ell^n$-modules, then $\sF$ is a $\dZ_\ell$-sheaf if and only if the
inverse systems (in $n$) $\sF_n/\ell^{k+1}\sF_n$ are essentially constant:
the inverse system of
$\sF_k'=\varprojlim_n\sF_n/\ell^{k+1}\sF_n$ then satisfies the conditions of
(\ref{II:2-1}), and
$\sF=\text{``}\varprojlim\text{''}\sF_k'$.

In what follows, we abandon the notation $\sF_n$ of (\ref{II:2-1}); if
$\sF$ is a $\dZ_\ell$-sheaf, the sheaf $\sF_n$ of (\ref{II:2-1}) will be
denoted by $\sF\otimes\dZ/\ell^{n+1}$.






\subsection{}\label{II:2-7}

In the abelian category of pro-sheaves, the subcategory of
$\dZ_\ell$-sheaves is stable under kernels, cokernels, and extensions. This is
clear for cokernels. For kernels, if $f:\sF\to\sG$ is a morphism of
$\dZ_\ell$-sheaves, it follows readily from
(\ref{II:2-4}, \ref{II:2-5}) and Artin--Rees that the inverse system of
$\ker(f:\sF\otimes\dZ/\ell^n\to\sG\otimes\dZ/\ell^n)$ satisfies the
criterion given in (\ref{II:2-6}) above. There is even an integer $r$ such
that for every $n$
\[
  \ker(f)\otimes\dZ/\ell^n = \ker(f:\sF\otimes\dZ/\ell^{n+r}\to\sG\otimes\dZ/\ell^{n+r})\otimes\dZ/\ell^n \text{.}
\]
The case of extensions is left to the reader.






\subsection{}\label{II:2-8}

A $\dZ_\ell$-sheaf $\sF$ is \emph{torsion-free} if
$\ker(\ell:\sF\to\sF) = 0$. Each $\sF\otimes\dZ/\ell^n$ is then flat over
$\dZ/\ell^n$ (i.e. its fibers are finite-rank free
$\dZ/\ell^n$-modules). It follows from (\ref{II:2-4}, \ref{II:2-5}) that for
every $\dZ_\ell$-sheaf $\sF$ there is an integer $n$ such that
$\sF/\ker(\ell^n)$ is torsion-free.






\subsection{}\label{II:2-9}

The abelian category of \emph{$\dQ_\ell$-sheaves} on $X$ is the quotient of
that of $\dZ_\ell$-sheaves by the subcategory of torsion
$\dZ_\ell$-sheaves. Equivalently, its objects are the $\dZ_\ell$-sheaves,
and, denoting by $\sF\otimes\dQ_\ell$ the $\dZ_\ell$-sheaf $\sF$, viewed as
an object of the category of $\dQ_\ell$-sheaves, one has
\[
  \hom(\sF\otimes\dQ_\ell,\sG\otimes\dQ_\ell) = \hom(\sF,\sG)\otimes_{\dZ_\ell}\dQ_\ell \text{.}
\]

The \emph{fiber} of a $\dQ_\ell$-sheaf $\sF\otimes\dQ_\ell$ at a geometric
point $x$ of $X$ is the finite-dimensional $\dQ_\ell$-vector space
$\sF_x\otimes_{\dZ_\ell}\dQ_\ell$.






\subsection{}\label{II:2-10}

Let $X$ be a separated scheme of finite type over an algebraically closed
field $k$, and let $\sF$ be a constructible $\dQ_\ell$-sheaf on $X$. Let
$\sF'$ be a constructible $\dZ_\ell$-sheaf such that
$\sF\sim \sF'\otimes_{\dZ_\ell}\dQ_\ell$. We shall see in
(\ref{II:4-11}) that
\[
  \left(\varprojlim \h_c^q(X,\sF'\otimes \dZ/\ell^n)\right)\otimes_{\dZ_\ell}\dQ_\ell
\]
is a finite-dimensional $\dQ_\ell$-vector space. It depends only on $\sF'$,
and is denoted by $\h_c^q(X,\sF)$. The proof does not use the fact that
$\ell$ is prime to the characteristic, but only that case concerns us.






\subsection{}\label{II:2-11}

In this exposé, all important calculations will be carried out at the level
of $\dZ/\ell^n$-sheaves, with passage to $\dQ_\ell$-sheaves taking place only
at the very end. In practice, however, it is often convenient to work
systematically with $\dQ_\ell$-sheaves. In Exposés V and VI of
\cite{sga5}, Jouanolou gives the formalism that makes this possible. One of
the essential results (\hyperref[VI:2-2]{VI 2.2}) is that if $f:X\to Y$ is a separated morphism
of finite type of Noetherian schemes, and $\sF$ a $\dZ_\ell$-sheaf on $Y$,
then, for every $q$, the pro-sheaf
$\text{``}\varprojlim\text{''}\eR^q f_! (\sF\otimes\dZ/\ell^n)$ is a
$\dZ_\ell$-sheaf on $Y$. He even proves a stronger regularity property (of
Artin--Rees type) for the inverse system of the
$\eR^q f_!(\sF\otimes\dZ/\ell^n)$.

Jouanolou even works in a more general setting, encompassing the very useful
case of $E_\lambda$-sheaves, for $E_\lambda$ a finite extension of
$\dQ_\ell$. They are defined as the $\dQ_\ell$-sheaves were, with
$\dQ_\ell$ replaced by $E_\lambda$, $\dZ_\ell$ by the ring $\sO_\lambda$ of
the valuation $\lambda$, and $\dZ/\ell^n$ by $\sO_\lambda/\pi^n$, where
$\pi$ is a uniformizer. It would amount to the same thing to define an
$E_\lambda$-sheaf as a $\dQ_\ell$-sheaf $\sF$ equipped with a homomorphism
$E_\lambda\to\operatorname{End}(\sF)$.









\section{Trace formula and \texorpdfstring{$L$}{L}-functions}\label{II:3}

\emph{We use the notation of \S\ref{II:1}, and $\ell$ is a prime number
$\ne p$.}

Let $X_0$ be a separated scheme of finite type over $\dF_q$ ($q=p^f$), and
let $\sF_0$ be a constructible $\dQ_\ell$-sheaf on $X_0$. Grothendieck's
cohomological interpretation of $L$-functions is the following theorem:






\begin{theorem_}\label{II:3-1}
$L(X_0,\sF_0) = \prod_i \det\left(1-F^* t^f, \h_c^i(X,\sF)\right)^{(-1)^{i+1}}$.
\end{theorem_}

It is obtained as a consequence of the trace formula:






\begin{theorem_}\label{II:3-2}
For every integer $n$,
$\sum_{x\in X^{F^n}} \tr\left({F^n}^*,\sF_x\right) = \sum_i (-1)^i \tr\left({F^n}^*,\h_c^i(X,\sF)\right)$.
\end{theorem_}






Let us briefly recall the classical method for deducing (\ref{II:3-1}) from
(\ref{II:3-2}). Set $T=t^f$; both sides are formal series in $T$ with
constant term $1$. Since $\dQ_\ell$ has characteristic $0$, it suffices to
prove that the logarithmic derivatives $T\frac{d}{dT}\log$ of the two sides
are equal.

To compute these logarithmic derivatives, we shall use the formula
\[
  T \frac{d}{dT}\log\det(1-f T)^{-1} = \sum_{n>0} \tr(f^n) T^n \text{,}
\]
% NOTE: the following lines are not typesetting correctly
whose proof (in a slightly more general setting) is recalled below. On the
left-hand side, one finds
\begin{align*}
  T\frac{d}{dT}\log L(X_0,\sF_0)
    &= \sum_{x\in |X_0|} \sum_n [k(x):\dF_q] \tr\left(F_x^n,\sF\right) T^{n\cdot[k(x):\dF_q]} \\
    &= \sum_{n>0} T^n \sum_{x\in X^{F^n}} \tr\left({F^n}^*,\sF_x\right)
\end{align*}
and on the right-hand side
\[
  \sum_{n>0} T^n \sum (-1)^i \tr\left({F^n}^*,\h_c^i(X,\sF)\right) \text{ ; }
\]
the terms are compared one by one.






\begin{proposition_}\label{II:3-3}
Let $A$ be a commutative ring, $M$ a projective $A$-module of finite type,
and $f$ an endomorphism of $M$. For every invertible formal series $Q$, set
$\frac{d}{dT}\log Q = Q^{-1}\frac{d}{dT}Q$. One has
\begin{equation*}\tag{3.3.1}\label{II:eq:3-3-1}
  T\frac{d}{dT}\log\det(1-f T)^{-1} = \sum_{n>0} \tr(f^n) T^n \text{.}
\end{equation*}
\end{proposition_}

Let $M'$ be an $A$-module such that $M\oplus M'$ is free of finite rank.
Replacing $(M,f)$ by $(M\oplus M',f\oplus 0)$ changes neither side of
(\ref{II:3-3}), and reduces us to the case where $M$ is free. For $M=A^d$,
\eqref{II:eq:3-3-1} is an algebraic identity in the coefficients $f_i^j$ of
$f$. By the principle of extension of algebraic identities, it suffices to
verify it for $A$ an algebraically closed field of characteristic $0$. Since
both sides are additive in $M$ in an extension
$0\to M'\to M\to M''\to 0$, it even suffices to treat only the case where
$M$ has rank $1$. If $f$ is multiplication by $a$,
\eqref{II:eq:3-3-1} then reduces to the formula
\[
  T\frac{d}{dT}\log\left((1-a T)^{-1}\right) = \sum_{n>0} a^n T^n \text{.}
\]






\begin{corollary_}\label{II:3-4}
Let $n$ be an integer, and let $f^{(n)}$ be the endomorphism
$(x_1,\dotsc,x_n)\mapsto \left(f(x_n),x_1,\dotsc,x_{n-1}\right)$ of $M^n$.
One has
\begin{equation*}\tag{3.4.1}\label{II:eq:3-4-1}
  \det(1-f T^n) = \det(1-f^{(n)} T) \text{.}
\end{equation*}
\end{corollary_}

Proceeding as in (\ref{II:3-3}), we reduce to assuming that $A$ has
characteristic $0$. It then suffices to verify that the logarithmic
derivatives $-T\frac{d}{dT}\log$ of the two sides are equal.
\begin{align*}
  -T\frac{d}{dT}\log\det(1-f T^n) &= - n T^n \frac{d}{dT^n} \log\det(1-f T^n) = n \sum_{m>0} \tr(f^m) T^{n m} \text{,} \\
  -T \frac{d}{dT}\log\det\left(1-f^{(n)} T\right) &= \sum_{m>0} \tr\left(f^{(n) m}\right) T^m \text{,}
\end{align*}
and one observes that
\begin{align*}
  \tr(f^{(n) m}) &= 0 \qquad \text{if $n\nmid m$} \\
  \tr(f^{(n) n m}) &= n \tr(f^m) \text{.}
\end{align*}






\subsection{Remark}\label{II:3-5}

For $A=\dQ_\ell$, \eqref{II:eq:3-4-1} is the special case
$X_0=\spec(\dF_{q^n})$ of (\ref{II:3-1}). Using this identity, it is easy to
verify directly that the right-hand side of (\ref{II:3-1}) is independent of
$q$ and of the structure morphism $X_0\to\spec(\dF_q)$.






\subsection{Remark}\label{II:3-6}

If $j:X_0\hookrightarrow \bar X_0$ is a compactification of $X_0$, one has
$L(X_0,\sF_0) = L(\bar X_0, j_!\sF_0)$ and
$\h_c^i(X,\sF) = \h^i(\bar X,j_!\sF)$. This reduces (\ref{II:3-1}) to the
proper case, and explains the use of cohomology with proper support. For $X$
proper, (\ref{II:3-2}) has the form of a Lefschetz trace formula. Note that
the local terms $\tr(F^{n *},\sF_x)$ ($x\in X^{F^n}$) depend only on the
fiber of $\sF$ at $x$ and carry no multiplicity. For $X$ and $\sF$ smooth,
this corresponds to the fact that the graph of $F$ is transverse to the
diagonal (since $d F = 0$).






\subsection{Remark}\label{II:3-7}

If the formalism of $\dQ_\ell$-sheaves is admitted (cf. \ref{II:2-11}; one
needs the Leray spectral sequence and the base-change theorem for the higher
direct images $\eR^q f_!$), it is easy to reduce the proof of
(\ref{II:3-1}, \ref{II:3-2}) to the case where $X_0$ is a smooth curve and
$\sF_0$ is lisse. This is clearly explained in \cite[\S 5]{gr95} (for
\ref{II:3-1}; \ref{II:3-2} is treated in the same way).






\subsection{}\label{II:3-8}

There are two methods for proving (\ref{II:3-2}).



\paragraph{Lefschetz--Verdier}
If $X_0$ is proper (cf. \ref{II:3-6}), the general Lefschetz--Verdier trace
formula expresses the right-hand side of (\ref{II:3-2}) as a sum of local
terms, one for each point of $X^{F^n}$. In the original version of
\cite{sga5}, this formula was proved only modulo resolution of singularities.
The reader will find an unconditional proof in the definitive version. In
the case of curves, to which one can reduce (\ref{II:3-7}), all the
ingredients of the proof were in fact already available.

To deduce (\ref{II:3-2}) from the Lefschetz--Verdier formula, one must be
able to compute its local terms. For a curve and the Frobenius endomorphism,
this had been done by Artin and Verdier \cite{ve67}.



\paragraph{Nielson--Wecken}
A method inspired by the work of Nielsen and Wecken makes it possible to
reduce (\ref{II:3-2}) to a special case proved by Weil; this is what will be
explained in the following sections.











\section{Reduction to theorems in \texorpdfstring{$\dZ/\ell^n$}{Z/l n}-cohomology}\label{II:4}

We shall prove (\ref{II:3-2}) by passage to the limit from an analogous
theorem in $\dZ/\ell^n$-cohomology. Its statement presents the following
difficulty. Let $X_0$ be a separated scheme of finite type over $\dF_q$, and
let $\sF_0$ be a flat constructible sheaf of $\dZ/\ell^n$-modules. Its fibers
are finite-rank free $\dZ/\ell^n$-modules, and the left-hand side of
(\ref{II:3-2}) therefore makes sense (it is an element of $\dZ/\ell^n$).
The $\h_c^i(X,\sF)$, on the other hand, will not in general be free, so the
right-hand side is not defined. Overcoming this difficulty is one of
Grothendieck's essential uses of the theory of derived categories.






\subsection{}\label{II:4-1}

For the definition of the derived category $\eD(\cA)$ of an abelian category
$\cA$, and that of the subcategories $\eD^+(\cA)$, $\eD^-(\cA)$,
$\eD^b(\cA)$, I refer to Verdier's text in this volume. For $\Lambda$ a ring,
write $\eD(\Lambda)$ instead of
$\eD(\text{category of left $\Lambda$-modules})$; for $X$ a site, write
$\eD(X,\Lambda)$ instead of
$\eD($category of sheaves of left $\Lambda$-modules$)$. Likewise for
$\eD^-$, etc\ldots

The sign $\eR$ (resp. $\mathsf{L}$) indicates a right-derived (resp.
left-derived) functor. For example, $\lotimes$ denotes the derived functor of
the tensor product: for $K$ in $\eD^-(\Lambda^\circ)$ and $L$ in
$\eD^-(\Lambda)$ ($\Lambda^\circ$ the opposite ring of $\Lambda$),
$K\lotimes L$ is computed as follows: take quasi-isomorphisms $K'\to K$,
$L'\to L$, with $K'$ and $L$ bounded above and one of them having flat
components, and take the simple complex $s(K'\otimes_\Lambda L)$ associated
with the double complex $K'\otimes_\Lambda L'$.






\subsection{}\label{II:4-2}

Even to treat only the case of $\dZ/\ell^n$-cohomology, we shall need the
theory of traces of endomorphisms of modules over rings that are not
necessarily commutative (finite-group algebras $\dZ/\ell^n[G]$). These
traces were introduced by Stallings and Hattori \cite{ba76}.

Let $\Lambda$ be a ring (with identity, but not necessarily commutative).
Denote by $\Lambda^\natural$ the quotient of the additive group of $\Lambda$
by the subgroup generated by the commutators $a b-b a$. If $f$ is an
endomorphism of the free left $\Lambda$-module $\Lambda^n$, with matrix
$f_i^j$, denote by $\tr(f)$ the image of $\sum_i f_i^i$ in
$\Lambda^\natural$. If $f$ and $g$ are two homomorphisms
$f:\Lambda^n \leftrightarrows \Lambda^m : g$, one immediately verifies that
$\tr(f g) = \tr(g f)$.

Let $f$ be an endomorphism of a projective $\Lambda$-module of finite type
$P$. Write $P$ as a direct factor of a free module $\Lambda^n$: this means
choosing a module $P'$ and an isomorphism
$\alpha:P\oplus P'\iso \Lambda^n$, or equivalently choosing a homomorphism
$a:P\to \Lambda^n$ and a retraction $b:\Lambda^n\to P$ (take
$P'=\ker(b)$). Let $f'=\alpha(f\oplus 0)\alpha^{-1} =af b$. Set
$\tr(f) = \tr(f')$. This element of $\Lambda^\natural$ depends only on $f$:
if another choice $c:P\leftrightarrows \Lambda^m:d$ is made, one has
$a=a(dc)(ba)$, whence
\[
 \tr(a f b) = \tr(adcbafb) = \tr(cbafbad) = \tr(cfd) \text{.}
\]

If $f$ is an endomorphism of a $\dZ/2$-graded projective
$\Lambda$-module of finite type, with components $f_i^j:P^j\to P^i$, set
\[
  \tr(f) = \tr(f_0^0) - \tr(f_1^1) \text{.}
\]
If there is a risk of ambiguity, we shall instead write $\tr(f;P)$. If the
homomorphisms $f:P\leftrightarrows Q:g$ are homogeneous, reduction to the
free case immediately verifies that
\begin{equation*}\tag{4.2.1}\label{II:eq:4-2-1}
  \tr(f g) = \tr\left((-1)^{\deg f \deg g} g f\right) \text{.}
\end{equation*}

If $f$ is an endomorphism of a filtered $\dZ/2$-graded $\Lambda$-module
$(P,F)$, with finite filtration compatible with the grading, and the
$\gr_F^i(P)$ are projective of finite type, then $P$ is projective of finite
type, and
\begin{equation*}\tag{4.2.2}\label{II:eq:4-2-2}
  \tr(f;P) = \sum \tr\left(f; \gr_F^i(P) \right)
\end{equation*}
(split the filtration to reduce to the case of a direct sum).






\subsection{}\label{II:4-3}

A $\dZ$-graded module will be regarded as $\dZ/2$-graded by the parity of
the degree. In particular, if $f$ is an endomorphism of a bounded complex of
projective $\Lambda$-modules of finite type, set
\begin{equation*}\tag{4.3.1}\label{II:eq:4-3-1}
  \tr(f) = \sum (-1)^i \tr(f^i) \text{.}
\end{equation*}

By \eqref{II:eq:4-2-1}, if $f$ is homotopic to zero, $\tr(f) = 0$:
\begin{equation*}\tag{4.3.2}\label{II:eq:4-3-2}
  f = d H + H d \rightarrow \tr(f) = 0 \text{.}
\end{equation*}

Let $\mathsf{K}_\text{parf}(\Lambda)$ be the category whose objects are
bounded complexes of projective $\Lambda$-modules of finite type and whose
arrows are morphisms of complexes taken up to homotopy. The functor
$\mathsf{K}_\text{parf}(\Lambda)\to\eD(\Lambda)$ is fully faithful. Denote its
essential image by $\mathsf{D}_\text{parf}(\Lambda)$. Formulas
\eqref{II:eq:4-3-1} and \eqref{II:eq:4-3-2} make it possible to define
$\tr(f)$ for $f$ an endomorphism of
$K\in\ob \mathsf{D}_\text{parf}(\Lambda)$ (and this trace depends only on the
isomorphism class of $(K,f)$).






\subsection{}\label{II:4-4}

Recall that the filtered derived category $\mathsf{DF}(\cA)$ of an abelian
category $\cA$ is the category obtained from that of complexes of objects of
$\cA$ with finite filtration by inverting the filtered quasi-isomorphisms
($=$ morphisms of filtered complexes inducing a quasi-isomorphism on the
associated graded). There are functors ``underlying complex'':
$\mathsf{DF}(\cA)\to\eD(\cA): (K,F)\mapsto K$, and ``$p$-th component of the
associated graded'':
$\mathsf{DF}(\cA)\to\eD(\cA): (K,F)\mapsto \gr_F^p(K)$.

Let $\mathsf{KF}_\text{parf}(\Lambda)$ be the category whose objects are
bounded filtered complexes $(K,F)$ with finite filtration, with the
$\gr_F^p(K^q)$ projective of finite type, and whose arrows are
filtration-preserving morphisms of complexes taken up to a
filtration-preserving homotopy. The functor
$\mathsf{KF}_\text{parf}(\Lambda)\to\mathsf{DF}(\Lambda)$ is fully faithful;
its essential image $\mathsf{DF}_\text{parf}(\Lambda)$ consists of the $(K,F)$ in
$\mathsf{DF}(\Lambda)$ such that the $\gr_F^p(K)$ belong to
$\mathsf{D}_\text{parf}(\Lambda)$ and are almost all zero; if $(K,F)$ belongs to
$\mathsf{DF}_\text{parf}(\Lambda)$, then $K$ belongs to
$\mathsf{D}_\text{parf}(\Lambda)$.

By \eqref{II:eq:4-2-2}, if $f$ is an endomorphism of
$(K,F)\in\ob\mathsf{DF}_\text{parf}(\Lambda)$, one has
\begin{equation*}\tag{4.4.1}\label{II:eq:4-4-1}
 \tr(f,K) = \sum_p \tr(f,\gr_F^p(K)) \text{.}
\end{equation*}






\subsection{}\label{II:4-5}

Let $\Lambda$ be a ring, and let $K\in\ob\eD^-(\Lambda)$. We say that $K$
has tor-dimension $\leqslant r$ if, for every right $\Lambda$-module $N$, the
hypertors $\h^i(N\lotimes_\Lambda K)$ vanish for $i<-r$. It is said to have
finite tor-dimension if this condition holds for some $r$.

The same terminology is used for a complex of sheaves
$K\in\ob\eD^-(X,\Lambda)$.






\begin{lemma}\label{II:4-5-1}
Let $\Lambda$ be a left Noetherian ring, and let
$K\in\ob\eD^-(\Lambda)$. Then $K$ belongs to
$\mathsf{D}_\text{parf}(\Lambda)$ if and only if it has finite tor-dimension
and the $\h^i(K)$ are of finite type.
\end{lemma}

After replacing $K$ by an isomorphic complex (in $\eD^-(\Lambda)$), we may
suppose that $K$ is bounded above and has finite-rank free components (cf.
\ref{II:4-7} below). If $K$ has tor-dimension $\leqslant r$, the
$\h^i(K)$ vanish for $i<-r$, $K^{-r}/\im(d)$ admits the free resolution
\[
  (\cdots \to K^{-r-1} \to K^{-r}) \to K^{-r}/\im(d) \text{,}
\]
and for every right $\Lambda$-module $N$,
$\h^{-r-k}(N\lotimes_\Lambda K) = \h^{-r-k}(N\otimes_\Lambda K)
= \tor_k(N,K^{-r}/\im d)$ for $k\geqslant 1$.

By hypothesis, these groups vanish; the module $K^{-r}/\im d$ is therefore
flat and finitely presented, i.e. projective of finite type. The quotient
\[
  0 \to K^{-r}/\im(d) \to K^{-r+1} \to K^{-r+2}\to \cdots
\]
of $K$ is quasi-isomorphic to $K$, and this shows that
$K\in\ob\mathsf{D}_\text{parf}(\Lambda)$.






\begin{prop-def_}\label{II:4-6}
Let $X$ be a Noetherian scheme, and let $\Lambda$ be a left Noetherian ring.
Denote by $\eD_\textnormal{ctf}^b(X,\Lambda)$ the subcategory of
$\eD^-(X,\Lambda)$ formed by the complexes $\sK$ satisfying the following
equivalent conditions:
\begin{enumerate}[\indent i)]
  \item $\sK$ is isomorphic (in $\eD^-(X,\Lambda)$) to a bounded complex of
    flat constructible sheaves of $\Lambda$-modules.
  \item $\sK$ has finite tor-dimension and the sheaves $\h^i(K)$ are
    constructible.
\end{enumerate}
\end{prop-def_}

The same argument as in (\ref{II:4-5}), starting from the following lemma.






\begin{lemma_}\label{II:4-7}
If a complex $\sK$ of sheaves of $\Lambda$-modules is such that the
$\h^i(\sK)$ are constructible and vanish for large $i$, there is a
quasi-isomorphism $\sK'\to \sK$, with $\sK'$ bounded above and having flat
constructible components.
\end{lemma_}

Construct $\sK'$ by proceeding from right to left. For an $m$ such that the
$\h^i(\sK)$ ($i\geqslant m$) vanish, take ${\sK'}^i = 0$ for
$i\geqslant m$. Next suppose that the ${\sK'}^i$ have already been
constructed for $i>n$:
\[\xymatrix{
  \cdots \ar[r]
    & \sK^{n-1} \ar[r]
    & \sK^n \ar[r]
    & \sK^{n+1} \ar[r]
    & \cdots \\
  & & & {\sK'}^{n+1} \ar[u] \ar[r]
    & \cdots \text{,}
}\]
with $\h^i(\sK')\iso \h^i(\sK)$ for $i>n+1$, and
$\ker(d)\to\h^{n+1}(\sK)$ surjective in degree $n+1$. Define sheaves $\sA$
and $\sB$ by the Cartesian diagram
\[\xymatrix{
  \sK^n \ar[r]
    & \sK^n/\im(d) \ar[r]
    & \ker(d) \\
  \sA \ar[u] \ar[r]^u
    & \sB \ar[u]\ar[r]
    & \ker(d) \ar[u] \text{.}
}\]

Then $\sB$ is constructible, $u$ is surjective, and to advance one step it
suffices to construct ${\sK'}^n$ flat and constructible and
$v:{\sK'}^n\to \sA$ such that $u v$ is surjective. The possibility of doing so
is guaranteed by the following lemma.






\begin{lemma_}\label{II:4-8}
Let $v:\sA\to \sB$ be an epimorphism of sheaves of $\Lambda$-modules, with
$\sB$ constructible. There is then $u:\sK\to \sA$, with $\sK$ flat and
constructible, such that $v u$ is an epimorphism.
\end{lemma_}

The sheaf $\sA$ is a quotient of a sum of sheaves of the form
$\varphi_! \Lambda$, for $\varphi:U\to X$ étale. Since $\sB$ is Noetherian
(\cite[IX.2.10]{sga4}), a finite sum of these already maps onto $\sB$, which
proves the lemma.






\begin{theorem_}\label{II:4-9}
Let $f:X\to Y$ be a separated morphism of finite type of Noetherian schemes.
If $\sK\in\ob\eD_\textnormal{ctf}^b(X,\Lambda)$, then
$\eR f_!\sK\in \ob\eD_\textnormal{ctf}^b(Y,\Lambda)$.
\end{theorem_}

Recall that if $j:X\hookrightarrow \bar X \xrightarrow{\bar f} Y$ is a
relative compactification of $X/Y$, then by definition
$\eR f_! = \eR \bar f_* \circ j_!$, where $\eR \bar f_*$ is the derived
functor of $\bar f_*$. This formula reduces the proof of (\ref{II:4-9}) to
the case where $f$ is proper. For $f$ proper, $f_*$ has finite cohomological
dimension; this makes it possible to define $\eR f_*$ on the entire derived
category, and also as a functor
\[
  \eR f_* : \eD^-(X,\Lambda) \to \eD^-(Y,\Lambda) \text{.}
\]
By composition, one likewise defines in general
\[
  \eR f_! : \eD^-(X,\Lambda) \to \eD^-(Y,\Lambda) \text{.}
\]

For every right $\Lambda$-module $N$, one has (proof given below)
\begin{equation*}\tag{4.9.1}\label{II:eq:4-9-1}
  N\lotimes \eR f_! \sK \iso \eR f_! (N\lotimes \sK) \text{.}
\end{equation*}

Let us deduce (\ref{II:4-9}) from \eqref{II:eq:4-9-1}.
\begin{enumerate}[\indent a)]
  \item The spectral sequence
    \[
      E_2^{p q} = \eR^p f_! \h^q(\sK) \rightarrow \h^{p+q} \eR f_! \sK \text{.}
    \]
    and the finiteness theorem for $\eR f_!$ show that
    $\eR f_!\sK\in \ob \eD^b(Y,\Lambda)$ and has constructible cohomology.
    It remains to prove finite tor-dimension (\ref{II:4-6}).
  \item If $\sK$ has tor-dimension $\leqslant -r$, the same is true of
    $\eR f_! \sK$: one has
    $\h^i(N\lotimes \eR f_! \sK) = \h^i(\eR f_!(N\lotimes \sK))$, and in the
    spectral sequence above for $N\lotimes \sK$, the $E_2^{p q}$ vanish for
    $q<r$, hence a fortiori for $p+q<r$.
\end{enumerate}

Let us prove \eqref{II:eq:4-9-1}, remaining at the level of complexes.
\begin{enumerate}[\indent a)]
  \item Using the defining formula $\eR f_! = \eR f_*\circ j_!$, we reduce to
    assuming that $f$ is proper.
  \item Represent $\sK$ by a complex bounded above whose components are
    acyclic for $f_*:\eR^p f_* \sK^q = 0$ for $p>0$. One then has
    $\eR f_* K \sim f_* \sK$. It is possible to work in this way with
    complexes bounded above only because $f_*$ has finite cohomological
    dimension.
  \item If $L$ is a free right $\Lambda$-module, one has
    $\eR^p f_*(L\otimes \sK^q) = L\otimes \eR^p f_* \sK^q$: this is trivial
    when $\Lambda$ is free of finite type, and the general case follows by
    passage to the limit. The sheaf $L\otimes \sK^q$ is acyclic for $f_*$,
    and $f_*(L\otimes \sK^q) = L\otimes f_* \sK^q$.
  \item Let $N_\bullet$ be a free resolution of $N$. One has
    $N\lotimes \sK\sim s(N_\bullet\otimes \sK)$ (the simple complex
    associated with the double complex $N_\bullet\otimes \sK$). By c),
    \[
      \eR f_*(N\lotimes \sK) \sim \eR f_* (s(N_\bullet\otimes \sK)) \sim f_*(s(N_\bullet\otimes \sK)) \sim s(N_\bullet\otimes f_* \sK) \sim N\lotimes \eR f_* \sK \text{,}
    \]
    which is the asserted formula.
\end{enumerate}

When $Y$ is the spectrum of a separably closed field, the functor $\eR f_!$
is identified with a functor denoted by
$\eR \Gamma_c : \eD_\text{ctf}^b(X,\Lambda)\to\eD_\text{parf}(\Lambda)$.






\begin{theorem_}\label{II:4-10}
With the notation of \S\ref{II:1}, let $X_0$ be a separated scheme of finite
type over $\dF_q$, and let $\Lambda$ be a Noetherian torsion ring annihilated
by an integer prime to $p$. Let
$\sK_0\in\ob\eD_\textnormal{ctf}^b(X_0,\Lambda)$. With the definition of
traces in (\ref{II:4-3}), one has for every $N$
\begin{equation*}\label{II:eq:4-10}
  \sum_{x\in X^{F^n}} \tr\left({F^n}^*,\sK_x\right) = \tr\left({F^n}^*,\eR\Gamma_c(X,\sK)\right) \text{.}
\end{equation*}
\end{theorem_}

In the remainder of this \S, we shall deduce \ref{II:3-2} from
\ref{II:4-10}, and prove \ref{II:2-10}.






\subsection{Proof of \texorpdfstring{\ref{II:2-10}}{2.10}}\label{II:4-11}

Let $X$ be a separated scheme of finite type over an algebraically closed
field $k$, and let $\sG$ be a $\dQ_\ell$-sheaf on $X$. By \ref{II:2-8}, it
can be written in the form $\sG=\sF\otimes\dQ_\ell$, with $\sF$ a
torsion-free $\dZ_\ell$-sheaf. Set
\[
  K_n = \eR \Gamma_c(X,\sF\otimes\dZ/\ell^{n+1})\in\eD(\dZ/\ell^{n+1}) \text{.}
\]
By \ref{II:4-9}, one has
$K_n\in\ob\eD_\text{parf}(\dZ/\ell^{n+1})$, and it follows from
\eqref{II:eq:4-9-1} that the natural arrow
\begin{equation*}\tag{4.11.1}\label{II:eq:4-11-1}
\xymatrix{
  K_n\lotimes_{\dZ/\ell^{n+1}} \dZ/\ell^n \ar[r] & K_{n-1} \text{,}
}
\end{equation*}
is an isomorphism. The problem of defining this arrow is discussed in
\ref{II:4-12}.

We may now invoke \cite[XV.3.3]{sga5} (p. 31, lines 1--2 from the end; this
passage is independent of the rest of \cite{sga5}) to realize each $K_n$ by
a bounded complex of finite-rank free $\dZ/\ell^{n+1}$-modules, and the
arrows \eqref{II:eq:4-11-1} by \emph{isomorphisms of complexes}
\[\xymatrix{
  K_n\otimes \dZ/\ell^n \ar[r]^-\sim
    & K_{n-1} \text{.}
}\]
Let $K$ be the complex $\varprojlim K_n$. It is a bounded complex of free
$\dZ_\ell$-modules, and
$K_n\simeq K\otimes_{\dZ_\ell}\dZ/\ell^{n+1}$. One has
\[
  \h^i(K) = \varprojlim \h^i(K_n) \text{,}
\]
(indeed, every inverse system of finite groups satisfies the Mittag--Leffler
condition), and $\h^i(K)$ is a $\dZ_\ell$-module of finite type. One has
\[
  \h_c^i(X,\sG) = \left(\varprojlim \h^i(K_n) \right)\otimes_{\dZ_\ell} \dQ_\ell \text{:}
\]
this is a finite-dimensional $\dQ_\ell$-vector space.






\subsection{The arrow \texorpdfstring{\eqref{II:eq:4-11-1}}{(4.11.1)}}\label{II:4-12}

To \emph{define} the arrow \eqref{II:eq:4-11-1}, one cannot proceed as in
\ref{II:4-9}, i.e. resolve $\dZ/\ell^n$ by a complex of free
$\dZ/\ell^{n+1}$-modules, since one wishes to construct an isomorphism in
$\eD_\text{parf}(\dZ/\ell^n)$. The desired arrow is a special case of the
following scalar-extension arrow.

(*) Let $\Lambda\to\Lambda'$ be a homomorphism of Noetherian torsion rings.
For $X$ a separated scheme of finite type over an algebraically closed $k$,
and $\sK\in\ob\eD_\text{ctf}(X,\Lambda)$, one has an isomorphism
\[\xymatrix{
  \eR\Gamma_c(X,\sK)\lotimes_\Lambda \Lambda' \ar[r]^-\sim & \eR \Gamma_c(X,\sK\lotimes_\Lambda \Lambda') \text{.}
}\]

More general arrows are constructed in \cite[XVII 4.2.12]{sga4}. The method
is as follows:
\begin{enumerate}[\indent a)]
  \item The definition of $\eR\Gamma_c$ and the formula
    $j_!(\sK\otimes_\Lambda\Lambda') = (j_! \sK)\otimes_\Lambda\Lambda'$ for
    an open immersion reduce us to the case where $X$ is proper.
  \item For $X$ proper, one must replace $\sK$ by a bounded complex whose
    components are acyclic for the functor $\Gamma$, and whose fiber at every
    geometric point (or merely every closed point) is homotopic to a complex
    of flat $\Lambda$-modules. For such a complex, one has
    $\Gamma(X,\sK)\sim \eR\Gamma(X,\sK)$ and
    $K\otimes_\Lambda\Lambda'\sim \sK\lotimes_\Lambda\Lambda'$; the arrow
    \eqref{II:eq:4-11-1} is defined as the composite
    \[
      \eR\Gamma(X,\sK)\lotimes_\Lambda \Lambda' \to \Gamma(X,\sK)\otimes_\Lambda\Lambda' \to \Gamma(X,\sK\otimes_\Lambda\Lambda') \xleftarrow\sim \eR\Gamma(X,\sK\lotimes_\Lambda\Lambda') \text{.}
    \]
  \item It remains to show that $\sK$ can be replaced by a complex of the
    desired kind \cite[XVII.4.2.10]{sga4}. First replace $\sK$ by a bounded
    complex with flat components (\ref{II:4-6}), then take a truncated
    canonical flasque resolution of it (this does not modify the homotopy
    type of the fiber complexes, and one truncates far enough for the
    acyclicity condition for $\Gamma$ to hold (i.e. beyond the cohomological
    dimension of $\Gamma$)).
\end{enumerate}






\subsection{Proof of \texorpdfstring{(\ref{II:3-2})}{(3.2)}}\label{II:4-13}

To simplify notation, we shall treat only the case $n=1$ of
(\ref{II:3-2}). The general case follows by replacing $F$ below by $F^n$.

Thus let $X_0$ be separated and of finite type over $\dF_q$, and let $\sG_0$
be a $\dQ_\ell$-sheaf on $X_0$. One has
$\sG_0=\sF_0\otimes\dQ_\ell$, for a suitable torsion-free
$\dZ_\ell$-sheaf $\sF_0$. Let
$K_n=\eR\Gamma_c(X,\sF\otimes\dZ/\ell^n)$. The Frobenius morphisms induce
morphisms
\[
  F^* : K_n\to K_n \qquad \text{(in $\eD_\text{parf}(\dZ/\ell^{n+1})$),}
\]
which are deduced from one another via the isomorphisms
\eqref{II:eq:4-11-1}.

As in (\ref{II:4-11}), realize the $K_n$ and the isomorphisms
\eqref{II:eq:4-11-1} at the level of complexes. The already cited passage
\cite[XV.3.3]{sga5} makes it possible to realize the morphisms $F^*$ by
endomorphisms of complexes, again denoted by $F^*$, which reduce to one
another. Again denote by $F^*$ their inverse limit and the endomorphisms
deduced from it. Since
$\h_c^i(X,\sG) = \h^i(K)\otimes\dQ_\ell = \h^i(K\otimes\dQ_\ell)$, one has
(with the sign convention of \ref{II:eq:4-3-1})
\[
  \tr\left(F^*, \h_c^\bullet(X,\sG)\right) = \tr\left(F^*,K^\bullet\otimes\dQ_\ell\right) = \varprojlim_n \tr\left(F^*,K_n^\bullet\right) = \varprojlim_n \tr\left(F^*,\eR\Gamma_c(X,\sF\otimes\dZ/\ell^{n+1})\right) \text{.}
\]
Apply (\ref{II:4-10}) to $\sF_0\otimes\dZ/\ell^{n+1}$ (viewed as a complex
concentrated in degree $0$), for $\Lambda=\dZ/\ell^{n+1}$. One finds that
\begin{align*}
  \tr\left(F^*,\eR\Gamma_c(X,\sF\otimes\dZ/\ell^{n+1})\right)
    &= \sum_{x\in X^F} \tr\left(F^*,\sF_x\otimes\dZ/\ell^{n+1}\right) \\
    &= \sum_{x\in X^F} \tr\left(F^*,\sG_x\right) \mod \ell^{n+1} \text{,}
\end{align*}
and (\ref{II:3-2}) follows by passage to the limit.











\section{The Nielsen--Wecken method}\label{II:5}

In this chapter, we prove two special cases of (\ref{II:4-10}). The general
case will be considered in the next chapter.






\begin{lemma_}\label{II:5-1}
Theorem (\ref{II:4-10}) is true for $X_0$ of dimension $0$.
\end{lemma_}

In dimension $0$, formula \eqref{II:eq:4-10} holds for every correspondence,
and not only for the iterates of a Frobenius. It is a special case of the
following statement (applied to $X(\dF)$, $K$, and the ${F^*}^n$).

(*) Let $X$ be a finite set, $F:X\to X$,
$\sK\in\ob\eD_\text{parf}(X,\Lambda)$, and $F^*:F^* \sK\to \sK$. One has
\[
  \sum_{X^F} \tr\left(F^*,\sK_x\right) = \tr\left(F^*,\Gamma(X,K)\right) \text{.}
\]

One has $\Gamma(X,\sK)=\bigoplus_{x\in X} \sK_x$, and the formula follows
from the fact that for every morphism
$u:\bigoplus_{x\in X} \sK_x\to \bigoplus_{x\in X} \sK_x$, with matrix
$u_x^y$, one has $\tr(u)=\sum_{x\in X} \tr\left(u_x^x\right)$.






\begin{lemma_}\label{II:5-2}
Let $\Lambda$ be a Noetherian torsion ring annihilated by a power of the prime
number $\ell\ne p$, let $X_0$ be a smooth projective connected curve over
$\dF_q$, let $j:U_0\hookrightarrow X_0$ be a dense open subscheme of $X_0$,
and let $\sG_0$ be a locally constant sheaf of projective
$\Lambda$-modules on $U_0$. One has
\[
  \sum_{x\in U^F} \tr\left(F^*,\sG_x\right) = \tr\left(F^*,\eR\Gamma_c(U,\sG)\right) \text{.}
\]
\end{lemma_}

More precisely, we shall deduce (\ref{II:5-2}) from the following theorem.






\begin{theorem_}\label{II:5-3}
Let $X$ be a nonsingular projective curve over an algebraically closed field
$k$, let $f$ be an endomorphism of $X$ with isolated fixed points, and let
$v(f)$ be the number of fixed points of $f$, each counted with its
multiplicity: $v(f)$ is the intersection number of the graph of $f$ with the
diagonal of $X\times X$. Then, for $\ell$ prime to the characteristic
exponent of $k$, one has
\[
  v(f) = \sum (-1)^i \tr\left(f^*,\h^i(X,\dQ_\ell)\right) \text{.}
\]
\end{theorem_}

This theorem is due to Weil (\cite{we48}, n$^\circ$ 43, 66, and 68), and Weil's
proof is reproduced in Lang \cite{la83} (VI \S3 combined with VII \S2,
Th. 3). It can also be deduced from the formalism of the cohomology class
associated with a cycle (\nameref{IV}, \ref{IV:3-7}).






\begin{corollary_}\label{II:5-4}
Let $j:U\hookrightarrow X$ be a dense open subscheme of $X$, such that
$f^{-1}(U)=U$. If the fixed points of $f$ in $X\setminus U$ have multiplicity
one, then the number $v_U(f)$ of fixed points of $f$ in $U$ (each counted
with its multiplicity) is
\[
  v_U(f) = \sum (-1)^i \tr\left(f^*,\h_c^i(U,\dQ_\ell)\right) \text{.}
\]
\end{corollary_}

Let $F$ be the finite set $X\setminus U$. The cohomology exact sequence
defined by the short exact sequence
$0\to j_! \dQ_\ell\to\dQ_\ell\to{\dQ_\ell}_F\to 0$,
\[\xymatrix{
  \ar[r]^-\partial
    & \h_c^i(U,\dQ_\ell) \ar[r]
    & \h^i(X,\dQ_\ell) \ar[r]
    & \h^i(F,\dQ_\ell) \ar[r]^-\partial
    &
}\]
gives, with the sign convention of (\ref{II:3-1}),
\[
  \tr\left(f^*,\h_c^\bullet(U,\dQ_\ell)\right)
    = \tr\left(f^*,\h^\bullet(X,\dQ_\ell)\right) - \tr\left(f^*,\h^\bullet(F,\dQ_\ell)\right)
    = v(f) - \tr(f^*,\dQ_\ell^F) \text{.}
\]
The trace of $f^*$ on $\dQ_\ell^F$ is the number $v_F(f)$ of fixed points of
$f$ in $F$; the multiplicity-one hypothesis therefore ensures that
\[
  \tr\left(f^*,\h_c^\bullet(U,\dQ_\ell)\right)
    = v(f) - v_F(f) = v_U(f) \text{,}
\]
as promised.

The proof of (\ref{II:5-2}) uses some lemmas concerning a special case of
noncommutative traces (\ref{II:4-2}), which we now develop. \emph{For
brevity, we shall say projective for projective of finite type}.






\subsection{}\label{II:5-5}

Let $\Lambda$ be a ring, and let $H$ be a finite group. The map
$\sum a_h \cdot a\mapsto $ class of $a_e$, from the group algebra
$\Lambda[H]$ to $\Lambda^\natural$, vanishes on the commutators $ab-ba$ of
$\Lambda[H]$. It therefore factors through
\[
  \varepsilon : \Lambda[H]^\natural \to \Lambda^\natural \text{.}
\]
If $F$ is an endomorphism of a projective $\Lambda[H]$-module $P$, set
\[
  \tr_\Lambda^H(F) = \varepsilon \tr_{\Lambda[H]} (F)
\]
where the trace of (\ref{II:4-2}) appears on the right-hand side. To avoid
ambiguity, we shall sometimes write this trace as $\tr_\Lambda^H(F,P)$.






\begin{proposition_}\label{II:5-6}
$\tr_\Lambda(F) = |H|\tr_\Lambda^H(F)$.
\end{proposition_}

The formal properties of traces reduce us to treating only the case
$P=\Lambda[H]$. The endomorphism $F$ is then right multiplication by an
element $\sum a_h\cdot h$ of $\Lambda[H]$, and
$\tr_\Lambda(F) = |H| a_e$, $\tr_\Lambda^H(F) = a_e$.






\subsection{}\label{II:5-7}

Let $A$ be a commutative ring, and let $\Lambda$ be an $A$-algebra.
Multiplication by an element of $A$ passes to the quotient to define an
$A$-module structure on $A^\natural$. If $H$ is a finite group, $P$ a
projective $A[H]$-module, and $M$ a $\Lambda[H]$-module projective as a
$\Lambda$-module, it is known that the $\Lambda$-module $P\otimes_A M$,
equipped with the diagonal action of $H$, is a projective
$\Lambda[H]$-module. Indeed:
\begin{enumerate}[\indent a)]
  \item It suffices to verify this for $P=A[H]$.
  \item The $\Lambda$-module $P\otimes_A M$, equipped with the action of $H$
    defined by its action on the first factor, is clearly a projective
    $\Lambda[H]$-module; denote it by $(P\otimes_A M)'$.
  \item The map $h\otimes m\mapsto h\otimes h^{-1} m$ induces an isomorphism
    \begin{equation*}\tag{5.7.1}\label{II:eq:5-7-1}
      A[H]\otimes_A M \iso (A[H]\otimes_A M)' \text{.}
    \end{equation*}
\end{enumerate}






\begin{proposition_}\label{II:5-8}
With the preceding notation, if $u$ is an endomorphism of $P$ and $v$ an
endomorphism of $M$, then
\[
  \tr_\Lambda^H(u\otimes v) = \tr_A^H(u)\tr_\Lambda(v) \text{.}
\]
\end{proposition_}

We reduce to assuming that $P=A[H]$; the endomorphism $u$ is then right
multiplication $m_x$ by an element $x=\sum a_h h$ of $A[H]$. Isomorphism
\eqref{II:eq:5-7-1} takes $u\otimes v$ to
$\sum a_h m_h\otimes h^{-1} v$, whose trace is
$a_e\cdot\tr_\Lambda(v)$.






\subsection{}\label{II:5-9}

Let $G$ be a group that is an extension of $\dZ$ by a finite group $H$
\[
  1 \to H \to G \to \dZ \to 1 \text{,}
\]
and let $G^+$ be the inverse image of $\dN$. Let $\Lambda$ be a ring, and let
$P$ be a $\Lambda[G^+]$-module projective as a $\Lambda[H]$-module. If
$g\in G^+$, denote by $Z_g$ the centralizer of $g$ in $H$; multiplication by
$g$ is an endomorphism of $P$, viewed as a $\Lambda[Z_g]$-module. Since $P$
is $\Lambda[Z_g]$-projective, a trace $\tr_\Lambda^{Z_g}(g)$ is defined. By
\ref{II:5-6}, one has
\begin{equation*}\tag{5.9.1}\label{II:eq:5-9-1}
  \tr_\Lambda(g) = |Z_g| \cdot \tr_\Lambda^{Z_g}(g) \text{.}
\end{equation*}

Suppose that $\Lambda$ is an algebra over a commutative ring $A$. For $P$ an
$A[G^+]$-module projective as an $A[H]$-module, and $M$ a
$\Lambda[G^+]$-module projective as a $\Lambda$-module, the
$\Lambda[G^+]$-module $P\otimes_A M$ (diagonal action of $G^+$) is projective
as a $\Lambda[H]$-module and, by \ref{II:5-8}, for every $g\in G^+$ one has
\begin{equation*}\tag{5.9.2}\label{II:eq:5-9-2}
  \tr_\Lambda^{Z_g}(g,P\otimes_A M) = \tr_A^{Z_g}(g,P) \cdot \tr_\Lambda(g,M) \text{.}
\end{equation*}






\subsection{}\label{II:5-10}

Let $P$ be a $\Lambda[G^+]$-module projective as a $\Lambda[H]$-module, and
let $P_H$ be the coinvariants of $H$ in $P$. This is a projective
$\Lambda$-module. The action of $g\in G^+$ passes to the quotient and defines
an endomorphism of $P_H$ that depends only on the image of $g$ in $\dN$. We
denote by $F$ the action of the elements $g\in G^+$ with image $1$.






\begin{proposition_}\label{II:5-11}
$\tr_\Lambda(F,P_H) = \sum_{g\mapsto 1}' \tr_\Lambda^{Z_g}(g,P)$, where
$\sum_{g\mapsto 1}'$ denotes a sum over the $H$-conjugacy classes of elements
of $G^+$ with image $1$ in $\dN$.
\end{proposition_}

After multiplication by $|H|$, this formula can be verified as follows:
\[
  |H|\cdot \tr_\Lambda(F,P_H)
    = \tr_\Lambda\left(\sum_{g\mapsto 1} g,P\right)
    = \sum_{g\mapsto 1}' \frac{|H|}{|Z_g|}\tr_\Lambda(g,P)
    = \sum_{g\mapsto 1}' |H|\cdot \tr_\Lambda^{Z_g}(g,P) \text{.}
\]

Choose an element $g\in G^+$ with image $1$ in $\dN$, and let $\sigma$ be
the automorphism of $\Lambda[H]$ induced by the automorphism
$\operatorname{ad} g$ of $H$. Giving a $\Lambda[G^+]$-module amounts to
giving a $\Lambda[H]$--module equipped with a $\sigma$-linear endomorphism
$\gamma$: let $g^n h$ act by $\gamma^n h$. In this language, formula
\ref{II:5-11} takes the form: for $P$ a projective $\Lambda[H]$-module and
$\gamma$ a $\sigma$-linear endomorphism of $P$, one has
\[
  \tr_\Lambda(\gamma,P_H) = \sum_h' \tr_\Lambda^{Z_g}(h\gamma,P)\text{,}
\]
where $\sum'$ is a sum over the $\operatorname{ad}g$-conjugacy classes in
$H$ (orbits of $H$ acting on itself by the maps
$x\mapsto h x \operatorname{ad} g(h)^{-1}$).

In this form, the properties of traces immediately reduce us to the case
$P=\Lambda[H]$. The endomorphism $\gamma$ then has the form
$x\mapsto \sigma(x) a$, with $a=\sum a_h h$ in $\Lambda[H]$. It suffices to
treat the universal case where the $a_h$ are indeterminates. In this case,
$\Lambda=\dZ\langle a_h\rangle_{h\in H}$ (noncommutative polynomials), and
$\Lambda^\natural$ is torsion-free: one obtains \ref{II:5-11} by the
argument at the beginning and division by $|H|$. One can also calculate
directly.






\subsubsection{Remark}\label{II:5-11-1}

Everything said above applies equally well to $\dZ/2$-graded modules, or to
complexes of modules (cf. \ref{II:4-2}, \ref{II:4-3}).






\subsection{}\label{II:5-12}

Retain the notation of \ref{II:5-2}, and let $A=\dZ/\ell^n$, with $n$ large
enough that $\Lambda$ is an $A$-algebra. Let $f:V_0\to U_0$ be a finite
\'etale connected covering of $U_0$, Galois with Galois group $H$, and such
that the sheaf $f^*\sG_0$ is constant. Denote its constant value
$\h^0(V_0,f^*\sG_0)$ by $M$; it is a $\Lambda[H]$-module projective as a
$\Lambda$-module.

The sheaf $f_* A$ on $U_0$, equipped with the natural action of $H$, is a
locally free sheaf (of rank one) of $A[H]$-modules. The complex
$\eR\Gamma_c(U,f_* A)$ is therefore an object of
$\eD_\text{parf}\left(A[H]\right)$. It is equipped with a Frobenius
endomorphism $F^*$.

For the natural action of $H$ on $f_* f^*\sG_0$, the trace morphism
$f_* f^*\sG_0\to\sG_0$ factors through an isomorphism
\[
  \left(f_* f^*\sG_0\right)_H \to \sG_0 \text{.}
\]
Moreover, one has $f_* f^*\sG_0=f_* A\otimes_A M$ (with diagonal action of
$H$). The isomorphism
\[\xymatrix{
  \sG_0
    & \ar[l]_-\sim (f_* A\otimes_A M)_H = (f_* A\otimes_A M)\otimes_{\Lambda[H]} \Lambda
}\]
derives (\ref{II:4-12}) to an isomorphism
\[
  \eR\Gamma_c(U,\sG) \sim \left(\eR\Gamma_c(U,f_* A)\otimes_A M\right)\otimes_{\Lambda[H]} \Lambda \text{.}
\]

The endomorphism $F^*$ of the left-hand side is deduced from the
endomorphism $F^*$ of $\eR\Gamma_c(U,f_* A)$. The complex of
$\Lambda[H]$-modules $\eR\Gamma_c(U,f_* A)$, equipped with $F^*$, can be
viewed as a complex of $\Lambda[G^+]$-modules, for $G^+=H\times \dN$.
Formulas \ref{II:5-10} and \eqref{II:eq:5-9-2}, amplified by
\ref{II:5-11-1}, therefore give
\[
  \tr\left(F^*,\eR\Gamma_c(U,\sG)\right) = \sum_h' \tr_A^{Z_h}\left(h F^*,\eR\Gamma_c(U,f_* A)\right)\cdot \tr_\Lambda(h,M)
\]
where $\sum'$ is a sum over the conjugacy classes in $H$. Moreover,
\[
  |Z_h|\cdot \tr_A^{Z_h}\left(h F^*,\eR\Gamma_c(U,f_* A)\right) = \tr_A\left((F h^{-1})^*,\eR\Gamma_c(V,A)\right) \text{.}
\]
This formula remains true for $A=\dZ/\ell^m$, with $m$ increasingly large.
Passing to the limit, one finds that
$\tr_A^{Z_h}\left(h F^*,\eR\Gamma_c(U,f_* A)\right)$ is equal to
\[
  \frac{1}{|Z_h|} \sum (-1)^i \tr\left((F h^{-1})^*,\h_c^i(V,\dQ_\ell)\right)
\]
(an element of $\dZ_\ell$), and
\begin{equation*}\tag{5.12.1}\label{II:eq:5-12-1}
\tr\left(F^*,\eR\Gamma_c(U,\sG)\right) = \sum_h' \frac{1}{|Z_h|} \tr\left((F h^{-1})^*,\h_c^\bullet(V,\dQ_\ell)\right)\cdot \tr_\Lambda(h,M) \text{.}
\end{equation*}






\subsection{}\label{II:5-13}

It remains to apply \ref{II:5-4} (note that if $\bar V$ is the nonsingular
projective completion of $V$, all the fixed points of $F h^{-1}$ have
multiplicity one). This can be done without any calculation:
\begin{enumerate}[\indent A.]
  \item If $U^F=\varnothing$,
    $\tr\left(F^*,\eR\Gamma_c(U,\sG)\right) = 0$. \end{enumerate}

Indeed, $F h^{-1}$ has no fixed point on $V$; \ref{II:5-4} shows that
$\sum (-1)^i \tr\left((F h^{-1})^*,\h_c^i(V,\dQ_\ell)\right)=0$, and one
applies \eqref{II:eq:5-12-1}.
\begin{enumerate}[\indent A.]
\setcounter{enumi}{1}
  \item In the general case, let $j:U'=U\setminus U^F\hookrightarrow U$.
    The filtration $j_! j^* \sG_0\subset \sG_0$ of $\sG_0$ induces a
    filtration of $\eR\Gamma_c(U,\sG)$ with successive quotients
    $\eR\Gamma_c(U',\sG)$ and $\eR\Gamma_c(U^F,\sG)$. By
    \eqref{II:eq:4-4-1} and \ref{II:5-1}, one has
    \[
      \tr\left(F^*,\eR\Gamma_c(U,\sG)\right) = \tr\left(F^*,\eR\Gamma_c(U',\sG)\right) + \sum_{x\in U^F} \tr(F^*,\sG_x) \text{,}
    \]
    and, since ${U'}^F=0$, the first term on the right-hand side is $0$.
\end{enumerate}











\section{End of the proof of \ref{II:4-10}}\label{II:6}

Let us prove \ref{II:4-10} under the additional hypotheses that the ring
$\Lambda$ is annihilated by a power of a prime number $\ell\ne p$, and that
$n=1$. The general case will follow: decompose $\Lambda$ into its
$\ell$-primary components (with $\ell$ running through the prime numbers),
and replace $\dF_q$ by $\dF_{q^n}$ and $X_0$ by
$X_0\otimes_{\dF_q}\dF_{q^n}$.

For $X_0$ separated and of finite type over $\dF_q$, and
$\sK_0\in\ob\eD_\text{ctf}(X_0,\Lambda)$, set
\[
  T'(X_0,\sK_0) = \sum_{x\in X^F} \tr(F^*,\sK_x)
\]
and
\[
  T''(X_0,\sK_0) = \tr(F^*,\eR\Gamma_c(X,\sK)) \text{.}
\]
We must prove the identity
\[
  (1;X_0,\sK_0) \qquad T'(X_0,\sK_0) = T''(X_0,\sK_0) \text{.}
\]
We begin by observing that $T'$ and $T''$ (denoted generically by $T$) obey
the same formalism:

(A) If $j:U_0\hookrightarrow X_0$ is an open subscheme of $X_0$, with
closed complement $Y_0$, one has
\[
  T(X_0,\sK_0) = T(U_0,\sK_0|U_0) + T(Y_0,\sK_0|Y_0) \text{.}
\]
This is clear for $T'$. For $T''$, observe that the filtration
$0\subset j_! j^* \sK_0\subset \sK_0$ of $\sK_0$ induces a filtration of
$\eR\Gamma_c(X,\sK)$ with successive quotients $\eR\Gamma_c(U,\sK)$ and
$\eR\Gamma_c(Y,\sK)$, and invoke \eqref{II:eq:4-4-1} (more precisely stated:
$\sK_0$, equipped with the said filtration, defines an object of
$\mathsf{DF}(X,\Lambda)$; applying $\eR\Gamma_c$ gives an object of
$\mathsf{DF}_\text{parf}(\Lambda)$ whose underlying object is
$\eR\Gamma_c(X,\sK)$ and whose associated graded is as stated).

(B) If $\sK_0$ is a bounded complex with constructible flat components,
\[
  T(X_0,\sK_0) = \sum (-1)^i T(X_0,\sK^i) \text{.}
\]
This is clear for $T'$; for $T''$, apply the argument above to the stupid
filtration of $\sK_0$ to obtain
\[
  T''(X_0,\sK_0) = \sum T''(X_0,\sK_0^i[-i]) = \sum (-1)^i T''(X_0,\sK_0^i) \text{.}
\]

(C) For a morphism $f:X_0\to Y_0$, one has
\[
  T''(X_0,\sK_0) = T''(Y_0,\eR f_! \sK_0)
\]
and the same identity holds for $T'$ if the fibers of $f$ at the rational
points of $Y_0$ satisfy
$\left(1;f^{-1}(y),\sK_0|f^{-1}(y)\right)$: use the fact that
\[
  T'(X_0,\sK_0) = \sum_{y\in Y^F} T'\left(f^{-1}(y), \sK_0|f^{-1}(y)\right)
\]
and
\[
  \eR\Gamma_c(Y, \eR f_! \sK) = \eR\Gamma_c(X,\sK) \text{.}
\]

Using this formalism and \ref{II:4-7}, the proof of \ref{II:4-10} is reduced
to the two special cases \ref{II:5-1} and \ref{II:5-2}.

% END INCLUDED components/sga4half_2_en.tex

% BEGIN INCLUDED components/sga4half_3_en.tex
% !TEX root = ../SGA4half_English_Reader_1_to_8.tex

\chapter{\texorpdfstring{$L$}{L}-functions modulo \texorpdfstring{$\ell^n$}{l n} and modulo \texorpdfstring{$p$}{p}}\label{III}

In this exposé, I prove a formula analogous to the main theorem of the report on
the trace formula (this volume; cited as ``Report'' below) for a coefficient
ring $A$ that is torsion prime to $p$, or a field of characteristic $p$.
Essential use will be made of \cite[XVII 5.5]{sga4}.










\section{Symmetric tensors}\label{III:1}





\subsection{}\label{III:1-1}

Let $A$ be a commutative ring. For a list of properties of the functors
$\Gamma^n:(\text{$A$-modules})\to (\text{$A$-modules})$, I refer to
\cite[XVII 5.5.1 and 2]{sga4}. We recall only that for $M$ flat (in fact, only
the case where $M$ is projective of finite type concerns us), one has
\[
  \Gamma^n M = \left(M^{\otimes n}\right)^{S_n} \qquad \text{(symmetric tensors of degree $n$).}
\]






\subsection{}\label{III:1-2}

Let $X$ be an $S$-scheme, let $X^n/S$ be the $n$-th fiber power of $X$, and
let $\pi$ be the projection from $X^n/S$ onto
$(X^n/S)/S_n=\operatorname{Sym}_S^n(X)$ (assumed to exist). If $\sG$ is a
sheaf of $A$-modules on $X$, the external tensor product of $n$ copies of
$\sG$, namely $\sG^{\boxtimes n}$, is an $S_n$-equivariant sheaf on $X^n/S$.
We shall have to use the \emph{external} functor $\Gamma^n$ of
\cite[XVII 5.5.7 to 9]{sga4}. We recall only that for $\sG$ a sheaf of flat
$A$-modules (in fact, only the constructible flat case concerns us), one has
\[
  \Gamma_\text{ext}^n(\sG) = \left(\pi_*\sG^{\boxtimes n}\right)^{S_n} \text{.}
\]






\begin{lemma_}\label{III:1-3}
For $S=\spec(k)$, with $k$ algebraically closed, $X$ a finite disjoint union of
copies of $S$, and $G$ a flat sheaf of $A$-modules on $X$, one has
\[
  \Gamma\left(\operatorname{Sym}_S^n(X),\Gamma_\text{ext}(\sG)\right) = \Gamma^n \Gamma(X,\sG)
\]
\end{lemma_}

Indeed,
\begin{align*}
  \Gamma\left(\operatorname{Sym}_S^n(X),\Gamma_\text{ext}^n(\sG)\right)
    &= \Gamma\left(\operatorname{Sym}_S^n(X),\pi_*\sG^{\boxtimes n}\right)^{S_n} \\
    &= \Gamma\left(X^n/S,\sG^{\boxtimes n}\right)^{S_n}
    = \left(\Gamma(X,\sG)^{\otimes n}\right)^{S_n}
    = \Gamma^n \Gamma(X,\sG) \text{.}
\end{align*}






\subsection{}\label{III:1-4}

We shall have to use the derived functors of the nonadditive functors
$\Gamma^n$ and $\Gamma_\text{ext}^n$.

The theory of such derived functors is due to Dold and Puppe
\cite{do61}. For more complete reminders than those given below, I refer to
\cite[XVII 5.5.3]{sga4}.

Let $\cA$ be an abelian category (or even an additive category in which every
idempotent endomorphism is the projection onto a direct factor). Denote by $N$
the normalization functor: $N:($cosimplicial complexes of objects of
$\cA)\to($differential complexes of objects of $\cA$, with $K^i$ for
$i<0)$. It is an equivalence of categories. Like its inverse, it takes
homotopic morphisms to homotopic morphisms. If $\Gamma$ is a functor
$\cA\to\cB$, and if $K$ is a complex of objects of $\cA$, in nonnegative
degrees ($K^i=0$ for $i<0$), we set $T K=N T N^{-1} K$.





\subsection{Example}\label{III:1-5}

If $K$ is concentrated in $M$ in degree $0$, $T K$ is concentrated in $T M$
in degree $0$.






\subsection{}\label{III:1-6}

If $K$ is a bounded complex, in nonnegative degrees, of projective
$A$-modules of finite type, the complex of $A$-modules $\Gamma^n K$ has the
same properties.

Denote by $\eD_\text{parf}^{\geqslant 0}(A)$ the subcategory of
$\eD_\text{parf}(A)$ (Report, \ref{II:4-3}) consisting of complexes of
tor-dimension $\leqslant 0$. One sees that $\Gamma^n$ induces a functor
\[
  \eL\Gamma^n : \eD_\text{parf}^{\geqslant 0}(A) \to \eD_\text{parf}^{\geqslant 0}(A) \text{.}
\]

For the analogous, but more complicated, definition of
$\eL\Gamma_\text{ext}^n$, I refer to \cite[XVII 5.5.14]{sga4}.






\subsection{}\label{III:1-7}

Let $K$ be a bounded complex of projective $A$-modules of finite type, and let
$u$ be an endomorphism of $K$. We denote by $\det(1-u t,K)$ the formal series
with constant term $1$
\[
  \det(1-u t,K) = \prod_i \det(1-u t,K^i)^{(-1)^i} \text{.}
\]

The verification of the following properties is left to the reader.


\subsubsection{}\label{III:1-7-1}

Let $F$ be a finite filtration stable under $u$, such that the
$\gr_F^p(K^q)$ are projective of finite type. Then
\[
  \det(1-u t,K) = \prod_p \det(1-u t,\gr_F^p(K)) \text{.}
\]


\subsubsection{}\label{III:1-7-2}

If $K$ is shifted $q$ places to the left, one has
\[
  \det(1-u t,K[q]) = \det(1-u t,K)^{(-1)^q} \text{.}
\]






\begin{proposition_}\label{III:1-8}
Let $u$ be an endomorphism of a bounded complex, in nonnegative degrees, of
projective $A$-modules of finite type. One has \
\begin{equation*}\tag{1.8.1}\label{III:eq:1-8-1}
  \det(1-u t,K)^{-1} = \sum_n \tr(u,\Gamma^n K) t^n \text{.}
\end{equation*}
\end{proposition_}

Denote the right-hand side of \eqref{III:eq:1-8-1} by $D(u,K)$.






\begin{lemma_}\label{III:1-9}
Let $u$ be an endomorphism of a short exact sequence of bounded complexes in
degrees $\geqslant 0$ of projective $A$-modules of finite type
\[
  0 \to K' \to K \to K'' \to 0 \text{.}
\]
One has $D(u,K)=D(u,K')\cdot D(u,K'')$.
\end{lemma_}

If $0\to M'\to M\to M''\to 0$ is a split short exact sequence of
$A$-modules, $\Gamma^n M$ has a natural filtration whose successive quotients
are the $\Gamma^i M'\otimes \Gamma^j M''$ ($i+j=n$). Consequently,
$\Gamma^n K$ has a natural filtration whose successive quotients are the
complexes $N(N^{-1} \Gamma^i K'\otimes N^{-1}\Gamma^j K'')$ ($i+j=n$),
canonically homotopic to the simple complexes
$s(\Gamma^i K'\otimes \Gamma^i K'')$ obtained from the double complexes
$\Gamma^i K'\otimes \Gamma^i K''$, and
\[
  \tr(u,\Gamma^n K) = \sum_{i+j=n} \tr(u,\Gamma^i K')\cdot \tr(u,\Gamma^j K'') \text{,}
\]
whence the lemma. From \ref{III:1-9}, one obtains by induction.






\begin{lemma_}\label{III:1-10}
Let $F$ be a finite filtration of $K$ stable under $u$, such that the
$\gr_F^p(K^q)$ are projective of finite type. Then
\[
  D(u,K) = \prod_p D\left(u,\gr_F^p(K)\right) \text{.}
\]
\end{lemma_}






\begin{lemma_}\label{III:1-11}
Let $u$ be an endomorphism of a projective $A$-module of finite type $M$, and
let $M[-q]$ be the complex concentrated in $M$ in degree $q$. One has
\[
  D\left(u,M[-q]\right) = \left(\sum \tr(u,\Gamma^n M) t^n\right)^{(-1)^q} \qquad \text{($q\geqslant 0$.)}
\]
\end{lemma_}






By \ref{III:1-5}, this is true for $q=0$. Proceeding by induction, it remains
to show that
\[
  D\left(u,M[-q]\right) \cdot D\left(u,M[-(q+1)]\right) = 1 \text{.}
\]
This formula follows from \ref{III:1-9}, applied to the complex $K$
concentrated in $M$ in degrees $q$ and $q+1$ (with
$d=\operatorname{id}$), and to its stupid filtration:
\[
  (\cdots 0 \to M \cdots) \to (\cdots M \to M) \to (\cdots M \to 0 \to 0 \cdots) \text{.}
\]
Indeed, $K$ is homotopic to $0$. In the derived category, one therefore has
$\Gamma^n(K)=\Gamma^n(0)$: $0$ for $n>0$ and $A$ for $n=0$, and $D(u,K)=1$.





\subsection{}\label{III:1-12}

Let us prove \ref{III:1-8}. By \ref{III:1-10}, applied to the ``stupid''
filtration of $K$ by the subcomplexes
\[
  0 \to \cdots \to 0 \to K^q \to K^{q+1} \to \cdots \text{,}
\]
one has
\[
  D(u,K) = \prod_q D\left(u,K^q[-q]\right) =_{\text{(\ref{III:1-10})}} \prod_q \left(\sum \tr\left(u,\Gamma^n(K^q)\right) t^n \right)^{(-1)^q} \text{.}
\]
It remains to prove that for $u$ an endomorphism of a projective module of
finite type $M$, one has
\begin{equation*}\tag{1.12.1}\label{III:eq:1-12-1}
  \det\left(1-u t,M\right)^{-1} = \sum \tr(u,\Gamma^n M) t^n \text{.}
\end{equation*}
If $u=0$, both sides are equal to $1$. By adjoining to $M$ a module $N$ on
which $u=0$, and applying \ref{III:1-9}, we reduce to the assumption that $M$
is free. Choose a basis for it. It is then a matter of proving algebraic
identities among the coordinates of $u$. The principle of extension of
algebraic identities reduces us to assuming that $A$ is an algebraically
closed field. The module $M$ then admits a filtration stable under $u$ with
successive quotients of rank $1$, and \ref{III:1-10} reduces us to the case
where $M$ has rank $1$. Formula \eqref{III:eq:1-12-1} then reduces to the
identity
\[
  \frac{1}{1-a t} = \sum a^n t^n \text{.}
\]





\begin{corollary_}\label{III:1-13}
Let $u$ and $u'$ be endomorphisms of bounded complexes $K$ and $K'$ of
projective $A$-modules of finite type. If, in the derived category, $K$,
equipped with $u$, is isomorphic to $K'$, equipped with $u'$, then
\[
  \det(1-u t,K) = \det(1-u' t,K') \text{.}
\]
\end{corollary_}

Applying \ref{III:1-7-2}, we reduce to assuming that $K$ and $K'$ are in
nonnegative degrees. One then observes that the right-hand side of
\ref{III:1-8} has the required invariance property.






\subsection{}\label{III:1-14}

This corollary makes it possible to define $\det(1-u t,K)$ for $u$ an
endomorphism of $K\in \ob\eD_\text{parf}(A)$. This construction is ad hoc; in
fact, for $v$ an automorphism of $L\in\ob\eD_\text{parf}(\Lambda)$, one can
define $\det(v)\in\Lambda^\times$, and $\det(1-u t,K)$ is obtained by taking
$v=1-u t$, $\Lambda=A\llbracket t\rrbracket$, and
$L=K\otimes_A\Lambda$.









\section{The theorem}\label{III:2}






\subsection{}\label{III:2-1}

We retain the notation of (Report, \S\ref{II:1}). For $A$ a commutative
Noetherian torsion ring, $X_0$ a separated scheme of finite type over
$\dF_q$, and $\sK_0\in\ob\eD_\text{ctf}(X_0;A)$, set
\[
  L(X_0,\sK_0) = \prod_{x\in |X_0|} \det\left(1-F_x^* t^{\deg(x)}, \sK\right)^{-1}
\]
(\ref{III:1-7}, \ref{III:1-13}, and Report, \ref{II:1-5}, \ref{II:1-6}).





\begin{theorem_}\label{III:2-2}
In each of the following two cases
\begin{enumerate}[\indent a)]
  \item $A$ is torsion prime to $p$;
  \item $A$ is reduced and of characteristic $p$
\end{enumerate}
one has
\begin{equation*}\tag{2.2.1}\label{III:eq:2-2-1}
  L(X_0,\sK_0) = \det\left(1-F^* t^f, \eR\Gamma_c(X,\sK)\right)^{-1} \text{.}
\end{equation*}
\end{theorem_}






\subsection{Remark}\label{III:2-3}

In case a), the method of Report, \S\ref{II:3} makes it possible to deduce
from the trace formula (Report \ref{II:4-10}) that the logarithmic derivatives
of the two sides of \eqref{III:eq:2-2-1} are equal. Since the ring $A$ is a
torsion ring, this does not suffice to prove \ref{III:2-2}.






\subsection{Remark}\label{III:2-4}

If $A$ is a field, one has
\[
  L(X_0,\sK_0) = \prod_i L\left(X_0,\underline\h^i(\sK_0)\right)^{(-1)^i}
\]
and the spectral sequence
$E_2^{p q} = \h_c^p(X,\underline\h^q(\sK))\Rightarrow \h_c^{p+q}(X,\sK)$ shows
that
\begin{align*}
  \det\left(1-F^* t, \eR\Gamma_c(X,\sK)\right)
    &= \prod_n \det\left(1-F^* t^f,\h_c^n(X,\sK)\right)^{(-1)^n} \\
    &= \prod_{p,q} \det\left(1-F^* t,\h_c^p(X,\underline\h^q(\sK))\right)^{(-1)^{p+1}} \text{.}
\end{align*}
This reduces \ref{III:2-2} (for $A$ a field) to the following special case.


\subsubsection{}\label{III:2-4-1}

If $\sG_0$ is a constructible sheaf of $A$-vector spaces, one has
\[
  L(X_0,\sG_0) = \prod_n \det\left(1-F^* t^f,\h_c^n(X,\sG)\right)^{(-1)^{n+1}} \text{.}
\]






\subsection{Remark}\label{III:2-5}

It is easy to reduce case b) to the case where $A$ is a field (or even a
finite field) of characteristic $p$; cf. \ref{III:4-2}.






\subsection{Remark}\label{III:2-6}

The special case of b) where $A=\dZ/p$, and where $X$ is reduced to the
constant sheaf $\dZ/p$ in degree $0$,
\[
  Z(X_0,t) = \prod_{x\in |X_0|} \left(1-t^{\deg(x)}\right)^{-1}
      \equiv \prod_n \det\left(1-F^* t^f,\h_c^n(X,\dZ/p)\right)^{(-1)^{n+1}} \pmod p
\]
was proved by N. Katz \cite[XXII 3.1]{sga7}. His method, entirely different
from the one followed here, starts from Dwork's $p$-adic theory of
$Z(X_0,t)$.






\subsection{Remark}\label{III:2-7}


Contrary to what I had rashly asserted to friends, in b) one cannot omit the
hypothesis ``$A$ reduced.'' For a counterexample, cf. \ref{III:4-5}.






\subsection{Outline of the proof of \ref{III:2-2}}\label{III:2-8}

A shift in degrees \ref{III:1-7-2} reduces us to assuming that $\sK_0$ is a
bounded complex, in nonnegative degrees, of constructible flat sheaves of
$A$-modules. One then has
$\eR\Gamma_c(X,\sK)\in \ob\eD_\text{parf}^{\geqslant 0}(A)$ (Report, proof of
\ref{II:4-9}, b). We shall use \ref{III:1-8} to expand both sides of
\eqref{III:eq:2-2-1} as power series in $T=t^f$. The equality of the
coefficients of $T^n$ will follow from \cite[XVII 5.5.21]{sga4} and a trace
formula on $\operatorname{Sym}^n(X)$: in case a), the formula (Report,
\ref{II:4-10}), and in case b), a formula to be proved in the next chapter.
It is first necessary to reduce to the case where the symmetric powers of
$X_0$ exist (for example, the case where $X$ is affine), using the
multiplicativity in $X_0$ of both sides of \eqref{III:eq:2-2-1}: for $U_0$ an
open subscheme of $X_0$, with closed complement $Y_0$, one has
$L(X_0,\sK_0)=L(U_0,\sK_0)\cdot L(Y_0,\sK_0)$, and the analogous formula for
the right-hand side follows from \ref{III:1-7-1} by the method of (Report
\S\ref{II:6}A).






\subsection{The case where $X_0$ is finite}\label{III:2-9}

The multiplicativity in $X_0$ of both sides of \eqref{III:eq:2-2-1} reduces
this case to the case where $X_0=\spec(\dF_{q^k})$. Returning to the
definitions, one sees that \eqref{III:eq:2-2-1} is then equivalent to Report
\ref{II:3-4}.






\subsection{The left-hand side}\label{III:2-10}

By \ref{III:1-8}, one has
\[
  L(X_0,\sK_0) = \prod_{x\in |X_0|} \sum_n \tr\left(F_x^*,\Gamma^n\sK\right) t^{m\deg(x)} \text{.}
\]
The coefficient of $T^n$ is therefore the sum, over those formal linear
combinations $\sum_{x\in |X_0|} n_x\cdot x$ ($n_x\geqslant 0$, with almost
all the $n_x$ zero) of elements of $|X_0|$ such that
$\sum n_x [k(x):\dF_q] = n$, of the products
\begin{equation*}\tag{2.10.1}\label{III:eq:2-10-1}
  \prod_x \tr\left(F_x^*,\Gamma^{n_x} \sK\right) \text{.}
\end{equation*}

For $x\in |X_0|$, denote by $(x)$ the $0$-cycle consisting of the
$[k(x):\dF_q]$ points of $X$ above $x$ (an orbit of $F$). The map
$\sum n_x\mapsto \sum n_x\cdot (x)$ identifies the formal linear combinations
above, with $\sum n_x [k(x):\dF_q]=n$, with the effective $0$-cycles of degree
$n$ on $X$ fixed under $F$, or equivalently with the fixed points of $F$ in
$\operatorname{Sym}^n(X)$. The coefficient of $T^n$ in $L(X_0,\sK_0)$ thus
appears as a sum of products \eqref{III:eq:2-10-1}, indexed by
$\operatorname{Sym}^n(X)^F$. More precisely,





\begin{lemma_}\label{III:2-11}
 $L(X_0,\sK_0) = \sum_n \sum_{x\in\operatorname{Sym}^n(X)^F} \tr\left(F_x^*,\eL\Gamma_\textnormal{ext}^n \sK\right) T^n$.
\end{lemma_}

For $Y_0$ an increasingly large finite subscheme of $X_0$, one has
\[
  L(X_0,\sK_0) = \lim L(Y_0,\sK_0) \text{,}
\]
i.e., for every $m$, there is a finite subscheme $Z_0\subset X_0$ such that
if $Y_0\supset Z_0$, then $L(X_0,\sK_0)$ and $L(Y_0,\sK_0|Y_0)$ are congruent
$\mod T^m$.
It suffices for $Z_0$ to contain the closed points of degree $\leqslant m$.
Since the same fact holds for the right-hand side, it suffices to prove
\ref{III:2-11} for $X_0$ finite. One then has
\begin{align*}
  L(X_0,\sK_0)
    &=_\text{\ref{III:2-9}} \det\left(1-F^* T,\Gamma(X,\sK)\right)^{-1} \\
    &=_\text{\ref{III:1-8}} \sum_n \tr\left(F^*,\Gamma^n \Gamma(X,\sK)\right) T^n \\
    &=_\text{\ref{III:1-3}} \sum_n \tr\left(F^*,\Gamma(\operatorname{Sym}^n(X),\Gamma_\text{ext}^n(\sK)\right) T^n \\
    &= \sum_n \sum_{x\in \operatorname{Sym}^n(X)^F} \tr\left(F_x^*,\Gamma_\text{ext}^n(\sK)\right) T^n
\end{align*}
(by the trace formula in the trivial case of a finite set).






\subsection{The right-hand side}\label{III:2-12}

By \ref{III:1-8} and \cite[XVII 5.5.12]{sga4}, one has
\begin{align*}
  \det\left(1-F^* T,\eR\Gamma_c(X,\sK)\right)^{-1}
    &= \sum_n \tr\left(F^*,\eL\Gamma^n\eR\Gamma_c(X,\sK)\right) T^n \\
    &= \sum_n \tr\left(F^*,\eR\Gamma_c(\operatorname{Sym}^n(X),\Gamma_\text{ext}^n\sK)\right) T^n \text{.}
\end{align*}






\subsection{}\label{III:2-13}

The theorem is therefore equivalent to the trace formulas
\[
  \tr\left(F^*,\eR\Gamma_c(\operatorname{Sym}^n(X,\Gamma_\text{ext}^n \sK)\right) = \sum_{x\in\operatorname{Sym}^n(X)^F} \tr\left(F_x^*,\Gamma_\text{ext}^n\sK\right) \text{.}
\]
In case a), they follow from Report \ref{II:4-10}. The analogous formula in
case b) will be proved in \S\ref{III:4}.










\section{Artin--Schreier theory}\label{III:3}






\subsection{}\label{III:3-1}

In what follows, a coherent sheaf on a scheme $S$ will always be regarded as
a sheaf on $S_\text{ét}$. Let $S$ be a scheme of characteristic $p>0$. If
$G$ is a locally constant sheaf of finite-rank $\dZ/p$-vector spaces,
$\sG=G\otimes_{\dZ/p}\sO$ is a locally free coherent sheaf on $S$. We shall
denote by $\Phi$ both the endomorphism $f\mapsto f^p$ of $\sO$ and its tensor
product with the identity of $G$: $\phi:\sG\to\sG$. Artin--Schreier theory
\cite[IX 3.5]{sga4} asserts that the sequence
\[\xymatrix{
  0 \ar[r]
    & \dZ/p \ar[r]
    & \sO \ar[r]^-{\Phi-1}
    & \sO \ar[r]
    & 0
}\]
is exact. Tensoring with $G$ gives an exact sequence
\[\xymatrix{
  0 \ar[r]
    & G \ar[r]
    & \sG \ar[r]^-{\Phi-1}
    & \sG \ar[r]
    & 0
}\]






\subsection{}\label{III:3-2}

Suppose that $S$ is an open subscheme of a Noetherian scheme $\bar S$. Let
$j:S\hookrightarrow \bar S$ be the inclusion morphism, and let $I$ be a
sheaf of ideals defining the closed complement $\bar S\setminus S$. We again
denote by $\Phi$ the extension of $\Phi$ to $j_*\sG$ by functoriality.






\begin{lemma_}\label{III:3-3}
There exist coherent extensions $\bar\sG\subset j_*\sG$ of $\sG$ such that
\begin{equation*}\tag{3.3.1}\label{III:eq:3-3-1}
  \Phi(\bar\sG)\subset I\cdot \bar\sG\text{.}
\end{equation*}
If $\bar\sG$ is such an extension, the sequence
\begin{equation*}\tag{3.3.2}\label{III:eq:3-3-2}
\xymatrix{
  0 \ar[r]
    & j_! G \ar[r]
    & \bar\sG \ar[r]^-{\Phi-1}
    & \bar\sG \ar[r]
    & 0
}
\end{equation*}
is exact.
\end{lemma_}
\begin{enumerate}[a)]
  \item \emph{Existence}: Let $\sG'\subset j_*\sG$ be a coherent extension
    of $\sG$ to $\bar S$. For every $n$, one has
    $\Phi(I^n \sG')\subset I^{p^n}\Phi(\sG')$. Let $n$ be large enough that
    $I^{n(p-1)-1}\Phi(\sG')\subset \sG'$, and set $\bar\sG=I^n\sG'$.
    One has $\Phi(\bar\sG)\subset I^{p n}\Phi(\sG')
    =I\cdot I^n\cdot I^{n(p-1)-1}\Phi(\sG')\subset I\cdot I^n\sG'
    = I\bar\sG$.
  \item \emph{Kernel of $\Phi-1$}: On $S$, the assertion follows from
    \ref{III:3-1}; it remains to verify that a section $x$ of
    $\ker(\Phi-1)$, on $U$ étale over $\bar S$, vanishes in a neighborhood of
    the inverse image of $\bar S\setminus S$. Indeed, since $x=\Phi(x)$, one
    has by \eqref{III:eq:3-3-1} that $x\in I\bar\sG$; moreover,
    $x\in I^n\bar\sG\Rightarrow x=\Phi(x)\in\Phi(I^n\bar\sG)\subset I^{n p}\bar\sG$, and $x$ lies in every $I^n\bar\sG$, hence is zero near
    $\bar S\setminus S$.
  \item \emph{Surjectivity of $\Phi-1$}: We may suppose $\bar S$ affine, and
    write $\sG$ as a quotient of $\sO^n$; $\Phi$ then lifts to a $p$-linear
    map $\widetilde\Phi:\sO^n\to\sO^n$, and it suffices to verify the
    surjectivity of $\widetilde\Phi-1$:
    \[\xymatrix{
      \sO^n \ar[r]^-{\widetilde \Phi-1} \ar[d]
        & \sO^n \ar[d] \\
      \sG \ar[r]^-{\Phi-1}
        & \sG
    }\]
    The sheaf $\sO^n$ is the sheaf of local sections of the group scheme
    $\dG_a^n$, and $\widetilde\Phi-1$ is induced by a homomorphism from
    $\dG_a^n$ to $\dG_a^n$. This homomorphism is étale. To verify that it is
    surjective, it suffices to check this after every base change
    $\spec(k)\to\bar S$ ($k$ algebraically closed), and over an algebraically
    closed $k$, every étale homomorphism between connected groups of the same
    dimension is surjective. Finally, an étale surjective morphism of
    $\bar S$-schemes induces an epimorphism between the corresponding sheaves
    on $\et{\bar S}$.
\end{enumerate}






\subsection{Frobenius and Frobenius}\label{III:3-4}

The absolute Frobenius morphism $F_\text{abs}:S\to S$ is the identity on the
underlying topological space of $S$, and $f\mapsto f^p$ on the structure
sheaf. For $U/S$ étale, one has a canonical isomorphism
$F_\text{abs}^* U=U$: the diagram
\[\xymatrix{
  U \ar[r]^-{F_\text{abs}} \ar[d]
    & S \ar[d] \\
  S \ar[r]^-{F_\text{abs}}
    & S
}\]
is Cartesian. The endomorphism of the ringed site
$(\et{\bar S},\sO)\to(\et{\bar S},\sO)$ defined by $F_\text{abs}$ can
therefore be described as follows:
\begin{center}
  $F_\text{abs}^*$ is the identity on $\et S$; the homomorphism
  $F_\text{abs}^*\sO=\sO\to\sO$ is $\Phi$.
\end{center}
With this description, for every étale sheaf $\sG$ on $S$, the Frobenius
correspondence $F_\text{abs}^*\sG\to\sG$ is the identity.

Giving a $p$-linear morphism of quasi-coherent sheaves
$u:\sM\to\sN$ on $S$ is equivalent to giving a linear morphism
$u':F_\text{abs}^*\sM\to\sN$. In particular, with the notation of
\ref{III:3-5}, the homomorphism $\Phi:\sG\to\sG$ corresponds to
$\Phi':F_\text{abs}^*\sG\to\sG$.






\begin{lemma_}\label{III:3-5}
With the notation of \ref{III:3-4}, the diagram
\[\xymatrix{
  F_\textnormal{abs}^* G \ar[r] \ar[d]
    & G \ar[d] \\
  F_\textnormal{abs}^*\sG \ar[r]^-{\Phi'}
    & \sG
}\]
is commutative.
\end{lemma_}

It suffices to verify this locally. We may therefore suppose $G$ constant, in
which case the lemma is evident.

The same theory applies to the powers of $F_\text{abs}$.






\subsection{}\label{III:3-6}

Suppose now that $\bar S$ is a scheme over $\dF_q$ ($q=p^f$), and in
accordance with the conventions of Report \S\ref{II:1}, write
$S_0,\bar S_0,G_0,\sG_0,\bar\sG_0$ in place of $S,\dots$. Lemma
\ref{III:3-5}, for $F_\text{abs}^F:S_0\to S_0$, gives a commutative diagram
\[\xymatrix{
  F^* G_0 \ar[r] \ar[d]
    & G_0 \ar[d] \\
  F^*\sG_0 \ar[r]^-{(\Phi^f)'} \ar[r]
    & \sG_0
}\qquad\text{, and then}\qquad
\xymatrix{
  F^*(j_! G_0) \ar[r] \ar[d]
    & j_!G_0 \ar[d] \\
  F^*\bar\sG_0 \ar[r]^-{(\Phi^f)'}
    & \sG_0
}\text{,}
\]
and then, by extension of scalars from $\dF_q$ to $\dF$ and passage to
cohomology,
\[\xymatrix{
  F^*(j_! G) \ar[r]^-{F_*} \ar[d]
    & j_!G \ar[d] \\
  F^*\bar\sG \ar[r]
    & \bar\sG
}\qquad\text{and}\qquad
\xymatrix{
  \h^\bullet(\bar S,j_! G) \ar[r]^-{F^*} \ar[d]
    & \h^\bullet(\bar S,j_!G) \ar[d] \\
  \h^\bullet(\bar S,\bar\sG) \ar[r]
    & \h^\bullet(\bar S,\bar\sG)
}\text{,}
\]
where the second row is obtained by extension of scalars from $\dF_q$ to
$\dF$ of the $\dF_q$-linear endomorphism $\Phi^f$ of
$\h^\bullet(\bar S_0,\bar G_0)$.

Artin--Schreier theory provides a long exact sequence
\[\xymatrix{
  \cdots \ar[r]
    & \h^i(\bar S,j_! G) \ar[r]
    & \h^i(\bar S,\bar\sG) \ar[r]^-{\Phi-1}
    & \h^i(\bar S,\bar\sG) \ar[r]
    & \cdots
}\]
where $\Phi$ is the $p$-linear extension to
$\h^\bullet(\bar S,\bar \sG)= \h^\bullet(\bar S_0,\bar\sG_0)\otimes_{\dF_q}\dF$
of the $p$-linear endomorphism $\Phi$ of
$\h^\bullet(\bar S_0,\bar\sG_0)$.






\subsection{}\label{III:3-7}

If $\bar S_0$ is proper over $\dF_q$, the $\h^i(\bar S,\bar\sG)$ are
finite-dimensional and $\Phi-1$ is surjective (cf. \ref{III:3-6}b). Thus
\[
  \h_c^i(S,G) = \ker\left(\Phi-1 : \h^i(\bar S,\bar G) \to \h^i(\bar S,\bar\sG)\right) \text{,}
\]
and $F^*$ is induced by the $\dF$-linear extension to
$\h^i(\bar S,\bar\sG) = \h^i(\bar S_0,\bar\sG_0)\otimes_{\dF_q}\dF$ of the
$\dF_q$-linear endomorphism $\Phi^f$ of $\h^i(\bar S_0m\bar\sG_0)$.

Applying Lemma \ref{III:3-8} below, one deduces the identity
\begin{equation*}\tag{3.7.1}\label{III:eq:3-7-1}
  \tr\left(F^*,\h_c^i(S,G)\right) = \tr\left(\Phi^f,\h^i(\bar S_0,\bar\sG_0)\right) \text{.}
\end{equation*}






\begin{lemma_}\label{III:3-8}
Let $\Phi$ be a $p$-linear endomorphism of a vector space $V$ over
$\dF_q$: $\phi(\lambda x) = \lambda^p\phi(x)$. Set $F=\Phi^f$ ($F$ is
linear), and again denote by $\Phi$ and $F$ the respective $p$-linear
extensions of $\Phi$ and $F$ to $V\otimes_{\dF_q}\dF$. Then $F$ stabilizes
$\ker(\Phi-1)\subset V\otimes_{\dF_q}\dF$ and
\[
  \tr\left(F,\ker(\Phi-1)\right) = \tr(F,V) = \tr\left(F,V\otimes_{\dF_q}\dF\right) \text{.}
\]
\end{lemma_}

On $V\otimes_{\dF_q}\dF$, one has $F\Phi = \Phi F$: the $p$-linear
endomorphism $F\Phi-\Phi F$ vanishes on $V$, hence everywhere. Write
$V=V'+V''$, with $F$ invertible on $V'$ and nilpotent on $V''$. The theory
of $p$-linear endomorphisms ensures that
\[
  \ker(\Phi-1)\otimes_{\dF_p}\dF \iso V'\otimes_{\dF_q}\dF \text{.}
\]
Since $\tr(F,V'') = 0$, one therefore has
\begin{align*}
  \tr\left(F,\ker(\Phi-1)\right)
    &= \tr\left(F,\ker(\Phi-1)\otimes_{\dF_q}\dF\right) = \tr\left(F,V'\otimes_{\dF_q}\dF\right) \\
    &= \tr\left(F,V\otimes_{\dF_q}\dF\right) = \tr(F,V) \text{.}
\end{align*}











\section{Trace formula modulo \texorpdfstring{$p$}{p}}\label{III:4}






\begin{theorem_}\label{III:4-1}
Let $X_0$ be a separated scheme of finite type over $\dF_q$, let $A$ be a
commutative reduced Noetherian ring of characteristic $p$, and let
$\sK_0\in\ob\eD_\textnormal{ctf}(X_0,A)$. One has
\begin{equation*}\tag{4.1.1}\label{III:eq:4-1-1}
  \sum_{x\in X^F} \tr\left(F_x^*,\sK_x\right) = \tr\left(F^*,\eR\Gamma_c(X,\sK)\right) \text{.}
\end{equation*}
\end{theorem_}






\subsection{Reduction to the case where \texorpdfstring{$A$}{A} is a finite field}\label{III:4-2}

%NOTE: in this paragraph and others, (3.1.1) seems to be referring to (4.1.1)
A standard passage-to-the-limit argument reduces us to the case where $A$ is
of finite type over $\dZ/p$. For every maximal ideal $\fm$ of $A$, the image
of each side of \eqref{III:eq:4-1-1} in $A/\fm$ is given by the same formula,
with $\sK_0$ replaced by
$\sK_0\lotimes_A A/\fm\in\ob\eD_\text{ctf}(X_0,A/\fm)$ (for the right-hand
side, cf. Report \ref{II:4-12}). Since the $A/\fm$ are finite, and the
intersection of the maximal ideals of $A$ is reduced to $0$, one obtains the
stated reduction.






\subsection{Reduction to the case where \texorpdfstring{$A=\dZ/p$}{A=Z/p}}\label{III:4-3}

Suppose that $A$ is a finite field, and denote by $T_A'(\sK_0)$ and
$T_A''(\sK_0)$ the two sides of \eqref{III:eq:4-1-1}. If, by restriction of
scalars, $\sK_0$ is regarded as a complex of sheaves of $\dZ/p$-modules, one
has
\[
  T_{\dZ/p}'(\sK_0) = \tr_{A/(\dZ/p)} T_A'(\sK_0)
\]
and likewise for $T''$. If \eqref{III:eq:4-1-1} holds for $A=\dZ/p$, one
therefore has
\begin{equation*}\tag{4.3.1}\label{III:eq:4-3-1}
  \tr_{A/(\dZ/p)} T_A'(\sK_0) = \tr_{A/(\dZ/p)} T_A''(\sK_0) \text{.}
\end{equation*}
I claim that, for every $\lambda\in A$, one even has
\begin{equation*}\tag{4.3.2}\label{III:eq:4-3-2}
  \tr_{A/(\dZ/p)}\left(\lambda\cdot T_A'(\sK_0)\right) = \tr_{A/(\dZ/p)}\left(\lambda\cdot T_A''(\sK_0)\right) \text{,}
\end{equation*}
and hence $T_A'(\sK_0) = T_A''(\sK_0)$. This is clear for $\lambda=0$. For
$\lambda$ invertible, denote by $A(\lambda)$ the inverse image on $X$ of the
free rank-one sheaf of $A$-modules on $\spec(\dF_q)$ for which
$F^*=\lambda$. Formula \eqref{III:eq:4-3-2} is precisely
\eqref{III:eq:4-3-1} for $A(\lambda)\otimes_A \sK_0$.






% originally theorem 4.1bis (i.e. theorem 4.1 again)
\begin{theorem_}\label{III:4-1bis}
Let $\bar S_0$ be a smooth projective curve, let
$j:S_0\hookrightarrow \bar S_0$ be a dense open subscheme of $\bar S_0$,
and let $G_0$ be a locally constant sheaf of finite-rank $\dZ/p$-vector
spaces on $S_0$. One has
\begin{equation*}\tag{4.4.1}\label{III:eq:4-4-1}
  \sum_{x\in S^F} \tr\left(F_x^*,G_x\right) = \sum_i (-1)^i \tr\left(F^*,\h_c^i(S,G)\right) \text{.}
\end{equation*}
\end{theorem_}
\begin{proof}
Apply \ref{III:3-3} to $(S_0,\bar S_0,G_0)$. The coherent extension
$\bar\sG_0$ of $\sG_0$ whose existence is guaranteed by \ref{III:3-3} is
locally free because it is torsion-free (we have
$\bar\sG_0\subset j_*\sG_0$) and $\bar S_0$ is regular of dimension one. The
$q$-linear endomorphism $\Phi^f$ of $\bar\sG_0$ can be interpreted as a
morphism
\[
  u : F^*\bar\sG_0 \to \bar\sG_0 \text{.}
\]
It factors through $I\bar G$, where $I$ defines the closed subset
$\bar S\setminus S$. Let $v:j_!G_0 \to \bar\sG_0$ be the unique extension of
the natural map $G_0 \to \sG_0 = G_0\otimes_{\dZ/p}\sO$. The diagram
\begin{equation*}\tag{4.4.2}\label{III:eq:4-4-2}
\xymatrix{
  F^* j_! G_0 \ar[r]^-{F^*} \ar[d]^-{F^* v}
    & j_! G_0 \ar[d]^-v \\
  F^*\bar \sG_0 \ar[r]^-u
    & \bar\sG_0
}
\end{equation*}
Finally, denote by $u F^*$ the composite map
\[\xymatrix{
  \h^\bullet(\bar S,\bar\sG) \ar[r]^-{F^*}
    & \h^\bullet(\bar S,F^*\bar\sG) \ar[r]^-u
    & \h^\bullet(\bar S,\bar\sG) \text{,}
}\]
one has by \eqref{III:eq:3-7-1}
\begin{equation*}\tag{4.4.3}\label{III:eq:4-4-3}
  \tr\left(F^*,\h_c^i(S,G)\right) = \tr\left(u F^*,\h^i(\bar S,\bar\sG)\right) \text{.}
\end{equation*}

%NOTE: trace down reference for Woods Hole seminar
For every closed point $i_x:x\hookrightarrow \bar S$ of $\bar S$, set
$\bar\sG_x = i_x^*\bar\sG$. The Lefschetz trace formula in coherent
cohomology (``Woods Hole fixed point formula''), applied to $F$ and $u$, says
that
\begin{equation*}\tag{4.4.4}\label{III:eq:4-4-4}
  \sum_i (-1)^i \tr\left(uF^*, \h^i(\bar S,\bar\sG)\right) = \sum_{x\in \bar S^F} \frac{\tr(u_x,\sG_x)}{\det(1-d F_x)} \text{.}
\end{equation*}
On the right-hand side,
\begin{enumerate}[\indent a)]
  \item since $d F=0$, the denominators are equal to 1;
  \item for $x\in \bar S^F\setminus S^F$, the endomorphism $u_x$ of
    $\bar\sG_x$ is zero, and therefore so is $\tr(u_x,\bar\sG_x)$;
  \item for $x\in S^F$, \eqref{III:eq:4-4-4} gives a contribution
    \[
      \tr(u_x,\bar\sG_x) = \tr(F_x^*,G_x) \text{.}
    \]
    Identity \eqref{III:eq:4-4-1} therefore follows from
    \eqref{III:eq:4-4-3} and \eqref{III:eq:4-4-4}.
\end{enumerate}
\end{proof}






\subsection{A counterexample}\label{III:4-5}

An example shows that the hypothesis that $A$ be reduced cannot be omitted
from \ref{III:4-1}.

Let $Y_0$ be the elliptic curve over $\dF_2$ for which the eigenvalues of
Frobenius are $\pm \sqrt{-2}$. It has $3$ rational points and is
supersingular. The involution $\sigma:x\mapsto -x$ therefore has only one
fixed point, the point $0$. The quotient of $Y_0$ by $\sigma$ is a projective
line; the quotient of $Y_0\setminus 0$ by $\sigma$ is therefore an affine
line $\dF_0^1$, and $Y_0\setminus 0$ is an unramified double covering of
$\dF_0^1$:
\[
  \pi:Y_0\setminus 0 \to \dA_0^1 \qquad \text{, $\pi=\pi\sigma$.}
\]

The sheaf $\pi_*\dZ/2$ is a free module of rank one over the algebra $A$ of
the two-element group $\{1,\sigma\}$ over $\dZ/2$. This algebra is isomorphic
to the algebra of numbers over $\dZ/2$.

The sheaf of $A$-modules $\pi_*\dZ/2$ is our example:
\begin{enumerate}[\indent a)]
  \item One has
    $\h_c^\bullet(\dA^1,\pi_*\dZ/2) = \h_c^\bullet(Y\setminus 0,\dZ/2)$.
    Since $Y_0$ is supersingular, one has $\h^i(Y,\dZ/2)=0$ for $i>0$, and
    the long exact sequence
    \[
      \cdots \to \h_c^i(Y\setminus 0,\dZ/2) \to \h^i(Y,\dZ/2) \to \h^i(\{0\},\dZ/2) \to \cdots
    \]
    shows that $\h_c^\bullet(Y\setminus 0,\dZ/2)=0$. Thus
    $\h_c^\bullet(\dA^1,\pi_*\dZ/2)=0$ and
    \[
      \tr_A\left(F^*,\eR\Gamma_c(\dA^1,\pi_*\dZ/2)\right) = 0 \text{.}
    \]
  \item Since $Y_0\setminus 0$ has two rational points, at one of the
    rational points of $\dA_1^0$, the Frobenius substitution is the identity,
    and at the other, it is $\sigma$. Thus
    \[
      \sum_{x\in {\dA^1}^F} \tr\left(F^*,\pi_*\dZ/2\right) = 1+\sigma \text{,}
    \]
    and $1+\sigma\ne 0$.
\end{enumerate}

The equation of the covering $\pi$ is $y^2-y=x^3$.

% END INCLUDED components/sga4half_3_en.tex

% BEGIN INCLUDED components/sga4half_4_en.tex
% !TEX root = ../SGA4half_English_Reader_1_to_8.tex

\chapter{The cohomology class associated with a cycle}\label{IV}
\blfootnote{by A.\ Grothendieck, written up by P.\ Deligne}





This exposé is inspired by notes of Grothendieck that formed a preliminary
version of \cite[IV]{sga5}. It defines the cohomology class of a cycle $X$ in
a separated smooth scheme of finite type over a field, and proves that
intersection corresponds to the cup product.

Chapter \ref{IV:1} contains some general formalities. In Chapter \ref{IV:2},
the class of a cycle is defined in several situations more general than the
one just mentioned. The principal compatibility not considered is that
between the direct image of a cycle and the trace morphism in cohomology. In
Chapter \ref{IV:3}, this formalism is used to deduce the Lefschetz trace
formula for an endomorphism with isolated fixed points of a proper smooth
scheme over an algebraically closed $k$---and for a Frobenius endomorphism of
a curve.

We adopt the following conventions:
\begin{enumerate}[\indent 1)]
  \item ``scheme'' means a separated Noetherian scheme (this is largely a
    hypothesis of convenience).
  \item In Chapter \ref{IV:2}, an integer $n$ is fixed, and $n$ is invertible
    on all schemes under consideration.
  \item In Chapter \ref{IV:3}, a prime number $\ell$ is fixed, and $\ell$ is
    invertible on all schemes under consideration. The cohomology used is
    always $\ell$-adic cohomology:
    \[
      \h^\bullet(X) = \dQ_\ell\otimes_{\dZ_\ell} \varprojlim \h^\bullet(X,\dZ/\ell^n) \text{.}
    \]
\end{enumerate}
Cohomology classes of cycles, trace morphisms, \ldots{} are defined by passage
to the limit from the case of finite coefficients $\dZ/\ell^n$ (cf.
\cite[VI]{sga5}).




















\section{Cohomology with support and cup products}\label{IV:1}

This section contains reminders from general topology, which the reader is
invited to consult only as the need arises.










\subsection{\texorpdfstring{$H^1$}{H1} and torsors}\label{IV:1-1}






\subsubsection{}\label{IV:1-1-1}

Let $\sF$ be an abelian sheaf on a site $X$. It is known that
$\h^1(X,\sF)$ classifies $\sF$-torsors on $X$. We normalize ($=$ choose the
sign of) the isomorphism $($set of isomorphism classes of
$\sF$-torsors$)\to \h^1(X,\sF)$ in such a way that for every exact sequence
$0\to \sF\xrightarrow\alpha\sG\xrightarrow\beta\sH\to 0$ and every
$h\in\h^0(X,\sH)$, the $\sF$-torsor $\beta^{-1}(h)\subset \sG$, on which
$\sF$ acts by $(f,x)\mapsto \alpha(f)+x$, has class $\partial h$.






\subsubsection{}\label{IV:1-1-2}

Let $\sP$ be an $\sF$-torsor. If $(U_i)$ is an open covering of $X$, and
$p_i$ a section of $\sP$ on $U_i$, one associates to $\sP$ the \v{C}ech
cocycle
\[
  p_{i j} = p_j - p_i \qquad \text{($p_{i j}\in \h^0(U_i\times U_j,\sF)$).}
\]
If, according to the usual rule, the map
$\operatorname{\check H}^\bullet(X,\sF)\to \h^\bullet(X,\sF)$ is defined so
as to be a morphism of $\delta$-functors, the image of
$(p_{i j})\in \operatorname{\check H}^1(X,\sF)$ in $\h^1(X,\sF)$ is the class
of $\sP$, normalized as in \ref{IV:1-1-1}.






\subsubsection{}\label{IV:1-1-3}

The definition of $\h^\bullet(X,\sF)$ is as follows: for $\sF^\bullet$ a
resolution of $\sF$ with acyclic components,
$\h^i(X,\sF) = \h^i\Gamma(X,\sF^\bullet)$. The $\delta$-functor structure is
obtained by associating to a short exact sequence of sheaves a short exact
sequence of resolutions that remains exact after application of the functor
$\Gamma$. If $\sF^\bullet$ is a resolution of $\sF$, the connecting
homomorphism $\partial$ associated with
\[\xymatrix{
  0 \ar[r]
    & \sF \ar[r]
    & \sF^0 \ar[r]^-d
    & \ker(d) \ar[r]
    & 0
}\]
induces the negative of the defining isomorphism
$\h^1(X,\sF) = \Gamma\left(X,\ker(d)\right)/d \Gamma(X,\sF^0)$.






\subsubsection{}\label{IV:1-1-4}

Let $U$ be an open part of $X$ (a subsheaf of the final sheaf), and let $D$
be the ``closed complement.'' It is known that $\h_D^1(X,\sF)$ classifies
the $\sF$-torsors on $X$ trivialized on $U$. For every short exact sequence
$0\to \sF\xrightarrow\alpha \sG\xrightarrow\beta\sH\to 0$, and every
section with support in $D$, $h\in\h_D^0(X,\sH)$, the torsor
$\beta^{-1}(h)$, trivialized on $U$ by the section $0$, has class
$\partial h$.

The long exact sequence of cohomology with support
\[\xymatrix{
  \cdots \ar[r]^-\partial
    & \h_D^i(X,\sF) \ar[r]
    & \h^i(X,\sF) \ar[r]
    & \h^i(U,\sF) \ar[r]^-\partial
    & \cdots
}\]
is defined from the sequence of functors
\[\xymatrix{
  0 \ar[r]
    & \Gamma_D \ar[r]
    & \Gamma \ar[r]
    & \Gamma(U,-) \ar[r]
    & 0
}\]
(exact on injective sheaves). For every section $f\in \h^0(U,\sF)$,
$\partial f\in \h_D^1(X,\sF)$ is the class of the trivial torsor $\sF$,
trivialized on $U$ by the section $f$.






\subsubsection{}\label{IV:1-1-5}

Let $j:U\hookrightarrow X$. If $\sF$ injects into $j_* j^* \sF$, the exact
sequence
\[\xymatrix{
  0 \ar[r]
  & \sF \ar[r]
  & j_* j^* \sF \ar[r]
  & j_* j^* \sF / \sF \ar[r]
  & 0
}\]
gives
$\partial:\h^0(j_* j^*\sF/\sF) = \h_D^0(j_* j^* \sF/\sF) \to \h_D^1(X,\sF)$.
The composite map
\[\xymatrix{
  \h^0(U,\sF) = \h^0(X,j_* j^*\sF) \ar[r]
    & \h^0(X,j_* j^*\sF / \sF) \ar[r]^-\partial
    & \h_D^1(X,\sF)
}\]
is the negative of the map considered in \ref{IV:1-1-4}.






\subsubsection{}\label{IV:1-1-6}

Recall that if $\sL$ is an invertible sheaf on a scheme $X$, the
$\dG_m$-torsor corresponding to it is the sheaf
$\operatorname{Isom}(\sO,\sL)$, on which $\dG_m$ acts by
$(\lambda,f)\mapsto f\circ (\lambda\cdot) = \lambda f$. Recall also that if
$D$ is a Cartier divisor on $X$, and $j:U\hookrightarrow X$ is the inclusion
of the complementary open subscheme, the invertible sheaf $\sO(D)$ is the
subsheaf of $j_*\sO_U$ consisting of local sections $s$ such that $s f$ lies
in $\sO_X$, where $f$ is a local equation for $D$.











\subsection{Cup products}\label{IV:1-2}

In this subsection, we develop some observations about cup products in
cohomology with support that will be used again in [\nameref{V}] to relate
Poincaré duality for curves and self-duality of the Jacobian.

\subsubsection{}\label{IV:1-2-1}

Let $X$ be a site, $Y$ a closed part of $X$, and $\sF,\sG$ two abelian
sheaves on $X$. For example: $X$ a scheme, $Y$ a closed subscheme, and
$\sF,\sG$ sheaves on $\et X$. Define a product
\begin{equation*}\tag{1.2.1.1}\label{IV:eq:1-2-1-1}
  \Gamma_Y(X,\sF)\otimes \Gamma(Y,\sG) \to \Gamma_Y(X,\sF\otimes\sG) \text{.}
\end{equation*}
The product of a section $s$ of $\sF$ with support in $Y$ by a section $t$ of
$\sG$ on $Y$ is obtained as follows: locally on $X$, $t$ is the restriction
to $Y$ of a section $t'$ of $\sG$ on $X$, and one forms the product
$s\otimes t'$. It has support in $Y$ and does not depend on the choice of
$t'$, which justifies the definition and permits it to be globalized.

Let $i:Y\hookrightarrow X$ be the inclusion morphism. The local analogue of
\eqref{IV:eq:1-2-1-1} is the product
\begin{equation*}\tag{1.2.1.2}\label{IV:eq:1-2-1-2}
  i^!\sF\otimes i^*\sG \to i^!(\sF\otimes\sG) \text{,}
\end{equation*}
from which \eqref{IV:eq:1-2-1-1} follows by applying $\Gamma(Y,-)$.






\subsubsection{}\label{IV:1-2-2}

Let us derive these arrows. Let $\sK$, $\sL$, and $\sM$ be in the derived
category, with a bilinear map $\sK\lotimes\sL\to\sM$. Deriving
\eqref{IV:eq:1-2-1-1} gives
\begin{equation*}\tag{1.2.2.1}\label{IV:eq:1-2-2-1}
  \eR\Gamma_Y(X,\sK) \lotimes \eR\Gamma(Y,\sL) \to \eR\Gamma_Y(X,\sM)
\end{equation*}
inducing
\begin{equation*}\tag{1.2.2.2}\label{IV:eq:1-2-2-2}
  \h_Y^i(X,\sK)\otimes \h^j(Y,\sL) \to \h_Y^{i+j}(X,\sM)
\end{equation*}
(the cup product). The local arrow \eqref{IV:eq:1-2-1-2} gives
\begin{equation*}\tag{1.2.2.3}\label{IV:eq:1-2-2-3}
  \eR i^! \sK\lotimes i^*\sL \to \eR i^!\sM \text{,}
\end{equation*}
from which \eqref{IV:eq:1-2-2-1} follows by applying $\eR\Gamma(Y,-)$.





\subsubsection{}\label{IV:1-2-3}

Above, I passed over in silence the conflicts caused by using right-derived
functors ($\eR\Gamma$) and left-derived functors ($\lotimes$) in the same
formula.
\begin{enumerate}[\indent a)]
  \item In the next subsection, all the right-derived functors under
    consideration will have finite cohomological dimension. This makes it
    possible to work systematically in the derived categories $\eD^-$. To
    have complexes that are both flat and flasque, one uses canonical
    flasque resolutions as in \cite[XVII]{sga4}.
  \item For a more general theory, one can no longer interpret bilinear maps
    from $\sK$ and $\sL$ to $\sM$ as morphisms from
    $\sK\lotimes\sL$ to $\sM$. For example,
    $\eR\Gamma_Y(X,\sK)\lotimes\eR\Gamma(Y,\sL)$ is not defined if both
    factors lie in $\eD^+$. One solution is to work in $\eD^+$, define
    \[
      \operatorname{Bil}(\sK,\sL;\sM) = \varinjlim \hom(\sK'\otimes\sL',\sM') \text{,}
    \]
    where the limit is taken over the quasi-isomorphisms $\sK'\iso \sK$,
    $\sL'\iso \sL$, $\sM\iso \sM'$ ($\sK'$, $\sL'$, $\sM'$ bounded
    below), where $\hom$ means ``morphism of complexes, up to homotopy,'' and
    systematically use such bilinear maps without ever mentioning
    $\otimes$.
\end{enumerate}





\subsubsection{Second theme}\label{IV:1-2-4}

Let $U$ be an open subscheme of $X$, and let $j:U\hookrightarrow X$ be the
inclusion morphism. For $\sK$ in the derived category on $U$, set
$\eR\Gamma_!(U,\sK)=\eR\Gamma(X,j_!\sK)$. For $\sK$, $\sL$, and $\sM$ on
$U$, with a bilinear map $\sK\lotimes\sL\to\sM$, we wish to define
\begin{equation*}\tag{1.2.4.1}\label{IV:eq:1-2-4-1}
  \eR\Gamma(U,\sK)\lotimes\eR\Gamma_!(U,\sL)\to\eR\Gamma_!(U,\sM) \text{.}
\end{equation*}
At the level of sheaves and their global sections, such a product follows
from the isomorphism
$j_*\sF\otimes j_!\sG\xleftarrow\sim j_!(\sF\otimes \sG)$, but care must be
taken because $\eR\Gamma_!$ is not in general the derived functor of
$\Gamma(X,j_!-)$.

We begin by defining
\[\xymatrix{
  \eR j_* \sK\lotimes j_! \sL
    & \ar[l]_-\sim j_!(\sK\lotimes \sL) \ar[r]
    & j_!\sM \text{.}
}\]
Applying $\eR\Gamma(X,-)$ gives
\[\xymatrix{
  \eR\Gamma(U,\sK)\lotimes\eR\Gamma_!(U,\sL) = \eR\Gamma(X,\eR j_*\sK)\lotimes \eR\Gamma(X,j_!\sL) \ar[r]
    & \eR\Gamma(X,j_! \sM) \text{.}
}\]

Here again, the conflict between left and right will be resolved in the next
subsection by working in $\eD^-$.






\subsubsection{Coda}\label{IV:1-2-5}

Let $j:U\hookrightarrow X$ be an open part, and let $i:Y\hookrightarrow U$ be
a closed part of $U$. Let $\bar Y$ be a closed subset of $X$ such that
$\bar Y\cap U=Y$ (for example, the complement of $U\setminus Y$). Let $\sK$
be on $U$. Set $\eR\Gamma_{Y!}(U,\sK)=\eR\Gamma_!(Y,\eR i^!\sK)$, where
$\eR\Gamma_!$ is relative to the inclusion of $Y$ into $\bar Y$. For every
open subset $V$ of $U$ containing $Y$, $\eR\Gamma_{Y!}(U,\sK)$ maps to
$\eR\Gamma_!(V,\sK)$: if $i$ also denotes the inclusion of $\bar Y$ into
$X$, $j$ that of $Y$ into $\bar Y$, and $k$ that of $V$ into $X$, one has
$\eR i^!\sK = j^*\eR i^! k_!(k^*\sK)$, whence a morphism
$j_!\eR i^!\sK \to \eR i^! k_!(k^*\sK)$. Applying
$\eR\Gamma(\bar Y,-)$ to it gives
$\eR\Gamma_{Y!}(U,\sK)\to \eR\Gamma_{\bar Y}(k_! k^* \sK) \to \eR\Gamma(k_! k^*\sK) = \eR\Gamma_!(V,\sK)$.

For $\sK,\sL,\sM$ on $U$, with a bilinear map
$\sK\lotimes\sL\to\sM$, we wish to define
\begin{equation*}\tag{1.2.5.1}\label{IV:eq:1-2-5-1}
  \h^n(U,\sK)\otimes \h_!^m(Y,\sL) \to \h_{Y!}^{n+m}(U,\sM)
\end{equation*}
(the last group itself mapping to $\h_!^{n+m}(V,\sM)$). In the derived
category, this means defining
\begin{equation*}\tag{1.2.5.2}\label{IV:eq:1-2-5-2}
  \eR\Gamma_Y(U,\sK) \lotimes \eR\Gamma_!(Y,\sL) \to \eR\Gamma_{Y!}(U,\sM) \text{.}
\end{equation*}
We identify $\eR\Gamma_Y$ with $\eR\Gamma(Y,\eR i^!\sK)$. The desired
product is then of type \eqref{IV:eq:1-2-4-1}, relative to the local product
\eqref{IV:eq:1-2-2-3} on $Y$:
$\eR i^! \sK\lotimes\sL \to \eR i^! \sM$.






\subsubsection{}\label{IV:1-2-6}

Above, we developed the absolute formalism. There is a parallel relative
formalism, with $\Gamma$ replaced by $f_*$ for $f$ a morphism $X\to S$.










\subsection{The Koszul rule}\label{IV:1-3}

Let $A$ be a commutative ring, and let $(V_i)_{i\in I}$ be a finite family of
graded (or $\dZ/2$-graded) $A$-modules. Recall the definition of the graded
tensor product $\bigotimes_{i\in I} V_i$, in the sense of the Koszul rule
(cf. \cite[XVII 1.1]{sga4}). For each total ordering $a$ on $I$, we shall
define a module $V(a)$. We shall also define a transitive system of
isomorphisms $\varphi_{a b}:V(b)\iso V(a)$, and the graded tensor product of
the $V_i$ will be the ``common value'' $\varprojlim V(a)$ of this system of
modules. We take
\begin{enumerate}[\indent a)]
  \item $V(a) = \bigotimes_{i\in I} |V_i|$ (the ordinary tensor product of
    the underlying ungraded modules of the $V_i$).
  \item if the $x_i\in V_i$ are homogeneous, we take
    $\varphi_{a b}(\otimes x_i) = (-1)^N \otimes x_i$, where $N$ is the sum
    of the $\deg(x_i)\deg(x_j)$ over the pairs $(i,j)$ such that
    $i<_a j$ and $i>_b j$.
\end{enumerate}






\subsubsection*{\hypertarget{sga4half.e04.example.1-3-1}{Example 1}}

Take $I=\{1,2\}$, and let $a$ be the ordering in which $1<2$, and $b$ the
ordering in which $1>2$. Let $v_i\in V_i$ be homogeneous, and denote by
$v_1\otimes v_2$ (resp. $v_2\otimes v_1$) the image in the graded tensor
product of the product of the $v_i$ in $V(a)$ (resp. $V(b)$). One has
\[
  v_1\otimes v_2 = (-1)^{\deg(v_1)\deg(v_2)} v_2\otimes v_1
\]
(the Koszul rule).






\subsubsection*{\hypertarget{sga4half.e04.example.1-3-2}{Example 2}}

If all the $V_i$ have degree $1$, there is a canonical isomorphism relating
the module underlying the graded tensor product and the ordinary tensor
product of the modules $|V_i|$ underlying the $V_i$:
\[
  \left|\bigotimes V_i\right| \simeq \bigotimes |V_i|\otimes_{\dZ}\bigwedge^{|I|}\dZ^I \text{.}
\]






\subsubsection*{\hypertarget{sga4half.e04.example.1-3-3}{Example 3}}

For $A$ a field, and $X_i$ a finite family of spaces, the K\"unneth formula
is written
$\h^\bullet\left(\prod X_i,A\right) = \bigotimes \h^\bullet(X_i,A)$.

Let $V^\vee$ be the graded dual of $V$. The canonical map
\begin{equation*}\tag{1.3.1}\label{IV:eq:1-3-1}
  V^\vee\otimes V = V\otimes V^\vee \to A
\end{equation*}
is $v'\otimes v\to v'(v)$.

Suppose now that $A$ is a field, and consider only finite-dimensional vector
spaces. The canonical isomorphism
\begin{equation*}\tag{1.3.2}\label{IV:eq:1-3-2}
  W\otimes V^\vee \to \hom(V,W)
\end{equation*}
is $w\otimes v'\mapsto \left(v\mapsto w\cdot v'(v)\right)$. Via this
isomorphism, the composition
$\hom(Y,Z)\otimes \hom(X,Y)\to \hom(X,Z)$ is identified with the morphism
induced by \eqref{IV:eq:1-3-1}:
$Z\otimes Y^\vee\otimes Y\otimes X^\vee \to Z\otimes X^\vee$.

The trace of $f:V\to V$ (zero for homogeneous $f$ of degree $\ne 0$) is
the image of $f$ under
\begin{equation*}\tag{1.3.4}\label{IV:eq:1-3-4}
  \tr:\hom(V,V) {\xleftarrow\sim}_\text{\eqref{IV:eq:1-3-2}} V\otimes V^\vee = V^\vee \otimes V \to_\text{\eqref{IV:eq:1-3-1}} A .
\end{equation*}

It is easily verified that, for $f$ of degree $0$,
\begin{equation*}\tag{1.3.5}\label{IV:eq:1-3-5}
  \tr(f,V) = \sum (-1)^i \tr(f,V^i) \text{.}
\end{equation*}

Expressing the commutativity of the two morphisms obtained by composing the
morphisms \eqref{IV:eq:1-3-1}
$V^\vee\otimes V\otimes W^\vee\otimes W\to k$, one finds that for homogeneous
$f:V\to W$ and $g:W\to V$, one has
\begin{equation*}\tag{1.3.6}\label{IV:eq:1-3-6}
  \tr(f g) = (-1)^{\deg(f)\deg(g)} \tr(g f) \text{.}
\end{equation*}




















\section{The cohomology class associated with a cycle}\label{IV:2}










\subsection{The class of a divisor}\label{IV:2-1}






\subsubsection{}\label{IV:2-1-1}

Let $D$ be a Cartier divisor in a scheme $X$. Away from $D$, the invertible
sheaf $\sO(D)$ is trivialized by the section $1$. The class $\cl(D)$ of $D$
in $\h_D^1(X,\dG_m)$ is the class of the corresponding $\dG_m$-torsor
trivialized on $X\setminus D$ (\ref{IV:1-1-6} and \ref{IV:1-1-4}).

Let $\partial:\h^i(X\setminus D,\dG_m)\to \h_D^i(X,\dG_m)$ be the morphism
of \ref{IV:1-1-4}. If $D$ admits a global equation $f$, multiplication by
$f$ is an isomorphism from $\sO(D)$, trivialized by $1$ on $X\setminus D$,
to $\sO$, trivialized by $f$ on $X\setminus D$. By \ref{IV:1-1-4}, one
therefore has
\begin{equation*}\tag{2.1.1}\label{IV:eq:2-1-1}
  \cl(D) = \partial f \text{.}
\end{equation*}

For every morphism $u:X'\to X$ such that $u^* D$ is again a Cartier divisor
(i.e., $u^{-1} D$ is disjoint from $\operatorname{Ass}(X')$), one has
$\cl(u^* D) = u^*\cl(D)$. If such functoriality were desired for every
morphism $u$, one would have to consider not Cartier divisors, but more
generally invertible sheaves equipped with a section.

Recall that the integer $n$ is henceforth assumed invertible on the schemes
under consideration. Let
$\partial:\h_D^i(X,\dG_m) \to \h_D^{i+1}(X,\dmu_n)$ be the boundary map for
the Kummer exact sequence $0 \to \dmu_n \to \dG_m \to \dG_m \to 0$.






\begin{definition}\label{IV:2-1-2}
The \emph{class $\cl_n(D)$} of $D$ in $\h_D^2(X,\dmu_n)$ is
$\partial \cl(D)$.
\end{definition}

When there is no risk of confusion, the subscript $n$ will be omitted.






\subsubsection{}\label{IV:2-1-3}

The diagram
\[\xymatrix{
  \h^0(X\setminus D,\dG_m) \ar[r]^-\partial \ar[d]^-\partial
    & \h_D^1(X,\dG_m) \ar[d]^-\partial \\
  \h^1(X\setminus D,\dmu_n) \ar[r]^-\partial
    & \h_D^2(X,\dmu_n)
}\]
is anticommutative. If $D$ admits a global equation $f$, $\cl_n(D)$ is
therefore the negative of the image under $\partial$ of the class in
$\h^1(X\setminus D,\dmu_n)$ of the $\dmu_n$-torsor of $n$-th roots of $f$.






\begin{proposition}\label{IV:2-1-4}
Let $i$ be the inclusion of $D$ into $X$. If $D$ and $X$ are regular, the
cohomology sheaves with support $\eR^p i^! \dmu_n$ are zero for $p=0,1$, and
$\eR^2 i^!\dmu_n = \underline{\dZ/n}$, generated by $\cl_n(D)$.
\end{proposition}

It suffices to prove that, for $X$ strictly local and $D$ defined by a regular
parameter, one has $\h_D^p(X,\dmu_n) = 0$ for $p=0,1$ and
$\h_D^2(X,\dmu_n) = \dZ/n$, generated by $\cl(D)$. Denoting reduced
cohomology by $\sim$, one has
$\widetilde\h^{p-1}(X\setminus D,\dmu_n)\iso \h_D^p(X,\dmu_n)$. The assertion
for $p=0,1$ says that $D$ does not disconnect $X$, and for $p=2$ it follows,
via \ref{IV:2-1-3}, from Abhyankar's lemma.

This is a partial analogue of the relative theorem
(\hyperref[I]{Étale cohomology}, \ref{I:5-3-4}). Grothendieck conjectures
that the $\eR^p i^! \dmu_n$ vanish for $p\ne 2$, at least for $X$ excellent
(purity conjecture), but this is known only in characteristic $0$
(\cite[XIX]{sga4}).






\begin{theorem}[Fundamental compatibility]\label{IV:2-1-5}
Let $X$ be a smooth curve over an algebraically closed field $k$, let $P$ be a
closed point of $X$, and let $\tr$ be the composite
$\h_P^2(X,\dmu_n) \to \h_c^2(X,\dmu_n) \xrightarrow{\tr} \dZ/n$. One has
\[
  \tr \cl(P) = 1 \text{.}
\]
\end{theorem}

Let $\bar X$ be the smooth projective curve completing $X$. The formula says
that the invertible sheaf $\sO(P)$ on $\bar X$ has degree $1$.











\subsection{Cohomological method}\label{IV:2-2}






\subsubsection{}\label{IV:2-2-1}

Let $X$ be a (Noetherian) scheme. Recall that a subscheme $Y$ of $X$ is said
to be a local complete intersection of codimension $c$ if, locally (on $Y$),
it is defined by a regular sequence of $c$ equations in $X$. For $X$ the
spectrum of a local ring $A$ with maximal ideal $\fm$, and $Y$ with ideal
$\fa$, this means that $\Ext^i(A/\fa,A) = 0$ for
$i<\dim(\fa/\fm\fa)=c$, and every sequence of elements of $\fa$ whose image
in $\fa/\fm\fa$ is a basis of $\fa/\fm\fa$ is a regular sequence of
equations for $Y$.






\subsubsection{}\label{IV:2-2-2}

Let $i:Y\hookrightarrow X$ be a local complete intersection of codimension
$c$. We propose to define a local fundamental class $\cl(Y)$ that is a
global section of the cohomology sheaf with support
$\eR^{2c}i^! \dZ/n(c)$ (recall that $\dZ/n(c)=\dmu_n^{\otimes c}$).

Locally, $Y$ is the intersection of a sequence of $c$ divisors $D_i$, and we
define $\cl(Y)$ as the cup product of the $\cl(D_i)$. Each $\cl(D_i)$ has
support in $D_i$; their product has support in $Y$. The fact that this product
is locally independent of the choice of the $D_i$ follows from
\ref{IV:2-2-3} below and the following invariance properties:
\begin{enumerate}[\indent a)]
  \item compatibility with localization;
  \item independence of the ordering of the $D_i$ (the $\cl(D_i)$ have
    degree $2$, which is even, so the cup product is commutative);
  \item the product depends only on the ``flag''
    $D_1\supset D_1\cap D_2\supset \cdots\supset Y$.
\end{enumerate}
To prove c), note the existence of a product (a variant of
\ref{IV:1-2-1})
\[
  \h_{D_1}^\bullet(X)\otimes \h_{D_1\cap D_2}^\bullet(D_1)\otimes \cdots \otimes \h_Y^\bullet(D_1\cap \cdots \cap D_{c-1}) \to \h_Y^\bullet(X) \text{;}
\]
the product of the $\cl(D_i)$ is also the product of the
($\cl(D_i)$ restricted to
$D_1\cap \cdots \cap D_{i-1})=(\cl(D_1\cap \cdots \cap D_i)$ in
$D_1\cap \cdots \cap D_{i-1})$.






\begin{lemma}\label{IV:2-2-3}
Let $A$ be a local ring with maximal ideal $\fm$, and let
$u=(u_1,\dots,u_c)$, $v=(v_1,\dots,v_c)$ be two regular sequences generating
the same ideal $\fa$. There is then a sequence $w_i$
($1\leqslant i\leqslant N$) of such sequences joining them, such that
$w_{i+1}$ is obtained from $w_i$ by a permutation or by changing only the
last element.
\end{lemma}

Since permutations are available, it would amount to the same thing to allow
one element, rather than the last element, to be changed. By
\ref{IV:2-2-1}, we are thus reduced to verifying that, in the vector space
$\fa/\fm\fa$, one can pass from one basis to another by a sequence of
permutations and changes of basis modifying only one vector. Indeed, the
linear group is generated by diagonal, elementary, and permutation matrices.






\subsubsection{}\label{IV:2-2-4}

The same methods make it possible to define a local fundamental class for
every $Y\subset X$ locally definable by $c$ equations (where fewer equations
suffice, the class is zero).






\subsubsection{}\label{IV:2-2-5}

One can pass from such a local fundamental class in
$\h^0\left(Y,\eR^{2 c} i^!\dZ/n(c)\right)$ to a global fundamental class in
$\h_Y^{2 c}\left(X,\dZ/n(c)\right)$ when semipurity results are available:






\begin{proposition}\label{IV:2-2-6}
\begin{enumerate}[(i)]
  \item If $\eR^p i^!\dZ/n=0$ for $p<2 c$, then
    \[
      \h_Y^{2 c}\left(X,\dZ/n(c)\right) \iso \h^0\left(Y,\eR^{2 c} i^! \dZ/n(c)\right) \text{.}
    \]
  \item Let $Z$ be a closed part of $Y$, with complement $V$ in $Y$, and let
    $k$ be the inclusion of $Z$ into $X$. If $\eR^p i^! \dZ/n = 0$ for
    $p\leqslant 2 c$, then
    \[
      \h_Y^{2 c}\left(X,\dZ/n(c)\right) \hookrightarrow \h_V^{2 c}\left(X,\dZ/n(c)\right) \text{.}
    \]
    If $\eR^p i^!\dZ/n=0$ for $p\leqslant 2 c+1$, this arrow is an
    isomorphism.
\end{enumerate}
\end{proposition}

The hypothesis $\eR^p i^!\dZ/n=0$ is equivalent to
$\eR^p i^!\dZ/n(c)=0$. With this understood, (i) follows from the spectral
sequence $\h^p(Y,\eR^q i^!)\Rightarrow \h_Y^{p+q}(X,-)$. If $k_1$ is the
inclusion of $Z$ into $Y$, one has $k=i k_1$, whence
$\eR k^! \dZ/n=\eR k_1^! \eR i^!$, and the long exact cohomology sequence for
$Z\subset Y$ gives a long exact sequence
\[\xymatrix{
  \cdots \ar[r]
    & \h^i\left(Z,\eR k^! \dZ/n(c)\right) \ar[r]
    & \h_Y^i\left(X,\dZ/n(c)\right) \ar[r]
    & \h_V^i\left(X,\dZ/n(c)\right) \ar[r]
    & \cdots
}\]
from which (ii) follows.






\subsubsection{Amplification}\label{IV:2-2-7}

Proposition \ref{IV:2-2-6} remains valid for $i$ any separated morphism of
finite type, and $2 c$ any integer (positive or negative), provided that
$\h_Y^{2 c}\left(X,\dZ/n(c)\right)$ is replaced there by
$\h^{2 c}\left(Y,\eR i^!\dZ/n(c)\right)$, and likewise for $\h_V$.

The following semipurity results follow from \cite[1.8, 1.10, 1.15]{sga2}.
We recall their proof.






\begin{theorem}\label{IV:2-2-8}
Let there be a morphism of $S$-schemes of finite type
\[\xymatrix{
  Y \ar[r]^-i \ar[dr]_-f
    & X \ar[d]^-g \\
  & S
}\]
Suppose that $X$ is smooth, purely of relative dimension $N$, and that $Y$
has fiberwise dimension $\leqslant d$. Let $c=N-d$.
\begin{enumerate}[\indent (i)]
  \item One has $\eR^p i^!\dZ/n=0$ for $p<2 c$; likewise with $\dZ/n$
    replaced by a sheaf $g^* \sF$.
  \item If, over a dense open subset $U$ of $S$, $Y$ has fiberwise dimension
    $<d$, one has $\eR^p i^!\dZ/n=0$ for $p\leqslant 2 c$. If in addition
    the complement $Z$ of $U$ does not locally disconnect $S$, one has
    $\eR^pi^!\dZ/n=0$ for $p\leqslant 2 c+1$.
\end{enumerate}
\end{theorem}

Since $\eR g^! \sF=g^*\sF(N)[2N]$, the transitivity formula
$\eR i^!\eR g^!=\eR f^!$ shows that
$\eR^{2 c+q} i^!(g^*\sF)=0\Leftrightarrow \eR^{-2 d+q}f^!(\sF)$. This reduces
us to studying $f$, i.e. to assuming that $X=S$. Assertion (i) is then
\cite[XVIII 3.17]{sga4}.

Let $Y'$ be the inverse image of $Z$ in $Y$:
\[\xymatrix{
  Y' \ar[r]^-v \ar[d]^-f
    & Y \ar[d]^-f \\
  Z \ar[r]^-u
    & S
}\]
By (i), the $\eR^p f^! \dZ/n$ have support in $Y'$ for
$p\leqslant -2 d+1$; and the spectral sequence
$\eR^a v^!\eR^b f^! \Rightarrow \eR^{a+b} (f v)^!$ shows that they coincide
with the $\eR^p(f v)^!\dZ/n=\eR^p(u f)^!\dZ/n$. Applying (i) to $Y'/Z$ and
to the spectral sequence
$\eR^a f^!\eR^b u^! \Rightarrow \eR^{a+b} (u f)^!$, one finds that (ii)
follows from the vanishing of $\eR^b u^! \dZ/n$ for $b=0$, or for $b=0$ and
$1$, as the case may be.






\subsubsection{}\label{IV:2-2-9}

Grothendieck conjectures the following absolute analogue of
\ref{IV:2-2-8} (semipurity conjecture---a consequence of the purity
conjecture): for $Y$ of codimension $\geqslant c$ in a regular $X$, one has
$\h_Y^i(X,\dZ/n)=0$ for $i<2 c$, at least if $X$ is excellent.






\subsubsection{}\label{IV:2-2-10}

We are now in a position to define the class of a cycle $Y$ of codimension
$c$ in $X$ smooth over a field $k$. Write $Y=\sum d_i Y_i$, where the $Y_i$
are reduced and irreducible. An open subset $U_i$ of $Y_i$, whose complement
has codimension $>c$, is then a local complete intersection in $X$. This
makes it possible to define the local fundamental class of $U_i$. By
\ref{IV:2-2-6} and \ref{IV:2-2-8}, it comes from a unique fundamental class
$\cl(Y_i)\in \h_{Y_i}^{2 c}\left(X,\dZ/n(c)\right)$, and we set
\[
  \cl(Y) = \sum d_i \cl(Y_i) \in \h_{|Y|}^{2 c}\left(X,\dZ/n(c)\right) \text{.}
\]






\subsubsection{}\label{IV:2-2-11}

In complex analytic geometry, a construction of Baum, Fulton, and
Mac Pherson \cite{bfm75} makes it possible to define without restriction the
cohomology class of a local complete intersection $Y\subset X$. Suppose $Y$
is purely of codimension $c$, and let $\sN$ be the vector bundle on $Y$ that
is the normal bundle of $Y$ in $X$. Its local sections are those of
$(\sI/\sI^2)^\vee$, where $\sI$ is the sheaf of ideals of $Y$. Recall that
the sheaf of $C^\infty$ functions on $X$ is defined locally, in terms of local
embeddings of $X$ into $\dC^n$, as the restriction to $X$ of the quotient of
the sheaf of $C^\infty$ functions on $\dC^n$ by the ideal generated by the
real and imaginary parts of the equations defining $X$. The bundle $\sN$
extends to a complex $C^\infty N$ vector bundle on a neighborhood $U$ of
$Y$ in $X$, and for $U$ sufficiently small, there exist sections $f$ of $N$
whose zero locus is $Y$ and such that, on $Y$, $d f:\sN\to N$ is the
identity. Two choices of $N$ and $f$ are homotopic on a sufficiently small
neighborhood of $Y$.

The cohomology class $\cl(U)$ of the section $0$ of $N$ (denoted
$z:U\to N$) is defined: $U$ is a local complete intersection in $N$, and
$\eR^i z^!\dZ = 0$ for $i\ne 2 c$, $\eR^{2 c} z^!\dZ = \dZ$. We have
$f^*:\h_U^\bullet(N,\dZ) \to \h_Y^\bullet(X,\dZ)$ and set
\[
  \cl(Y) = f^*\cl(U) \text{.}
\]

The same construction works whenever an analytic subspace $Y$ of $X$ has
the following normal structure: a locally free sheaf $\sC$ of rank $c$ on
$Y_\text{red}$, and an epimorphism
$\sC\to \sI/\sI\cdot \sI_\text{red}$.










\subsection{Homological method}\label{IV:2-3}






\subsubsection{}\label{IV:2-3-1}

Let $f:Y\to S$ be a flat morphism of finite type with fibers of dimension
$\leqslant d$. In \cite[XVIII 2.9]{sga4}, we defined a trace morphism
\begin{equation*}\tag{2.3.1.1}\label{IV:eq:2-3-1-1}
  \tr_f:\eR^{2 d} f_! \dZ/n(d) \to \dZ/n \text{.}
\end{equation*}
One has $\eR^i f_!\dZ/n(d)=0$ for $i>2 d$, and
$\eR^i f^!\dZ/n=0$ for $i<-2 d$. This, together with the adjunction between
$\eR f_!$ and $\eR f^!$, gives isomorphisms
\begin{equation*}\tag{2.3.1.2}\label{IV:eq:2-3-1-2}
\begin{aligned}
  \hom\left(\eR^{2 d}f_!\dZ/n(d),\dZ/n\right)
    &= \hom\left(\eR f_!\dZ/n(d),\dZ/n[-2 d]\right) \\
    &= \hom\left(\dZ/n,\eR f^!\dZ/n(-d)[-2 d]\right) \\
    &= \h^0\left(Y,\eR^{-2 d}f^!\dZ/n\right) \text{.}
\end{aligned}
\end{equation*}
At least when $S$ is a point, the image of $\tr_f$ in the last two groups
deserves the name fundamental class of $Y$ (in homology).

We again denote by $\tr_f$ the image of $\tr_f$ in the second group, and the
morphisms deduced from it by functoriality. For example, for $Y/S$ proper,
the morphism
\begin{equation*}\tag{2.3.1.3}\label{IV:eq:2-3-1-3}
\xymatrix{
  \h^i\left(Y,\dZ/n(d)\right) = \h^i\left(S,\eR f_!\dZ/n(d)\right) \ar[r]
    & \h^{i-2 d}(S,\dZ/n) \text{.}
}
\end{equation*}

Suppose $Y$ is contained in $X$, smooth over $S$ and purely of relative
dimension $N$, and set $c=N-d$.
\[\xymatrix{
  Y \ar@{^{(}->}[r]^-i \ar[dr]_-f
    & X \ar[d] \\
  & S
}\]
One has $\eR f^! = \eR i^!\eR g^!$, and
$\eR g^! \dZ/n=\dZ/n(-N)[-2 N]$. The fundamental class of $Y$ is identified
with an element of
$\h^0\left(Y,\eR i^!\dZ/n(c)[2 c]\right)=\h_Y^{2 c}\left(X,\dZ/n(c)\right)$,
the class $\cl(Y)$ of $Y$ in $X$. We shall see below that it depends only on
$Y\subset X$, not on the projection of $X$ onto $S$, and that for $Y$ a local
complete intersection, it induces the local class of the preceding
subsection.

Let us make explicit the isomorphism
\[
  \h_Y^{2 c}\left(X,\dZ/n(c)\right) \iso \hom\left(\eR^{2 d}f_!\dZ/n(d),\dZ/n\right)
\]
that takes $\cl(Y)$ to $\tr_f$: via the isomorphisms
\begin{align*}
  \h_Y^{2 c}\left(X,\dZ/n(c)\right)
    &= \h_Y^{2 c}\left(\eR i^!\dZ/n(c)\right) \\
    &= \hom_Y\left(\dZ/n(d)[2 d],\eR i^!\dZ/n(N)[2 N]\right) \\
    &= \hom_X\left(\dZ/n(d)[2 d]_Y,\dZ/n(N)[2 N]\right) \text{,}
\end{align*}
it is the top row of
\[\xymatrix{
  \h_Y^{2 c} \ar[r] \ar[dr]_-{(3)}
    & \hom\left(\eR f_!\dZ/n(d)[2 d],\eR f_! \eR i^!\dZ/n(N)[2 N]\right) \ar[r]^-{(1)} \ar[d]^-{\tr_i} \ar@{}[dr]|-{(2)}
    & \hom\left(\eR f_!\dZ/n(d)[2d],\dZ/n\right) \ar@{=}[d] \\
  & \hom\left(\eR g_!\dZ/n(d)[2 d]_Y,\eR g_!\dZ/n(N)[2 N]\right) \ar[r]^-{\tr_g}
    & \hom\left(\eR g_!\dZ/n(d)[2 d]_Y,\dZ/n\right) \text{,}
}\]
where (1) is the adjunction arrow $\eR f_!\eR f^!\to \text{id}$. The
commutativity of (2) expresses the fact that the isomorphism
$\eR f^!=\eR i^!\eR g^!$ is defined by adjunction. Finally, the arrow
deduced from (3)
\[
  \h_Y^{2c}\left(X,\dZ/n(c)\right) \otimes \eR^{2 d} f_!\dZ/n(d) \to \eR^{2 N} g_!\dZ/n(N)
\]
is interpreted as a cup product (cf. \ref{IV:1-2}).






\begin{definition}\label{IV:2-3-2}
The class $\cl(Y)\in\h_Y^{2 c}\left(X,\dZ/n(c)\right)$ is characterized by
the property that, for every local section $u$ of
$\eR^{2 d}f_! \dZ/n(d)$, one has
\[
  \tr_f(u) = \tr_g\left(\cl(Y)\smallsmile u\right) \text{.}
\]
\end{definition}

Because $\cl(Y)$ already has an image $\tr_f$ in
$\hom\left(\eR f_!\dZ/n(d)[2 d],\dZ/n\right)$, formula
\ref{IV:2-3-2} holds for the arrows deduced from $\tr_f$ by functoriality.
For example, for $X$ proper over $S$ and
$u\in \h^\bullet(Y,\dZ/n)$, the formula holds in
$\h^\bullet\left(S,\dZ/n(-d)\right)$. For
$v\in \h^\bullet(X,\dZ/n)$, it gives
\[
  \tr_f(i^* v) = \tr_g\left(\cl(Y)\smallsmile v\right) \text{,}
\]
where the cup product can be computed in $\h^\bullet(X,\dZ/n)$.






\subsubsection{}\label{IV:2-3-3}

We shall define the trace morphisms, and hence the class $\cl(Y)$, under more
general hypotheses. Consider a diagram
\[\xymatrix{
  Y \ar@{^{(}->}[r] \ar[dr]_-f
    & X \ar[d]^-g \\
  & S
}\]
with $g$ smooth and purely of relative dimension $N$, and $Y$ closed in $X$
with fiberwise dimension $\leqslant d$. Equip $Y$ in addition with a
``weight'' $\sK$ of the following kind:
$\sK\in\eD_{\text{parf},X}^b$ is a bounded complex of sheaves of
$\sO$-modules on $X$, of finite $\tor$-dimension over $S$ (or over $X$,
which amounts to the same thing), whose cohomology sheaves are coherent and
have support in $Y$. We propose to define a morphism
$\tr_{f,\sK}:\eR^{2 d}f_! \dZ/n(d) \to \dZ/n$. For $Y$ flat over $S$ and
$\sK=\sO_Y$, this will be the preceding trace morphism. In general, it
depends only on the lengths of $\sK$ at the generic points $y$ of the
irreducible components of
$Y$ ($\operatorname{lg}_y(\sK) = \sum (-1)^i \operatorname{lg} \h^i(\sK)_y$).
% NOTE: notation is not clear, is it a calligraphic H?
Finally, $\cl(Y,\sK)$ will denote the class with support in $Y$ such that
\begin{equation*}\tag{2.3.3.1}\label{IV:eq:2-3-3-1}
  \tr_g\left(\cl(Y,\sK)\smallsmile u\right) = \tr_{f,\sK}(u) \text{.}
\end{equation*}






\subsubsection{}\label{IV:2-3-4}

The construction of $\tr_{f,\sK}$ is parallel to that of
\cite[XVIII.2]{sga4}; we shall indicate only its main outlines.

\paragraph{A}
$d=0$ ($f$ quasi-finite). Let $x$ be a geometric point of $Y$, $s$ its image
in $S$, $\sK_{(x)}$ the inverse image of $\sK$ on the strict localization of
$X$ at $x$, $\sK_{(x)s} = \sK_{(x)}\lotimes_{\sO_{S,s}^h} k(s)$ its inverse
image on the geometric fiber at $s$, and
$n(x)=\sum (-1)^i \dim_{k(s)} \h^i(\sK_{(x)s})$. The function
$x\mapsto n(x)$ is a weighting of $f$, and we take the corresponding trace
morphism \cite[XVII.6.2.5]{sga4}. Formation of the weighting $n(x)$ and of
$\tr_{f,\sK}$ is compatible with every base change.

\paragraph{B}
General case. If $u:X\to \dA_S^d$ is such that $u i$ is quasi-finite, set
$\tr_{f,\sK} = \tr_{\dA_S^d}\circ \tr_{u i,\sK}$. This trace morphism is
clearly compatible with every base change $S'/S$. To prove that it does not
depend on $u$, we may therefore suppose that $S$ is the spectrum of an
algebraically closed field $k$. In that case, let $Y_i$ be the irreducible
components of $Y$, and let $\operatorname{lg}_i$ be the length of $\sK$ at
the generic point of $Y_i$. If $\alpha_i$ is the inclusion of
$(Y_i)_\text{red}$ into $Y$, one has
\begin{equation*}\tag{2.3.4}\label{IV:eq:2-3-4}
  \tr_{f,\sK} = \sum \operatorname{lg}_i \tr_{Y_{i,\text{red}/S}} \alpha_i^\ast
\end{equation*}
(this formula follows from the analogous and easy formula for $\tr_{u i}$).

This defines $\tr_{f,\sK}$ locally on $Y$; one then proceeds as in
\cite[XVIII.2.9]{sga4}.






\begin{lemma}\label{IV:2-3-5}
If $g=g'g'':X\xrightarrow{g''} S'\xrightarrow{g'} S$, with $g'$ and $g''$
smooth and purely of relative dimensions $N'$ and $N''$, and if $Y$ still
has fiberwise codimension $\geqslant c$ over $S'$, then the class
$\operatorname{cl}(Y,\sK)$ is the same whether computed in terms of $g$ or
of $g''$.
\end{lemma}

Let $f'=g'' i$ and $d'=d-N'$. The formula
$\tr_g=\tr_{g'}\tr_{g''}$ (\cite[XVIII.2.9]{sga4} Var 3) ensures the
commutativity of
\[\xymatrix{
  \h_Y^{2 c}\left(X,\dZ/n(c)\right) \ar[d] \\
  \hom_{S'}\left(\eR f_!\dZ/n(d')[2 d'],\eR g_!''\dZ/n(N'')[2 N'']\right) \ar[r] \ar[d]^-{\eR g_!'(N')}
    & \hom_{S'}\left(\eR^{2 d'} f_!' \dZ/n,\dZ/n\right) \ar[d]^-{\tr_{g'}\circ \eR^{2 N'} g_!'} \\
  \hom_S\left(\eR f_!\dZ/n(d)[2d],\eR g_!'\dZ/n(N)[2 N]\right) \ar[r]
    & \hom_S\left(\eR^{2 d} f_!\dZ/n,\dZ/n\right)
}\]
while the analogous and easily verified formula
$\tr_{f,\sK}=\tr_{g'}\tr_{f',\sK}$ ensures that the image of
$\tr_{f',\sK}$ is $\tr_{f,\sK}$. The assertion follows.






\begin{lemma}\label{IV:2-3-6}
If $Y$ is a divisor in $X$, the class of \ref{IV:2-3-2} coincides with the
class of \ref{IV:2-1-2}.
\end{lemma}

By semipurity (\ref{IV:2-2-6}(i) and \ref{IV:2-2-8}(i)), the problem is local
on $Y$; using \ref{IV:2-3-5} to replace $S$ by a suitable $S'$, this allows
us to suppose that $N=1$. Using semipurity once more
(\ref{IV:2-2-6}(ii) and \ref{IV:2-2-8}(ii)), it suffices to prove
\ref{IV:2-3-6} over the generic points of $S$. An étale localization then
reduces us to assuming that $S$ is the spectrum of an algebraically closed
field $k$. Since the classes of \ref{IV:2-3-2} and \ref{IV:2-1-2} are each
additive in $Y$, this reduces us to the case where $Y$ is a closed point on a
smooth curve over $k$. The fact that the characteristic property
\ref{IV:2-3-2} is satisfied is then the fundamental compatibility
\ref{IV:2-1-5}.






\begin{lemma}\label{IV:2-3-7}
The formation of $\cl(Y,\sK)$ is compatible with every base change $S'/S$.
\end{lemma}

This follows from the same assertion for the trace morphisms.






\begin{theorem}\label{IV:2-3-8}
(i) $\cl(Y,\sK)$ depends only on $Y\subset X$ and on the lengths of $\sK$ at
the generic points of $Y$. This class is additive in $\sK$. For
$\sK=\sO_Y$ and $Y$ a local complete intersection, it induces the local
class of \ref{IV:2-2-2}.

(ii) Let there be a commutative diagram
\[\xymatrix{
  X' \ar[r]^-u \ar[d]
    & X \ar[d] \\
  S' \ar[r]
    & S
}\]
with $X'$ smooth over $S'$. If $u^{-1}(Y)$ still has fiberwise codimension
$\geqslant c$, then
\[
  \cl\left(u^{-1}(Y),\eL u^\ast\sK\right) = u^\ast \cl(Y,\sK) \text{.}
\]

(iii) Let $Y'\subset X$ have fiberwise codimension $\geqslant c'$, let
$\sK'$ be a weight on $Y'$, and suppose that $Y\cap Y'$ has fiberwise
codimension $\geqslant c+c'$. Then
\[
\cl\left(Y\cap Y',\sK\lotimes \sK'\right) = \cl(Y,\sK) \smallsmile \cl(Y',\sK') \text{.}
\]
\end{theorem}



% The following paragraphs were originally numbered

% (2.3.8.1)
\paragraph{(A)}
Given $Y\subset X/S$ as above, it follows from semipurity
(\ref{IV:2-2-6}, \ref{IV:2-2-8}) that to verify that two classes in
$\h_Y^{2 c}\left(X,\dZ/n(c)\right)$ coincide, it suffices to verify this
locally at the generic points of $Y$, or even after every base change
$s\to S$ ($s$ a geometric generic point of $S$), at the generic points of
$Y$.

% (2.3.8.2)
\paragraph{(B) Proof of (ii) for $S=S'$ and $u$ a closed immersion}
By localization (A), we may suppose that $X'$ is the inverse image of the
section $0$ under a smooth morphism $v:X\to \dA_S^{N'}$
\[\xymatrix{
  X' \ar[r]^-u \ar[d]
    & X \ar[d]^-v \\
  S \ar[r]^-0
    & \dA_S^{N'} \ar[r]^-\pi
    & S
}\]
Apply \ref{IV:2-3-5} to $g=\pi v$ and \ref{IV:2-3-7} to the base change
$S\to \dA_S^{N'}$.

% (2.3.8.3)
\paragraph{(C) In a product situation:}
$X=X'\times_S X''$, $Y=Y'\times_S Y''$, $\sK=\sK'\lotimes_S \sK''$, one has
$\cl(Y,\sK) = \cl(Y',\sK')\smallsmile \cl(Y'',\sK'')$ (external cup product,
i.e. $\operatorname{pr}_1^\ast \cl(Y',\sK')\smallsmile \operatorname{pr}_2^\ast \cl(Y'',\sK'')$).

The functorial properties of $\tr_g$ \cite[XVIII.2.12]{sga4} imply the
commutativity of
\small
\[\xymatrix@=0.5cm{
  \h_{X'}^{2c'}\otimes \h_{X''}^{2 c''}\ar[r] \ar[d]
    & \hom_S\left(\eR f_!' \dZ/n(d')[2 d'],\eR g_!'\dZ/n(N')[2 N']\right)\otimes \cdots \ar[r] \ar[d]
    & \hom_S\left(\eR^{2 d'} f_!' \dZ/n,\dZ/n\right) \otimes \cdots \ar[d] \\
  \h_X^{2 c} \ar[r]
    & \hom_S\left(\eR f_!\dZ/n(d)[2 d],\eR g_!\dZ/n(N)[2 N]\right) \ar[r]
    & \hom_S\left(\eR^{2 d}f_!\dZ/n,\dZ/n\right)
}\]
\normalsize
and one verifies that $\tr_f$ is the image of
$\tr_{f'}\otimes \tr_{f''}$, whence the assertion.

\paragraph{Proof of (ii)}
(\ref{IV:2-3-7}) and a preliminary base change reduce us to assuming that
$S=S'$. Factor $u$ through its graph:
\[\xymatrix{
  X' \ar[r]^-{(\text{id},u)}
    & X'\times X \ar[r]^-{\operatorname{pr}_2}
    & X \text{.}
}\]
This leads us to treat $\operatorname{pr}_2$ and $(\text{id},u)$ separately:
the former is covered by (C) with $Y'=X'$, $\sK'=\sO_{X'}$ (with
$\cl(Y')=1\in \h^0(X',\dZ/n)$), and the latter by (B).

\paragraph{Proof of (iii)}
Let $\delta:X\to X\times_S X$ be the diagonal inclusion. We use the formula
$\delta^\ast(Y_1\times_S Y_2) = Y_1\cap Y_2$ to reduce to (B) and (C).

\paragraph{Proof of (i)}
The base change $S_\text{red}\to S$ replaces $X$ by $X_\text{red}$; this
reduces us to assuming $X$ and $S$ reduced. Localization near the generic
points of $Y$ then reduces us to the case where $Y_\text{red}$ is irreducible
and the complete intersection of $c$ divisors $D_i$. The same localization,
together with the definition of $\tr$ in the quasi-finite case, then shows
that
$\cl(Y,\sK)=\operatorname{lg}\cl(Y_\text{red},\sO_{Y_\text{red}})$, where
$\operatorname{lg}$ is the length of $\sK$ at the generic point of $Y$.
By (iii), one therefore has
\[
  \cl(Y,\sK) = \operatorname{lg} \prod_i \cl(D_i,\sO_{D_i})
\]
and the conclusion follows from \ref{IV:2-3-6}.






\subsubsection{Remark}\label{IV:2-3-9}

For $S$ the spectrum of a field, with the class of a cycle defined as in
\ref{IV:2-2-10}, assertion (iii) says that if two cycles $Y'$ and $Y''$
intersect with the correct dimension, then
\[
  \cl(Y'\cap Y'') = \cl(Y')\smallsmile \cl(Y'')\text{,}
\]
provided that the multiplicities of the components of $Y'\cap Y''$ are
computed as alternating sums of $\tor$.






\subsubsection{Remark}\label{IV:2-3-10}

Let $X/\spec(k)$, with $k$ algebraically closed. If two cycles of codimension
$c$ in $X$ are algebraically equivalent, the existence of the relative
theory above implies that the images in $\h^{2 c}(X,\dZ/n(c))$ of their
classes are the fibers, at two points of a suitable connected $S$, of a
section over $S$ of the constant sheaf $\h^{2 c}(X,\dZ/n(c))$ (the sheaf
$\eR^{2 c} f_\ast \dZ/n(c)$ for
$f=\operatorname{pr}_2:X\times S\to S$): two algebraically equivalent cycles
have the same class in $\h^{2 c}(X,\dZ/n(c))$.





















\section{Application: the Lefschetz trace formula in the proper smooth case}\label{IV:3}






\subsection{}\label{IV:3-1}

Let $X$ and $Y$ be proper smooth algebraic varieties over an algebraically
closed field $k$, purely of dimensions $N$ and $M$. Via the K\"unneth
isomorphism $\h^\bullet(X\times Y)=\h^\bullet(X)\otimes \h^\bullet(Y)$, the
trace morphism $\h^\bullet(X\times Y)(M)\to \h^\bullet(X)$ is simply
$\h^\bullet(X)\otimes ($the trace morphism
$\h^\bullet(Y)(M)\to \dQ_\ell)$. Since the morphism is homogeneous of even
degree, there is no sign problem. We denote it by $\int_Y$.






\subsection{}\label{IV:3-2}

A class $\eta\in \h^{2 q}(X\times Y)(q)$ defines a morphism of degree
$2(q-M)$, $\eta_{XY}:\h^\bullet(Y)(M-q)\to \h^\bullet(X)$, by the formula
\[
  \beta \mapsto \int_Y \eta\cdot \operatorname{pr}_2^\ast \beta \text{.}
\]






\begin{proposition_}\label{IV:3-3}
Let $p$ and $q$ be two integers, with $p+q=N+M$, and let
$\varepsilon\in \h^{2 p}(X\times Y)(p)$,
$\eta\in \h^{2 q}(X\times Y)(q)$. With the trace of
$\eta_{X Y}\varepsilon_{Y X}:\h^\bullet(X)\to \h^\bullet(X)$ understood in
the sense of \ref{IV:1-3}, one has
\[
  \tr_{X\times Y}(\eta\cdot \varepsilon) = \tr\left(\eta_{X Y}\varepsilon_{Y X},\h^\bullet(X)\right) \text{.}
\]
\end{proposition_}

To exorcise the signs, it is helpful to retain from the grading of
$\h^\bullet$ only the underlying $\dZ/2$-grading. To avoid carrying Tate
twists along, also fix an isomorphism $\dQ_\ell(1)=\dQ_\ell$. Let $\alpha_Y$
be the homomorphism $\h^\bullet(Y)\to \h^\bullet(Y)^\vee$ (an isomorphism)
that makes the diagram
\[\xymatrix{
  \h^\bullet(Y)\otimes \h^\bullet(Y) \ar[r] \ar[d]^-{\alpha_Y\otimes 1}
    & \h^\bullet(Y) \ar[r]^-{\tr}
    & \dQ_\ell \ar@{=}[d] \\
  \h^\bullet(Y)^\vee \otimes \h^\bullet(Y) \ar[rr]^-{\text{(\ref{IV:eq:1-3-1})}}
    & & \dQ_\ell \text{.}
}\]
commute. The map $\eta_{X Y}$ is the image of $\eta$ under
\[\xymatrix{
  \h^\bullet(X\times Y) = \h^\bullet(X)\otimes \h^\bullet(Y) \ar[r]^-{\alpha_Y}
    & \h^\bullet(X) \otimes \h^\bullet(Y)^\vee \ar@{=}[r]^-{\text{(\ref{IV:eq:1-3-2})}}
    & \hom\left(\h^\bullet(Y),\h^\bullet(X)\right) \text{,}
}\]
and likewise for $\varepsilon$. Definition \ref{IV:eq:1-3-4} of the trace
then reduces \ref{IV:3-3} to the formula
$\tr_{X\times Y}=\tr_X\otimes \tr_Y$.






\subsection{Remark}\label{IV:3-4}

The proof of \ref{IV:3-3} does not use the fact that $\alpha_X$ and
$\alpha_Y$ are isomorphisms. It is valid without assuming $X$ and $Y$
smooth. Outside the smooth case, however, it is difficult to find
\emph{cohomology} classes to which one would want to apply \ref{IV:3-3}.






\subsection{Remark}\label{IV:3-5}

If $\varepsilon,\eta$ are the cohomology classes of algebraic cycles $A$ and
$B$ on $X\times Y$, of codimensions $p$ and $q$, and $A\cap B$ has dimension
$0$, then $\tr_{X\times Y}(\eta\varepsilon)$ is the intersection
multiplicity $A\cdot B$, computed by an alternating sum of $\tor$
(\ref{IV:2-3-9}).






\begin{proposition_}\label{IV:3-6}
For $q=M$ and $\eta=\cl(B)$ the class of a subscheme
$B\subset X\times Y$, finite over $X$, the map $\eta_{X Y}$ is the composite
\[\xymatrix{
  \eta_{X Y} : \h^\bullet(Y) \ar[r]^-{\operatorname{pr}_2^\ast}
    & \h^\bullet(B) \ar[r]^-{\tr_{B/X}}
    & \h^\bullet(X) \text{.}
}\]
\end{proposition_}

Indeed, by \ref{IV:3-2} and \eqref{IV:eq:2-3-3-1} for $\sK=\sO_B$, one has
\[
  \eta_{X Y}(\beta) = \tr_{X\times Y/X}\left(\cl(B)\cdot \operatorname{pr}_2^\ast\beta\right) = \tr_{B/X}\left(\operatorname{pr}_2^\ast \beta\right) \text{.}
\]

The most important case is that in which $B$ is the graph of a map
$f:X\to Y$. In this case, \ref{IV:3-6} shows that $\eta_{X Y}=f^\ast$.






\begin{corollary_}\label{IV:3-7}
Let $f$ be an endomorphism of $X$ proper and smooth over an algebraically
closed field $k$. Suppose that the fixed points of $f$ are isolated. Then the
trace $\sum (-1)^i \tr(f^\ast,\h^i(X))$ is the number of fixed points of $f$,
each counted with its multiplicity.
\end{corollary_}

In \ref{IV:3-5}, take $A=$ the graph of the identity and $B=$ the graph of
$f$, and apply \ref{IV:3-3}, \ref{IV:3-6}.






\begin{corollary_}\label{IV:3-8}
Let $X$ be a curve over an algebraically closed $k$, obtained by extension of
scalars from $X_0$ over $\dF_q$, and let $f$ be the Frobenius morphism. Then
$\sum (-1)^i \tr\left(f^\ast,\h_c^i(X)\right)$ is the number of fixed
points of $f$.
\end{corollary_}

One can replace $X$ by $X_\text{red}$, and hence suppose that $X$ is reduced.
Let $U$ be the open subset of $X$ on which $X$ is smooth of dimension $1$,
and let $\bar U$ be the smooth projective completion of $U$. The long exact
cohomology sequences for $(X,U)$ and $(\bar U,U)$ give
\[
  \tr\left(f^\ast,\h_c^\bullet(X)\right) = \tr\left(f^\ast,\h_c^\bullet(U)\right) + \tr\left(f^\ast,\h_c^\bullet(X\setminus U)\right)
\]
and
\[
  \tr\left(f^\ast,\h_c^\bullet(\bar U)\right) = \tr\left(f^\ast,\h_c^\bullet(U)\right) + \tr\left(f^\ast,\h_c^\bullet(\bar U\setminus U)\right) \text{.}
\]
The same formulas hold for the numbers of fixed points. Since formula
\ref{IV:3-8} is clear for a scheme of dimension $0$, this reduces us to
proving \ref{IV:3-8} for $\bar U$. This case follows from \ref{IV:3-7}. The
fact that $d f=0$ guarantees that all fixed points have multiplicity one.






\subsection{Remark}\label{IV:3-9}

For $X$ of dimension $\leqslant 1$ over an algebraically closed $k$, and
$f:X\to X$ such that $d f=0$, it is not difficult to verify that
\ref{IV:3-8} is still valid.

% END INCLUDED components/sga4half_4_en.tex

% BEGIN INCLUDED components/sga4half_5_en.tex
% !TEX root = ../SGA4half_English_Reader_1_to_8.tex

\chapter{Duality}\label{V}

This exposé contains several theorems and compatibility results, all related
to Poincaré duality. In Section \ref{V:1}, we give the local biduality theorem
in dimension $1$ \cite[I.5.1]{sga5}, together with some computations of
duals. In Section \ref{V:2}, we give a very economical proof of Poincaré
duality for curves, taught to me by M.\ Artin. In Section \ref{V:3}, we prove
a compatibility linking two definitions of the pairing giving rise to
Poincaré duality for curves: one by cup product, the other by autoduality of
the Jacobian. Finally, in Section \ref{V:4}, we prove the compatibility
announced in the title.

Throughout the exposé, all schemes are noetherian and separated, and $n$ is
an integer invertible on every scheme under consideration.


\section{Local biduality in dimension \texorpdfstring{$1$}{1}}\label{V:1}


\subsection{}\label{V:1-1}

Let $S$ be a regular scheme, purely of dimension $1$. We shall show that the
complex with value $\dZ/n$ concentrated in degree $0$ is dualizing; that is,
for $\sK\in \ob\eD_c^b(S,\dZ/n)$ (where $(-)_c$ means constructible), if we
put $D\sK=\rHom(\sK,\dZ/n)$, then $D\sK$ is again
constructible with bounded cohomology, and the canonical morphism $\alpha$
from $\sK$ to $D D\sK$ is an isomorphism.

For $\sK$ in $\eD^-$ and arbitrary $\sL$, we have
\[
  \hom(\sK\lotimes \sL,\dZ/n) = \hom(\sL,\Hom(\sK,\dZ/n)) \text{.}
\]
The morphism from $\sK$ to $DD\sK$ is defined when $D\sK$ lies in $\eD^-$:
if $\beta:D\sK\lotimes\sK\to \dZ/n$ is the canonical pairing, it is the
morphism defined by the pairing
$\sK\lotimes D\sK=D\sK\lotimes \sK \xrightarrow\beta \dZ/n$.


\subsection{}\label{V:1-2}

Let $f:X\to S$ be a separated morphism of finite type. Put
$\sK_X=\eR f^!\dZ/n$. For $\sK\in \ob\eD^-(X,\dZ/n)$, put
$D\sK = \rHom(\sK,\sK_X)$. The adjunction between $\eR f_!$ and $\eR f^!$
gives
\[\xymatrix{
  \eR f_\ast \sK\lotimes \eR f_\ast D\sK \ar[r]
    & \eR f_! (\sK\lotimes D\sK) \ar[r]
    & \eR f_! \eR f^! \dZ/n \ar[r]
    & \dZ/n \text{.}
}\]
The symmetry of this description shows that, when $f$ is proper, the diagram
\begin{equation*}\tag{1.2.1}\label{V:eq:1-2-1}
\xymatrix{
  \eR f_\ast \sK \ar[r] \ar[d]
    & D D \eR f_\ast \sK \ar[d]^-\wr \\
  \eR f_\ast D D \sK \ar[r]^\sim
    & D\eR f_\ast D\sK
}
\end{equation*}
is commutative.

If $f$ is the inclusion of a closed point, we know that
$\eR f^!\dZ/n=\dZ/n(-1)[-2]$---that is, up to a twist and a shift, $\dZ/n$.
Over the spectrum of a field this complex is dualizing (Pontryagin duality
for $\dZ/n$-modules). By \eqref{V:eq:1-2-1}, we therefore have
$\sK\iso D D\sK$ for $\sK$ of the form $\eR f_!\sL$, and hence whenever the
support of $\underline\h^\bullet(\sK)$ is finite.


\begin{theorem_}\label{V:1-3}
Let $j:U\hookrightarrow S$ be a dense open subscheme of $S$, and let $\sF$
be a locally constant constructible sheaf of $\dZ/n$-modules on $U$. Then
$\eD j_\ast \sF = j_\ast D \sF$; that is,
$\Hom(j_\ast \sF,\dZ/n) = j_\ast\Hom(\sF,\dZ/n)$ and
$\Ext^i(j_\ast\sF,\dZ/n) = 0$ for $i>0$.
\end{theorem_}

On $U$, $\Ext^i(j_\ast\sF,\dZ/n)=0$ for $i>0$, since $\sF$ is locally
constant (and $\dZ/n$ is an injective $\dZ/n$-module), while for $i=0$ it is
the dual $\sF^\vee$ of $\sF$. One checks that
$\hom(j_\ast\sF,\dZ/n)=j_\ast\sF^\vee$, and it remains to verify the
vanishing of the $\Ext^i$ for $i>0$ at the points of $S\setminus U$. The
problem is local at these points. We are thus reduced to assuming that $S$
is a strictly local trait and that $U$ consists of its generic point $\eta$.
Let $I=\gal(\bar\eta/\eta)$. The sheaf $\sF$ is identified with the Galois
module $\sF_{\bar\eta}$, and the special fiber of $j_\ast \sF$ with
$\sF_{\bar\eta}^I$.

Let $i$ be the inclusion of the closed point $s$, and apply $D$ to the exact
sequences
\[\xymatrix{
  0 \ar[r]
    & j_! \sF \ar[r]
    & j_\ast \sF \ar[r]
    & i_\ast \sF_{\bar\eta}^I \ar[r]
    & 0 \\
  0 \ar[r]
    & j_! \sF_{\bar\eta}^I \ar[r] \ar[u]
    & \sF_{\bar\eta}^I \ar[r] \ar[u]
    & i_\ast \sF_{\bar\eta}^I \ar[r] \ar[u]
    & 0 \text{.}
}\]

We obtain a morphism of triangles
\[\xymatrix{
  i_\ast\left(\sF_{\bar\eta}^I\right)^\vee (-1)[-2] \ar[r] \ar@{=}[d]
    & D j_\ast\sF \ar[r] \ar[d]
    & \eR j_\ast \sF^\vee \ar[d] \\
  i_\ast\left(\sF_{\bar\eta}^I\right)^\vee(-1)[-2] \ar[r]
    & \left(\sF_{\bar\eta}^I\right)^\vee \ar[r]
    & \eR j_\ast \sF_{\bar\eta}^I \text{.}
}\]
The long exact sequence arising from the first row gives the vanishing of the
$\Ext^i$ for $i>2$, and, on taking the fiber at $s$, we find
\[\xymatrix@=0.7cm{
  0 \ar[r]
    & \left(\underline\h^1(D j_\ast\sF)\right)_s \ar[r]
    & \h^1\left(I,\sF_{\bar\eta}^\vee\right) \ar[r]^-\partial \ar[d]
    & \left(\sF_{\bar\eta}^I\right)^\vee(-1) \ar[r] \ar@{=}[d]
    & \left(\underline\h^2 D j_\ast\sF\right)_s \ar[r]
    & 0 \\
  & 0 \ar[r]
    & \h^1\left(I,(\sF_{\bar\eta}^I)^\vee\right) \ar[r]^-\partial
    & \left(\sF_{\bar\eta}^I\right)^\vee(-1) \ar[r]
    & 0 \text{.}
}\]
Since
$\left(\sF_{\bar\eta}^I\right)^\vee = \left(\sF_{\bar\eta}^\vee\right)_I$,
putting $M=\sF_{\bar\eta}^\vee$, it remains only to check that
\[\xymatrix{
  \h^1(I,M) \ar[r]^-\sim
    & \h^1(I,M_I) \text{.}
}\]
If $p$ is the residual characteristic exponent, then $I$ is an extension of
a group isomorphic to $\widehat\dZ_{p'} = \varprojlim_{(m,p)=1} \dZ/m$ by a
$p$-group $P$. Since the order of $P$ is prime to the order of $M$, we have
$\h^i(P,M)=0$ for $i>0$ and $(M_I)^P=(M^P)_I$. This allows us to replace $M$
by $M^P$ and $I$ by $I/P$. Finally, there is a functorial isomorphism
$\h^1(\widehat\dZ_{p'},M)\sim ($coinvariants of $\widehat\dZ_{p'}$ in $M)$,
which proves the theorem.


\begin{theorem_}\label{V:1-4}
For $\sK\in \ob\eD_c^b(S,\dZ/n)$, one has $\sK\iso D D \sK$.
\end{theorem_}

By dévissage, we reduce to the case in which $\sK$ is a constructible sheaf
$\sF$ concentrated in degree $0$, and $\sF$ either has finite support
(\ref{V:1-2}) or is of the form $j_\ast\sF_1$ as in \ref{V:1-3}. In the
second case, \ref{V:1-3} reduces us to local biduality for $\sF$ locally
constant on $U$.


\section{Poincaré duality for curves, after M.\ Artin}\label{V:2}


\subsection{}\label{V:2-1}

Let $X$ be a smooth projective curve over an algebraically closed field $k$.
Put $\sK_X=\dZ/n(1)[2]$ and, for $\sK\in\ob \eD_c^b(X,\dZ/n)$,
$D\sK=\rHom(\sK,\sK_X)$. For $M$ a $\dZ/n$-module, also put
$D M = \hom(M,\dZ/n)$, and similarly for complexes of modules. The trace
morphism $\h^0(X,\sK_X) \to \dZ/n$, or
$\eR\Gamma(X,\sK_X)\to \dZ/n$, defines a pairing
$\eR\Gamma(X,\sK)\lotimes\eR\Gamma(X,D\sK) \to \dZ/n$. Poincaré duality
between the cohomology and compactly supported cohomology of an open
subscheme $j:U\hookrightarrow X$ of the curve $X$ says that, for
$\sK=j_!\dZ/n$, this pairing identifies each factor with the dual of the
other.

I give below a proof, communicated to me by M.\ Artin, that for
$\sK\in\ob\eD_c^b(X,\dZ/n)$ this pairing is always perfect; that is, it
defines an isomorphism
\begin{equation*}\tag{2.1.1}\label{V:eq:2-1-1}
\xymatrix{
  \eR\Gamma(X,\sK) \ar[r]^-\sim
    & D \eR\Gamma(X,D\sK) \text{.}
}
\end{equation*}

For every constructible sheaf $\sF$, put
\[
  '\h^i(X,\sF) = \text{the $\dZ/n$-dual of $\h^{-i}(X,D\sF)$.}
\]
The assertion that \eqref{V:eq:2-1-1} is an isomorphism is equivalent to the
following theorem.


\begin{theorem_}\label{V:2-2}
For a constructible sheaf $\sF$ of $\dZ/n$-modules, one has
\begin{equation*}\tag{2.2.1}\label{V:eq:2-2-1}
\xymatrix{
  \h^i(X,\sF) \ar[r]^-\sim
    & '\h^i(X,\sF) \text{.}
}
\end{equation*}
\end{theorem_}


\begin{lemma_}\label{V:2-3}
Let $f:X\to Y$ be a generically étale morphism between smooth curves over
$k$. Then $\sK_X=f^\ast\sK_Y=\eR f^!\sK_Y$, with $\tr_f$ as the adjunction
morphism $\eR f_!\eR f^!\sK_Y\to \sK_Y$.
\end{lemma_}

% NOTE: \rhom or \rHom ? originally \rhom
We reduce to checking that, for $f$ finite,
$f_\ast \eR f^!\sK_Y = \rHom(f_\ast\dZ/n,\sK_Y)$ is $f_\ast\dZ/n$, with
$\tr_f$ as the adjunction morphism. The vanishing of the $\Ext^i$ for $i>0$
follows from \ref{V:1-3}, and the assertion follows.

More explicitly, $\eR f^!$ is the derived functor of the functor
$f^!:f^!\sF(U)=\hom(f_!\dZ_U,\sF)$, and \ref{V:1-3} implies the vanishing of
$\eR^i f^!\dZ/n$ for $i>0$.


\begin{corollary_}\label{V:2-4}
Let $f:X'\to X$ be a generically étale morphism of smooth projective curves
over $k$. Then $'\h^i(X,f_\ast \sF)='\h^i(X',\sF)$; this isomorphism is
functorial, and the diagram
\[\xymatrix{
  \h^i(X,f_\ast\sF) \ar[r] \ar@{=}[d]
    & '\h^i(X,f_\ast\sF) \ar@{=}[d] \\
  \h^i(X',\sF) \ar[r]
    & '\h^i(X',\sF)
}\]
is commutative.
\end{corollary_}

The isomorphism is given by \ref{V:1-2}: $D f_\ast\sF=f_\ast D\sF$.


\subsection{}\label{V:2-5}

We prove \ref{V:2-2}. If $\sF$ is $\dZ/n$ at one point, extended by $0$,
then $\h^\bullet(D\sF) = \ext^\bullet(\sF,\sK_X)=\dZ/n$ in degree $0$, and
in the pairing $\h^\bullet(\sF)\otimes\h^\bullet(D\sF) \to \dZ/n$, one has
$1\otimes 1\mapsto 1$ (\hyperref[IV]{Cycle}, \ref{IV:2-1-5}): the duality is
perfect. The case in which $\sF$ has finite support is treated in the same
way.

For arbitrary constructible $\sF$, let $\sA$ be the subsheaf of its sections
with finite support, and let $\sG=\sF/\sA$. The exact sequence
\begin{equation*}\label{V:eq:2-5-1}\tag{2.5.1}
\xymatrix{
  0 \ar[r]
    & \sA \ar[r]
    & \sF \ar[r]
    & \sG \ar[r]
    & 0
}
\end{equation*}
gives a diagram
\begin{equation*}\tag{2.5.2}\label{V:eq:2-5-2}
\xymatrix@=.5cm{
  & & 0 \ar[r]
    & \h^0(X,\sA) \ar[r] \ar[d]^-\wr
    & \h^0(X,\sF) \ar[r] \ar[d]
    & \h^0(X,\sG) \ar[r] \ar[d]
    & 0 \\
  0 \ar[r]
  & '\h^{-1}(X,\sF) \ar[r]
  & '\h^{-1}(X,\sG) \ar[r]
  & '\h^0(X,\sA) \ar[r]
  & '\h^0(X,\sF) \ar[r]
  & '\h^0(X,\sG) \ar[r]
  & 0
}
\end{equation*}
and compatible isomorphisms $\h^i(X,\sF)\iso \h^i(X,\sG)$,
$'\h^i(X,\sF) \iso '\h^i(X,\sG)$ for $i\ne 0,-1$. If $j:U\to X$ is an
open subscheme on which $\sF$ is locally constant, there is an exact sequence
\begin{equation*}\tag{2.5.3}\label{V:eq:2-5-3}
\xymatrix{
  0 \ar[r]
    & \sG \ar[r]
    & j_\ast j^\ast \sF \ar[r]
    & \sB \ar[r]
    & 0
}
\end{equation*}
with $\sB$ of finite support. Since
$'\h^i(X,j_\ast j^\ast\sF) = D \h^{2-i} (X,j_\ast(j^\ast\sF)^\vee(1))$
by \ref{V:1-3}, one has $'\h^i=0$ for $i<0$; then
\eqref{V:eq:2-5-2} and the long exact sequence defined by
\eqref{V:eq:2-5-3} show that $'\h^i$ always vanishes for $i<0$.

If $\sF=\dZ/n$, the theorem is true for $i=0$ and $2$. This expresses the
fact that, for connected $X$, the trace morphism
$\h^2(X,\dmu_n) \iso \dZ/n$ is an isomorphism. By \ref{V:2-4}, the theorem
remains true for $i=0$ and $\sF=f_\ast\dZ/n$.

For arbitrary $\sF$ again, with $\sG$ as above, there exists
$f:X'\to X$ as in \ref{V:2-4} such that $\sG$ embeds into
$f_\ast\dZ/n$:
\[\xymatrix{
  0 \ar[r]
    & \sG \ar[r]
    & f_\ast\dZ/n \ar[r]
    & \sQ \ar[r]
    & 0 \text{.}
}\]
After replacing $X'$ by an $X''/X'$, we may assume that
$\h^i(X,\sF) \to \h^i(X,f_\ast\dZ/n) = \h^i(X',\dZ/n)$ is zero for $i>0$.
Consider then
\small
\[\xymatrix@=.3cm{
  0 \ar[r]
    & \h^0(X,\sG) \ar[r] \ar[d]
    & \h^0(X',\dZ/n) \ar[r] \ar[d]^-\wr
    & \h^0(X,\sQ) \ar[r] \ar[d]
    & \h^1(X,\sG) \ar[r]^-0 \ar[d]
    & \h^1(X',\dZ/n) \ar[r] \ar[d]
    & \h^1(X,\sQ) \ar[r] \ar[d]
    & \cdots \\
  0 \ar[r]
    & '\h^0(X,\sG) \ar[r]
    & '\h^0(X',\dZ/n) \ar[r]
    & '\h^0(X,\sQ) \ar[r]
    & '\h^1(X,\sG) \ar[r]
    & '\h^1(X',\dZ/n) \ar[r]
    & '\h^1(X,\sQ) \ar[r]
    & \cdots
}\]
\normalsize

The isomorphism at $\h^0(X')$ gives
$\h^0(X,\sG)\hookrightarrow '\h^0(X,\sG)$, and the same injectivity for
$\sF$ by \eqref{V:eq:2-5-2}. Applied to $\sQ$, this gives
$\h^1(X,\sG)\hookrightarrow '\h^1(X,\sG)$. We apply this to
$\h^1(X',\dZ/n)$. This group has the same order as
$'\h^1(X',\dZ/n)=D(\h^1(X',\dZ/n)(1))$, so we obtain an isomorphism, and
$\h^1(X,\sG) \iso '\h^1(X,\sG)$. The same holds for $\sF$. For $i=2$, we
already know that $\h^2(X',\dZ/n) \iso '\h^2(X',\dZ/n)$, and proceed in the
same way, then stop for lack of combatants.


\section{Poincaré duality for curves}\label{V:3}

In this section we prove the compatibility
(\hyperref[I]{Arcata} \ref{I:eq:6-2-3-3}).


\subsection{}\label{V:3-1}

We use the notation of (\hyperref[I]{Arcata} \ref{I:6-2-3}): $\bar X$ is a
smooth, connected, projective curve over an algebraically closed field $k$,
$D$ is a reduced divisor on $\bar X$, and $0$ is a point of
$X=\bar X\setminus D$,
\[\xymatrix{
  X \ar@{^{(}->}[r]^-j
    & \bar X
    & \ar[l]_-i D \text{,}
}\]
and $\pic_D(\bar X)=\h^1(\bar X,_D\dG_m)$, where $_D\dG_m$ is the sheaf of
sections of $\dG_m$ congruent to $1\mod D$. Recall that
$\h_c^1(X,\dmu_n) = \pic_D^0(\bar X)_n$, and that $x\to \sO(x)$ defines a
canonical map $f:X\to \pic_D(\bar X)$. Put
$f_0(x)=f(x)-f(0):X\to \pic_D^0(\bar X)$.

Again write $j$ for the inclusion of $X\times X$ in $X\times \bar X$. The
diagonal $\Delta$ of $X\times X$ is closed in $X\times \bar X$; it defines a
class in
$\h_\Delta^2(X\times X,\dmu_n)=\h_\Delta^2(X\times \bar X,j_!\dmu_n)$. Let
$c$ denote its image in $\h^2(X\times \bar X,j_!\dmu_n)$.

The K\"unneth formula gives
\[
  \h^\bullet(X\times \bar X,j_!\dmu_n) = \h^\bullet(X,\dZ/n)\otimes \h^\bullet(\bar X,j_!\dmu_n) = \h^\bullet(X,\dZ/n)\otimes \h_c^\bullet(X,\dmu_n) \text{.}
\]
We shall compute the $(1,1)$-component of $c$,
\[
  c^{1,1} \in \h^1(X,\dZ/n)\otimes \h_c^1(X,\dmu_n) = \h^1(X,\h_c^1(X,\dmu_n)) = \h^1\left(X,\pic_D^0(\bar X)_n\right) \text{.}
\]

Consider the exact sequence
\[\xymatrix{
  0 \ar[r]
    & \pic_D^0(\bar X)_n \ar[r]
    & \pic_D^0(\bar X) \ar[r]^n
    & \pic_D^0(\bar X) \ar[r]
    & 0
}\]
This makes $\pic_D^0(\bar X)$ a torsor under $\pic_D^0(\bar X)_n$ over
$\pic_D^0(\bar X)$. Let $u$ be the class, in
$\h^1(X,\h_c^1(X,\dmu_n))=\h^1(X,\pic_D^0(\bar X)_n)$, of its inverse image
under $f_0$.


\begin{proposition_}\label{V:3-2}
With the preceding notation, $c^{1,1}=-u$.
\end{proposition_}

Let $e$ be the difference between the images in
$\h^2(X\times \bar X,j_!\dmu_n)$ of the cohomology classes of $\Delta$ and
$X\times\{0\}$. Since the latter is of K\"unneth type $(0,2)$, one has
$c^{1,1}=e^{1,1}$. Consider the Leray spectral sequence for
$\pr_1:X\times \bar X\to X$ and the sheaf $j_!\dmu_n$:
\[
  E_2^{pq} = \h^p(X,\eR^q {\pr_1}_\ast j_!\dmu_n) = \h^p(X,\h_c^q(X,\dmu_n)) \Rightarrow \h^{p+q}(X\times \bar X,j_!\dmu_n) \text{.}
\]
It degenerates: it is the tensor product of $\h^\bullet(X)$ with the
(trivial) Leray spectral sequence for $\bar X\to \spec(k)$ and the sheaf
$j_!\dmu_n$. Fiber by fiber, the divisor $\Delta-X\times\{0\}$ has degree
$0$, so the image of $e$ in $E^{02}$ is zero. We may therefore speak of its
image $\bar e$ in $E^{1,1}$, and we must show that $\bar e = -u$.

Let $\sG$ on $X\times\bar X$ be the subsheaf of $\dG_m$ consisting of local
sections whose restriction to the closed subscheme $X\times D$ is $1$.
There is again an exact sequence
\begin{equation*}\tag{3.2.1}\label{V:eq:3-2-1}
\xymatrix{
  0 \ar[r]
    & j_!\dmu_n \ar[r]
    & \sG \ar[r]
    & \sG \ar[r]
    & 0
}
\end{equation*}
and $e$ is the image under $\partial$ of the class
$e_1\in \h^1(X\times\bar X,\sG)$ of the invertible sheaf
$\sO(\Delta-X\times\{0\})$, trivialized by $1$ on $X\times D$.

Applying $\eR\pr_{1\ast}$ to the distinguished triangle defined by
\eqref{V:eq:3-2-1} gives a distinguished triangle
\begin{equation*}\tag{3.2.2}\label{V:eq:3-2-2}
\xymatrix{
  & \ar[r]^-\partial
    & \eR \pr_{1\ast} j_!\dmu_n \ar[r]
    & \eR \pr_{1\ast}\sG \ar[r]
    & \eR \pr_{1\ast}\sG \ar[r]^-\partial
    &
}
\end{equation*}
and the diagram
\[\xymatrix{
  \h^1(X\times \bar X,\sG) \ar[r]^-\partial \ar@{=}[d]^-{(1)}
    & \h^2(X\times \bar X,j_!\dmu_n) \ar@{=}[d] \\
  \h^1(X,\eR\pr_{1\ast} \sG) \ar[r]^-\partial
    & \h^2(X,\eR\pr_{1\ast} j_!\dmu_n)
}\]
is commutative. Our problem is thus to identify the image in
$E^{11}=\h^1(X,\h_c^1(X,\dmu_n))$ (the image in $E^{02}$ being zero) of the
image under
$\partial_{\text{\eqref{V:eq:3-2-2}}}:\h^1(X,\eR\pr_{1\ast}\sG) \to \h^2(X,\eR\pr_{1\ast} j_!\dmu_n)$
of the class, still denoted by $e_1$, corresponding to $e_1$ under (1).

The sheaf $\eR^1\pr_{1\ast}\sG$ is defined by the group scheme over $X$
obtained by pulling back the algebraic group $\pic_D(\bar X)$ over
$\spec(k)$, and the image of $e_1$ in
$\h^0(X,\eR^1\pr_{1\ast}\sG) = \hom(X,\pic_D(\bar X))$ is $f_0$.

Represent the triangle \eqref{V:eq:3-2-2} by a short exact sequence of
complexes of sheaves:
\[\xymatrix{
  0 \ar[r]
    & \sA \ar[r]
    & \sB \ar[r]
    & \sC \ar[r]
    & 0 \text{.}
}\]
Let $\sB'$ be the subsheaf of
\[
  \tau_{\leqslant 1} \sB : \cdots \sB^0 \to \ker(d) \to 0 \cdots
\]
obtained by replacing the degree-one term $\ker(d)$ by the inverse image, under
$\ker(d)\to \underline\h^1(\sB) = \pic_D(\bar X)$ over $X$, of
$\pic_D^0(\bar X)$ over $X$. Let $\sA'=\sA\cap \sB'$ and let $\sC'$ be the
image of $\sB'$. We have $\sA'=\tau_{\leqslant 1}(\sA)$, and $\sC'$ is
obtained from $\sC$ in the same way as $\sB'$ from $\sB$. Hence there is a
commutative diagram of short exact sequences
\begin{equation*}\tag{3.2.3}\label{V:eq:3-2-3}
\xymatrix{
  0 \ar[r]
    & \sA \ar[r]
    & \sB \ar[r]
    & \sC \ar[r]
    & 0 \\
  0 \ar[r]
    & \tau_{\leqslant 1} \sA \ar[r] \ar[u] \ar[d]
    & \sB' \ar[r] \ar[u] \ar[d]
    & \sC' \ar[r] \ar[u] \ar[d]
    & 0 \\
  0 \ar[r]
    & \h_c^1(X,\dmu_n)[-1] \ar[r]
    & \pic_D^0(\bar X)[-1] \ar[r]
    & \pic_D^0(\bar X)[-1] \ar[r]
    & 0
}
\end{equation*}
where the last row is to be viewed as an exact sequence of complexes of
sheaves concentrated in degree $1$ on $X$.
\begin{enumerate}[a)]
  \item $e\in \hh^2(X,\sA)$ comes from
    $\widetilde e\in \hh^2(X,\tau_{\leqslant 1}\sA)$ (because its image in
    $E^{02}$ is zero). We seek the image $\bar e$ of $\widetilde e$ in
    $\hh^2(X,\h_c^1(X,\dmu_n)[-1]) = \h^1(X,\h_c^1(X,\dmu_n))$.
  \item $e_1\in \hh^1(X,\sC)$ comes from
    $\widetilde e_1\in \hh^1(X,\sC')$, because its image $f_0$ in
    $\h^0(X,\pic_D(\bar X))$ lies in $\h^0(X,\pic_D^0(\bar X))$. We may
    take $\widetilde e=\partial \widetilde e_1$.
  \item Thus $\bar e = \partial f_0$, where $\partial$ is defined by
    \eqref{V:eq:3-2-2}, and where we use the usual isomorphisms
    $\hh^2(X,\sF[-1]) = \h^1(X,\sF)$. These isomorphisms anticommute with
    $\partial$, so that $\bar e=-\partial f_0$, for $\partial$ defined by
    the exact sequence
    \[\xymatrix{
      0 \ar[r]
        & \h_c^1(X,\dmu_n) \ar[r]
        & \pic_D^0(\bar X) \ar[r]
        & \pic_D^0(\bar X) \ar[r]
        & 0
    }\]
    The conclusion follows from (\hyperref[IV]{Cycles}, \ref{IV:1-1-1}).
\end{enumerate}


\subsection{}\label{V:3-3}

For every homomorphism $\varphi:\h_c^1(X,\dmu_n)\to \dZ/n$, let
$\varphi(u)\in \h^1(X,\dZ/n)$ be the image of $u$ under
$\varphi:\h^1(X,\h_c^1(X,\dmu_n)) \to \h^1(X,\dZ/n)$. If
$x\in \h_c^1(X,\dmu_n)$, the cup product $\varphi(u)\smallsmile x$ lies in
$\h_c^2(X,\dmu_n)$. We shall prove the following.


\begin{proposition_}\label{V:3-4}
$\tr(\varphi(u)\smallsmile x) = \varphi(x)$.
\end{proposition_}

This identity follows, after applying $\varphi$, from
\begin{equation*}\tag{3.4.1}\label{V:eq:3-4-1}
  \tr(u\smallsmile x) = x
\end{equation*}
where $\tr$ now denotes the map
\[\xymatrix{
  \tr:\h_c^2(X,\h_c^1(X,\dmu_n)\otimes\dmu_n) \ar[r]
    & \h_c^1(X,\dmu_n) \text{,}
}\]
so that $x\mapsto \tr(u\smallsmile x)$ is the composite
\small
\begin{equation*}\tag{3.4.2}\label{V:eq:3-4-2}
\begin{aligned}
&\xymatrix{
  \h_c^1(X,\dmu_n) \ar[r]^-{u\otimes x}
    & \h^1(X,\h_c^1(X,\dmu_n)) \otimes \h_c^1(X,\dmu_n)
    & \ar[l]_-\sim \h^1(X,\dZ/n)\otimes \h_c^1(X,\dmu_n) \otimes \h_c^1(X,\dmu_n)} \\ &\xymatrix{
    \ar[r]^-{(1)}
      & \h_c^2(X,\dmu_n) \otimes \h_c^1(X,\dmu_n) \ar[r]^-{\tr\otimes\operatorname{id}}
      & \h_c^1(X,\dmu_n)
}
\end{aligned}
\end{equation*}
\normalsize
where (1) is given by
$a\otimes b\otimes c\mapsto (a\smallsmile c)\otimes b$.

We express this morphism in terms of the trace morphism
$\tr:\h_c^3(X\times X,\dmu_n^{\otimes 2}) \to \h_c^2(X,\dmu_n)$ relative to
the second projection. For $x\in \h_c^3$, $\tr(x)$ is computed as follows:
among its K\"unneth components, retain only the component of type $(2,1)$,
and apply $\tr:\h_c^2(X,\dmu_n) \to \dZ/n$.

Let $u'$ be the image of $u$ under the map
\begin{align*}
&\xymatrix{
  \h^1(X,\h_c^1(X,\dmu_n))
    & \ar[l]_-\sim \h^1(X,\dZ/n) \otimes \h_c^1(X,\dmu_n) = \h^1(X,\dZ/n)\otimes \h^1(\bar X,j_!\dmu_n)
} \\
&\xymatrix{
  \ar[r]
    & \h^2(X\times \bar X,\dZ/n\boxtimes j_!\dmu_n) \text{.}
}
\end{align*}
In view of the relation between the K\"unneth formula and cup product, one
has
\[
  \tr(u\smallsmile x) = -\tr(u'\smallsmile \pr_1^\ast x) \text{.}
\]
The second cup product is the one in
(\hyperref[IV]{Cycle} \ref{IV:1-2-4}):
\[\xymatrix{
  \h^2(X\times \bar X,\dZ/n\boxtimes j_!\dmu_n) \otimes \h^1(\bar X\times \bar X,j_!\dmu_n\boxtimes \dZ/n) \ar[r]
    & \h^3(\bar X\times \bar X,j_!\dmu_n \boxtimes j_!\dmu_n) \text{.}
}\]
The minus sign comes from permuting two symbols of degree $1$ in (1). Since
$c\smallsmile\pr_1^\ast x$ and $c^{1,1}\smallsmile \pr_1^\ast x$ have the
same $(2,1)$-component in the K\"unneth decomposition, and since
$u'=-c^{1,1}$, one has
\[
  \tr(u\smallsmile x) = \tr(c\smallsmile \pr_1^\ast x) \text{.}
\]
The cup product on the right can now be interpreted as the cup product
(\hyperref[IV]{Cycle} \ref{IV:1-2-5}) of
$\cl(\Delta)\in \h_\Delta^2(X\times X,\dmu_n)$ with
$\pr_1^\ast x\in \h_c^1(\Delta,\dmu_n)$. On $\Delta$, one has
$\pr_1=\pr_2$, whence
\[
  \tr(c\smallsmile \pr_1^\ast x) = \tr(\cl(\Delta)\smallsmile \pr_1^\ast x) = \tr(\cl(\Delta)\smallsmile \pr_2^\ast x) = \tr(\cl(\Delta))\smallsmile x = x \text{.}
\]


\section{Compatibility with SGA 4 XVIII 3.1.10.3}\label{V:4}

In \cite[XVIII]{sga4}, overcome by discouragement, I did not verify the
compatibility announced in the title. Although it turns out to be no more
than smoke, I feel bound to remedy that omission here. For $f:X\to S$
separated and of finite type, the point was to verify that the adjunction
morphism $\eR f_!\eR f^!\sL\to \sL$ is compatible with every étale
localization $k:V\to S$. Using $k$ and $f$ for several arrows, as below, we
must verify commutativity:
\[\xymatrix{
  X_V \ar[r]^-k \ar[d]^-f
    & X \ar[d]^-f \\
  V \ar[r]^-k
    & S
} \qquad
\xymatrix{
  k^\ast \eR f_! \eR f^!\sL \ar[r]^-{k^\ast(\text{adj})} \ar[d]^-\sim
    & k^\ast \sL \ar@{=}[d] \\
  \eR f_! \eR f^! k^\ast \sL \ar[r]^-{\text{adj}}
    & k^\ast \sL \text{.}
}\]
At the level of complexes, the required commutativity is written (loc. cit.)
\[\xymatrix{
  k^\ast f_!^\bullet f^!_\bullet \sL \ar[r] \ar[d]^-{(1)}
    & k^\ast \sL \ar@{=}[d] \\
  f_!^\bullet f^!_\bullet k^\ast \sL \ar[r]
    & k^\ast \sL \text{.}
}\]
It is checked component by component, reducing the problem to a
compatibility for $\sL$ a sheaf and for each pair of adjoint functors
$(f_i^!,f_!^i)$, denoted simply by $f^!$ and $f_!$. The arrow (1) is defined
using the isomorphism $(2):k^\ast f_!=f_! k^\ast$, and a morphism
$(3):k^\ast f^! \to f^! k^\ast$ deduced from it. In loc. cit., one first
introduces a morphism $(4):k_! f_!\to f_! k_!$, deduced from $(2)$ by the
adjunction of $k^\ast=k^!$ and
$k_!:k_! f_!\to k_! f_! k^\ast k_!\to k_! k^\ast f_! k_! \to f_! k_!$, and
then deduces $(3)$ from $(4)$ by adjunction. Equivalently, one may deduce
$(3)$ from $(2)$ by the adjunction of $f_!$ and
$f^!:k^\ast \to f^! f_! k^\ast f^!\to f^! k^\ast f_! f^! \to f^! k^\ast$.

Writing $k$ for $k^\ast$ and omitting $\sL$, the required commutativity is
that of the outer boundary of
\[\xymatrix@=.5cm{
  k f_! f^! \ar[r] \ar@{=}[d]
    & f_! k f^! \ar[r]
    & f_! f^! f_! k f^! \ar[r]
    & f_! f^! k f_! f^! \ar@{=}[d] \\
  k f_! f^! \ar[d]
    & & & \ar[lll] f_! f^! k f_! f^! \ar[d] \\
  k
    & & & \ar[lll] f_! f^! k \text{.}
}\]
In the first row, we find a homomorphism
$k f_! f^! \to f_! f^! k f_! f^!$, defined because $k f_! f^!$ is of the
form $f_! X$ (with $X=k f^!$). It admits the adjunction morphism in the
second row as a retraction, and the commutativity of the lower square is the
functoriality of this adjunction morphism.

% END INCLUDED components/sga4half_5_en.tex

% BEGIN INCLUDED components/sga4half_6_en.tex
% !TEX root = ../SGA4half_English_Reader_1_to_8.tex

\chapter{Applications of the trace formula to trigonometric sums}\label{VI}
\chaptermark{Applications to trigonometric sums}










In this exposé, I explain how the trace formula can be used to calculate 
or study various trigonometric sums and how, together with the 
Weil conjecture, it can be used to bound them.

The first two sections provide a ``user's guide'' to these tools. 
\hyperref[VI:3]{Section 3} presents, in cohomological language, Weil's results on sums in 
$1$ variable. Sections $4$ through $6$ give a 
detailed study of Gauss and Jacobi sums -- including 
Weil's old and recent results on Hecke characters 
defined by Jacobi sums. In Section $7$, we study a 
generalization of Kloosterman sums to several variables. Finally, 
Section $8$ gives some indications of other uses that have been, or can be, 
made of these methods. 










\section*{Notations}\label{VI:0}





\subsection{}\label{VI:0-1}

We use the notation of \hyperref[II]{Report}, Section \ref{II:1}. We 
will often have to consider a finite extension $\dF_{q^n}\subset \dF$ of 
$\dF_q$. A subscript $0$ will denote an object over $\dF_q$, and a subscript $1$ 
an object over $\dF_{q^n}$. Replacing a subscript $0$ by a subscript $1$ (resp. 
omitting the subscript) means extending scalars to $\dF_{q^n}$ (resp. 
to $\dF$).





\subsection{}\label{VI:0-2}

We denote by $\ell$ a prime number $\ne p$. We freely use the 
language of $\dQ_\ell$-sheaves, as well as that of $E_\lambda$-sheaves, 
where $E_\lambda$ is a finite extension of $\dQ_\ell$ (cf. \hyperref[II]{Report}, 
Section \ref{II:2}, especially \ref{II:2-11}). 





\subsection{}\label{VI:0-3}

Let $H^\bullet$ be a graded vector space. If $T$ is an endomorphism of 
$H^\bullet$, we put (cf. \hyperref[IV]{Cycle}, \ref{IV:eq:1-3-5}) 
\[
  \tr(T,H^\bullet) = \sum (-1)^i \tr(T,H^i) \text{.}
\]










\section{Principles}\label{VI:1}





\subsection{}\label{VI:1-1}

Let $X_0$ be a separated scheme of finite type over $\dF_q$, let $E_\lambda$ be a 
finite extension of $\dQ_\ell$, and let $\sF_0$ be an $E_\lambda$-sheaf on $X_0$. The 
trace formula says that 
\begin{equation*}\tag{1.1.1}\label{VI:eq:1-1-1}
  \sum_{x\in X^F} \tr(F_x^\ast,\sF) = \sum (-1)^i \tr\left(F^\ast,\h_c^i(X,\sF)\right) \text{.}
\end{equation*}

With the notation \ref{VI:0-3}, the right-hand side is simply written 
$\tr(F^\ast,\h_c^\bullet(X,\sF))$. 

This trace formula for $E_\lambda$-sheaves can either be 
deduced from \hyperref[II]{Report} \ref{II:4-10} by passage to the limit (cf. 
\hyperref[II]{Report} \ref{II:4-11} through \ref{II:4-13}), or be deduced 
from the trace formula for $\dQ_\ell$-sheaves (\hyperref[II]{Report} 
\ref{II:3-2}) by the method of (\hyperref[III]{Functions $L$ mod. $\ell^n$}, 
\ref{III:4-3}). 

We shall interpret various trigonometric sums as the left-hand side of 
\eqref{VI:eq:1-1-1}, for suitable $X_0$ and $\sF_0$. 





% NOTE: in this section, following the source (but not convention in the 
% LaTeX version, torsors are denoted in roman (not script) letters
\subsection{}\label{VI:1-2}

Let $A$ be a finite group. An \emph{$A$-torsor} on a scheme $X$ is a 
sheaf $T$ on $X$, equipped with a right action of $A$, which, locally 
(for the étale topology) on $X$, is isomorphic to the constant sheaf $A$ on 
which $A$ acts by right translation. 

If $\tau:A\to B$ is a homomorphism and $T$ is an $A$-torsor, there exists, 
up to unique isomorphism, a unique $B$-torsor $\tau(T)$, equipped with a 
morphism of sheaves $\tau:T\to \tau(T)$ such that 
\begin{equation*}\tag{1.2.1}\label{VI:eq:1-2-1}
  \tau(t a) = \tau(t) \tau(a)\text{.}
\end{equation*}

If $T$ is an $A$-torsor on $X$ and $\rho$ is a linear representation 
of $A$, $\rho:A\to \gl(V)$ (where $V$ is a finite-dimensional vector space over 
$E_\lambda$), there exists, up to unique isomorphism, a unique 
$E_\lambda$-sheaf $\sF$, smooth of rank $\dim(V)$, equipped with a morphism of 
sheaves 
\begin{equation*}\tag{1.2.2}\label{VI:eq:1-2-2}
  \rho:T\to \underline{\operatorname{Isom}}(V,\sF) \text{.}
\end{equation*}
such that $\rho(t a) = \rho(t) \rho(a)$. We denote it by $\rho(T)$. For 
$\tau:A\to B$ and $\rho$ a representation of $B$, we have a canonical 
isomorphism $\rho(\tau(T)) = (\rho\tau)(T)$. 

In this paragraph (except the appendix), we only consider the case where $A$ 
is commutative and where $V=E_\lambda$: $\rho$ is a character 
$A\to E_\lambda^\times$, and $\rho(T)$ is a smooth $E_\lambda$-sheaf of rank 
one. The morphism \eqref{VI:eq:1-2-2} is interpreted as an everywhere 
nonzero morphism 
\begin{equation*}\tag{1.2.3}\label{VI:eq:1-2-3}
  \rho:T\to \rho(T) 
\end{equation*}
such that $\rho(t a) = \rho(a)\rho(t)$. 





\subsection{}\label{VI:1-3}

Let $S$ be a scheme, let $G$ be a commutative group scheme over $S$, and let $G'$ 
be an extension of commutative group schemes over $S$
\[\xymatrix{
  0 \ar[r] 
    & A \ar[r] 
    & G' \ar[r]^-\pi 
    & G \ar[r] 
    & 0 \text{.}
}\]
The sheaf $T$ of local sections of $\pi$ is then an $A$-torsor on 
$G$. 

If $X$ is a scheme over $S$ and $f,g$ are two $S$-morphisms from $X$ to $G$, 
the fact that $T$ is defined by an extension gives a canonical isomorphism 
\begin{equation*}\tag{1.3.1}\label{VI:eq:1-3-1}
  (f+g)^\ast T = f^\ast T + g^\ast T \text{.}
\end{equation*}
On the left, $f+g$ is the sum of $f$ and $g$ with respect to the group law on $G$; 
on the right, $+$ denotes a sum of torsors. 

If $f:X\to G$ factors through $G'$, $f^\ast T$ is trivial (and a 
factorization gives a trivialization); up to (non-unique) isomorphism, 
the $A$-torsor $f^\ast T$ depends only on the image of $f$ in 
$\hom_S(X,G)/\pi\hom_S(X,G')$: it is given by the morphism $\partial$ in 
the exact sequence 
\[\xymatrix{
  \hom(X,G') \ar[r] 
    & \hom(X,G) \ar[r]^-\partial 
    & \h^1(X,A) \text{.}
}\]

We will apply these constructions to the exact Kummer sequences 
$0 \to \dmu_n \to \dG_m \to \dG_m \to 0$, and to Lang torsors. 





\subsection{The Lang torsor}\label{VI:1-4}

Let $G_0$ be a connected commutative algebraic group over $\dF_q$ (for the 
noncommutative case, see the appendix). The group law is written 
\emph{multiplicatively}. Lang's isogeny 
\[
\fL:G_0 \to G_0: x\mapsto F x\cdot x^{-1}
\]
is étale; its image, an open subgroup of $G_0$, can only be 
$G_0$ itself; its kernel is the finite group $G_0(\dF_q)$. The \emph{Lang 
torsor} $L$ is the $G_0(\dF_q)$-torsor on $G_0$ defined by the exact sequence 
\begin{equation*}\tag{1.4.1}\label{VI:eq:1-4-1}
\xymatrix{
  0 \ar[r] 
    & G_0(\dF_q) \ar[r] 
    & G_0 \ar[r]^-\fL 
    & G_0 \ar[r] 
    & 0 \text{.}
}
\end{equation*}

\paragraph{Notation}
We also write $L$ and $\fL$ (or $L_{(q)}$ and $\fL_{(q)}$) for the objects obtained 
from $L$ and $\fL$ by extending the base field. When necessary, we 
specify them with a subscript $0$ or $1$.


\paragraph{Examples}
For $G_0=\dG_a$ or $\dG_m$, the sequences \eqref{VI:eq:1-4-1} are written 
\begin{align*}\tag{1.4.2}\label{VI:eq:1-4-2}
&\xymatrix{
  0 \ar[r] 
    & \dF_q \ar[r] 
    & \dG_a \ar[r]^-{x^q-x} 
    & \dG_a \ar[r] 
    & 0 
} \\
\tag{1.4.3}\label{VI:eq:1-4-3}
&\xymatrix{
  0 \ar[r] 
    & \dmu_{q-1} \ar[r] 
    & \dG_m \ar[r]^-{x^{q-1}} 
    & \dG_m \ar[r] 
    & 0
}
\end{align*}





\subsection{}\label{VI:1-5}

Let us compute the endomorphism $F^\ast$ of the fiber of the Lang torsor at 
$\gamma\in G^F=G_0(\dF_q)$. If $g\in \fL^{-1}(\gamma)\subset G$, we have 
$F g = F g\cdot g^{-1}\cdot g = \gamma g$. Therefore (\hyperref[II]{Report}, 
\ref{II:1-2}) 
\begin{equation*}\tag{1.5.1}\label{VI:eq:1-5-1}
  \text{on $L_0(G_0)_\gamma\simeq \fL^{-1}(\gamma)$, $F^\ast$ is $g\mapsto g\gamma^{-1}$.}
\end{equation*}





\subsection{}\label{VI:1-6}

On $\dF_{q^n}$, the identity $F_{(q^n)}=F_{(q)}^n$ between endomorphisms of 
$G_1$ implies that $\fL_{(q^n)}=\fL_{(q)}\circ \prod_{i=0}^{n-1} F_{(q)}^i$. 
On $G_0(\dF_{q^n})$, $F_{(q)}^i$ acts as the element 
$x\mapsto x^{q^i}$ of $\gal(\dF_{q^n}/\dF_q)$. On $G_0(\dF_{q^n})$, 
$\prod F_{(q)}^i$ is therefore the composite of the norm 
$G_0(\dF_{q^n})\to G_0(\dF_q)$ and the inclusion 
$G_0(\dF_q)\subset G_0(\dF_{q^n})$: the diagram 
\begin{equation*}\tag{1.6.1}\label{VI:eq:1-6-1}
\xymatrix{
  0 \ar[r] 
    & G_0(\dF_{q^n}) \ar[r] \ar[d]^-N 
    & G_1 \ar[r]^-{\fL_{(q^n)}} \ar[d]^-{\prod_{i=0}^{n-1} F_{(q)}^i} 
    & G_1 \ar[r] \ar@{=}[d] 
    & 0 \\
  0 \ar[r] 
    & G_0(\dF_q) \ar[r] 
    & G_1 \ar[r]^-{\fL_{(q)}} 
    & G_1 \ar[r] 
    & 0
}
\end{equation*}
is commutative. In the language of torsors, this gives a canonical isomorphism 
between $G_0(\dF_q)$-torsors on $G_1$ 
\begin{equation*}\tag{1.6.2}\label{VI:eq:1-6-2}
  N L_{(q^n)} = L_{(q)} \text{.}
\end{equation*}





\begin{definition_}\label{VI:1-7}
Let $f_0:X_0\to G_0$ be a morphism and 
$\chi:G_0(\dF_q)\to E_\lambda^\times$ a character. We set 
$\sF(\chi,f_0) = \chi^{-1}(f_0^\ast (L_0(G_0))) = f_0^\ast\chi^{-1}(L_0(G_0))$. 
\end{definition_}

We have the following functoriality properties:
\begin{enumerate}[a)]
  \item $\sF(\chi,f_0)$ is bimultiplicative in $\chi$ and $f_0$: 
    \begin{align*}\tag{1.7.1}\label{VI:eq:1-7-1}
      \sF(\chi,f_0'\cdot f_0'') &= \sF(\chi,f_0')\otimes \sF(\chi,f_0'') \text{,} \\
      \tag{1.7.2}\label{VI:eq:1-7-2}
      \sF(\chi'\chi'',f_0) &= \sF(\chi',f_0)\otimes \sF(\chi'',f_0) \text{.}
    \end{align*}
    This follows from the definition of $\chi(T)$, for $T$ a torsor, combined 
    with \eqref{VI:eq:1-3-1} for \eqref{VI:eq:1-7-1}. 
  \item For $g_0:Y_0\to X_0$ a morphism, we have 
    \begin{equation*}\tag{1.7.3}\label{VI:eq:1-7-3}
      \sF(\chi,f_0\circ g_0) = g_0^\ast \sF(\chi,f_0) \text{.}
    \end{equation*}
    In particular, $\sF(\chi,f_0) = f_0^\ast \sF(\chi)$, where $\sF(\chi)$ 
    denotes $\sF(\chi,\operatorname{id}_{G_0})$. 
  \item For $u_0:G_0\to H_0$ a morphism, and 
    $\chi:H_0(\dF_q) \to E_\lambda^\times$, we have 
    \begin{equation*}\tag{1.7.4}\label{VI:eq:1-7-4}
      \sF(\chi,u_0 f_0) = \sF(\chi u_0,f_0) \text{.}
    \end{equation*}
  \item For $G_0=\prod_{i\in I} G_0^i$, with $\chi$ having coordinates $\chi_i$ 
    ($i\in I$) and $f_0$ having coordinates $f_0^i$, it follows from 
    \eqref{VI:eq:1-7-1} through \eqref{VI:eq:1-7-4} that 
    \begin{equation*}\tag{1.7.5}\label{VI:eq:1-7-5}
      \sF(\chi,f_0) = \bigotimes_{i\in I} \sF(\chi_i,f_0^i) \text{.}
    \end{equation*}    
\end{enumerate}

Note the $\chi^{-1}$ in Definition \ref{VI:1-7}. It ensures that, for 
$x\in X^F$, on $\sF(\chi,f_0)_x$ one has 
\begin{equation*}\tag{1.7.6}\label{VI:eq:1-7-6}
  F_x^\ast = \chi f_0(x)
\end{equation*}
(use the fact that the fiber at $x$ of the morphism \eqref{VI:eq:1-2-2}, for 
$\rho=\chi^{-1}$, commutes with $F_x^\ast$, together with \eqref{VI:eq:1-5-1}). 

If $\sF(\chi,f_0)_1$ is the sheaf obtained from $\sF(\chi,f_0)$ by extension 
from the base field to $\dF_{q^n}$, we deduce from \eqref{VI:eq:1-6-2} that 
\begin{equation*}\tag{1.7.7}\label{VI:eq:1-7-7}
  \sF(\chi,f_0)_1 = \sF(\chi\circ N,f_1) \text{.}
\end{equation*}





\subsection{Abuse of notations}\label{VI:1-8}

(i) If $\Xi$ denotes the composite map 
$\chi f_0:X_0(\dF_q) \to E_\lambda^\times$, we will sometimes write $\sF(\Xi)$ 
instead of $\sF(\chi,f_0)$. By \eqref{VI:eq:1-7-1} through \eqref{VI:eq:1-7-5}, 
there is little risk of ambiguity. For example:
\begin{enumerate}[a)]
  \item we write $\sF(\chi)$ for $\sF(\chi,\operatorname{id}_{G_0})$ 
    (notation already used in \ref{VI:1-7}b)); 
  \item we write $\sF(\chi f_0)$ for $\sF(\chi,f_0)$;
  \item with the notation of \ref{VI:1-7}d), we write 
    $\sF(\prod \chi_i f_0^i)$ for $\sF(\chi,f_0)$. 
\end{enumerate}

With this notation, \eqref{VI:eq:1-7-4} shows that the notation 
$\sF(\chi u_0 f_0)$ is unambiguous, \eqref{VI:eq:1-7-1} \eqref{VI:eq:1-7-2} 
\eqref{VI:eq:1-7-5} express the multiplicativity of $\sF(\Xi)$ with respect to $\Xi$, 
and \eqref{VI:eq:1-7-6} becomes $F_x^\ast = \Xi(x)$ on $\sF(\Xi)_x$.

(ii) If $X$ is a scheme over an extension $k$ of $\dF_q$, and $f$ is a 
morphism of $X$ into $G_0\otimes_{\dF_q} k$, we shall again denote by $\sF(\chi,f)$, 
$\sF(\chi f)$, or $\sF(\chi)$ (for $f$ an inclusion) the inverse image under 
$f$ of the $E_\lambda$-sheaf obtained from $\chi^{-1}(L_0(G))$ by extension of 
scalars from $\dF_q$ to $k$. If $f_0$ is the composite 
$X\to G_0\otimes_{\dF_q} k \to G_0$, this is again the $\sF(\chi,f_0)$ of 
\ref{VI:1-7}. If $k$ is a finite extension of $\dF_q$, we have 
$\sF(\chi f) = \sF(\chi N_{k/\dF_q} f)$. 

Applying \eqref{VI:eq:1-1-1} to \eqref{VI:eq:1-7-6} and \eqref{VI:eq:1-7-7}, 
we find: 





% NOTE: originally a 'scholium'
\begin{theorem_}\label{VI:1-9}
Let $S_0$ be a separated scheme of finite type over $\dF_q$, let $G_0$ be a 
connected commutative algebraic group over $\dF_q$, let $f_0:S_0\to G_0$ be a morphism, and let 
$\chi:G_0(\dF_q) \to E_\lambda^\times$. We have 
\begin{equation*}\tag{1.9.1}\label{VI:eq:1-9-1}
  \sum_{s\in S_0(\dF_q)} \chi f_0(s) = \tr(F^\ast,\h_c^\bullet(S,\sF(\chi,f_0))) 
\end{equation*}
and, for any integer $n\geqslant 1$,
\begin{equation*}\tag{1.9.2}\label{VI:eq:1-9-2}
  \sum_{s\in S_0(\dF_{q^n})} \chi\circ N_{\dF_{q^n}/\dF_q} f_0(s) = \tr({F^\ast}^n, \h_c^\bullet\left(S,\sF(\chi,f_0))\right) \text{.}
\end{equation*}
\end{theorem_}





\addtocounter{subsubsection}{2}
\subsubsection{Remark}\label{VI:1-9-3}

Take $G_0$ to be a product. With the notation of \ref{VI:1-7}d) and 
\ref{VI:1-8}c), the formula \eqref{VI:eq:1-9-2} becomes 
\[
  \sum_{s\in S_0(\dF_{q^n})} \prod_i \chi_i\circ N_{\dF_{q^n}/\dF_q} (f_0^i(s)) = \tr\Big({F^\ast}^n, \h_c^\bullet\Big(S,\sF\Big(\prod_i \chi_i f_0^i\Big)\Big)\Big) \text{.}
\]





\subsubsection{Remark}\label{VI:1-9-4}

Take $G_0=\{e\}$. Formula \eqref{VI:eq:1-9-1} becomes 
\[
  \left|S_0(\dF_q)\right| = \sum_{s\in S_0(\dF_q)} 1 = \tr(F^\ast,\h_c^\bullet(S,\dQ_\ell)) \text{.}
\]

It is rare to be able to compute the right-hand side of 
\eqref{VI:eq:1-9-1} explicitly. Here is an amusing example, with $G_0=\{e\}$. 





\subsection{Example}\label{VI:1-10}

A nonsingular projective quadric of odd dimension over \texorpdfstring{$\dF_q$}{Fq} has the same number of rational points as projective space of the same dimension over \texorpdfstring{$\dF_q$}{Fq}.

If, over $\dC$, $X$ is a nonsingular quadric hypersurface in 
the projective space $\dP^{2 N}$, and $Y$ a hyperplane, we know that for 
$i\leqslant 2\dim(X)=2\dim(Y)$, the inclusions 
$X\hookrightarrow \dP^{2 N}\hookleftarrow Y$ induce isomorphisms (in 
ordinary cohomology) 
\[\xymatrix{
  \h^i(X,\dQ) 
    & \ar[l]_-\sim \h^i(\dP^{2 N},\dQ) \ar[r]^-\sim 
    & \h^i(Y,\dQ) \text{.}
}\]

By specialization, it follows that if $X_0'$ is a nonsingular quadric 
hypersurface in projective space $\dP_0^{2 N}$ over $\dF_q$, 
and $Y_0'$ a hyperplane, the inclusions 
$X_0'\hookrightarrow \dP_0^{2 N}\hookleftarrow Y_0'$ induce 
isomorphisms 
\[\xymatrix{
  \h^i(X',\dQ_\ell) 
    & \ar[l]_-\sim \h^i(\dP^{2 N},\dQ_\ell) \ar[r]^-\sim 
    & \h^i(Y',\dQ_\ell) \text{.}
}\]
These isomorphisms commute with $F^\ast$, and the trace formula applies. 





\subsection{}\label{VI:1-11}

The trace formula can also be used to understand the classical formulas 
giving the number of rational points of linear groups over 
finite fields, and of certain homogeneous spaces (see \hyperref[VI:8]{Section 8}). 





\subsection{}\label{VI:1-12}

There is a dictionary that translates various kinds of classical 
manipulations of trigonometric sums into cohomological terms. This dictionary 
will be given in \hyperref[VI:2]{Section 2}. Since the cohomological statement is more 
``geometric,'' it can sometimes be applied 
to situations where the classical argument only applies after a 
extension of the finite base field. Here is an example (a special case of a 
theorem of Kazhdan).

\begin{theorem_}[Kazhdan]\label{VI:1-13}
Let $X_0$ be a scheme over $\dF_q$. Suppose there is an action $\rho$ of 
$\dG_a$ on $X$, and a morphism $f:X\to Y$ of schemes over $\dF$, making 
$X$ a $\dG_a$-torsor over $Y$. Then the number of rational points of $X_0$ 
is divisible by $q$.
\end{theorem_}

First suppose that $\rho$, $Y$, and $f$ are defined over $\dF_q$. 


\paragraph{Classical argument}
The fibers of $f:X_0(\dF_q)\to Y_0(\dF_q)$ all have $q$ elements: 
$|X_0(\dF_q)| = q\cdot |Y_0(\dF_q)|$, hence the divisibility. 


\paragraph{Cohomological translation}
The higher direct image sheaves $\eR^i f_!\dQ_\ell$ are 
\[
  \eR^i f_!\dQ_\ell = \begin{cases}
                       \dQ_\ell(-1) & \text{if $i = 2$} \\
                       0            & \text{for $i\ne 2$}
                     \end{cases}
\]
The Leray spectral sequence for $f$ therefore degenerates, giving an isomorphism 
\[
  \h_c^i(X,\dQ_\ell) = \h_c^{i-2}(Y,\dQ_\ell)(-1) \text{.}
\]
To understand the effect of a Tate twist on the eigenvalues of 
Frobenius, it is most convenient to use the Galois point of view 
(\hyperref[II]{Report} \ref{II:1-8}) and note that the Frobenius substitution 
$\varphi$ acts on $\dQ_\ell(1)$ by multiplication by $q$. The 
eigenvalues of $F^\ast$ on $\h_c^p(X,\dQ_\ell)$ are therefore obtained by 
multiplying by $q$ the eigenvalues of $F^\ast$ acting on $\h_c^{i-2}(Y,\dQ_\ell)$. We 
know that the latter are algebraic integers \cite[XXI 5.2.2]{sga7}; hence 

\begin{equation*}\tag{$*$}\label{VI:eq:*}
  \substack{\displaystyle\text{the eigenvalues of $F^\ast$ acting on $\h_c^i(X,\dQ_\ell)$ are} \\
  \displaystyle\text{algebraic integers divisible by $q$.}}
\end{equation*}


\paragraph{Descent}
Return to the hypotheses of the theorem. For a suitable $n$, we may 
suppose that $\rho$, $Y$, and $f$ are defined over $\dF_{q^n}$. Since the 
Frobenius morphism relative to $\dF_{q^n}$ is the $n$-th power of the one 
relative to $\dF_q$, statement \eqref{VI:eq:*} for the scheme over 
$\dF_{q^n}$ obtained from $X_0$ by extension of scalars gives: 
\begin{equation*}\tag{$**$}\label{VI:eq:**}
  \substack{\displaystyle\text{the $n$-th powers of the eigenvalues of $F^\ast$ acting on} \\ \displaystyle\text{$\h_c^i(X,\dQ_\ell)$ are algebraic integers divisible by $q^n$.}}
\end{equation*}

This assertion implies \eqref{VI:eq:*}. The number of rational points of 
$X_0$, namely $\sum (-1)^i \tr(F^\ast,\h^i(X,\dQ_\ell))$, is therefore an 
algebraic integer divisible by $q$. It follows that, as a rational integer, 
it is divisible by $q$. 





\subsection{}\label{VI:1-14}

Formula \ref{VI:1-9-3} controls the dependence on $n$ of the trigonometric 
sum on the left-hand side. It implies identities between a 
trigonometric sum and those which are deduced from it by ``extension of 
scalars.'' The most famous of these identities is that of 
Hasse--Davenport. Let $\chi:\dF_q^\times \to \dC^\times$ and 
$\psi:\dF_q\to \dC^\times$ be characters, with $\psi$ nontrivial, and 
define the Gauss sum $\tau(\chi,\psi)$ by 
\begin{equation*}\tag{1.14.1}\label{VI:eq:1-14-1}
  \tau(\chi,\psi) = - \sum_{x\in \dF_q^\times} \psi(x) \chi^{-1}(x) \text{.}
\end{equation*}





\begin{theorem_}[Hasse-Davenport]\label{VI:1-15}
We have 
\begin{equation*}\tag{1.15.1}\label{VI:eq:1-15-1}
  \tau\left(\chi\circ N_{\dF_{q^n}/\dF_q},\psi\circ \tr_{\dF_{q^n}/\dF_q}\right) = \tau(\chi,\psi)^n \text{.}
\end{equation*}
\end{theorem_}

Let $E\subset \dC$ be a number field containing the values of $\chi$ and 
$\psi$, and let $E_\lambda$ be the completion of $E$ at a place of residue 
characteristic $\ell\ne p$. Identity \eqref{VI:eq:1-15-1} is an identity in 
$E$. Proving it in $\dC$ is equivalent to proving it in $E_\lambda$. We 
shall prove it in $E_\lambda$, regarding $\chi$ and $\psi$ as taking values 
in $E_\lambda^\times$.

Apply \ref{VI:1-9-3} with $X_0=\dG_m$ over $\dF_q$ and 
$G_0=\dG_m\times \dG_a$. The groups $\h_c^i(\dG_m,\sF(\chi^{-1}\psi))$ vanish for 
$i\ne 1$, and $\h_c^1$ has dimension $1$ (\ref{VI:4-2}). By 
\ref{VI:1-9-3}, 
$\tau(\chi\circ N_{\dF_{q^n}/\dF_q},\psi\circ \tr_{\dF_{q^n}/\dF_q})$ is 
the unique eigenvalue of $(F^\ast)^n$ on $\h_c^1$, and \eqref{VI:eq:1-15-1} 
follows. 

Similar examples are given in Weil \cite[App V]{we74}. 





\subsection{}\label{VI:1-16}

An $E_\lambda$-sheaf $\sF_0$ on $X_0$ is said to be \emph{pointwise pure of 
weight $n$} if for any closed point $x\in |X_0|$, the eigenvalues of 
$F_x^\ast$ (\hyperref[II]{Report} \ref{II:1-2}) are algebraic numbers 
all of whose complex conjugates have absolute value 
$q_x^{n/2}$, where $q_x$ is the number of elements of $k(x)$. The 
following theorem will be proved in \cite{de80}. 





\begin{theorem_}\label{VI:1-17}
If $\sF_0$ is pointwise pure of weight $n$, then for any eigenvalue $\alpha$ 
of the endomorphism $F^\ast$ of $\h_c^i(X,\sF)$, there exists an integer 
$m\leqslant n+i$ such that the complex conjugates of $\alpha$ are all of 
absolute value $q^{m/2}$.
\end{theorem_}

The sheaves considered in \ref{VI:1-9} have weight $0$. The theorem 
therefore gives, for the sum \ref{VI:eq:1-9-1} (or rather, for all its 
complex conjugates), the bound 
\begin{equation*}\tag{1.17.1}\label{VI:eq:1-17-1}
  \left|\sum_{s\in S_0(\dF_q)} \prod_i \chi_i(f^i(s)) \right| \leqslant \sum_i \dim \h_c^i\left(S,\sF\left(\prod \chi_i f^i\right)\right)\cdot q^{i/2} \text{.}
\end{equation*}





\subsection{Remarks}\label{VI:1-18}

\begin{enumerate}[a)]
  \item If $\dim S=n$, for any $E_\lambda$-sheaf $\sF$ on $S$, we have 
    \[
      \h_c^i(S,\sF) = 0 \qquad \text{for $i\notin [0,2n]$.}
    \]
  \item The group $\h_c^0(X,\sF)$ is the group of global sections of $\sF$ 
    on $S$ whose support is proper. If $S$ is connected, not proper and 
    $\sF$ is smooth, this group is zero; the same holds if $S$ is connected, 
    $\sF$ is smooth of rank one, and is nonconstant.
  \item If $S$ is smooth and purely of dimension $n$, and $\sF$ is smooth, 
    Poincaré duality says that the vector spaces 
    $\h_c^i(X,\sF)$ and $\h^{2n-i}(X,\sF^\vee)(n)$ are dual to each other 
    (for the effect of a Tate twist on the eigenvalues of Frobenius, 
    cf. the $2$nd part of the proof of \ref{VI:1-13}).
  \item If $\dim S=n$, Poincaré duality allows us to compute 
    $\h_c^{2n}(S,\sF)$ as follows: take an open subset $U$ of $S_\text{red}$ 
    whose complement has dimension $<n$; require this open subset to be smooth and 
    purely of dimension $n$, with $\sF$ smooth on it. We then have 
    $\h_c^{2n}(U,\sF) \iso \h_c^{2n}(S,\sF)$, because in the long exact sequence 
    of cohomology, $\h_c^i(S\setminus U,\sF) = 0$ for $i=2n,2n-1$. 
    Poincaré duality therefore gives 
    \[
      \h_c^{2n}(S,\sF) \iso \left(\h^0(U,\sF^\vee)(n)\right)^\vee \text{.}
    \]
    
    Let $U$ be connected, and let $u$ be a geometric point of $U$. The 
    ``local system'' $\sF$ corresponds to a representation of 
    $\pi_1(U,u)$ on $\sF_u$, and $\h^0(U,\sF^\vee)$ is 
    $(\sF_u^\vee)^{\pi_1(U,u)}=((\sF_u)_{\pi_1(U,u)})^\vee$ (dual of 
    coinvariants). We therefore have 
    \[
      \h_c^{2n}(S,\sF) \simeq (\sF_u)_{\pi_1(U,u)}(-n) \text{.}
    \]
  \item These remarks often allow the calculation of $b_0$ and $b_{2n}$. 
    Computing the other $b_i$ can be difficult. If $b_{2n}=0$ and 
    $b_{2n-1}=0$, the bound \eqref{VI:eq:1-17-1} is a Lang--Weil bound, 
    with, for large $q$, a gain of $\sqrt q$ over the 
    trivial $O(q^n)$ bound. One should bear in mind, however, that 
    \eqref{VI:eq:1-17-1} does not formally imply the trivial bound 
    \[
      \left|\sum_{s\in S_0(\dF_q)} \prod_i \chi_i(f^i(s))\right| \leqslant \left| S_0(\dF_q)\right| \text{,}
    \]
    so it may be useful to take this into account (cf. the proof of \ref{VI:3-8}). 
\end{enumerate}

Below I explain a method for proving that $b_i=0$ for $i\ne n$. 
When it applies, it gives an estimate of order $O(q^{n/2})$. The constant 
implicit in $O$ is $b_n=(-1)^n \chi(S,\sF(\prod \chi_i f^i))$, which may 
be difficult to compute. 





\begin{proposition_}\label{VI:1-19}
Let $\bar X$ be a proper scheme over an algebraically closed field $k$, 
let $j:X\hookrightarrow \bar X$ be an open immersion with domain $X$, and let $\sF$ be an $E_\lambda$-sheaf 
on $X$. Assume that 
\begin{enumerate}[\indent a)]
  \item $(j_\ast\sF)_x = 0$ for $x\in \bar X\setminus X$, i.e. 
    $j_!\sF\iso j_\ast\sF$; 
  \item $\eR^i j_\ast \sF = 0$ for $i>0$.
\end{enumerate}
Then the maps 
\[
  \h_c^i(X,\sF) \to \h^i(X,\sF)
\]
are isomorphisms. 
\end{proposition_}

Assumptions a) and b) mean that $j_! \sF \iso \eR j_\ast \sF$, hence 
\[
  \h_c^i(X,\sF) = \h^i(\bar X,j_!\sF) \iso \hh^i(\bar X,\eR j_\ast \sF) = \h^i(X,\sF) \text{.}
\]
In other words, the Leray spectral sequence for $j$ 
\[
  E_2^{pq} = \h^p(\bar X,\eR^q j_\ast \sF) \Rightarrow \h^{p+q}(X,\sF) 
\]
degenerates, by b) into isomorphisms 
\[
  \h^i(\bar X,j_\ast\sF) \iso \h^i(X,\sF) \text{,}
\]
while, by a) 
\[
  \h_c^i(X,\sF) = \h^i(\bar X,j_!\sF) \iso \h^i(\bar X,j_\ast\sF) \text{.}
\]





\subsubsection{Example}\label{VI:1-19-1}

If $X$ is a curve, or if $\bar X$ is smooth, $X$ is the complement of a 
normal-crossings divisor, and $\sF$ is tamely ramified, 
hypothesis a) of \ref{VI:1-19} (no invariants under local monodromy) 
implies hypothesis b). 





\begin{proposition_}\label{VI:1-20}
Let $X_0$ be a smooth affine scheme, purely of dimension $n$, over $\dF_q$, and let 
$\sF_0$ be a smooth $E_\lambda$-sheaf, pointwise pure of weight $m$, on $X_0$. We 
assume that the maps $\h_c^i(X,\sF) \to \h^i(X,\sF)$ are 
isomorphisms (cf. \ref{VI:1-17}). Then, $\h_c^i(X,\sF) = 0$ for $i\ne n$, and 
the complex conjugates of the eigenvalues of $F^\ast$ on $\h_c^n(X,\sF)$ 
are of absolute value $q^{(m+n)/2}$. 
\end{proposition_}

Poincaré duality pairs $\h_c^i(X,\sF)$ (resp. 
$\h_c^i(X,\sF^\vee)$) with $\h^{2n-i}(X,\sF^\vee(n))$ (resp. 
$\h^{2n-i}(X,\sF(n))$). According to \cite[XIV 3.2]{sga4}, the ordinary $\h^i$ groups 
vanish for $i>n$. By duality, the $\h_c^i$ are zero for $i<n$; since 
$\h_c^i(X,\sF) \iso \h^i(X,\sF)$, these groups are zero for $i\ne n$. 
Apply \ref{VI:1-17} to $\sF_0$ and to $\sF_0^\vee(n)$ (of weight 
$-m-2n$). We find that the complex conjugates $\alpha$ of the eigenvalues 
of $F^\ast$ on $\h_c^n(X,\sF) = \h_c^n(X,\sF^\vee(n))^\vee$ satisfy 
\begin{align*}
  |\alpha| &\leqslant q^{(m+n)/2} && \text{and} \\
  |\alpha^{-1}| &\leqslant q^{((-m-2n)+n)/2} = q^{-(m+n)/2} \text{,}
\end{align*}
hence the assertion. 





\subsection{Remark}\label{VI:1-21}

The last argument shows that if $X_0$ is a smooth separated scheme over 
$\dF_q$, if $\sF_0$ is a smooth $E_\lambda$-sheaf pointwise pure of weight 
$m$ on $X_0$, and if $\h_c^i(X,\sF) \hookrightarrow \h^i(X,\sF)$, then the 
complex conjugates of the eigenvalues of $F^\ast$ on $\h_c^i(X,\sF)$ are 
of absolute value $q^{(m+i)/2}$. 





\subsection*{Appendix. The Lang isogeny in the noncommutative case}




\subsection{}\label{VI:1-22}

Let $G_0$ be a connected algebraic group over $\dF_q$ and let $\fL$ be the morphism 
$x\mapsto F x\cdot x^{-1}$ from $G_0$ to itself. It is a morphism of 
$G_0$-homogeneous spaces if, on the source, $G_0$ acts by left translation 
$x\mapsto (y\mapsto x y)$, and on the target by 
$x\mapsto (y\mapsto F x\cdot y\cdot x^{-1})$). We have $\fL(xy) = \fL(x)$ for 
$y\in G_0(\dF_q)$, and $\fL$ induces an isomorphism $G_0/G_0(\dF_q)\iso G_0$. 

As in the commutative case, $\fL$ therefore makes $G_0$ a $G_0(\dF_q)$-torsor 
on $G_0$, the \emph{Lang torsor} $L$. If 
$\rho:G_0(\dF_q) \to \gl(n,E_\lambda)$ is a linear representation of 
$G_0(\dF_q)$, we denote by $\sF(\rho)$ the $E_\lambda$-sheaf $\rho(L)$ 
(\ref{VI:1-2}). 





\subsection{}\label{VI:1-23}

Let $\gamma\in G^F=G_0(\dF_q)$, and compute the endomorphism $F^\ast$ of 
$\sF(\rho)_\gamma$. Let $g\in G$ be such that $\fL(g)=\gamma$, and choose the frame 
$\rho(g)\in \operatorname{Isom}(E_\lambda^n,\sF(\rho)_\gamma)$. For 
$e\in E_\lambda^n$, we have 
\[
  F(\rho(g)(e)) = \rho(F g)(e) = \rho(\gamma g)(e) = \rho(g\cdot g^{-1}\gamma g)(e) \text{.}
\]
We will see that $g^{-1} \gamma g\in G_0(\dF_q)$, hence 
\[
  F(\rho(g)(e)) = \rho(g) \rho(g^{-1}\gamma g)(e)
\]
and 
\begin{equation*}\tag{1.23.1}\label{VI:eq:1-23-1}
  \rho(g)^{-1} (F_\gamma^\ast)^{-1} \rho(g) = \rho(g^{-1} \gamma g) \text{:}
\end{equation*}
$\rho(g)$ identifies the inverse of $F^\ast$ at $\gamma$ with the automorphism 
$\rho(g^{-1} \gamma g)$ of $E_\lambda^n$. It remains to understand 
$g^{-1} \gamma g$.

\begin{lemma_}\label{VI:1-24}
\begin{enumerate}[(i)]
  \item If $F g\cdot g^{-1} \in G_0(\dF_q)$, we have 
    $g^{-1} (F g\cdot g^{-1}) g = g^{-1} F g\in G_0(\dF_q)$.
  \item Let $\gamma\in G_0(\dF_q)$, and $g$ such that 
    $F g\cdot g^{-1} = \gamma$. The conjugation class of $\gamma'=g^{-1} Fg$ 
    in $G_0(\dF_q)$ only depends on that of $\gamma$. 
  \item The map induced by $\gamma\mapsto \gamma'$ on the set 
    $G_0(\dF_q)^\natural$ of conjugacy classes of $G_0(\dF_q)$ 
    is bijective. Its inverse is the corresponding map for the opposite 
    group. 
\end{enumerate}
\end{lemma_}

(i) If $\gamma= F g\cdot g^{-1}\in G_0(\dF_q)$, the identities $F g=\gamma g$ 
and 
$F\gamma=\gamma$ give $F(g^{-1} F g) = (F g)^{-1} F F g = g^{-1} \gamma^{-1} F(\gamma g) = g^{-1} \gamma^{-1} F \gamma F g = g^{-1} F g$, 
so that $g^{-1} F g \in G_0(\dF_q)$. 

(ii) For $\alpha,\gamma\in G_0(\dF_q)$, the $g$ such that 
$F g\cdot g^{-1} = \gamma$ form a right coset under $G_0(\dF_q)$, and 
if $F g\cdot g^{-1} = \gamma$, we have 
$F(\alpha g \alpha^{-1})(\alpha g \alpha^{-1})^{-1} = \alpha \gamma\alpha^{-1}$. 
As $\gamma$ ranges over a conjugacy class in $G_0(\dF_q)$, the $g$ 
such that $F g\cdot g^{-1} = \gamma$ therefore range over a double coset under 
$G_0(\dF_q)$, and $\gamma'= g^{-1} F g$ ranges over a conjugacy class. 

(iii) Finally, the inverse of $F g\cdot g^{-1}\mapsto g^{-1} F g$ is 
$g^{-1} F g\mapsto F g\cdot g^{-1}$. 





\subsection{}\label{VI:1-25}

Formula \eqref{VI:eq:1-23-1} can be read succinctly: $F^\ast$ at 
$\gamma$ is $\rho(\gamma')^{-1}$. If $\gamma$ belongs to the identity component 
of its centralizer, we can take $g$ in $Z(\gamma)$. We then have 
$\gamma=\gamma'$. If this condition is not met, $\gamma$ and $\gamma'$ 
need not be conjugate. 





\subsection{An example}\label{VI:1-26}

For $G_0=\operatorname{SL}(2)$, $q$ odd, 
$\alpha\in \dF_q^\times\setminus {\dF_q^\times}^2$, and $\bar\gamma$ is the 
conjugacy class of $-1(1+N)$, with $N^2=0$, then $\bar\gamma'$ is the 
conjugacy class of $-1\cdot (1+\alpha N)$, so that $\bar\gamma'\ne \bar\gamma$ 
if $N\ne 0$. 










\section{Dictionary}\label{VI:2}

Elementary manipulations of trigonometric sums can often 
be seen, via the trace formula, as counterparts of statements 
in cohomology. In this section, I list such statements and their 
counterparts; they carry the same number, with a $\ast$ appended to 
the cohomological statement. 





\subsection{}\label{VI:2-1}

``A sum of algebraic integers is an algebraic integer'' has as its 
cohomological analogue the following theorem, proved in 
\cite[XXI 5.2.2]{sga7} (see the use of \ref{VI:2-1}$^\ast$ in 
\ref{VI:1-13}). 





\begin{theorem*}[2.1*]\label{VI:2-1*}
Let $X_0$ be a separated scheme of finite type over $\dF_q$ and let $\sF_0$ be an 
$E_\lambda$-sheaf on $X_0$. If, for all $x\in |X_0|$, the 
eigenvalues of $F_x^\ast$ on $\sF_0$ are algebraic integers, then the 
eigenvalues of $F^\ast$ on $\h_c^i(X,\sF)$ are algebraic integers. 
\end{theorem*}





\subsection{}\label{VI:2-2}

Other integrality results are proved in 
\cite[XXI, 5.2.2 et 5.4]{sga7}: if $\dim(X_0)\leqslant n$, and $\alpha$ is 
an eigenvalue of $F^\ast$ on $\h_c^i(X,\sF)$, we have: 
\begin{enumerate}[a)]
  \item under the assumptions of \ref{VI:2-1}$^\ast$, and if $i>n$, then 
    $q^{i-n}$ divides $\alpha$;
  \item if the inverses of the eigenvalues of $F_x^\ast$ are algebraic 
    integers, then $\alpha^{-1}$ is integral away from $p$. More 
    precisely, $q^{\inf(n,i)}\alpha^{-1}$ is an algebraic integer. 
\end{enumerate}

For the sheaves considered in \ref{VI:1-9}, the eigenvalues of 
$F_x^\ast$ are roots of unity, so the assumptions of 
the theorem and of b) above are satisfied. 





\subsection{}\label{VI:2-3}

Let $f:A\to B$ be a map from one finite set to another, and 
$\varepsilon$ a function on $A$. We have 
\begin{equation*}\tag{2.3.1}\label{VI:eq:2-3-1}
  \sum_{a\in A} \varepsilon(a) = \sum_{b\in B} \sum_{f(a)=b} \varepsilon(a) \text{.}
\end{equation*}





\subsection*{2.3*}\label{VI:2-3*}

Let $f:X\to Y$ be a morphism of separated schemes of finite type over an 
algebraically closed field $k$, and let $\sF$ be an $E_\lambda$-sheaf on $X$. The 
Leray spectral sequence for $f$ in cohomology with compact support is 
\begin{align*}\tag{2.3.1*}\label{VI:eq:2-3-1*}
  E_2^{p q} = \h_c^p(Y,\eR^q f_! \sF) \Rightarrow \h_c^{p+q}(X,\sF) \text{.}
\end{align*}

Let $f_0:X_0 \to Y_0$ be a morphism of separated schemes of finite type over 
$\dF_q$ and let $\sF_0$ be an $E_\lambda$-sheaf on $X_0$. Between 
\eqref{VI:eq:2-3-1} and \eqref{VI:eq:2-3-1*} there is the following compatibility. 

\begin{enumerate}[a)]
  \item The sheaf $\eR^q f_! \sF$ is obtained by extending scalars from 
    $\dF_q$ to $\dF$ in the sheaf $\eR^q f_{0!} \sF_0$ on $Y_0$. As 
    such, it is equipped with a Frobenius correspondence 
    $F^\ast:F^\ast\eR^q f_! \sF \to \eR^q f_! \sF$. It is obtained by applying 
    the functoriality of $\eR f_!$ to the Frobenius correspondence of $\sF$: 
    namely the composite 
    \[\xymatrix{
      F^\ast \eR^q f_! \sF \ar[r]^-\sim 
        & \eR^q f_! F^\ast \sF \ar[r]^-\sim 
        & \eR^q f_! \sF \text{.}
    }\]
  \item The trace formula for $\sF_0$ reads 
    \begin{align*}\tag{1}\label{VI:eq:1}
      \sum_{x\in X_0(\dF_q)} \tr(F_x^\ast,\sF_0) = \tr(F^\ast,\h_c^\bullet(X,\sF)) \text{;}
    \end{align*}
    while that for $\eR^q f_! \sF$ reads 
    \begin{align*}\tag{2$_q$}\label{VI:eq:2_q} 
      \sum_{y\in Y_0(\dF_q)} \tr(F_y^\ast,\eR^q f_! \sF) = \tr(F^\ast, \h_c^\bullet(Y,\eR^q f_! \sF)) \text{.}
    \end{align*}
\end{enumerate}

We have, for $y\in Y(\dF)$, $(\eR^q f_! \sF)_y = \h_c^q(f^{-1}(y),\sF)$, and if 
$y\in Y_0(\dF_q)$, 
\[
  \tr(F_y^\ast,\eR^q f_! \sF) = \tr(F^\ast,\h_c^q(f^{-1}(y),\sF)) \text{,}
\]
hence by the trace formula on the fiber 
\[
  \sum_{\substack{x\in X_0(\dF_q) \\ f(x) = y}} \tr(F_x^\ast,\sF_0) = \sum (-1)^q \tr(F_y^\ast,\eR^q f_! \sF) \text{.}
\]
Taking the alternating sum of the identities \eqref{VI:eq:2_q}, we therefore find 
\begin{align*}\tag{2}\label{VI:eq:2}
  \sum_{y\in Y_0(\dF_q)} \sum_{\substack{x \in X_0(\dF_q) \\ f(x) = y}} \tr(F_x^\ast,\sF_0) = \tr(F^\ast,\h_c^\bullet(Y,\eR^\bullet f_! \sF)) \text{.}
\end{align*}

Equality of the left-hand sides of \eqref{VI:eq:1} and \eqref{VI:eq:2} 
follows from \eqref{VI:eq:2-3-1}; equality of the right-hand sides follows from 
\eqref{VI:eq:2-3-1*} (taking into account that $F^\ast$ is an endomorphism of 
the entire spectral sequence). 





\subsection{}\label{VI:2-4}

Let $(A_i)_{i\in I}$ be a finite family of finite sets, 
$A=\prod_{i\in I} A_i$, $\varepsilon_i$ a function on $A_i$, and 
$\varepsilon$ the function $a\mapsto \prod \varepsilon_i(a_i)$ on $A$. We have 
\begin{align*}\tag{2.4.1}\label{VI:eq:2-4-1}
  \sum_{a\in A} \varepsilon(a) = \prod_{i\in I} \sum_{a_i\in A_i} \varepsilon_i(a_i) \text{.}
\end{align*}





\subsection*{2.4*}\label{VI:2-4*}

Let $(X_i)_{i\in I}$ be a finite family of separated schemes of finite type 
over an algebraically closed field $k$, let $X=\prod_{i\in I} X_i$, let $\sF_i$ be an 
$E_\lambda$-sheaf on $X_i$ and $\sF$ the external tensor product 
$\bigboxtimes_{i\in I}\sF_i = \bigotimes_{i\in I} \pr_i^\ast(\sF_i)$ 
of the $\sF_i$. 
The K\"unneth formula gives 
\begin{align*}\tag{2.4.1*}\label{VI:eq:2-4-1*}
  \h_c^\bullet(X,\sF) = \bigotimes_{i\in I} \h_c^\bullet(X_i,\sF_i) \text{.}
\end{align*}

For the isomorphism \eqref{VI:eq:2-4-1*} to be canonical and independent of 
the choice of a total ordering of $I$, the tensor product on the right must be 
understood as a graded tensor product, with the Koszul sign rule 
(\hyperref[IV]{Cycle}, \ref{IV:1-3}). 





\subsection{}\label{VI:2-5}

Let $A$ be a finite set, let $B$ be a subset of $A$, and let $\varepsilon$ be a 
function on $A$. We have 
\begin{align*}\tag{2.5.1}\label{VI:2-5-1}
  \sum_{a\in A} \varepsilon(a) = \sum_{a\in A\setminus B}\varepsilon(a) + \sum_{a\in B} \varepsilon(a) \text{.}
\end{align*}





\subsection*{2.5*}\label{VI:2-5*}

Let $X$ be a separated scheme of finite type over an algebraically closed field $k$, 
let $Y$ be a closed subscheme, and let $\sF$ be an $E_\lambda$-sheaf on $X$. We have 
a long exact cohomology sequence 
\begin{align*}\tag{2.5.1*}\label{VI:eq:2-5-1*}
\xymatrix{
  \cdots \ar[r]^-\partial 
    & \h_c^i(X\setminus Y,\sF) \ar[r] 
    & \h_c^i(X,\sF) \ar[r] 
    & \h_c^i(Y,\sF) \ar[r]^-\partial 
    & \cdots 
}
\end{align*}

More generally, given a finite filtration of $X$ by closed subsets 
$X_p$ ($p\in \dZ$; assume that $X_p\supset X_{p+1}$, that $X_p=X$ 
for $p$ sufficiently small, and that $X_p=\varnothing$ for $p$ sufficiently large), there is a 
spectral sequence 
\begin{align*}\tag{2.5.2*}\label{VI:eq:2-5-2*}
  E_1^{p q} = \h_c^{p+q}(X_p\setminus X_{p+1},\sF) \Rightarrow \h_c^{p+q}(X,\sF) \text{.}
\end{align*}





\subsection{}\label{VI:2-6}

Let $A$ be a finite set, $(A_i)_{i\in I}$ a finite covering of $A$ and 
$\varepsilon$ a function on $A$. For $J\subset I$, denote by $A_J$ 
the intersection of $A_j$ ($j\in J$). We have 
\begin{align*}\tag{2.6.1}\label{VI:eq:2-6-1}
  \sum_{J\subset I} (-1)^{|J|} \sum_{a\in A_J} \varepsilon(a) = 0 \text{.}
\end{align*}





\subsection*{2.6*}\label{VI:2-6*}

Let $X$ be a separated scheme of finite type over an algebraically closed field $k$, 
$(X_i)_{i\in I}$ a closed (resp. open) finite covering of $X$ by 
subschemes, and $\sF$ an $E_\lambda$-sheaf on $X$. For $J\subset I$, 
let $X_J$ be the intersection of the $X_j$ ($j\in J$). We have the respective 
Leray spectral sequences 
\begin{align*}\tag{2.6.1*}\label{VI:eq:2-6-1*}
  E_1^{p q} &= \bigoplus_{|J|=p+1>0} \h_c^q(X_J,\sF) \otimes_\dZ \textstyle\bigwedge^{|J|} \dZ^J \Rightarrow \h_c^{p+q}(X,\sF) \\ 
  \tag{2.6.2*}\label{VI:eq:2-6-2*}
  E_1^{p q} &= \bigoplus_{|J|=1-p>0} \h_c^q(X_J,\sF) \otimes_\dZ \textstyle\bigwedge^{|J|} \dZ^J \Rightarrow \h_c^{p+q}(X,\sF)
\end{align*}
If $J=\varnothing$ were not excluded, there would also be spectral 
sequences converging to $0$. The factors $\bigwedge^{|J|} \dZ^J$ are included 
to avoid choosing a total order on $I$. 





\subsection{}\label{VI:2-7}

Let $A$ be a finite commutative group, and $\chi:A \to E_\lambda^\times$ a 
nontrivial character. We have 
\begin{align*}\tag{2.7.1}\label{VI:eq:2-7-1}
  \sum_{a\in A} \chi(a) = 0 \text{.}
\end{align*}
We will prove the cohomological analogue of \eqref{VI:eq:2-7-1} by transposing 
the following proof: if $x\in A$, $a\mapsto x a$ is a permutation of $A$, and 
\begin{align*}
  \sum_{a\in A}\chi(a) = \sum_{a\in A} \chi(x a) &= \chi(x) \sum_{a\in A} \chi(a) && \text{, whence} \\
  (\chi(x)-1)\sum_{a\in A}\chi(a) &= 0 \text{,}
\end{align*}
and by hypothesis there exists $x$ such that $\chi(x)-1\ne 0$. 





\begin{theorem*}[2.7*]\label{VI:2-7*}
Let $G_0$ be a connected commutative algebraic group over $\dF_q$ and 
$\chi:G_0(\dF_q) \to E_\lambda^\times$ a nontrivial character. We have 
\begin{align*}\tag{2.7.1*}\label{VI:eq:2-7-1*}
  \h_c^\bullet(G,\sF(\chi)) = 0 \text{.}
\end{align*}
\end{theorem*}

For $x$ a rational point of $G$, let $t_x$ be the translation $t_x(g) = x g$ 
of $G$. The formula $\fL t_x = t_{\fL(x)} \fL$ expresses that $(t_x,t_{\fL(x)})$ 
is an automorphism of the diagram $G\xrightarrow\fL G$ ($G$ equipped with the torsor of 
Lang). Let $\rho(g)$ be the automorphism of $(G,\sF(\chi))$ induced by it. 
For $g\in G_0(\dF_q)$, i.e. for $\fL(g) = e$, it is the identity on $G$, and 
multiplication by $\chi(g)^{-1}$ on $\sF(\chi)$. 

Let $\rho_H(g)$ denote the automorphism of $\h_c^\bullet(G,\sF(\chi))$ induced by 
$\rho(g)$. A homotopy argument (Lemma \ref{VI:2-8} below) shows that 
$\rho_H(g) = \rho_H(e)$, and hence is the identity. On the other hand, for 
$g\in G_0(\dF_q)$, $\rho_H(g)$ is multiplication by 
$\chi^{-1}(g)$: on $\h^\bullet(G,\sF(\chi))$, multiplication by 
$(\chi^{-1}(g)-1)$ is zero. Taking $g$ such that $\chi(g)\ne 1$, we obtain 
\ref{VI:2-7*}. 





\begin{lemma_}\label{VI:2-8}
Let $X$ and $Y$ be two schemes over an algebraically closed field $k$, with $X$ 
separated and of finite type and $Y$ connected. Let $\sF$ be a sheaf on $X$ and 
let $(\rho,\varepsilon)$ be a family of endomorphisms of $(X,\sF)$ parametrized 
by $Y$: 
\begin{align*} 
  \rho:Y\times_k X &\to Y\times_k X && \text{is a $Y$-morphism, and} \\
  \varepsilon:\rho^\ast \pr_2^\ast \sF &\to \pr_2^\ast\sF && \text{a morphism of sheaves.} 
\end{align*}
Assume that $\rho$ is proper. For $y\in Y(k)$, let $\rho_H(y)^\ast$ denote 
the endomorphism of $\h_c^\bullet(X,\sF)$ induced by $\rho_y:X\to X$ and 
$\varepsilon_y:\rho_y^\ast\sF\to \sF$. Then, $\rho_H(y)^\ast$ is 
independent of $y$. 
\end{lemma_}

Indeed, $\eR^p \pr_{1!} \pr_2^\ast\sF$ is the 
constant sheaf on $Y$ with fiber $\h_c^p(X,\sF)$, and $\rho_H(y)^\ast$ is the 
fiber at $y$ of the endomorphism 
\[\xymatrix{
  \eR^p \pr_{1!} \pr_2^\ast \sF \ar[r]^-{\rho^\ast} 
    & \eR^p \pr_{1!} \rho^\ast \pr_2^\ast \sF \ar[r]^-\varepsilon 
    & \eR^p \pr_{1!} \pr_2^\ast \sF 
}\]
of this sheaf. 





\subsection{Remark}\label{VI:2-9}

Let $G_0$ be a connected algebraic group over $\dF_q$ and 
$\rho:G_0(\dF_q) \to \operatorname{GL}(V)$ a linear representation 
such that $V^{G_0(\dF_q)}=0$. Then again $\h_c^\bullet(G,\rho(L))=0$. 










\section{Sums in one variable}\label{VI:3}





\subsection{}\label{VI:3-1}

Weil was the first to apply ``cohomological'' methods to 
the study of trigonometric sums in one variable; since he had 
proved the analogue of the Riemann hypothesis for Artin's $L$ functions 
over function fields, these methods gave him excellent 
estimates (of order $O(\text{square root of the number of terms})$) and an 
understanding of how these sums behave under ``extension of the base field.'' 

The substance of this section is a presentation, in cohomological language, of 
his results. 





\subsection{}\label{VI:3-2}

Cohomological methods here lead us to compute the groups 
$\h_c^i(X,\sF)$, where $\sF_0$ is an $E_\lambda$-sheaf on a curve $X_0$ over 
$\dF_q$. These groups vanish for $i\ne 0,1,2$, and for 
$i=0,2$, they have a simple interpretation (\ref{VI:1-8} a,b,c). Furthermore, the 
Euler-Poincaré characteristic 
$\chi_c(\sF) = \sum (-1)^i \dim \h_c^i(X,\sF)$ (which is also equal to 
$\chi(\sF) = \sum (-1)^i \dim \h^i(X,\sF)$) can be computed in terms of 
$X$, the rank of $\sF$ at the generic points of $X$, and the 
ramification of $\sF$. The essential result is the following. 

Let $\bar X$ be a smooth connected projective curve of genus $g$ over an 
algebraically closed field $k$, let $X$ be the dense open subset of $\bar X$ 
complementary to a finite set $S$ of points, and let $\sF$ be a 
smooth $E_\lambda$-sheaf on $X$. Let $\chi(X)=2-2 g-\# S$ be the 
Euler--Poincaré characteristic of $X$, and let $\operatorname{rg}(\sF)$ be the 
rank of $\sF$. For each point $s\in S$, we define an integer 
$\swan_s(\sF)$, the Swan conductor, measuring the wild ramification of 
$\sF$, and 
\begin{align*}\tag{3.2.1}\label{VI:eq:3-2-1}
  \chi_c(\sF) = \operatorname{rg}(\sF) \cdot \chi(X) - \sum_{s\in S} \swan_s(\sF) 
\end{align*}
\cite{ra65}. 





\subsection{}\label{VI:3-3}

Let $\bar X$, $j:X\hookrightarrow\bar X$ and $\sF$ as in \ref{VI:3-2}. The 
Poincaré duality theorem takes the particularly convenient form: 
$\h^i(\bar X,j_\ast \sF)$ and $\h^{2-i}(\bar X,j_\ast(\sF^\vee(1)))$ are 
dual to each other (\hyperref[V]{Duality}, \ref{V:1-3}). 





\subsection{}\label{VI:3-4}

Take $k=\dF$, and suppose that $X$, $\bar X$, and $\sF$ arise from 
$X_0$, $\bar X_0$, and $\sF_0$ over $\dF_q$. If $\sF_0$ becomes trivial on a 
finite covering of $X_0$, Weil proved that the complex conjugates of 
the eigenvalues of $F^\ast$ on $\h^i(\bar X,j_\ast\sF)$ have absolute value 
$q^{i/2}$. Such sheaves $\sF_0$ correspond to Artin $L$-functions 
over the function field of $\bar X_0$. 





\subsection*{3.5}\label{VI:3-5_} % originally one of two 3.5s

Under the same hypotheses, or more generally if $E_\lambda$ is the 
completion at a place $\lambda$ of a number field $E$ and if $\sF_0$ 
belongs to an infinite compatible system of $v$-adic representations 
($v$ a place of $E$), we can compute 
\[
  \prod_i \det(-F^\ast,\h^i(X,\sF))^{(-1)^i} 
\]
from \emph{local} information on $\sF_0$. This is explained in 
\cite{de73}. For trivial $\sF_0$ on a finite covering of $X_0$, it is 
the expression of the constant in the functional equation of an Artin 
$L$-function as a product of local constants. 





\subsection{Example}\label{VI:3-5}

Let $\psi$ be the character 
$\exp\left(\frac{2\pi i}{p} \tr_{\dF_q/\dF_p}\right)$ of $\dF_q$; we have 
$\psi(a^p-a)=1$. Let $X_0$ be a smooth projective geometrically 
irreducible curve of genus $g$ over $\dF_q$, and let $f$ be a rational function on 
$X_0$, i.e. a morphism $f:X_0 \to \dP^1$, not identically equal to 
$\infty$. We are interested in the sum 
\[
  S_f' = \sum_{\substack{x\in X_0(\dF_q) \\ f(x)\ne \infty}} \psi(f(x)) 
  \text{.}
\]
This sum is zero for $f$ of the form $g^p-g$. This suggests modifying 
the sum $S_f'$ as follows: 
\begin{enumerate}[\indent a)]
  \item for any closed point $x$ of $X_0$, set $v_x(f) = $ the order of the 
    pole of $f$ at $x$ if $f(x)=\infty$, and $v_x(f)=0$ otherwise; also set 
    $v_x^\ast(f) = \inf v_x(f+g^p-g)$ (the infimum being over $g$). 
  \item If $v_x^\ast(f)=0$, $x\in X_0(\dF_q)$, and $f+g^p-g$ is 
    regular at $x$, put $\psi(f(x)) = \psi((f+g^p-g)(x))$. Finally set 
    \begin{equation*}\tag{3.5.1}\label{VI:eq:3-5-1}
      S_f = \sum'_{x\in X_0(\dF_q)} \psi(f(x)) \text{,}
    \end{equation*}
    where $(-)'$ indicates that the sum runs over those $x$ for which 
    $v_x^\ast(f)=0$. 
\end{enumerate}

We have $S_f=S_{f+g^p-g}$. We shall verify that if $f$ is not of the 
form $g^p-g+c^{t e}$, then 
\begin{equation*}\tag{3.5.2}\label{VI:eq:3-5-2}
  |S_f|\leqslant \left(2 g-2+\sum_{v_x^\ast(f)\ne 0} [k(x):\dF_q] (1+v_x^\ast(f))\right) q^{1/2} \text{.}
\end{equation*}

We begin by passing from complex coefficients to $\ell$-adic coefficients, replacing $\psi$ by 
a character of the form $\psi_0\circ \tr_{\dF_q/\dF_p}$, with 
$\psi_0:\dF_p\hookrightarrow E_\lambda^\times$. For $j:U_0\hookrightarrow X_0$ 
the inclusion of the open where $f\ne\infty$, we then prove that 
\begin{equation*}\tag{3.5.3}\label{VI:eq:3-5-3}
  S_f = \sum (-1)^i \tr(F^\ast,\h^i(X,j_\ast \sF(\psi f))) \text{.} 
\end{equation*}

The trace formula reduces this statement to the following (for 
$x\in X_0(\dF_q)$): 
\begin{enumerate}[\indent a)]
  \item If $v_x^\ast(f)=0$, then $F_x^\ast$ on $j_\ast \sF(\psi f)$ is 
    $\psi(f(x))$; 
  \item If $v_x^\ast(f)\ne 0$, then $j_\ast \sF(\psi f)$ is ramified at $x$: 
    $(j_\ast \sF(\psi f))_{\bar x}=0$. 
\end{enumerate}

Formula a) follows from \eqref{VI:eq:1-7-6} if $v_x(f)=0$, and the general 
case reduces to this one by replacing $f$ by $f+g^p-g$: up to isomorphism, this 
does not change $\sF(\psi f)_0$ on the open where $f$ and $g$ are regular 
(\ref{VI:1-3}). 

Formula b) follows from 
\begin{equation*}\tag{3.5.4}\label{VI:eq:3-5-4}
  \swan_x \sF(f,\psi) = v_x^\ast(f) \text{.} 
\end{equation*}
After reduction to the case $v_x(f)=v_x^\ast(f)$ (hence $p\nmid v_x(f)$) 
and extension of scalars to $\dF$, this formula can be found in 
\cite[4.4]{se61}.

Finally, we check that the sheaf $\sF(\psi f)$ on $U$ is nonconstant if 
$f$ is not of the form $g^p-g+c^{t e}$: since $\psi_0$ is injective, the 
triviality of $\sF(\psi f)=\sF(\psi_0 f)$ (\ref{VI:1-8}.ii) is equivalent to 
that of the $\dF_p$-torsor on $U$ defined by $T^p-T-f=0$. If it were trivial, 
the torsor on $U_0$ defined by the same equation would be the inverse image of an 
$\dF_q$-torsor on $\spec(\dF_q)$, with equation $T^p-T-\lambda=0$, and the 
$\dF_p$-torsor equation $T^p-T-(f-\lambda)=0$ on $U_0$ would be trivial 
and $f$ would therefore be of the form $g^p-g+\lambda$. The 
$\h^i(\bar X,j_\ast \sF(\psi f))$ therefore vanish for $i\ne 1$, and formula 
\eqref{VI:eq:3-5-3} reduces to 
\begin{equation*}\tag{3.5.5}\label{VI:eq:3-5-5}
  S_f = -\tr(F^\ast,\h^1(\bar X,j_\ast \sF(\psi f))) 
\end{equation*}
and, by \eqref{VI:eq:3-2-1} and \eqref{VI:eq:3-5-4}, the dimension of $\h^1$ is 
$2 g-2+\sum_{v_x^\ast(f)>0} [k(x):\dF_q](1+v_x^\ast(f))$; moreover, $-S_f$ 
is the sum of this many eigenvalues of $F^\ast$, each having complex 
absolute value $q^{1/2}$. 





\subsection{}\label{VI:3-6}

Suppose that, for an automorphism $\sigma$ of $X_0$, we have 
\[
  f(\sigma x) = - f(x) \text{.} 
\]
The sum $S_f$ is then real. If $\sigma$ is an involution, Poincaré 
duality allows us to say a little more: set 
$\sF_0=j_\ast \sF(\psi f)$, $\sG_0 = j_\ast \sF(\psi(-f))$, we have 
\begin{enumerate}[\indent a)]
  \item $\sF_0$ and $\sG_0$ are in duality, 
  \item $\sigma^\ast \sF_0 \iso \sG_0$ and $\sigma^\ast \sG_0 \iso \sF_0$, for 
    natural isomorphisms such that the composite 
    $\sF_0=(\sigma^2)^\ast \sF_0 \iso \sigma^\ast \sG_0 \iso \sF_0$ 
    is the identity. 
  \item the pairing $\sF_0\otimes \sG_0 \to E_\lambda$ satisfies 
    $\sigma^\ast(f\cdot g) = \sigma^\ast(f) \cdot \sigma^\ast (g)$. 
\end{enumerate} 

Passing to cohomology, we find that $\h^1(X,\sF)$ and $\h^1(X,\sG)$ are in 
perfect duality with values in $E_\lambda(1)$, and that the bilinear form 
$\alpha\cdot \sigma^\ast\beta$ on $\h^1(X,\sF)$ is \emph{alternating}: 
$\sigma^\ast$ acts trivially on $E_\lambda(1)$, and 
\[
  \alpha\cdot \sigma^\ast \beta = \sigma^\ast(\alpha\cdot \sigma^\ast \beta) = \sigma^\ast \alpha\cdot \beta = -\beta\cdot \sigma^\ast \alpha \text{.} 
\]
We conclude that $\h^1(X,\sF)$ is of even dimension (we also check 
    easily that each nonzero $v_x^\ast(f)$ is odd) and that the 
eigenvalues of $F^\ast$ on $\h^1(X,\sF)$ occur in pairs $\alpha$ and 
$q/\alpha$. 





\subsection{Example}\label{VI:3-7}

This applies to the Kloosterman sums 
$\sum_{x\in \dF_q^\times} \psi\left(x+\frac a x\right)$ (take 
$X_0=\dP^1$, $f=x+\frac a x$, $\sigma x=-x$; the poles of $f$ are $0$ and 
$\infty$, and at each of them $v_x^\ast(f)=1$; the dimension of $\h^1$ is 
$-2+2+2=2$). We therefore have (for $a\ne 0$) 
\[
  \sum_{\substack{x y=a \\ x,y\in \dF_{q^n}}} \exp\left(\frac{2\pi i}{p} \tr_{\dF_{q^n}/\dF_p}(x+y)\right) = -(\alpha^n+\bar\alpha^n) \qquad \text{, }\alpha\bar \alpha = q \text{.} 
\]
This result is due to L.\ Carlitz \cite{ca69}. 

We now give an application of a Lang--Weil method. 





\begin{proposition_}\label{VI:3-8}
Let $P\in \dF_q[X_1,\dots,X_n]$ be a polynomial in $n$ variables, of degree 
$d$, which is not of the form $Q^p-Q+C^{te}$. Again setting 
$\psi(x) = \exp\left(\frac{2\pi i}{p} \tr_{\dF_q/\dF_p}(x)\right)$, we have 
\[
  \left|\sum_{x_i\in \dF_q} \psi P(x_1,\dots,x_n)\right| \leqslant (d+1) q^{n-\frac 1 2} \text{.}
\]
\end{proposition_}

For the left-hand side we have the trivial estimate $q^n$. It is therefore enough to 
prove \ref{VI:3-8} when $d-1<q^{1/2}$; assume therefore that 
$d<q+1$, and let $P_d$ be the homogeneous part of degree $d$ of $P$. Recall: 





\begin{lemma_}\label{VI:3-9}
A hypersurface of degree $d$ in $\dP^r$ ($r\geqslant 1$) cannot contain 
all rational points over $\dF_q$ unless $d\geqslant q+1$. 
\end{lemma_}

We argue by induction: if there exists a rational hyperplane not 
entirely contained in the hypersurface (this is not the case for $r=1$), we 
apply the induction hypothesis to the intersection of the hypersurface with this 
hyperplane. Otherwise, the degree $d$ is $\geqslant$ the number of rational 
hyperplanes, $\geqslant q+1$. 

Let us now distinguish two cases. 


\paragraph{Case 1: $p\nmid d$.}
Applying the lemma to $P_d$, we see that, after making a linear change 
of variables if necessary, we may assume that $P_d(1,0,\ldots,0)\ne 0$, i.e. 
that the coefficient of $X_1^d$ in $P$ is nonzero. A polynomial in $1$ 
variable 
\[
  S(X) = \sum_{i=0}^d a_i x^i \text{,} 
\]
with $a_d\ne 0$ ($p\nmid d$) is never of the form $Q^p-Q+C^{t e}$, and 
the \eqref{VI:eq:3-5-2} estimate is reduced to 
\[
  \left| \sum \psi(S(x)) \right| \leqslant (d-1) q^{1/2} \text{.} 
\]
Applying this estimate to the partial sums obtained by varying 
only $x_1$, we find 
\[
  \left| \sum \psi P(x_1,\dots,x_n) \right| = \left| \sum_{x_2,\dots,x_n} \sum_{x_1} \psi P(x_1,\dots,x_n)\right| \leqslant q^{n-1} (d-1) q^{1/2} 
\]
as promised. 


\paragraph{Case 2: $p\mid d$ ($d>0$).}
If $P_d$ is a $p$-th power, replacing $P$ by $P-(P_d-P_d^{1/p})$ 
does not change the sum under consideration and lowers the degree: this case 
is disposed of by induction on $d$. Otherwise, the differential of 
$P_d$ is not identically zero; applying \ref{VI:3-9}, one can assume, 
after a linear change of variables if necessary, that it is nonzero 
at the point $(1,0,\dots,0)$. If we write 
\[
  P_d = \sum_{i=0}^d X_1^{d-i} S_i(X_2,\dots,X_n) \text{,} 
\]
this means that the linear form $S_1$ is not identically zero. 

The form $S_0$ is constant, and by replacing $P$ with 
$P-(S_0 X_1^d -S_0^{1/p} X_1^{d/p})$, we can assume that it is zero. 
Finally, denote by $-\lambda$ the coefficient of $X_1^{d-1}$ in $P$. For fixed 
$x_2,\dots,x_n$, we have 
\[
  P(X_1,x_2,\dots,x_n) = (S_1(x_2,\dots,x_n)-\lambda) X_1^{d-1} + \text{terms of lower degree in $X_1$} 
\]
and if $S_1(x_2,\dots,x_n)\ne \lambda$, we therefore have 
\[
  \left| \sum_{x_1} \psi P(x_1,\dots,x_n)\right| \leqslant (d-2) q^{1/2} \text{.} 
\]
In total, 
\begin{align*}
  \left|\sum \psi P(x_1,\dots,x_n)\right| 
    &\leqslant \left|\sum_{S_1=\lambda} \psi P(x_1,\dots,x_n)\right| + \sum_{S_1(X_2,\dots,X_n)\ne\lambda} \left| \sum_{x_1} \psi P(x_1,\dots,x_n)\right| \\
    &\leqslant q^{n-1} + (q^{n-1}-q^{n-2})(d-2) q^{1/2} \\
    &< (d-2) q^{n-1/2} + q^{n-1} \\
    &< (d-1) q^{n-1/2} \text{.}
\end{align*}

Another result of this nature is given by R.A.\ Smith 
\cite{sm70}. 










\section{Gauss and Jacobi sums}\label{VI:4}





\subsection{}\label{VI:4-1}

Let $k$ be a finite field of characteristic $p$, let $\chi$ be a character of 
$k^\times$, and let $\psi$ be a nontrivial character of the additive group of $k$. 
We take the following as the definition of Gauss sums: 
\begin{equation*}\tag{4.1.1}\label{VI:eq:4-1-1}
  \tau(\chi,\psi) = -\sum_{x\in k^\times} \psi(x) \chi^{-1}(x) 
\end{equation*}
(note the sign). Classically, $\chi$ and $\psi$ are complex-valued. 
We take them to have values in $E_\lambda^\times$ (see the proof of 
\ref{VI:1-15}). Regard $-\tau(\chi,\psi)$ as a sum over the rational 
points of the scheme $\dG_m$ over $k$. By Section \ref{VI:1}, 
we have $-\tau(\chi,\psi) = \tr(F^\ast,\h_c^\ast(\dG_m,\sF(\psi \chi^{-1}))$. Moreover,

\begin{proposition_}\label{VI:4-2}
The cohomology of $\dG_m$ with coefficients in $\sF(\psi\chi^{-1})$ satisfies 
\begin{enumerate}[(i)]
  \item $\h_c^i=0$ for $i\ne 1$, and $\dim \h_c^1=1$. 
  \item $F^\ast$, acting on $\h_c^1$, is multiplication by 
    $\tau(\chi,\psi)$. 
  \item If $\chi$ is nontrivial, we have $\h_c^\bullet \iso \h^\bullet$. 
\end{enumerate}
\end{proposition_}
\begin{proof}[Proof]
For any $n$ prime to $p$, denote by $K_n$ the $\dmu_n$-torsor on $\dG_m$ 
defined by the exact K\"ummer sequence 
$0 \to \dmu_n \to \dG_m \xrightarrow{x^n} \dG_m \to 0$. If $k$ has $q$ 
elements, then over $k$ one has $\dmu_{q-1}=k^\times$, and the Lang torsor on 
$\dG_m/k$ is $K_{q-1}$. Hence $\sF(\chi^{-1})=\chi(K_{q-1})$. 
Assertion \ref{VI:4-2}(ii) follows from (i,iii), which in turn follow from 
the following geometric statement, where $\sF(\psi)$ denotes the sheaf 
on $\dG_m/\dF$ obtained by extending scalars from $k$ to $\dF$ in 
$\sF(\psi)$ on $\dG_m/k$ (\ref{VI:1-8}.b). 
\end{proof}





\begin{proposition_}\label{VI:4-3}
Let $\chi:\dmu_n \to E_\lambda^\times$. The cohomology of $\dG_m$ with 
coefficients in $\sF(\psi)\otimes \chi(K_n)$ satisfies 
\begin{enumerate}[\indent (i)]
  \item $\h_c^i=0$ for $i\ne 1$, and $\dim \h_c^1{} = 1$. 
  \item If $\chi$ is nontrivial, we have $\h_c^\bullet{}\iso \h^\bullet{}$. 
\end{enumerate}
\end{proposition_}
\begin{proof}[Proof]
The sheaf $\sF(\psi)$ is the restriction to $\dG_m$ of a sheaf locally 
constant on $\dG_a$, wildly ramified at infinity, with Swan conductor 
$1$. The sheaf $\chi(K_n)$ is constant if $\chi=1$; if $\chi\ne 1$, it 
is tamely ramified at $0$ and $\infty$. 

The sheaf $\sF(\psi)\otimes\chi(K_n)$ is therefore ramified at $\infty$; 
applying \ref{VI:1-18}(b,c), we find that 
$\h_c^i(\dG_m,\sF(\psi)\otimes \chi(K_n))=0$ for $i\ne 1$. If $\chi\ne 1$, it 
is ramified at $0$ and $\infty$, and (ii) follows from \ref{VI:1-19}.a. 

The Swan conductors are $0$ at $0$ and $1$ at $\infty$. By 
\eqref{VI:eq:3-2-1}, the Euler--Poincaré characteristic is therefore $-1$, 
which completes the proof. 
\end{proof}





\subsection{Remark}\label{VI:4-4}

If $\chi$ is nontrivial, \ref{VI:4-2}(iii) and Poincaré duality 
show that $\h_c^1(\dG_m,\sF(\psi\chi^{-1}))$ and 
$\h_c^1(\dG_m,\sF(\psi^{-1}\chi))$ are in duality (duality with values 
in $E_\lambda(-1)$). We therefore have 
$\tau(\chi,\psi)\cdot \tau(\chi^{-1},\psi^{-1}) = q$, i.e. 
$|\tau(\chi,\psi)|=q^{1/2}$. 





\subsection{}\label{VI:4-5}

If $k$ is an extension of degree $N$ of $\dF_q$, we can also regard 
\eqref{VI:eq:4-1-1} as a sum in $N$ variables over $\dF_q$. More 
generally, let $k$ be an étale algebra over $\dF_q$, of degree $N$ over 
$\dF_q$. It is a product of fields $k_i$ of degrees $N_i$, and we set 
\begin{equation*}\tag{4.5.1}\label{VI:eq:4-5-1}
  \varepsilon(k) = (-1)^{\sum (N_i+1)} \text{.} 
\end{equation*}
This is the sign of the permutation of $S=\hom_{\dF_q}(k,\dF)$ induced by the 
Frobenius substitution $\varphi\in \gal(\dF/\dF_q)$. 

Let $\psi:\dF_q\to E_\lambda^\times$ be a nontrivial character, and let 
$\chi:k^\times \to E_\lambda^\times$. We set 
\begin{equation*}\tag{4.5.2}\label{VI:eq:4-5-2}
  \tau_{\dF_q}(\chi,\psi) = (-1)^N \sum_{x\in k^\times} \psi \tr_{k/\dF_q}(x) \cdot \chi^{-1}(x) \text{.} 
\end{equation*}

If $\chi$ has coordinates $\chi_i:k_i^\times \to E_\lambda^\times$, we 
have the elementary identity 
\begin{equation*}\tag{4.5.3}\label{VI:eq:4-5-3}
  \tau_{\dF_q}(\chi,\psi) = \varepsilon(k) \prod \tau(\chi_i,\psi\circ \tr_{k_i/\dF_q}) \text{.}
\end{equation*}

After some preliminaries, we give in \ref{VI:4-10} a 
cohomological interpretation of the sums \eqref{VI:eq:4-5-2}. In 
\ref{VI:4-12}, we will interpret the Hasse--Davenport identity as a 
twisted form (cf. \ref{VI:1-12}) of the following special case of 
\eqref{VI:eq:4-5-3}: for $k=\dF_q^N$, and $\chi$ a character of 
$\dF_q^\times$, we have 
$\tau_{\dF_q}(\chi\circ N_{k/\dF_q},\psi)=\tau(\chi,\psi)^N$. 





\subsection{}\label{VI:4-6}

Recall that for $M$ a finitely generated projective module over a ring $A$, the 
functor $\spec(B)\mapsto M\otimes_A B$ from affine schemes over $\spec(A)$ 
to $\mathsf{Ens}$ is represented by the affine scheme $\dV(M^\vee)$ 
(EGA notation), the spectrum of $\operatorname{Sym}_A^\bullet(M^\vee)$. If 
$M=A^n$, this is the standard affine space of dimension $n$. 

Suppose that $M$ is a unital $A$-algebra, and let 
$V=\dV(M^\vee)$. By definition, for any $B$ equipped with an $A$-algebra structure, the set 
$V(B)$ of points of $V$ with coordinates in $B$ is $M\otimes_A B$. Since the 
functor $B\mapsto M\otimes_A B$ takes values in unital rings, 
$V$ is a unital ring scheme over $\spec(A)$. The 
norm and trace morphisms $M\otimes_A B\to B$ are functorial in $B$; they 
therefore correspond to morphisms of schemes $N$ and $T$ from $V$ to 
$\dG_a$. We shall consider the following schemes obtained from $V$: 
\begin{enumerate}[a)]
  \item $V^\ast$ is the open subscheme of invertible elements of $V$ (a group 
    scheme under $\cdot$);
  \item $W$ is the hyperplane of equation $T=0$ and $W^\ast=W\cap V^\ast$; 
  \item $P$ is the hyperplane at infinity $V\setminus \{0\}/\dG_m$ of the affine 
    space $V$ over $\spec(A)$. If $Q$ is the hyperplane at infinity of $W$, 
    $W^\ast/\dG_m$ and $V^\ast/\dG_m$ are open subsets of the projective spaces 
    $Q$ and $P$ over $\spec(A)$. 
    \begin{equation*}\tag{4.6.1}\label{VI:eq:4-6-1}
    \xymatrix{
      W^\ast \ar@{^{(}->}[r] \ar[d]^-\pi 
        & V^\ast \ar[d]^-\pi \\
      W^\ast/\dG_m \ar@{^{(}->}[r] 
        & V^\ast/\dG_m 
    }\qquad\subset\qquad
    \xymatrix{
      W\setminus \{0\} \ar@{^{(}->}[r] \ar[d]^-\pi 
        & V\setminus \{0\} \ar[d]^-\pi \\
      Q \ar@{^{(}->}[r] 
        & P \text{.}
    }
    \end{equation*}
\end{enumerate}

If $M=A^I$, $I$ a finite set, then $V\simeq \dG_a^I$, $V^\ast=\dG_m^I$, 
$V^\ast/\dG_m$ is the torus $\dG_m^I/(\dG_m\text{ diagonal})$, $N$ and $T$ 
are given by $\prod x_i$ and $\sum x_i$, and $Q$ is therefore the projective 
hyperplane with equation $\sum x_i=0$ in $P$. 

The case of interest is that in which $M$ is finite étale over $A$. The 
previous description then applies locally to $\spec(A)$ (for the 
étale topology) and we can write $V^\ast=\dG_m^I$, for $I$ a locally 
constant sheaf of finite sets on $\spec(A)$. If, for example, $A$ is a 
field with algebraic closure $\bar A$, and $M$ is a separable $A$-algebra, 
we put $I=\hom_A(M,\bar A)$, and over $\bar A$ we have 
$V\sim \dG_a^I$. Via this isomorphism, $N$ and $T$ are given by $\prod x_i$ and 
$\sum x_i$. 





\subsection{The Kummer torsor}\label{VI:4-7}

Let $S$ be a scheme, let $n$ be an integer invertible on $S$, and let $G$ be a 
commutative group scheme with connected fibers (written multiplicatively) over $S$. 
The sequence 
\[\xymatrix{
  0 \ar[r] 
    & G_n \ar[r] 
    & G \ar[r]^-{x^n} 
    & G \ar[r] 
    & 0 
}\]
is exact. It defines a $G_n$-torsor $K_n(G)$ (or simply $K_n$) on 
$G$. For $\chi:G_n \to E_\lambda^\times$ a homomorphism from the étale 
$S$-sheaf $G_n$ to the constant sheaf $E_\lambda^\times$, we denote by 
$\sK_n(\chi)$ the $E_\lambda$-sheaf $\chi^{-1}(K_n)$. 

The $K_n$ torsors form an inverse system of torsors under the 
projective system of groups $G_n$ (transition morphisms 
$x\mapsto x^d:G_{n d}\to G_n$). This expresses the commutativity of the diagrams 
\[\xymatrix{
  0 \ar[r] 
    & G_{n d} \ar[r] \ar[d]^-{x^d} 
    & G \ar[r]^-{x^{n d}} \ar[d]^-{x^d} 
    & G \ar[r] \ar@{=}[d] 
    & 0 \\
  0 \ar[r] 
    & G_n \ar[r] 
    & G \ar[r]^-{x^n} 
    & G \ar[r] 
    & 0 \text{.} 
}\]
We therefore have $\sK_n(\chi) = \sK_{n d}(\chi\circ x^d)$. 





\subsection{}\label{VI:4-8}

Apply the construction of \ref{VI:4-6} with $A=\dF_q$ and with $M=k$ an 
étale algebra over $\dF_q$. In accordance with our general conventions, we denote 
by a subscript $0$ the $\dF_q$-schemes denoted $V,V^\ast,\ldots$ in 
\ref{VI:4-6}. The same letters without a subscript denote the schemes over 
$\dF$ obtained by extension of scalars. 

We put $I=\hom(k,\dF)$; then $V\sim \dG_a^I$ and $T$ is given by 
$(x_i)\mapsto \sum x_i$. 
\begin{equation*}\tag{4.8.1}\label{VI:eq:4-8-1}
  \sF(\psi\circ \tr_{k/\dF_q}) = T^\ast \sF(\psi) = \bigotimes_i \pr_i^\ast \sF(\psi) \text{.} 
\end{equation*}
Since $V^\ast\sim \dG_m^I$, a character $\chi$ of $V_n^\ast\sim \dmu_n^I$, 
with values in $E_\lambda^\times$, amounts to a family 
$(\chi_i)_{i\in I}$ of characters of $\dmu_n$ indexed by $I$. It is 
defined over $\dF_q$ if, for all $\sigma\in \gal(\dF/\dF_q)$, we have 
$\chi_{\sigma_i} = \chi_i\circ \sigma^{-1}$. It then defines a 
$E_\lambda$-sheaf $\sK_n((\chi_i)_{i\in I})$ on $V_0^\ast$. If the product 
of the $\chi_i$ is trivial, $\chi$ factors through a character of 
$(V_0^\ast/\dG_m)_n$ and $\sK_n((\chi_i)_{i\in I})$ is the inverse image 
of a sheaf, denoted in the same way, on $V_0^\ast/\dG_m$. This construction 
is related to the Lang torsor as follows. 





\begin{lemma_}\label{VI:4-9}
\begin{enumerate}[(i)]
  \item For $n$ sufficiently divisible, we have a commutative diagram 
    \[\xymatrix{
      0 \ar[r] 
        & (V_0^\ast)_n \ar[r] \ar[d]^-\tau 
        & V_0^\ast \ar[r]^-{x^n} \ar[d]^-\tau 
        & V_0^\ast \ar[r] \ar@{=}[d] 
        & 0 \\
      0 \ar[r] 
        & k^\times \ar[r] 
        & V_0^\ast \ar[r]^-\fL 
        & V_0^\ast \ar[r] 
        & 0 \text{;} 
    }\]
    if $\chi$ is a character $\chi:k^\times \to E_\lambda^\times$, we have 
    $\sF(\chi) = \sK_n(\chi\circ\tau)$. 
  \item For sufficiently divisible $n$, $\tau$ identifies $k^\times$ with the coinvariants 
    of $\gal(\dF/\dF_q)$ acting on $V_n^\ast\sim \dmu_n^I$. 
  \item For $k$ a field, $N=[k:\dF_q]$, $\omega\in I$ an embedding of $k$ 
    in $\dF$, and $n=q^N-1$, the index component $\omega$ of 
    $\chi\circ\tau$ is $\chi\circ \omega^{-1}$. 
  \item If $\chi$ is nontrivial on each factor of $k$, the 
    $(\chi\circ \tau)_i$ are all nontrivial. 
\end{enumerate}
\end{lemma_}

It is enough to prove the lemma when $k$ is a field of degree $N$ over 
$\dF_q$. In this case, $n$ is ``sufficiently divisible'' if $q^N-1\mid n$. 
Choose an embedding $\omega$ of $k$ into $\dF$; then $I$ is identified 
with $\dZ/N$: the element $i\in \dZ/N$ corresponds to $\omega_i=\omega^{q^i}$. Via 
the isomorphism $V_0^\ast(\dF) = {\dF^\times}^I$, we have 
\begin{enumerate}[\indent a)]
  \item $x\in k^\times$ corresponds to $(\omega_i(x))\in {\dF^\times}^I$; 
  \item $F((x_i)_{i\in \dZ/N})=(x_{i-1}^q)_{i\in \dZ/N}$. 
\end{enumerate}

A character $\chi=(\chi_i)$ of $V_n^\ast\sim \dmu_n^I$ is defined over 
$\dF_q$ if $\chi_{-i}(x) =\chi_0(x^{q^i})$ ($i\in \dZ$). If $q^N-1\mid n$, there are 
$q^N-1$ such characters: $\chi_0$ factors through $\dmu_{q^N-1}$ and 
determines the $\chi_i$. To prove (ii), it is therefore sufficient to verify (i) 
and (iii) (or its corollary (iv)) which ensures that $\chi\mapsto \chi\circ \tau$ 
is injective. 

Let us prove (i) and (iii) for $n=q^N-1$. Write the group of 
endomorphisms of $V^\ast\sim \dG_m^I$ additively; in particular, denote by $n$ 
the $x\mapsto x^n$ operator. If $\alpha$ is the cyclic permutation operator 
$(x_i)\mapsto (x_{i-1})$, we have 
$q^N-1=(q\alpha)^N-1 = (q\alpha-1)((q\alpha)^{N-1}+\cdots + 1)$. This 
determines $\tau$: we have $\tau((x_i)) = \prod_{0\leqslant j<N} x_{i-j}^{q^j}$, 
and the induced map from $(V_0^\ast)_n$ to $k^\times$ is [as in a) above] 
$(x_i)\mapsto \prod x_{-i}^{q^i}$; hence (iii) follows. 





\begin{proposition_}\label{VI:4-10}
The cohomology of $V^\ast$ with coefficients in 
$\sF(\chi^{-1}\cdot \psi\tr_{k/\dF_q})$ satisfies 
\begin{enumerate}[\indent (i)]
  \item $\h_c^i = 0$ for $i\ne N$, and $\dim \h_c^N=1$. 
  \item On $\h_c^N$, $F^\ast$ is multiplication by 
    $\tau_{\dF_q}(\chi,\psi)$. 
  \item If $\chi$ is nontrivial on each factor of $k$, we have 
    $\h_c^\bullet \iso \h^\bullet$. 
\end{enumerate}
\end{proposition_}

Applying \ref{VI:4-9}(i), over $\dF$, if $\chi\circ \tau=(\chi_i)_{i\in I}$, 
we have $\sF(\chi) = \sK_n((\chi_i)_{i\in I})$. Points (i) and (iii) 
therefore follow from the following more geometric statement (for (iii), 
apply \ref{VI:4-9}(iv)), and (ii) follows from the trace formula: 





\begin{proposition_}\label{VI:4-11}
Let $(\chi_i)_{i\in I}$ be a family of characters of $\dmu_n(\dF)$. The 
cohomology of $V^\ast\simeq \dG_m^I$ with coefficients in 
$\sK_n((\chi_i)_{i\in I})\otimes \sF(\psi\tr_{k/\dF_q})$ satisfies 
\begin{enumerate}[\indent (i)]
  \item $\h_c^i{}=0$ for $i\ne N$, and $\dim \h_c^N{} = 1$. 
  \item If the $\chi_i$ are all nontrivial, we have 
    $\h_c^\bullet{}\iso \h^\bullet$. 
\end{enumerate}
\end{proposition_}

We have $V^\ast\sim \dG_m^I$, 
$\sK((\chi_i)) = \bigotimes_i \pr_i^\ast \chi_i(K_n(\dG_m))$ and 
$\sF(\psi\circ \tr_{k/\dF_q})=T^\ast \sF(\psi) = \bigotimes_i \pr_i^\ast \sF(\psi)$. 
The K\"unneth formula then applies, and \ref{VI:4-11} follows from 
\ref{VI:4-3}. 





\subsection{}\label{VI:4-12}

From \ref{VI:2-4*}* and (\hyperref[IV]{Cycle}, \ref{IV:1-3} Example 2), it also follows 
that the group of permutations $\sigma$ of $I$ such that 
$\chi_i = \chi_{\sigma i}$ ($i\in I$) acts on $\h_c^N$ by multiplication by 
the sign $\varepsilon(\sigma)$. We deduce a second proof 
of the identity of Hasse-Davenport. 

If $\chi$ is a character of $\dF_q^\times$, the fact that, over $\dF$, $N$ 
is given by $(x_i)\mapsto \prod x_i$ gives 
$\sF(\chi\circ N)=N^\ast\sF(\chi) = \bigotimes_i \pr_i^\ast\sF(\chi)$ 
and the K\"unneth formula gives 
\[
  \h_c^\bullet(V^\ast,\sF(\chi^{-1}\circ N,\psi\circ \tr)) \sim \h_c^\bullet(\dG_m,\sF(\chi^{-1}\psi))^{\otimes I} \text{,} 
\]
where the tensor product on the right-hand side is taken in the sense of 
\ref{VI:2-4*}*. Since $\h_c^1(\dG_m,\sF(\psi\chi^{-1}))$ is of dimension 
$1$, its tensor power $\otimes I$, in the ordinary sense, depends only on 
the cardinality $n$ of $I$. Applying (\hyperref[IV]{Cycle}, \ref{IV:1-3} Example 
2), we obtain a canonical isomorphism 
\begin{equation*}\tag{4.12.1}\label{VI:eq:4-12-1}
  \h_c^N\left(V^\ast,\sF(\chi^{-1}\circ N)\otimes T^\ast \sF(\psi)\right) \sim \h_c^1(\dG_m,\sF(\psi\chi^{-1})) \otimes_\dZ \textstyle \bigwedge^N \dZ^I \text{.}
\end{equation*}
To compute the action of $F^\ast$, it is most convenient to adopt the Galois 
point of view and observe that, since the isomorphism \eqref{VI:eq:4-12-1} is 
canonical, it is compatible with the transport-of-structure action of 
$\gal(\dF/\dF_q)$. If $\varepsilon(k)$ is the sign of the permutation 
$\varphi$ of $I$, we find that 
\begin{align*} 
  \tau_{\dF_q}(\chi\circ N_{k/\dF_q},\psi) 
    &= \tr\left(\varphi^{-1},\h_c^N\left(V^\ast,N^\ast\sF(\chi^{-1})\otimes T^\ast \sF(\psi)\right)\right) \\ 
    &= \varepsilon(k) \tr\left(\varphi^{-1},\h_c^1(\dG_m,\sF(\chi^{-1}\psi))\right)^N \\
    &= \varepsilon(k) \tau(\chi,\psi)^N \text{,} 
\end{align*}
that is, 
\begin{equation*}\tag{4.12.2}\label{VI:eq:4-12-2}
  \tau_{\dF_q}\left(\chi\circ N_{k/\dF_q},\psi\right) = \varepsilon(k) \tau(\chi,\psi)^N \text{.} 
\end{equation*}

If $k$ is a field, $\varphi$ is a cyclic permutation of $I$, 
$\varepsilon(k) = (-1)^{N+1}$, and we obtain the Hasse--Davenport identity: 
\[
  \tau\left(\chi\circ N_{k/\dF_q},\psi\circ \tr_{k/\dF_q}\right) = (-1)^{N+1} \tau_{\dF_q}\left(\chi\circ N_{k/\dF_q},\psi\right) = \tau(\chi,\psi)^N \text{.} 
\]





\begin{lemma_}\label{VI:4-13}
For $\chi\circ \tau=(\chi_i)_{i\in I}$ as in \ref{VI:4-9}(i), the
following conditions are equivalent 
\begin{enumerate}[\indent (i)]
  \item $\chi| \dF_q^\times$ is trivial. 
  \item The product of $\chi_i$ is trivial. 
  \item $\sF(\chi) = \sK_n((\chi_i))$ is the inverse image of a unique 
    sheaf on $V^\ast/\dG_m$. 
  \item The inverse image of $\sF(\chi)$ on $\dG_m$ (embedded in 
    $V^\ast$ by the structural morphism $\dF_q \to k$) is trivial. 
\end{enumerate}
\end{lemma_}

(i)$\Rightarrow$(ii). We regard $\chi$ as a character of 
$V^\ast/\dG_m(\dF_q)=k^\times/\dF_q^\times$, and we take $\sF(\chi)$ on 
$V^\ast/\dG_m$. 

(ii)$\Rightarrow$(iii). Likewise, we regard $(\chi_i)$ as a character 
of $(V^\ast/\dG_m)_n$. The uniqueness in (iii) follows from the fact that $V^\ast$ is 
a fibration with connected fibers over $V^\ast/\dG_m$, and (iii)$\Rightarrow$(iv)  
is trivial. 

(iv)$\Rightarrow$(i), (ii). This inverse image is 
$\sF(\chi|\dF_q^\times)$ and $\sK_n(\prod \chi_i)$. 





\subsection{}\label{VI:4-14}

Let $k$ be an étale algebra over $\dF_q$, of 
dimension $N+1$, and let $\chi$ be a nontrivial character of $k^\times$, trivial on 
$\dF_q^\times$. The Jacobi sum $J(\chi)$ is defined by 
\begin{equation*}\tag{4.14.1}\label{VI:eq:4-14-1}
  J(\chi) = (-1)^{N-1} \sum_{\substack{x\in k^\times/\dF_q^\times \\ \tr(x)=0}} \chi^{-1}(x) \text{.} 
\end{equation*}

We have the following identity between Gauss sums and Jacobi sums. 





\begin{proposition_}\label{VI:4-15}
For $\chi$ as above and $\psi$ a nontrivial additive character of 
$\dF_q$, we have 
\[
  q J(\chi) = \tau_{\dF_q}(\chi,\psi) \text{.} 
\]
\end{proposition_}

In the special case $k=\dF_q^n$, formula \eqref{VI:eq:4-5-3} rewrites this as 
\begin{equation*}\tag{4.15.1}\label{VI:eq:4-15-1}
  q J(\chi) = \prod_i \tau(\chi_i,\psi) \text{,} 
\end{equation*}
or equivalently as 
\begin{align*}
  \chi_0(-1) J(\chi) &= (-1)^{N-1} \sum_{\substack{x_1,\dots,x_N\in \dF_q^\times \\ \sum x_i = 1}} \prod_{i=1}^N \chi_i^{-1}(x_i) = \tau(\chi_0^{-1},\psi)^{-1} \prod_{i=1}^N \tau(\chi_i,\psi) \\
  \chi_0^{-1} &= \prod_{i=1}^N \chi_i 
\end{align*}





\begin{proof}[Proof of \ref{VI:4-15}]
\[
  \tau_{\dF_q}(\chi,\psi) = (-1)^{N+1} \sum_{x\in k^\times} \chi(x)^{-1} \psi \tr(x) = (-1)^{N+1} \sum_{x\in k^\times/\dF_q^\times} \chi(x)^{-1} \sum_{\lambda\in \dF_q^\times} \psi(\lambda \tr x) \text{.}
\]
The sum $\sum_{\lambda\in \dF_q^\times} \psi(\lambda \tr x)$ equals $q-1$ if 
$\tr(x)=0$, and $-1$ if $\tr(x)\ne 0$. Therefore, 
\[
  \tau_{\dF_q}(\chi,\psi) = (-1)^{N+1} \left(q\sum_{\substack{x\in k^\times/\dF_q^\times \\ \tr(x)=0}} \chi(x)^{-1} - \sum_{x\in k^\times/\dF_q^\times} \chi(x)^{-1}\right) \text{.}
\]
The second term on the right-hand side is zero, since it is the sum of the values 
of a character, and \ref{VI:4-15} follows. 
\end{proof}

Proposition \ref{VI:4-15} therefore has the following cohomological counterpart: 





\begin{proposition_}\label{VI:4-16}
The cohomology of $W^\ast/\dG_m$ (\ref{VI:4-6}, \ref{VI:4-8}) with coefficients 
in $\sF(\chi^{-1})$ satisfies 
\begin{enumerate}[\indent (i)]
  \item $\h_c^i=0$ for $i\ne N-1$, and $\dim \h_c^{N-1}{}=1$. 
  \item On $\h_c^{N-1}$, $F^\ast$ is multiplication by $J(\chi)$. 
  \item $\h_c^{N-1}$ is canonically isomorphic to 
    $\h_c^{N+1}\left(V^\ast,\sF(\chi^{-1},\psi\circ \tr_{k/\dF_q})\right)(1)$. 
\end{enumerate}
\end{proposition_}

Assertion (ii) follows from (i) and the trace formula; (i) and (iii) 
follow from \ref{VI:4-11} and the following more geometric statement. 





\begin{proposition_}\label{VI:4-17}
Let $(\chi_i)_{i\in I}$ be a family of characters of 
$\dmu_n(\dF)$ that are not all trivial and whose product is $1$. Then 
$\h_c^{i-1}\left(W^\ast/\dG_m,\sK_n((\chi_i)_{i\in I})\right)$ is 
canonically isomorphic to 
\newline
$\h_c^{i+1}\left(V^\ast,\sK_n((\chi_i)_{i\in I})\otimes \sF(\psi\tr_{k/\dF_q})\right)(1)$. 
\end{proposition_}

The first line of \ref{VI:4-15}, namely the projection formula, gives ($\pi$ as in 
\eqref{VI:eq:4-6-1}). 
\begin{equation*}\tag{4.17.1}\label{VI:eq:4-17-1}
  \eR^i \pi_!\left(\sK_n((\chi_i)_{i\in I})\otimes \sF(\psi \tr_{k/\dF_q})\right) = \sK_n((\chi_i)_{i\in I}) \otimes \eR^i \pi_! \sF(\psi\tr_{k/\dF_q}) 
\end{equation*}
(on $V^\ast/\dG_m$). Let us compute the sheaves 
$\eR^i\pi_! \sF(\psi\tr_{k/\dF_q}) = \eR^i \pi_! T^\ast \sF(\psi)$. Let 
$v_0:\widetilde V_0 \to V_0$ be the blow-up of $V_0$ at $\{0\}$. In the 
diagram 
\begin{equation*}\tag{4.17.2}\label{VI:eq:4-17-2}
\xymatrix{
  V_0\setminus \{0\} \ar@{^{(}->}[r] \ar[dr]_-\pi 
    & \widetilde V_0 \ar[r]^-{v_0} \ar[d]^-{\bar\pi} 
    & V_0 \\
  & P_0 
}
\end{equation*}
$\pi$ is a fibration whose fibers are punctured affine lines, $\widetilde V_0$ 
is the corresponding line bundle, and $Z_0=v_0^{-1}(0)$ is its section $0$. 
The $T^\ast\sF(\psi)$ sheaf on $V_0\setminus \{0\}$ extends to 
$(T v_0)^\ast \sF(\psi)$ on $\widetilde V_0$, and the restriction of this sheaf 
to $Z_0$ is the constant sheaf $E_\lambda$, because $T_0 v_0$ vanishes on 
$Z_0$. The long exact sequence for cohomology with compact supports is therefore 
\begin{equation*}\tag{4.17.3}\label{VI:eq:4-17-3}
\xymatrix{
  \ar[r] 
    & \eR^i \pi_! T^\ast \sF(\psi) \ar[r] 
    & \eR^i \bar\pi_! (T v_0)^\ast \sF(\psi) \ar[r] 
    & (E_\lambda\text{ for $i=0$, $0$ otherwise}) \ar[r] 
    & 
}
\end{equation*}

The inverse image $\bar\pi^{-1}(Q_0)$ is the blow-up $\widetilde W_0$ of 
$W_0$ at $0$; $T_0 v_0$ vanishes on $\widetilde W_0$; hence 
$(T v_0)^\ast \sF(\psi)=E_\lambda$ on $\bar\pi^{-1}(Q_0)$ and, since $\bar\pi$ is an 
affine-line fibration, 
\begin{equation*}\tag{4.17.4}\label{VI:eq:4-17-4}
  \eR^i \bar\pi_! (T v_0)^\ast \sF(\psi)|Q = 
    \begin{cases}
      0 & \text{for $i\ne 2$} \\
      E_\lambda(-1) & \text{for $i=2$.} 
    \end{cases} 
\end{equation*}

If $x\in P$, $x\notin Q$, the line $D=\pi^{-1}(x)$ is mapped 
isomorphically onto $\dG_a$ by $T v_0$; therefore (by 
\ref{VI:2-7}*) 
\begin{align} \notag 
  \h_c^\bullet(D,(T v_0)^\ast \sF(\psi)) &= \h_c^\bullet(\dG_a,\sF(\psi)) = 0 && \text{and} \\
  \tag{4.17.5}\label{VI:eq:4-17-5}
  \eR^i \bar\pi_! (T v_0)^\ast \sF(\psi) 
    &= \begin{cases} 
         0 & \text{for $i\ne 2$} \\
         E_\lambda(-1)_Q & \text{for $i=2$} 
       \end{cases} 
\end{align}

Combining \eqref{VI:eq:4-17-3} and \eqref{VI:eq:4-17-5}, we finally find 
\begin{equation*}\tag{4.17.6}\label{VI:eq:4-17-6}
  \eR^i \pi_! T^\ast\sF(\psi) = 
    \begin{cases}
      0 & \text{for $i\ne 1,2$} \\
      E_\lambda & \text{for $i=1$} \\
      E_\lambda(-1)_Q & \text{for $i=2$.} 
    \end{cases}
\end{equation*}





\subsection{}\label{VI:4-18}

Let us compute the compactly supported cohomology of the sheaf 
$\sK_n((\chi_i))\otimes T^\ast \sF(\psi)$ on $V^\ast$ using the 
Leray spectral sequence of $\pi:V^\ast\to V^\ast/\dG_m$. Applying 
\eqref{VI:eq:4-17-1} and \eqref{VI:eq:4-17-6}, we obtain as initial terms 
$E_2^{p,1} = \h_c^p(V^\ast/\dG_m,\sK_n(\chi_i)) = 0$ (because $(\chi_i)\ne 1$) 
and $E_2^{p,2} = \h_c^p(W^\ast/\dG_m,\sK(\chi_i))(-1)$. 

The spectral sequence reduces to an isomorphism, and \ref{VI:4-17} follows. 





\subsection{}\label{VI:4-19}

The sheaves $\sK_n((\chi_i)_{i\in I})$ make sense over any field, 
or even over any base scheme. This allows us to generalize 
\ref{VI:4-16}(i). With the notation of \ref{VI:4-6}, suppose that $M$ is finite 
étale of constant rank $N$ over $A$, and let $n$ be invertible in $A$. 
Let $a$ be the projection of $W^\ast/\dG_m$ to $\spec(A)$. Let also $\chi$ be 
a character (morphism of sheaves) 
$\chi:(V^\ast/\dG_m)_n \to E_\lambda^\times$, and $\sK_n(\chi)$ the sheaf 
corresponding. Locally for the étale topology, we can regard $\chi$ 
as a family of characters $(\chi_i)_{i\in I}$ of $\dmu_n$, whose product is 
trivial.

\begin{proposition_}\label{VI:4-20}
\begin{enumerate}[(i)]
  \item If $\chi$ is everywhere nontrivial, we have 
    $\eR^i a_! (\sK_n(\chi)) = 0$ for $i\ne N-2$, and 
    $\eR^{N-2} a_!(\sK_n(\chi))$ is smooth of rank $1$. 
  \item An automorphism $\sigma$ of $M$ which preserves $\chi$ acts on this 
    $\eR^{N-2} a_!$ by multiplication by the sign $\varepsilon(\sigma)$ 
    of $\sigma$, seen as a permutation of $I$. 
  \item If the $\chi_i$ are all nontrivial, we have $\eR a_! \iso \eR a_\ast$. 
\end{enumerate}
\end{proposition_}
\begin{proof}[Proof]
The scheme $W^\ast/\dG_m$ is the complement of a relative normal-crossings 
divisor in the scheme $Q$, which is proper and smooth over $\spec(A)$, and 
$\sK_n(\chi)$ is locally constant on $W^\ast/\dG_m$ and is tamely ramified 
along the boundary. Consequently, the $\eR^i a_!$ and 
$\eR^i a_\ast$ are smooth sheaves, compatible with every change of 
base. A standard argument then reduces us to the case in which $A$ is a 
finite field, and (i) follows from \ref{VI:4-17} and \ref{VI:4-11}(i). 
To prove (ii), use the compatibility of the action of $\sigma$ with 
the isomorphism \ref{VI:4-17}, together with \ref{VI:4-12}. For (iii), note that if the 
$\chi_i$ are all nontrivial, then $\sK_n(\chi)$ on 
$W^\ast/\dG_m\subset Q$ is ramified along every divisor at infinity, 
and apply \ref{VI:1-19-1}. 
\end{proof}










\section{Hecke characters}\label{VI:5}





\subsection{}\label{VI:5-1}

Let $F$ be a number field (of finite degree over $\dQ$) and let $k$ be a field of 
characteristic $0$. Here are several equivalent descriptions of an 
\emph{algebraic} homomorphism $\alpha:F^\times \to k^\times$. 
\begin{enumerate}[(i)]
  \item Let $\{e_i\}$ be a basis of $F$ over $\dQ$. The homomorphism $\alpha$ is 
    algebraic if it is given by a formula 
    $\alpha(\sum x^i e_i) = A(x^i)$, $A\in k(X^i)$. 
\end{enumerate}

This means that $\alpha$ agrees on $F^\times$ with a rational map defined over $k$ 
from $\res_{F/\dQ}(\dG_m)$ to $\dG_m$. By the Zariski density of 
$F^\times$, and because a rational homomorphism is defined everywhere, such 
a map is a homomorphism of group schemes: 
\begin{enumerate}[(i)]
\setcounter{enumi}{1}
  \item $\alpha$ is induced by a homomorphism of group schemes over $k$ 
    \[
      \res_{F/\dQ}(\dG_m)\otimes_\dQ k \to \dG_m \text{.} 
    \] 
\end{enumerate}

If $k$ is a number field, the adjunction property of restriction of 
scalars $R$ shows that this is equivalent to 
\begin{enumerate}[(i)]
\setcounter{enumi}{2}
  \item ($k$ a number field) $\alpha$ is induced by a homomorphism of 
    group schemes over $\dQ$: $\res_{F/\dQ}(\dG_m) \to \res_{k/\dQ}(\dG_m)$. 
\end{enumerate}

Let $\bar k$ be an algebraic closure of $k$, and let $I=\hom(F,\bar k)$. Over 
$\bar k$, the character group of $\res_{F/\dQ}(\dG_m)$ is $\dZ^I$, with basis 
the embeddings of $F$ into $\bar k$. The characters defined over $k$ 
are those invariant under $\gal(\bar k/k)$; they can be described either as 
those sending $F^\times$ into $k^\times\subset \bar k$, or in terms of 
orbits of $\gal(\bar k/k)$ on $I$, which in turn correspond to the factors of 
$F\otimes k$. 
\begin{enumerate}[(i)]
\setcounter{enumi}{3}
  \item $\alpha$ is of the form 
    $\alpha=\prod_{\omega\in I} \omega^{n_\omega}$. The allowed families of exponents 
    $(n_\omega)$ are those such that $n_\omega=n_{\omega'}$ for 
    $\omega$ and $\omega'$ in the same orbit of $\gal(\bar k/k)$. Equivalently, these 
    are those such that $\prod_\omega \omega(x)^{n_\omega}\in k$ for every 
    $x\in F$. 
  \item Write $F\otimes k=\prod_{j\in J} F_j$, where the $F_j$ are fields. 
    Then $\alpha$ can be written $\alpha=\prod N_{F_j/k}^{m_j}$. 
\end{enumerate}

In the special case where $k$ contains a normal closure of $F$, these 
expressions become: $\alpha=\prod \omega^{n_\omega}$, where $\omega$ 
runs over the $[F:\dQ]$ embeddings of $F$ into $k$. 





\subsection{}\label{VI:5-2}

Let $k$ be a number field, and let $\alpha$ be an algebraic homomorphism 
from $F^\times$ to $k^\times$. There is then a unique 
homomorphism, also denoted $\alpha$, of the group $I(F)$ 
of fractional ideals of $F$ to that of $k$, such that $\alpha((x))=(\alpha(x))$. 
Uniqueness follows from the fact that every ideal $a$ has a power that is 
principal, and that $I(k)$ is torsion-free. To prove 
existence, we use for example \ref{VI:5-1}(v): we have 
$\alpha(a) = \prod N_{F_j/k}^{m_j}((a))$. 





\subsection{}\label{VI:5-3}

Let us recall the definition of algebraic Hecke characters, also called 
Weil Hecke characters (or gr\"ossencharaktere) of type $A_0$. Let $F$ be 
a number field, let $\fm$ be an ideal of $F$ (i.e., of the ring of integers of 
$F$), let $I_\fm$ be the group of fractional ideals of $F$ prime to $\fm$, 
and let $k$ be a field of characteristic $0$. A homomorphism 
$\chi:I_\fm \to k^\times$ is an \emph{algebraic Hecke character} (of 
conductor $\leqslant \fm$) if there exists an algebraic homomorphism 
$\chi_\text{alg}:F^\times \to k^\times$ satisfying 
\begin{equation*}\tag{$*$}\label{VI:eq:5-3-1}
\text{For $x\in F^\times$, prime to $\fm$, totally positive, and 
$\equiv 1\pmod\fm$, one has $\chi((x))=\chi_\text{alg}(x)$.}
\end{equation*}

By the density of the set of $x$ in \eqref{VI:eq:5-3-1}, $\chi_\text{alg}$ 
is entirely determined by $\chi$. This is the \emph{algebraic part} 
of $\chi$. If $\chi((x))=\chi_\text{alg}(x)$ for $x$ totally positive and 
$\equiv 1\pmod{\fm'}$, $\chi$ extends to a Hecke character of 
conductor $\leqslant \inf(\fm,\fm')$: $I_{\fm+\fm'}\to k^\times$. We 
identify Hecke characters that coincide on their common domain of 
definition, and we call \emph{conductor} of $\chi$ the smallest $\fm'$ 
such that $\chi$ is of conductor $\leqslant \fm'$. 





\subsection{Remark}\label{VI:5-4}

If an algebraic Hecke character $\chi$ takes its values in a 
subfield $k'$ of $k$, it is already a Hecke character with values in 
$k'$: it suffices to observe that $\chi_\text{alg}:\res_{F/\dQ}(\dG_m) \to \dG_m$ 
is already defined over $k'$, and this follows from the Zariski density 
in $\res_{F/\dQ}(\dG_m)$ of the set of totally positive $x$ 
$\equiv 1\pmod\fm$. We can always take for $k'$ a subfield of $k$ of 
finite degree over $\dQ$. 





\subsection{}\label{VI:5-5}

If $\varepsilon$ is a totally positive unit $\equiv 1\pmod\fm$, we have 
$\chi_\text{alg}(\varepsilon) = \chi((\varepsilon)) = 1$. The homomorphism 
$\chi_\text{alg}$ therefore factors through the quotient $T_\fm$ of 
$\res_{F/\dQ}(\dG_m)$ by the Zariski closure of the group 
$E_\fm\subset F^\times$ of totally positive units $\equiv 1\pmod\fm$. 

According to Serre \cite[II.3]{se68}, if $k$ is an algebraic closure of 
$\dQ$, and $\fm$ is large enough, the characters $\prod \omega^{n_\omega}$ 
of $\res_{F/\dQ}(\dG_m)$ of the form $\chi_\text{alg}$ are characterized 
as follows: there must exist an integer $N$, the \emph{weight} of $\chi$ (or 
$\chi_\text{alg}$), such that for any element $\sigma$ of $\gal(k/\dQ)$ 
conjugate to complex conjugation, we have $n_\omega+n_{\sigma \omega}=N$. If 
$F_1$ is the largest subfield of $F$ which is a quadratic extension 
that is totally imaginary over a totally real field $F_1'$, this amounts to 
  saying that $\chi_\text{alg}$ is of the form $\chi_1\circ N_{F/F_1}$, and that 
  $\chi_1|F_1'=(N_{F_1'/\dQ})^N$. 

If $\chi$ is of weight $N$, for any ideal $\fa$ of $F$, $\chi(\fa)$ is a 
algebraic number all of whose conjugates have absolute value 
$N(\fa)^{N/2}$ \cite[II.3 prop.2]{se68}. 





\subsection{}\label{VI:5-6}

For a fixed number field $k$, additional conditions are 
imposed on $\chi_\text{alg}$. For any ideal $\fa$ of $F$ prime to the 
conductor, we have in fact 
\begin{equation*}\tag{5.6.1}\label{VI:eq:5-6-1}
  \chi_\text{alg}(\fa) = (\chi(\fa)) \text{,} 
\end{equation*}
so that $\chi_\text{alg}(\fa)$ is principal. Since the group of 
  ideal classes of $k$ is torsion, it is enough to prove the $n$-th power 
of \eqref{VI:eq:5-6-1}, for a suitable $n\ne 0$; this allows us 
to replace $\fa$ by $\fa^n$, therefore to assume $\fa=(x)$, with $x$ totally 
positive $\equiv 1\pmod\fm$. All that remains is to use the definitions. 

Together with \ref{VI:5-5}, this shows that $\chi_\text{alg}$ determines the 
valuations of $\chi(\fa)$ at all places of $k$. 





\subsection{}\label{VI:5-7}

Serre's group $S_\fm$ can be characterized as the 
group of multiplicative type whose group of characters (on any 
field) is the group of algebraic Hecke characters of conductor 
$\leqslant \fm$ (see \cite[II 2.1 et 2.2]{se68}). Its relation with 
$\ell$-adic representations is explained in \cite[II 2.3]{se68}.





\begin{theorem_}[{\cite{se68}}]\label{VI:5-8}
Let $\chi$ be an algebraic Hecke character of $F$ with values in $E_\lambda$, for 
$E_\lambda$ a finite extension of $\dQ_\ell$. There is then a unique 
  homomorphism $\chi_\lambda:\gal(\bar F/F)^\textnormal{ab} \to E_\lambda^\times$, such 
that 
\begin{enumerate}[\indent (i)]
  \item $\chi_\lambda$ is unramified outside the conductor $\ff$ of 
    $\chi$ and outside $\ell$. 
  \item For $\fp$ a prime ideal of $F$ prime to $\ff$ and to $\ell$, 
    and $F_\fp\in \gal(\bar F/F)^\textnormal{ab}$ the geometric Frobenius 
    at $\fp$, we have 
    \[
      \chi_\lambda(F_\fp) = \chi(\fp) \text{.}
    \] 
\end{enumerate}
\end{theorem_}





\subsection{}\label{VI:5-9}

At the end of this section, we give a criterion for a 
$\lambda$-adic representation to arise from a Hecke character, 
and in the next section we apply it to Jacobi sums. We 
thereby recover, in a somewhat generalized form, Weil's results 
\cite{we52,we74-2}, with the second part of Weil's proof 
replaced by a $\ell$-adic cohomology argument. 





\begin{theorem_}\label{VI:5-10}
Let $F$ and $k$ be two number fields, let $\lambda$ be a place of $k$ of 
residue characteristic $\ell$, and let 
$\chi_\lambda:\gal(\bar F/F)^\textnormal{ab}\to k_\lambda^\times$ be a 
homomorphism, unramified outside $\ell$ and a finite set $S$ of 
places. We assume that there exists a set $T$ of places (containing $S$ and the 
places above $\ell$), of density $0$, and an algebraic homomorphism 
$\chi_0:F^\times \to k^\times$ such that 
\begin{enumerate}[(i)]
  \item For $\bar k$ an algebraic closure of $k$, 
    $\chi_0:F^\times \to k^\times \hookrightarrow \bar k^\times$ is of the type 
    considered in \ref{VI:5-5}, of weight $N$; 
  \item For $\fp$ prime $\notin T$, $\chi_\lambda(F_\fp)$ lies in 
    $k^\times$, $(\chi_\lambda(F_\fp))=\chi_0(\fp)$, and all complex conjugates 
    of $\chi_\lambda(F_\fp)$ have absolute value $(N \fp)^{N/2}$. 
\end{enumerate}

Then $\chi_\lambda$ is defined (\ref{VI:5-8}) by an algebraic Hecke character 
of $F$ with values in $k$, whose algebraic part is $\chi_0$. 
\end{theorem_}

Let $\chi'$ be an algebraic Hecke character with values in a 
finite Galois extension $k'$ of $k$, of algebraic part $\chi_0$ 
(\ref{VI:5-5}). After enlarging $T$ if necessary, we may suppose that the conductor of 
$\chi'$ is supported in $T$. Let $\lambda'$ be a place of $k'$ above 
$\lambda$, and let $\chi_{\lambda'}$ be the composite 
$\gal(\bar F/F)^\text{ab}\to k_\lambda^\times \to {k'}_{\lambda'}^\times$, and 
$\chi_{\lambda'}':\gal(\bar F/F)^\text{ab}\to {k'}_{\lambda'}^\times$ defined 
by $\chi'$ (\ref{VI:5-8}). Set 
$\varepsilon = \chi_{\lambda'} \cdot {\chi'}_{\lambda'}^{-1}$.

Hypothesis (ii), together with \ref{VI:5-5}, \ref{VI:5-6}, and \ref{VI:5-8}, ensures that 
for $\fp\notin T$, we have $\varepsilon(F_\fp)\in {k'}^\times$, and that in every 
completion of $k'$ this number has absolute value $1$. It follows that 
$\varepsilon(F_\fp)$ belongs to the finite group $\dmu$ of roots of unity 
in $k'$. By the Chebotarev density theorem and the density-zero hypothesis 
on $T$, the $F_\fp$ ($\fp\notin T$) are dense in 
  $\gal(\bar F/F)^\text{ab}$, so that $\varepsilon(\sigma)\in \dmu$ for all 
$\sigma\in \gal(\bar F/F)^\text{ab}$: the character $\varepsilon$ is a 
  finite order character $\gal(\bar F/F)^\text{ab}\to \dmu$. By 
class field theory, it corresponds to a finite-order character 
of the idèle class group of $F$, i.e. to an algebraic Hecke character 
with trivial algebraic part. After correcting $\chi'$ by this 
character, we are reduced to the case in which 
$\chi_{\lambda'}=\chi_{\lambda'}'$, and it remains to show that $\chi'$, 
a priori taking values in $k'$, in fact takes values in $k$. If 
$\sigma\in \gal(k'/k)$, $\chi'$ and ${\chi'}^\sigma$ coincide on every prime 
ideal $\fp\notin T$, since 
$\chi'(\fp) = \chi'_{\lambda'}(F_\fp) = \chi_{\lambda'}(F_\fp) \in k^\times$. 
The characters $\chi_{\lambda'}'$, and ${\chi_{\lambda'}'}^\sigma$, coincide 
on a dense subset of $\gal(\bar F/F)^\text{ab}$ and hence are equal. Thus 
$\chi'={\chi'}^\sigma$, so $\chi'$ takes values in $k$. 










\section{Hecke characters defined by Jacobi sums}\label{VI:6}

The idea of this section is as follows. Consider a Gauss sum 
$\tau_{\dF_q}(\chi,\psi)$. Both this sum and the corresponding sheaf 
$\sF(\chi^{-1}\psi)$ are tied in two ways to the characteristic $p$ and to the 
base field $\dF_q$: 
\begin{enumerate}[\indent a)]
  \item because the additive character $\psi$ and the sheaf 
    $\sf(\psi)$ occur in it; 
  \item because $\sF(\chi)$ is defined from the exact Lang sequence 
    \[\xymatrix{
      0 \ar[r] 
        & k^\times \ar[r] 
        & V^\ast \ar[r] 
        & V^\ast \ar[r] 
        & 0 \text{.} 
    }\]
\end{enumerate}
These present two difficulties if we wish to lift the corresponding 
cohomological constructions to a number field $F$ and find 
$\ell$-adic representations of $\gal(\bar F/F)$ whose Frobenius eigenvalues 
are Gauss sums. 

If $\chi$ is trivial on $\dF_q^\times$, the sum $\tau_{\dF_q}(\chi,\psi)$ 
is independent of $\psi$ and is $q$ times a Jacobi sum. The latter is 
related to the cohomology of $W^\ast/\dG_m$, with coefficients in a sheaf 
obtained from $\chi$: thus $\psi$, and with it difficulty a), disappear. To 
resolve b), it suffices to describe the sheaves in question by means of 
exact K\"ummer sequences (which make sense in every characteristic). This is 
possible by \ref{VI:4-9}. 

Once the $\ell$-adic representations have been found, we use 
\ref{VI:5-10} and Stickelberger's theorem to show that they 
come from algebraic Hecke characters. 





\subsection{}\label{VI:6-1}

If $E$ is a field with algebraic closure $\bar E$, and $\lambda$ is 
a character, of order dividing $n$, of $\mu_d(\bar E)$, we denote by $_n\lambda$ the 
character of $\mu_n(\bar E)$ such that $\lambda$ and $_n \lambda$ induce the 
same character of $\widehat\dZ(1)_{\bar E} = \varinjlim \mu_n(\bar E)$. The 
same notation is used when $\lambda$ is a character of $\widehat\dZ(1)_{\bar E}$.

Let $F$ and $k$ be two number fields, $\bar F$ an algebraic closure of 
$F$, $I$ a finite set equipped with an action of $\gal(\bar F/F)$ and 
$\lambda=(\lambda_i)_{i\in I}$ a family of characters 
$\lambda_i:\widehat\dZ(1)_{\bar F} \to k^\times$. It is assumed that 
$\lambda_i\ne 1$ ($i\in I$), that the product of $\lambda_i$ is trivial and that 
the family $\lambda$ is defined over $F$, i.e. that for every 
$\sigma\in \gal(\bar F/F)$, we have 
$\lambda_{\sigma(i)}=\lambda_i\circ \sigma^{-1}$. 

The Galois set $I$ is the set of homomorphisms into $\bar F$ of an 
algebra $E$ separable over $F$. If $C$ is the set of orbits of 
$\gal(\bar F/F)$ in $I$, we have a decomposition of $E$ as a product of fields 
$E=\prod_{\alpha\in C} E_\alpha$, where $\alpha$ is identified with the set of 
embeddings of $E_\alpha$ in $\bar F$. If the characters $\lambda_i$, for 
$i\in \alpha$, are of order $d_\alpha$, the field $E_\alpha$ contains the 
$d_\alpha$-th roots of $1$ and there exists 
$\lambda_\alpha:\mu_{d_\alpha}(E_\alpha) \hookrightarrow k^\times$ such that, if 
$i\in \alpha$ corresponds to the embedding $\sigma_i$ of $E_\alpha$ in $\bar F$, 
then $_{d_\alpha} \lambda_i = \lambda_\alpha \sigma_i^{-1}$. 





\subsection{}\label{VI:6-2}

Let $n>0$ be a multiple of the $d_\alpha$, let $\Sigma$ be the set of places of $F$ 
at which $E/F$ ramifies or which divide $n$, let $\mu$ be a place of $k$ of 
residue characteristic $\ell$, and let $A$ be the ring of elements of $F$ 
integral outside $\Sigma$ and $\ell$. The integral closure $M$ of $A$ 
is finite étale of rank $N=|I|$ over $A$. With the notation of \ref{VI:4-6}, 
\ref{VI:4-8}, $_n \lambda_i$ define 
\[
  _n\lambda:(V^\ast/\dG_m)_n \to k^\times \subset k_\mu^\times \text{.} 
\]

Let $a$ be the projection of $W^\ast/\dG_m$ onto $\spec(A)$. By 
\ref{VI:4-20}, the sheaf $\eR^{N-2} a_! \sK_n(_n \lambda)$ is smooth of rank 
one; it defines an $\ell$-adic representation 
\[
  j_\mu[\lambda]:\gal(\bar F/F)^\text{ab} \to k_\mu^\times 
\]
  (or simply $j_\mu$) unramified outside $\Sigma$ and $\ell$. 





\subsection{}\label{VI:6-3}

Let us compute $j_\mu(F_\fp)$. By \ref{VI:4-20}, it suffices to do so 
after reduction modulo $\fp$. Put $f=A/\fp$, $e=M/\fp  M$, and let 
$\bar f$ be the algebraic closure of $f$ determined by a place of $\bar F$ 
above $\fp$. We still have $I=\hom_f(e,\bar f)$. After reduction, 
$_n\lambda$ still has a description as in \ref{VI:6-1}: 
\begin{enumerate}[a)]
  \item Let $D$ be the set of orbits of $\gal(\bar f/f)$ in $I$. The 
    decomposition of $e$ as a product of fields has the form 
    $e=\prod_{\beta\in D} e_\beta$, where $\sigma_i$ ($i\in \beta$) 
    are identified with the embeddings of $e_\beta$ into $\bar f$. For $\beta\in D$, 
    denote by $\alpha(\beta)$ the element of $C$ containing $\beta$, and put 
    $d_\beta=d_{\alpha(\beta)}$. Let 
    \[\xymatrix{
      \lambda_\beta : \mu_{d_\beta}(e_\beta) 
        & \ar[l]_-\sim \mu_{d_\beta}(E_{\alpha(\beta)}) \ar[r]^-{\lambda_{\alpha(\beta)}} 
        & k^\times \text{.} 
    }\]
  \item If $i\in \beta$ corresponds to the embedding $\sigma_i$ of $e_\beta$ into 
    $\bar f$, the reduction 
    $_{d_\beta} \bar\lambda_i:\mu_{d_\beta}(\bar f) \to k^\times$ 
    of $_{d_\beta} \lambda_i$ is 
    $\lambda_\beta \circ \sigma_i^{-1}$. 
\end{enumerate}

Let $q_\beta$ be the number of elements of $e_\beta$, 
  $\chi_\beta={}_{q_\beta-1}\lambda_\beta$, and 
$\chi:e^\times \to k^\times$ the character with coordinates $\chi_\beta$. On the 
affine space over $A/\fp$ defined by $e$ (\ref{VI:4-6}), or rather on the torus 
$e_\beta^\times$, we have isomorphisms 
  $\sK_n((_n\bar\lambda_i)_{i\in\beta}) = \sK_{d_\beta}((\lambda_\beta \circ \sigma_i^{-1})_{i\in\beta}) =\sK_{q_\beta-1}((\chi_\beta\circ \sigma_i^{-1})_{i\in\beta}) = \sF(\chi_\beta)$ (\ref{VI:4-9}). 
  On $e^\times$, we therefore have $\sK_n((_n\bar\lambda_i)_{i\in I}) = \sF(\chi)$, 
and this isomorphism descends to the reduction modulo $\fp$ of $V^\ast/\dG_m$. 
It follows that





\begin{proposition_}\label{VI:6-4}
With the notations of \ref{VI:6-2} and \ref{VI:6-3}, $j_\mu(F_\fp)$ is the 
Jacobi sum $J(\chi)$. 
\end{proposition_}

Let us rewrite the formulas linking $(\lambda_i)_{i\in I}$ to 
  $(\chi_\beta)_{\beta\in D}$: for $i\in \beta\subset \alpha$, 
\begin{align*}
  _{d_\alpha} \lambda_i &= \lambda_\alpha\circ \sigma_i^{-1} && \lambda_\alpha:\mu_{d_\alpha}(E_\alpha) \to k^\times \\
  \lambda_\beta &= \text{reduction of $\lambda_\alpha$} && :\mu_{d_\alpha}(e_\beta) \to k^\times \\
  \chi_\beta &= {}_{q_\beta-1}\lambda_\beta: x\mapsto \lambda_\beta(x^{(q_\beta-1)/d_\alpha}) && :e_\beta^\times \to k^\times 
\end{align*}





\begin{theorem_}\label{VI:6-5}
The representation $j_\mu$ is defined by an algebraic Hecke character 
of $F$ with values in $k^\times$. 
\end{theorem_}

We shall verify this using \ref{VI:5-10}. Here we follow Weil's 
proof \cite{we52,we74-2}, using Stickelberger's theorem, 
and I shall give only a few indications. 





\subsection{}\label{VI:6-6}

Given $j:\gal(\bar F/F) \to k_\lambda^\times$ and a finite extension 
$F'/F$, one verifies using \ref{VI:5-10} that $j$ is defined by a 
Hecke character if and only if $j|\gal(\bar F/F')$ is. This 
remark allows us, in \ref{VI:6-5}, to reduce to the case where Galois acts 
trivially on $I$ and where $F$ contains the $n$-th roots of 
unity. We may then reduce to the case $F=\dQ(\sqrt[n]{1})$ and 
take $k=F$. In this case, each $_n \lambda_i$ is defined by 
$a_i\in \dZ/n$: it is $x\mapsto x^{a_i}:\mu_n \to k^\times$. We have 
$\sum a_i=0$, and $a_i\ne 0$. 





\subsection{}\label{VI:6-7}

In a number field, the places of degree greater than $1$ always form a 
set of density $0$. By applying \ref{VI:5-10}, we may therefore disregard 
them and use Stickelberger only for a prime field. Let 
$p$ be prime, let $0<a<p-1$, and let $\widetilde x:\dF_p^\times \to \dQ_p^\times$ be the 
``multiplicative lift'' character; put 
  $g=-\sum_{x\in \dF_p^\times} \widetilde x^{-a} \zeta^x\in \dQ_p(\zeta)$, for 
$\zeta$ a $p$-th root of $1$. Again denote by $x$ the integer in $[0,p-1]$ 
  lifting $x$, and expand $\zeta^x=(1+\pi)^x$ by the binomial formula. 
  We find $g=-\sum_{i=0}^{p-1} A_i\pi^i$ with $A_i\in \dZ_p$, whose reduction 
$\mod p$ is the sum $\sum_{x\in \dF_p^\times} x^{-a} \binom x i$. For $b$ not 
divisible by $p-1$, we have $\sum_{x\in \dF_p^\times} x^b=0$. If $i<a$, 
$\binom x i$ is a polynomial in $x$ of degree $i<a$, whence 
$A_i\equiv 0\pmod p$. Likewise, $A_a\equiv \frac{p-1}{a!} \pmod p$, and 
$g\sim \frac{\pi^a}{a!}$: the valuation of $g$ is given by 
$v(g)/v(p) = \frac{a}{p-1}$. 





\subsection{}\label{VI:6-8}

Let $\sigma_i$ ($i\in (\dZ/n)^\times$) be the element of 
$\gal(\dQ(\sqrt[n]{1})/\dQ)$ such that $\sigma_i(\zeta)=\zeta^i$ for $\zeta^n=1$, 
and write the group of algebraic homomorphisms 
$\dQ(\sqrt[n]{1})^\times \to \dQ(\sqrt[n]{1})^\times$ additively. For $\lambda$ as in 
\ref{VI:6-6}, defined by a family 
$(a_i)_{i\in I}$ of nonzero integers $\mod n$ whose sum is zero, 
\ref{VI:5-10}, \eqref{VI:eq:4-15-1} and \ref{VI:6-7} 
show that $j_\mu[\lambda]$ is defined by an algebraic Hecke character 
$j[a]$ and that, if $N$ denotes the norm, its algebraic part is 
given by 
\begin{equation*}\tag{6.8.1}\label{VI:eq:6-8-1}
  (N j[a])_\text{alg} = \sum_{i\in (\dZ/n)^\times} \sum_{j\in I} \left\langle \frac{i a_j}{n} \right\rangle \sigma_i^{-1} 
\end{equation*}
  where $\langle \cdot \rangle$ is the fractional part. 





\subsection{}\label{VI:6-9}

The most interesting case is that in which $F\subset \dQ(\sqrt[n]{1})$ and 
the family of $_n \lambda_i$ is constructed as follows: 
\begin{enumerate}[a)]
  \item Take a finite family $(A_j)_{j\in J}$ of orbits of 
    $\gal(\dQ(\sqrt[n]{1})/F)=H\subset (\dZ/n)^\times$ in $\dZ/n$. Assume that 
    $A_j\ne\{0\}$ and that $\sum_{j\in J} \sum_{a\in A_j} a=0$. 
  \item Take $I=\sqcup_{j\in J} A_j$; Galois acts on each $A_j$. 
  \item Take $\bar F\supset \dQ(\sqrt[n]{1})$, $k=\dQ(\sqrt[n]{1})$ and, if 
    $i\in I$ corresponds to $a\in A_j\subset (\dZ/n)^\times$, 
    $_n\lambda_i(x)=x^a$. 
\end{enumerate}

Using \ref{VI:4-20}, one verifies that a general $j_\mu$ (relative to 
$F_1$) is obtained from a $j_\mu$ as above (for 
$F\subset \dQ(\sqrt[n]{1})$, with $F_1$ an extension of $F$) by restriction to 
$\gal(\bar F/F_1)$ and multiplication by a character of order two. 

For $a\in \gal(k/\dQ)=(\dZ/n)^\times$, $a$ transforms the family of $A_j$ 
into the family of $a A_j$. If $a\in H$, the family of $A_j$ is therefore preserved, 
and the $j_\mu(F_\fp)$ are invariant under $H$: 





\begin{proposition_}\label{VI:6-10}
Under the assumptions of \ref{VI:6-9}, $j_\mu$ is defined by a character 
that is an algebraic Hecke character of $F$ with values in $F$. 
\end{proposition_}

Let us again take $F=\dQ(\sqrt[n]{1})$ and consider the Hecke characters 
$j[a]$ defined by a family $(a_i)_{i\in I}$, $i\in \dZ/n$, 
$a_i\ne 0$, $\sum a_i=0$. The following theorem was obtained 
independently by Vishik. 





\begin{theorem_}\label{VI:6-11}
Every algebraic Hecke character of $\dQ(\sqrt[n]{1})$ has a power 
in the group generated by the $j[a]$ and the norm. 
\end{theorem_}

Since we work ``up to torsion,'' only the algebraic part of the 
Hecke characters under consideration matters. Applying \eqref{VI:eq:6-8-1} to 
$j[a]$, with $a=$ an element of $\dZ/n$ repeated $n$ times, we 
are reduced to the following lemma.





\begin{lemma_}\label{VI:6-12}
Write $\sum x_j \sigma_j$ for the elements of the group algebra 
  $\dQ[(\dZ/n)^\times]$, and let $X$ be the subgroup consisting of those for which 
$x_j+x_{-j}=C^{\underline{t e}}$. Then $X$ is generated over $\dQ$ by 
$N=\sum \sigma_j$ and by the elements 
  $g_a=\sum \left\langle \frac{i a}{n}\right\rangle \sigma_i^{-1}$ 
($a\in \dZ/n$, $a\ne 0$). 
\end{lemma_}

Both $X$ and $Y=\langle N,(g_a)\rangle$ are ideals of the group 
algebra, and it suffices to check that they are annihilated by the same 
complex characters $\chi$ of $(\dZ/n)^\times$. This amounts to saying that a 
odd character $(\chi(j)=-\chi(-j)$) never annihilates all the $g_a$. If 
$\chi$ is primitive, we have 
\[
  \chi(g_1) = \sum \frac i n \chi^{-1}(i) = -L(\chi,0)\ne 0 \text{.} 
\]
In the general case, if $\chi$ comes from a primitive character of 
  $(\dZ/d)^\times$, $\chi(g_{n/d})$ reduces to an analogous sum over 
  $(\dZ/d)^\times$, hence the lemma. 










\section{Generalized Kloosterman sums}\label{VI:7}





\subsection{}\label{VI:7-1}

Let $\dF_q$ be a finite field with $q$ elements, and let $\psi$ be a nontrivial 
additive character. In this section, we undertake a cohomological study of 
generalized Kloosterman sums 
\begin{equation*}\tag{7.1.1}\label{VI:eq:7-1-1}
  \kloos_{n,a} = \sum_{x_1\dotsm x_n = a} \psi(x_1+\cdots + x_n) \qquad (n\geqslant 1) \text{.} 
\end{equation*}
For $a=0$, this sum is elementary (see \ref{VI:7-7}): 
\begin{equation*}\tag{7.1.2}\label{VI:eq:7-1-2}
  \kloos_{n,0} = (-1)^{n-1} \text{.} 
\end{equation*}
The interesting case is $a\ne 0$. The $x_i$ then belong to 
$\dF_q^\times$. In this case we shall prove the bound 
\begin{equation*}\tag{7.1.3}\label{VI:eq:7-1-3}
  \left|\kloos_{n,a}\right| \leqslant n q^{\frac{n-1}{2}} \text{.} 
\end{equation*}

Our essential tools will be the cohomological counterparts of the 
following identities: 
\begin{align*}\tag{7.1.4}\label{VI:eq:7-1-4}
  \kloos_{n,a} &= \sum_{x_1} \psi(x_1) \sum_{x_2\dotsm x_n = a/x_1} \psi(x_2+\cdots + x_n) \\ \notag
    &= \sum_{x\in \dF_q^\times} \psi(x) \kloos_{n-1,a/x} \qquad (a\ne 0,n\geqslant 2) \\ \tag{7.1.5}\label{VI:eq:7-1-5}
  \sum_a \kloos_{n,a} &= \sum_a \sum_{x_1 \dotsm x_n = a} \psi\left(\sum x_i\right) = \sum_{x_1,\dots,x_n} \psi\left(\sum x_i\right) = 0 
\end{align*}
For $\chi$ a nontrivial character of $\dF_q^\times$, 
\begin{equation*}\tag{7.1.6}\label{VI:eq:7-1-6}
  \sum \chi(a) \kloos_{n,a} = \sum_{x_1,\dots,x_n} \chi\left(\prod x_i\right) \cdot \psi\left(\sum x_i\right) = (-\tau(\chi^{-1},\psi))^n 
\end{equation*}
(Gauss sum). 





\subsection{}\label{VI:7-2}

Let $k$ be an étale algebra of degree $n$ over $\dF_q$, and let $\varepsilon(k)$ 
be as in \eqref{VI:eq:4-5-1}. Put 
\[
  \kloos_{k,a} = \sum_{N_{k/\dF_q}(x)=a} \psi \tr(x) \text{.} 
\]
We likewise have 
\begin{align*}\tag{7.2.1}\label{VI:eq:7-2-1}
  \kloos_{k,0} &= (-1)^{n-1} \cdot \varepsilon(k) \\ \tag{7.2.2}\label{VI:eq:7-2-2}
  \sum_a \kloos_{k,a} &= 0 \\ \tag{7.2.3}\label{VI:eq:7-2-3}
  \sum_a \chi(a) \kloos_{k,a} &= (-1)^n \tau_{\dF_q}\left(\chi^{-1}\circ N_{k/\dF_q},\psi\right) \text{.} 
\end{align*}

By Fourier inversion on $\dF_q^\times$, we deduce an expression of 
$\kloos_{k,a}$ in terms of Gauss sums: 
\begin{align*} 
  \kloos_{k,a} &= \frac{1}{q-1} \sum_\chi \chi(a) \sum_{b\in \dF_q^\times} \chi^{-1}(b) \kloos_{k,b} \\ \tag{7.2.4}\label{VI:eq:7-2-4}
  (-1)^n \kloos_{k,a} &= \frac{1}{q-1} \left(\varepsilon(k) + \sum_{\chi\ne 1} \chi(a) \tau_{\dF_q}\left(\chi\circ N_{k/\dF_q},\psi\right) \right) \text{.} 
\end{align*}

The Hasse--Davenport identity now makes it possible to compare $\kloos_{k,a}$ 
with $\kloos_{n,a} = \kloos_{\dF_{q^n},a}$: 
\begin{equation*}\tag{7.2.5}\label{VI:eq:7-2-5}
  \kloos_{k,a} = \varepsilon(k) \kloos_{n,a} \text{.} 
\end{equation*}

I know of no cohomological interpretation of \eqref{VI:eq:7-2-4}, but 
I shall give one for \eqref{VI:eq:7-2-5}. Let us return to the sums $\kloos_{n,a}$. 





\subsection{}\label{VI:7-3}

Let $\dF$ be an algebraic closure of $\dF_q$ and, for $a\in \dF$, let 
$V_a^{n-1}$, or simply $V_a$, be the hypersurface in $\dA^n$ with equation 
$x_1 \dotsm x_n=a$. As usual, we shall regard $\psi$ as taking 
$\ell$-adic rather than complex values. With the notation of 
\ref{VI:1-8}(ii), \eqref{VI:eq:7-1-3} follows from 





\begin{theorem_}\label{VI:7-4}
The cohomology of $V_a^{n-1}$ with coefficients in $\sF(\psi(\sum x_i))$ satisfies 
\begin{enumerate}[\indent (i)]
  \item $\h_c^i=0$ for $i\ne n-1$; 
  \item $\h_c^\bullet{} \iso \h^\bullet{}$; 
  \item for $a\ne 0$, $\dim \h_c^{n-1}{}=n$; 
  \item for $a=0$, $\h_c^{n-1}{}$ is canonically isomorphic to $E_\lambda$. 
\end{enumerate}
\end{theorem_}





\subsection{}\label{VI:7-5}

As usual, this theorem also controls the dependence of 
$\kloos_{n,a}$ on $q$: for each $a\in \dF_q^\times$, there are $n$ Frobenius 
eigenvalues $\alpha_1,\dots,\alpha_n$, of complex absolute value 
$q^{\frac{n-1}{2}}$, such that 
\[
  \sum_{\substack{x_1\dotsm x_n = a \\ x_i\in \dF_{q^m}}} \psi\circ \tr\left(\sum x_i\right) = (-1)^{n-1} \left(\alpha_1^m + \cdots + \alpha_n^m \right) \text{.} 
\]

For even $n$, $x\mapsto -x$ is an involution of $V_a$, which transforms 
$\sF(\psi)$ into its dual. Arguing as in \ref{VI:3-6}, we deduce a 
canonical nondegenerate alternating form on $\h_c^{n-1}(V_a,\sF(\psi))$, 
with values in $\dQ_\ell(-(n-1))$, and hence the following. 





\begin{corollary_}\label{VI:7-6}
With the notations of \ref{VI:7-5}, if $n$ is even, the $\alpha_i$ group together 
into $\frac n 2$ pairs of roots $\alpha,\bar\alpha$ whose product is 
$q^{n-1}$. 
\end{corollary_}





\subsection{}\label{VI:7-7}

Let us verify the case $a=0$ of \ref{VI:7-5}. The hypersurface $V_0$ is the union 
of the coordinate hyperplanes $x_i=0$. Let us compute the Leray spectral sequence 
\eqref{VI:eq:2-6-1*} for this covering of $V_0$, for cohomology with or 
without supports and with coefficients in $\sF(\psi)$. By \ref{VI:2-7*}*, all 
the initial terms vanish except the cohomology of the intersection $\{0\}$ of 
all the coordinate hyperplanes. What remains is an isomorphism 
$\h_c^0{}\iso \h^0{}=E_\lambda$ which justifies \ref{VI:7-4}. 

We shall prove \ref{VI:7-4} and \ref{VI:7-8} below by simultaneous induction 
on $n$, using the spectral sequence of which 
\eqref{VI:eq:7-1-4} is the counterpart. This requires control, as $a$ varies, 
of $\kloos_{n,a}$. 

Let $\pi:\dA^n \to \dA$ be the map ``product of the coordinates,'' and let 
$\sigma$ be the sum of the coordinates. The hypersurfaces $V_a$ are the fibers of 
$\pi$.





\begin{theorem_}\label{VI:7-8}
\begin{enumerate}[(i)]
  \item The sheaf $\eR^{n-1}\pi_! \sF(\psi\sigma)$ is smooth of rank $n$ on 
    $\dA^1\setminus \{0\}$. 
  \item Its extension by $0$ to $\dP^1$ is the direct image of its 
    restriction to $\dA^1\setminus \{0\}$. 
  \item At $0$, the monodromy is unipotent, with a single Jordan block.
  \item At $\infty$, wild inertia acts without fixed points $\ne 0$, and the 
    Swan conductor is $1$. 
  \item We have $\eR^i \pi_! \sF(\psi\sigma) \iso \eR^i \pi_\ast \sF(\psi\sigma)$ 
    (zero for $i\ne n-1$).
\end{enumerate}
\end{theorem_}
\begin{proof}[Proof]
that \ref{VI:7-4} (for a given value of $n$) implies \ref{VI:7-8} 
(for the same value of $n$). 

The fiber of $\eR^i \pi_! \sF(\psi\sigma)$ at $a$ is 
$\h_c^i(V_a,\sF(\psi\sigma))$, and on the dense open subset of $\dA^1$ where the 
  $\eR^i \pi_!\sF(\psi\sigma)$ are locally constant, the fiber of 
$\eR^i \pi_\ast \sF(\psi\sigma)$ is $\h^i(V_a,\sF(\psi\sigma))$ 
(\hyperref[VII]{Finiteness\ Theorem} \ref{VII:2-1}). The hypothesis \ref{VI:7-4}($n$) 
therefore gives us the following. 





\begin{lemma_}\label{VI:7-9}
\begin{enumerate}[(i)]
  \item $\eR^i \pi_! \sF(\psi\sigma)=0$ for $i\ne n-1$. 
  \item For $i=n-1$, the fiber of this sheaf at any point $a\ne 0$ has 
    constant rank $n$. At $0$, it has rank $1$. 
  \item On a dense open $U$, we have 
    $\eR^i \pi_! \sF(\psi\sigma) \iso \eR^i \pi_\ast \sF(\psi\sigma)$. 
\end{enumerate}
\end{lemma_}





\subsection{}\label{VI:7-10}

We shall now use the Leray spectral sequence of $\pi$ for 
cohomology with or without compact supports, with coefficients in $\sF(\psi\sigma)$. 
The abutment is known: 
$\h_c^\bullet(\dA^n,\sF(\psi\sigma))=\h^\bullet(\dA^n,\sF(\psi\sigma)) = 0$ 
(\ref{VI:2-7*}*). For cohomology with compact supports, \ref{VI:7-9}(i) ensures that 
$E_2^{pq}=0$ for $q\ne n-1$. The terms $E_2$ are therefore all zero: 
\begin{equation*}\tag{7.10.1}\label{VI:eq:7-10-1}
  \h_c^p(\dA^1,\eR^{n-1} \pi_!\sF(\psi\sigma)) = 0 \text{.} 
\end{equation*}

  For $p=0$, this means that $\eR^{n-1}\pi_! \sF(\psi\sigma)$ has no sections 
with point support. In view of the constancy of rank in \ref{VI:7-9}(ii), we obtain 
\begin{equation*}\tag{7.10.2}\label{VI:eq:7-10-2}
  \begin{array}{c}
    \text{The sheaf $\eR^{n-1} \pi_! \sF(\psi\sigma)$ on $\dA^1$ is smooth on $\dA^1\setminus \{0\}$, and is } \\ 
    \text{a subsheaf of the direct image $\sG$ of its restriction to $\dA^1\setminus \{0\}$.}
  \end{array}
\end{equation*}

The long exact cohomology sequence of 
\[\xymatrix{
  0 \ar[r] 
    & \eR^{n-1} \pi_! \sF(\psi\sigma) \ar[r] 
    & \sG \ar[r] 
    & \sQ \ar[r] 
    & 0 \text{.} 
}\]
where the support of $\sQ$ is concentrated at $0$, gives 
\[\xymatrix{
  \h_c^0(\dA^1,\sG) \ar[r] 
    & \sQ_0 \ar[r] 
    & \h_c^1(\dA^1,\eR^{n-1} \pi_! \sF(\psi\sigma)) \text{.} 
}\]
Since the outer terms vanish, $\sQ=0$ and 
\begin{equation*}\tag{7.10.3}\label{VI:eq:7-10-3}
  \text{$\eR^{n-1} \pi_! \sF(\psi\sigma)$ is the direct image of its restriction to $\dA^1\setminus \{0\}$.}
\end{equation*}

Let $j$ be the inclusion of $\dA^1$ into $\dP^1$. The ``derived-category 
version'' of the Leray spectral sequences is 
\begin{align*}
  \hh^\bullet(\dP^1,j_! \eR\pi_! \sF(\psi\sigma)) &= \hh_c^\bullet(\dA^1,\eR \pi_! \sF(\psi\sigma)) = \h_c^\bullet(\dA^n,\sF(\psi\sigma)) = 0 \\
  \hh^\bullet(\dP^1,\eR j_\ast \eR \pi_\ast \sF(\psi\sigma)) &= \hh^\bullet(\dA^1,\eR \pi_\ast \sF(\psi\sigma)) = \h^\bullet(\dA^n,\sF(\psi\sigma)) = 0 \text{.}
\end{align*}
Let $\Delta$ be the mapping cylinder of 
$j_! \eR\pi_! \sF(\psi\sigma) \to \eR j_\ast \eR \pi_\ast \sF(\psi\sigma)$. 
According to \ref{VI:7-9}(iii), the cohomology sheaves of this complex of 
sheaves have finite support. On the other hand, the long exact sequence in 
cohomology of the triangle 
$(j_!\eR \pi_! \sF(\psi\sigma),\eR j_\ast \eR \pi_\ast \sF(\psi\sigma),\Delta)$ 
shows that $\hh^\bullet(\dP^1,\Delta)=0$. We have 
$\hh^\bullet(\dP^1,\Delta) = \h^0(\dP^1,\sH^\bullet(\Delta))$, and finally 
$\Delta=0$: 
\[\xymatrix{
  j_! \eR\pi_! \sF(\psi\sigma) \ar[r]^-\sim 
    & \eR j_\ast \eR \pi_\ast \sF(\psi\sigma) \text{.} 
}\]

In particular, $\eR\pi_! \sF(\psi\sigma)\iso \eR\pi_\ast \sF(\psi\sigma)$ and, 
since these complexes of sheaves have only one nonzero cohomology sheaf, 
\begin{equation*}\tag{7.10.4}\label{VI:eq:7-10-4}
\xymatrix{
  j_! \eR \pi_! \sF(\psi\sigma) \ar[r]^-\sim 
    & j_\ast \eR \pi_! \sF(\psi\sigma) \text{.} 
}
\end{equation*}
This completes the proof of (i)(ii)(v). 





\subsection{}\label{VI:7-11}

Let $\swan_0$ and $\swan_\infty$ be the Swan conductors at $0$ and at 
$\infty$ of $\eR^{n-1} \pi_! \sF(\psi\sigma)$. We have 
$\chi_c(\eR^{n-1} \pi_! \sF(\psi\sigma))=0$ \eqref{VI:eq:7-10-1} and the formula 
of Euler--Poincaré \eqref{VI:eq:3-2-1} reduces to 
\begin{equation*}\tag{7.11.1}\label{VI:eq:7-11-1}
  \swan_0 + \swan_\infty =1 \text{:} 
\end{equation*}
one of these conductors is $0$ (tame ramification), and the other is $1$. 





\subsection{}\label{VI:7-12}

Let $\sG$ be a sheaf of rank $1$ on $\dG_m$, of Kummer type 
\ref{VI:4-7} ($\sG=\sK_n(\chi)$) and nonconstant ($\chi\ne 1$). The sheaf 
$\sF(\psi\sigma)\otimes \pi^\ast \sG$ on 
$V^\ast=\pi^{-1}(\dG_m) = \dA^n\setminus \pi^{-1}(0)$ is then of the type 
encountered in the study of Gauss sums, and \ref{VI:4-11}(ii) gives us 
\begin{equation*}\tag{7.12.1}\label{VI:eq:7-12-1}
\xymatrix{
  \h_c^\bullet(V^\ast,\pi^\ast\sG\otimes \sF(\psi\sigma)) \ar[r]^-\sim 
    & \h^\bullet(V^\ast,\pi^\ast \sG\otimes \sF(\psi\sigma)) \text{.} 
}
\end{equation*}
The Leray spectral sequence of $\pi$ transforms this statement into 
\begin{equation*}\tag{7.12.2}\label{VI:eq:7-12-2}
\xymatrix{
  \h_c^\bullet(\dG_m,\sG\otimes \eR^{n-1} \pi_! \sF(\psi\sigma)) \ar[r]^-\sim 
    & \h^\bullet(\dG_m,\sG\otimes \eR^{n-1} \pi_\ast \sF(\psi\sigma)) \text{.} 
}
\end{equation*}

Let $i$ be the inclusion of $\dG_m$ into $\dP^1$, and let $\Delta$ be the mapping cylinder 
of 
  $\delta:i_!(\sG\otimes \eR^{n-1}\pi_! \sF(\psi\sigma))\to \eR i_\ast (\sG\otimes \eR^{n-1} \pi_! \sF(\psi\sigma))$
(where $\eR^{n-1} \pi_! = \eR^{n-1}\pi_\ast$). Reasoning as in 
\ref{VI:7-10}, we deduce from \eqref{VI:eq:7-12-2} that $\Delta=0$. In 
particular 
\begin{equation*}\tag{7.12.3}\label{VI:eq:7-12-3}
  \text{At $0$ and at $\infty$, inertia acts without fixed points $\ne 0$ on $\sG\otimes \eR^{n-1}\pi_! \sF(\psi\sigma)$.}
\end{equation*}

At $\infty$, this holds even if $\sG$ is trivial \eqref{VI:eq:7-10-4}, 
so wild inertia acts without a fixed point. In particular, the 
ramification is wild: by \ref{VI:7-11}, we have $\swan_\infty=1$, and the 
ramification at $0$ is tame. 

At $0$, \eqref{VI:eq:7-12-3} forces the ramification -- which is tame -- of 
$\eR^{n-1} \pi_! \sF(\psi\sigma)$ to be unipotent. By 
\eqref{VI:eq:7-10-3} and \ref{VI:7-9}(ii), there is a single Jordan block. This 
completes the proof of \ref{VI:7-8}($n$). 





\subsection{}\label{VI:7-13}

For $n=1$, \ref{VI:7-4} is trivial. To complete the proof of \ref{VI:7-4} and 
\ref{VI:7-8}, we still have to prove that \ref{VI:7-8}, for a given value 
of $n$, implies \ref{VI:7-4} for $n+1$. The case $a=0$ having been 
  treated (\ref{VI:7-7}), we assume that $a\ne 0$. We shall write $x_0,\dots,x_n$ for 
the coordinates of $\dA^{n+1}$. Let $g=x_0:V_a\to \dG_m$, $\tau$ 
the involution $x\mapsto a x^{-1}$ of $\dG_m$, and let $\pi$, as above, be the 
product of the coordinates: $\dA^{n+1} \to \dA^1$, with $\sigma$ their sum. We 
shall denote by $\sF(\psi\sigma)$ the sheaf $\sF(\psi(\sum x_i))$ both on $\dA^{n+1}$ 
and on $\dA^n$. 





\subsection{}\label{VI:7-14}

Let us write down the Leray spectral sequence for $g$. Since the sheaf $\sF(\psi\sigma)$ on 
$\dA^{n+1}$ is the external tensor product of the sheaves $\sF(\psi)$ on 
$\dA^1$ and $\sF(\psi\sigma)$ on $\dA^n$, we find (cf. \eqref{VI:eq:7-1-4}). 
\begin{align*}
  'E_2^{p q} &= \h_c^p(\dG_m,\sF(\psi)\otimes \tau^\ast \eR^q \pi_! \sF(\psi\sigma)) \Rightarrow \h_c^{p+q}(V_a,\sF(\psi\sigma)) \\
  ''E_2^{p q} &= \h^p(\dG_m,\sF(\psi)\otimes \tau^\ast \eR^q \pi_\ast \sF(\psi\sigma)) \Rightarrow \h^{p+q}(V_a,\sF(\psi\sigma)) \text{.} 
\end{align*}
By hypothesis, 
$\eR^q \pi_! \sF(\psi\sigma) \iso \eR^q \pi_\ast \sF(\psi\sigma)$, and this 
sheaf is zero for $q\ne n-1$. The sheaf $\sF(\psi)$ on $\dG_m$ is 
wildly ramified at $\infty$ and unramified at $0$, while 
$\tau^\ast \eR^q \pi_! \sF(\psi\sigma)$ is wildly ramified at $0$ (without 
invariants under wild inertia, and with Swan conductor $1$) and tamely 
ramified at $\infty$. Their tensor product is therefore 
\begin{enumerate}[\indent a)]
  \item wildly ramified at $0$ and $\infty$, without invariants under wild 
    inertia; 
  \item of Swan conductor $1$ at $0$ and $n$ (the rank of 
    $\eR^q \pi_! \sF(\psi\sigma))$ at $\infty$. 
\end{enumerate}
We have $'E_2^{p q} \iso ''E_2^{p q}$, and $'E_2^{p q}=0$ except for $p=1$ and 
$q=n-1$. The Euler-Poincaré formula finally gives 
\[
  \dim{'E_2^{1,n-1}} = 1+n \text{;} 
\]
this completes the proof. 
\end{proof}





\subsection{Remark}\label{VI:7-15}

The isomorphism obtained in \ref{VI:7-14} 
\begin{equation*}\tag{7.15.1}\label{VI:eq:7-15-1}
  \h_c^n(V_a^n,\sF(\psi\sigma)) = \h_c^1(\dG_m,\sF(\psi)\otimes \tau^\ast \eR^{n-1} \pi_! \sF(\psi\sigma)) 
\end{equation*}
allows us, with the notation of \ref{VI:7-5}, to calculate the product $\det(F)$ 
of the Frobenius eigenvalues $\alpha_i$. The formalism of \cite{de73} 
applies, since $\sF(\psi)$ and $\eR^{n-1} \pi_! \sF(\psi\sigma)$ 
belong to infinite compatible systems of $\ell$-adic 
representations. This reduces us to local problems, which can be 
simplified by noting that (cf. \cite[9.5]{de73}) 
\begin{enumerate}[\indent a)]
  \item the global constants for 
    $\sF(\psi)$ and $\tau^\ast \eR^{n-1} \pi_! \sF(\psi\sigma)$ are easy to 
    calculated easily, since the cohomology of these sheaves is elementary. 
  \item at every point, for either of these sheaves, the 
    corresponding $\ell$-adic representation of the decomposition group 
    has unramified semisimplification. 
\end{enumerate}

The result, for the sum with $n$ variables, is that 
\begin{equation*}\tag{7.15.2}\label{VI:eq:7-15-2}
  \det\left(F,\h_c^{n-1}(V_a^{n-1},\sF(\psi))\right) = q^{\frac{n(n-1)}{2}} \text{.} 
\end{equation*}





\subsection{Remark}\label{VI:7-16}

For $p=2$, the sheaf $\sf(\psi)$ is orthogonal. For $a\ne 0$ and $n$ 
odd, we deduce from \ref{VI:7-4}(ii) and Poincaré duality a 
symmetric bilinear form on $\h_c^{n-1}(V_a^{n-1},\sF(\psi))$ with 
values in $\bar\dQ_\ell(1-n)$. With respect to this form, $F$ is an 
orthogonal similitude of multiplier $q^{n-1}$, and $q^{\frac{1-n}{2}} F$ 
belongs to the special orthogonal group, by \eqref{VI:eq:7-15-2}. One 
of the eigenvalues of $F$ is therefore $q^{\frac{n-1}{2}}$, and the others 
occur in pairs of roots $\alpha$, $\bar\alpha$ whose product is $q^{n-1}$. 

For $n=3$, L.\ Carlitz \cite{ca69-2} obtained a more precise result, 
equivalent to the following proposition. 





\begin{proposition_}\label{VI:7-17}
For $p=2$, and $\psi(x)=(-1)^{\tr_{\dF_q/\dF_2}(x)}$, we have an isomorphism 
\[
  \h_c^2(V_a^2,\sF(\psi\sigma)) = \operatorname{Sym}^2 \h_c^1(V_a^1,\sF(\psi\sigma)) \text{.} 
\]
\end{proposition_}

The second tensor power of $\h_c^1(V_a^1,\sF(\psi\sigma))$ is 
$\h_c^2(V_a^1\times V_a^1,\sF(\psi\sigma)\boxtimes \sF(\psi\sigma))$ 
(K\"unneth), with the symmetry $x\otimes y\mapsto y\otimes x$ represented 
geometrically by $-\tau^\ast$, where $\tau$ is the automorphism 
$(x,y)\mapsto (y,x)$ of $V_a\times V_a$. If $X'$ is the quotient of 
$V_a^1\times V_a^1$ by $\{\operatorname{id}{},\tau\}$ and $\pi$ is the 
projection of $V_a^1\times V_a^1$ onto $X'$, the group 
$\operatorname{Sym}^2 \h_c^1(V_a^1,\sF(\psi\sigma))$ is therefore identified with the 
cohomology of $X'$ with coefficients in the part anti-invariant under $\tau$ of 
$\pi_\ast(\sF(\psi\sigma)\boxtimes \sF(\psi\sigma))$. Denote by $s$ the function 
on $X'$ such that $s\pi =\sigma\pr_1+\sigma\pr_2$. We have 
$\sF(\psi\sigma)\boxtimes \sF(\psi\sigma) = \sF(\psi(\sigma\pr_1+\sigma\pr_2)) = \sF(\psi s \pi) = \pi^\ast \sF(\psi s)$ 
(an isomorphism compatible with $\tau$). On the image $\pi(\Delta)$ of the 
diagonal, the desired sheaf of anti-invariants is therefore zero, while on 
the complement $X=X'\setminus \pi(\Delta)$ it is the tensor product of 
$\sF(\psi s)$ by $\varepsilon(P)$, where $\varepsilon$ is the unique character 
of order $2$ of $\{1,\tau\}$, and $P$ the $\{1,\tau\}$-torsor on $X$ given by the 
double covering induced by $V_a^1\times V_a^1$. We have 
\begin{equation*}\tag{7.17.1}\label{VI:eq:7-17-1}
  \operatorname{Sym}^2 \h_c^1(V_a^1,\sF(\psi\sigma)) = \h_c^2(X,\sF(\psi s)\otimes \varepsilon(P)) \text{.}
\end{equation*}

Let us compute $X$, $s$ and $\varepsilon(P)$. On $V_a^1\times V_a^1$, put 
$x_1=x_1\circ\pr_1$, $x_2=x_2\circ\pr_2$, $x_1'=x_1\circ \pr_2$, 
$x_2'=x_2\circ \pr_1$. We have $x_1 x_2=x_1'x_2'=a$. Identify functions on 
$X'$ with functions on $X$ invariant under $\tau$. Let $s_1=x_1+x_1'$, 
$s_2=x_2+x_2'$, $p_1=x_1 x_1'$, $p_2=x_2 x_2'$. We have $s_1 p_2 = a s_2$, 
$s_2 p_1 = a s_1$ and $X$ is defined in $X'$ by $s_1\ne 0$, $s_2\ne 0$. On 
$X$, we therefore have $p_1 = a s_1 s_2^{-1}$, and $(s_1,s_2)$ identifies $X$ with 
$\dG_m\times \dG_m$. On $X$, the covering $P$ has equation 
$T^2-s_1 T+a s_1 s_2^{-1}$ (take $T$ to be the coordinate $x_1$). If 
$T=s_1 t$, this is rewritten $t^2-t+a s_1^{-1} s_2^{-1}$. We therefore have 
$\varepsilon(P) = \sF(\psi(a s_1^{-1} s_2^{-1}))$. Since $s=s_1+s_2$, we have 
$\sF(\psi s)\otimes \varepsilon(P) = \sF(\psi(s_1+s_2+ a s_1^{-1} s_2^{-1}))$ 
and the right-hand side of \eqref{VI:eq:7-17-1} is identified with the left-hand side of 
\ref{VI:7-17}. 


\paragraph{Remark}

$\h_c^1(V_a^1,\sF(\psi\sigma))$ is also the first cohomology group of 
the elliptic curve obtained by completing the affine curve with equations 
\begin{align*}
  y^2-y &= x_1+x_2 \\
  x_1 x_2 &= a \text{;} 
\end{align*}
its modular invariant is $j=a^{-2}$. 





\subsection{Remark}\label{VI:7-18}

For $\chi$ a multiplicative character of $\dF_q^\times$, and $a\ne 0$, the 
method of \ref{VI:7-14} applies to the study of the sum 
\[
  S = \sum_{x_0\dotsm x_n = a} \chi(x_0) \psi(x_0+\cdots + x_n) 
\]
(replace the sheaf $\sF(\psi)$ on $\dG_m$ by $\sF(\chi\cdot \psi)$). We 
again find a cohomology group of dimension $n+1$, and 
$|S|\leqslant (n+1)q^{n/2}$. Such sums were considered by 
Salie and Mordell \cite{mo73}.

One may hope that our methods make it possible to study the sums 
\[
  S_{k,\chi,a} = \sum_{N(x)=a} \chi(x) \psi \tr(x) 
\]
($k$ an étale algebra over $\dF_q$, $a\in \dF_q^\times$, $\chi$ a character of 
$k^\times$). The analog of \ref{VI:7-4} (for $a\ne 0$) should be true. 





\subsection{}\label{VI:7-19}

The symmetric group $S_n$ acts on $\dA^n$, preserving $V_a$ and 
the map $\sigma$, and hence also the sheaf $\sf(\psi\sigma)$, the inverse 
image under $\sigma$ of the sheaf $\sf(\psi)$ on $\dG_a$. It therefore acts 
on $\h_c^{n-1}(V_a,\sF(\psi\sigma))$. Formula \eqref{VI:eq:7-2-5} admits 
the following cohomological interpretation.





\begin{proposition_}\label{VI:7-20}
The action of $S_n$ on $\h_c^{n-1}(V_a,\sF(\psi\sigma))$ is multiplication 
by the sign character. 
\end{proposition_}

It suffices to show that a transposition $\tau$ acts by multiplication 
by $-1$. Since $\tau^2=1$, the only possible eigenvalues of $\tau$ are 
$\pm 1$. Let $V_a/\tau$ be the quotient of $V_a$ for 
$\{\operatorname{id},\tau\}$. The map $\sigma$, and hence the sheaf 
$\sf(\psi\sigma)$, descend to the quotient, and the subspace of 
  $\h_c^{n-1}(V_a,\sF(\psi\sigma))$ fixed by 
$\tau$ is $\h_c^{n-1}(V_a/\tau,\sF(\psi\sigma))$. We must prove that 
this cohomology is zero. 

Suppose $a\ne 0$. For $\tau=(1,2)$ the quotient $V_a/\tau$ is identified 
with $\dG_a\times \dG_m^{n-2}$, the quotient map being 
$(x_1,x_2,x_3,\dots,x_n)\mapsto (x_1+x_2,x_3,\dots,x_n)$. Indeed, $x_1$ and 
$x_2$ are determined up to order by $x_1+x_2$ and 
$x_1x_2 = a/x_3\dotsm x_n$. In the coordinates $(t,x_3,\dots,x_n)$, the 
sheaf takes the form 
$\sF(\psi(t+x_3+\cdots + x_n)) = \sF(\psi(t))\otimes \sF(\psi(x_3+\cdots + x_n))$. 
  The K\"unneth formula and $\h_c^\bullet(\dG_a,\sF(\psi)) = 0$ therefore give 
the required vanishing. 

For $a=0$, we can deduce \ref{VI:7-20} from calculation \ref{VI:7-7}. 





\subsection{}\label{VI:7-21}

The sums $\kloos_{n,a}$ in this section were studied by 
S.\ Sperber, in his thesis, and in \cite{sp77} by methods 
$p$-adic methods inspired by Dwork and by Bombieri's article \cite{bo66}. 
For $p\ne 2$, these methods determine the dependence on $q$ in \ref{VI:7-5}, 
the functional equation $\alpha_i = q^{n-1} \beta_i^{-1}$ between the 
Frobenius eigenvalues for the sums $\kloos_{n,a}$ and 
$\kloos_{n,b}$ if $b=(-1)^n a$ (see \ref{VI:7-6}), and the value \ref{VI:7-15} 
of the product of the roots. For $p\geqslant n+3$, they also determine the 
$p$-adic valuations of the roots. The result is very beautiful: if we arrange 
the roots $\alpha_i$ in a suitable order, with 
  $0\leqslant i<n$, $\alpha_i q^{-i}$ are units. The 
restrictions on $p$ are probably not essential. 










\section{Other applications}\label{VI:8}





\subsection{}\label{VI:8-1}

The identity \ref{VI:1-9-4} can sometimes be used to calculate 
the number of rational points of an algebraic variety over $\dF_q$. In 
most cases where an exact result has been obtained, cohomology 
is expressed in terms of algebraic cycles; this is the case for rational 
surfaces, quadrics, smooth intersections of two quadrics of 
odd dimension, flag varieties\ldots. This is not, however, 
always the case; in Lusztig's beautiful article \cite{lu76}, algebraic cycles 
do not appear. 





\subsection{}\label{VI:8-2}

For reductive groups, cohomology can be calculated by lifting to 
characteristic $0$: if $B$ is a Borel subgroup of $G$, $G$ is a 
$B$-torsor over $G/B$; the cohomology of the smooth projective variety 
$G/B$ is invariant under specialization, and we use the Leray spectral sequence 
of $G\to G/B$ to prove the same result for $G$. It can be expressed 
conveniently in terms of \emph{the} maximal torus $T$ of $G$ and the action 
\emph{of the} Weyl group $W$ on $T$ (for the definition \emph{of the} maximal 
torus, cf. \cite[VIII \S 2 Rem.2]{bo05} or \cite[p.105]{dl76}). It is 
the exterior algebra on its primitive part, which agrees up to a shift 
with the quotient $I(\h^\bullet(B G,\dQ_\ell))$ of the cohomology of 
$B G$ formed by the ``indecomposable elements,'' and 
$\h^\bullet(B G,\dQ_\ell)=\h^\bullet(B T,\dQ_\ell)^W$. If $X$ is the group of 
characters of $T$, we have altogether 
\[
  \h^\bullet(G,\dQ_\ell) = \bigwedge \left(I\left(\operatorname{Sym}^\bullet(X\otimes \dQ_\ell(-1))^W\right)[-1]\right) \text{.}
\]
If $G$ is defined over a finite field $\dF_q$, $\gal(\dF/\dF_q)$ acts on $X$ 
and this isomorphism is Galois-compatible. Passing to 
cohomology with compact supports, we recover the classical formulas for the 
number of rational points of $G$. If $F^\ast:X\to X$ is defined by the 
Frobenius morphism $F:T\to T$, and $\deg_G(F)$ is the degree 
$q^{\dim G}$ of $F:G\to G$, we have 
\begin{equation*}\tag{8.2.1}\label{VI:eq:8-2-1}
  \left|G_0(\dF_q)\right| = \deg_G(F) \cdot \det\left(1-{F^\ast}^{-1},I(\operatorname{Sym}^\bullet(X\otimes \dQ)^W)\right) \text{.}
\end{equation*}

The same formula applies to the Ree and Suzuki groups; these groups are 
of the form $G^F$ (reductive $G$, $F^2$ a Frobenius), and it suffices to 
verify that the Lefschetz trace formula holds for $F$ acting 
on $G$. If $B$ is a Borel subgroup stable by $F$, we verify it 
directly for $F$ acting on $B$; for the action on $G/B$, which is proper and 
smooth, we can invoke the general theorems. 





\subsection{}\label{VI:8-3}

The trace formula was used by Deligne--Lusztig \cite{dl76}, 
Kazhdan \cite{ka77} and Springer \cite{sp76} in the study of 
complex representations of finite groups $G_0(\dF_q)$, for $G_0$ 
reductive over $\dF_q$. The work of Kazhdan and Springer contains 
admirable examples of how the trace formula permits ``analytic 
continuation'' from a split situation to a nonsplit situation 
(see \ref{VI:1-13}).

% END INCLUDED components/sga4half_6_en.tex

% BEGIN INCLUDED components/sga4half_7_en.tex
% !TEX root = ../SGA4half_English_Reader_1_to_8.tex


\chapter{Finiteness theorems in \texorpdfstring{$\ell$}{l}-adic cohomology}\label{VII}











\section{Statement of the theorems}\label{VII:1}

Throughout this exposé, $S$ will be a Noetherian scheme and $A$ a left
Noetherian torsion ring annihilated by an integer invertible on $S$. Our main
result is the following.






\begin{theorem_}\label{VII:1-1}
Suppose that $S$ is regular of dimension $0$ or $1$. Let $f:X\to Y$ be a
morphism of $S$-schemes of finite type, and let $\sF$ be a constructible sheaf
of left $A$-modules on $X$. Then the sheaves $\eR^i f_\ast \sF$ are
constructible.
\end{theorem_}

The proof will be given in Section \ref{VII:2} and in \ref{VII:3-10}.






\subsection{Remark}\label{VII:1-2}

In characteristic $0$, this result is less general than
\cite[XIX paragraphe 5]{sga4}, which proves the conclusion of the theorem for
every morphism of finite type of excellent schemes of characteristic $0$.
The proof of \cite[XIX]{sga4} uses, on the one hand, resolution of
singularities and, on the other, the assumption that the schemes have equal
characteristic (in order to deduce the ``purity theorem'' from resolution).






\subsection{Remark}\label{VII:1-3}

We shall say that a complex $K\in\ob\eD(X,A)$ is \emph{constructible} if its
cohomology sheaves are constructible, and shall use a subscript $c$ for the
subcategory of $\eD(X,A)$ (or $\eD^+$, $\eD^-$, $\eD^b$) consisting of
constructible complexes. In the language of derived categories, which we
shall use freely, \ref{VII:1-1} says that
$\eR f_\ast:\eD^+(X,A) \to \eD^+(Y,A)$ sends $\eD_c^+$ into $\eD_c^+$.
More explicitly:
\begin{enumerate}[\indent a)]
  \item For $K$ concentrated in a sheaf $\sF$ in degree $0$, one has
    $\sH^i\eR f_\ast K = \eR^i f_\ast \sF$; the derived statement therefore
    implies the theorem.
  \item Conversely, one invokes the spectral sequence
    \[
      E_1^{p q} = \eR^q f_\ast \sH^p(K) \Rightarrow \sH^{p+q} \eR f_\ast K \text{.}
    \]
\end{enumerate}

Speaking the language of derived categories has the advantage of replacing,
by the simple transitivity formula $\eR(f g)_\ast = \eR f_\ast \eR g_\ast$,
the Leray spectral sequence
$E_2^{p q} = \eR^p f_\ast \eR^q g_\ast \Rightarrow \eR^{p+q} (f g)_\ast$.

Under the hypotheses of \ref{VII:1-1}, it follows from \cite{sga4} that
$\eR f_\ast$ has finite cohomological dimension. This makes it possible above
to replace $\eD^+$ by $\eD$, $\eD^-$, or $\eD^b$.






\subsection{Ideas of the proof}\label{VII:1-4}
\begin{enumerate}[\indent a)]
  \item Poincaré duality treats the case where $X$ is smooth, $\sF$ is
    locally constant, and $Y=S$. The dimension hypothesis is seen na\"ively
    when calculating the $\rHom$ on $S$.
  \item Factor $f$ as $g j$: $g$ proper and $j$ an open immersion. One has
    $\eR f_\ast = \eR g_\ast \eR j_\ast$, and the finiteness theorem for
    proper morphisms controls $\eR g_\ast$. By dévissage, one reduces to
    assuming that $X$ is smooth, $\sF$ is locally constant, $f$ is an open
    immersion, and $Y$ is proper over $S$.
  \item Naively, an induction on $\dim X$ (with a change of $S$) makes it
    possible to suppose that the theorem is true outside a part of $Y$ finite
    over $S$. Denoting the structure morphisms by $a:X\to S$ and $b:Y\to S$,
    one then shows that if the theorem were false for $\sF$ and
    $\eR f_\ast$, it would also be false for $\sF$ and
    $\eR b_\ast \eR f_\ast = \eR a_\ast$: a contradiction.
\end{enumerate}






\begin{corollary_}\label{VII:1-5}
Under the hypotheses of the theorem, the categories $\eD_c(X,A)$ and
$\eD_c(Y,A)$ are taken into one another by the $4$ operations
$\eR f_\ast$, $\eR f_!$, $f^\ast$, $\eR f^!$.
\end{corollary_}

The case of $\eR f_\ast$ is treated above, that of $\eR f_!$ in
\cite[XVII 5.3]{sga4}, and that of $f^\ast$ is clear. It remains to consider
$\eR f^!$. The problem is local, which makes it possible to suppose that $f$
factors as $X\xrightarrow i Z \xrightarrow g Y$ ($i$ a closed immersion and
$g$ smooth, purely of relative dimension $n$). Poincaré duality
$\eR g^! K \iso g^\ast K(n)[2 n]$ and the transitivity
$\eR f^! = \eR i^! \eR g^!$ reduce us to proving \ref{VII:1-5} for $i$.
For every $L\in \eD(Z,A)$, if $j$ is the inclusion into $Z$ of
$U=Z\setminus X$, then $i_\ast \eR i^! L$ is the mapping cylinder of
$L\to \eR j_\ast j^\ast L$, and \ref{VII:1-5} for $i$ follows from
\ref{VII:1-1} for $j$.






\begin{corollary_}\label{VII:1-6}
Let $X$ be as in the theorem, and suppose that $A$ is commutative. Then, if
$\sF$ and $\sG$ are constructible sheaves of $A$-modules on $X$, the
$\Ext_A^i(\sF,\sG)$ are constructible: $\rHom$ sends
$\eD_c^-(X,A)\times \eD_c^+(X,A)$ into $\eD_c^+(X,A)$.
\end{corollary_}

By dévissage of $\sF$, we reduce to assuming that $\sF$ has the form
$j_!\sF_1$ ($j:Y\hookrightarrow X$ a locally closed immersion and $\sF_1$
locally constant on $Y$). One then has
\[
  \rHom(j_! \sF_1,\sG) = \eR j_\ast \Hom(\sF_1,\eR j^! \sG) \text{:}
\]
by \ref{VII:1-3}, we are reduced to the case where $\sF$ is locally
constant---or even constant, since the problem is local. For $\sF$ locally
constant and constructible, one has
$\Hom(\sF,\sG)_x=\hom(\sF_x,\sG_x)$, and likewise for $\rHom$. For $\sF$
constant, the $\Ext^i$ can therefore be calculated using a projective
resolution of finite type of its constant value, and the $\Ext^i$ are
constructible if $\sG$ is---whence the corollary.






\subsection{Remark}\label{VII:1-7}

Since $\eR f_\ast$ has finite cohomological dimension, it takes complexes of
finite tor-dimension (resp. $\leqslant d$) to complexes of finite
tor-dimension (resp. $\leqslant d$) \cite[XVII 5.2.11]{sga4}.
\cite[XVII 5.2.10]{sga4} and the proof of \ref{VII:1-5} then show that finite
tor-dimension is stable under the $4$ operations $\eR f_\ast$,
$\eR f_!$, $f^\ast$, $\eR f^!$. A variant of the argument of
\ref{VII:1-6} (apply dévissage to $K$ along a partition of $X$ so that the
cohomology sheaves become locally constant, then localize so as to replace $K$
by a finite complex of locally free sheaves of finite type) then shows that,
for $A$ commutative, $\rHom$ induces
\[
  \rHom:\eD_\text{ctf}^b(X,A)\times \eD_\text{tf}^b(X,A) \to \eD_\text{tf}^+(X,A)
\]
(t.f.: finite tor-dimension).






\subsection{}\label{VII:1-8}

Under the hypotheses of the theorem, the same method allows us to prove a
local biduality theorem $K\iso D D K$ (Section \ref{VII:4}). We also prove a
finiteness theorem for vanishing-cycle sheaves. For arbitrary $S$, one still
obtains a ``generic'' theorem:






\begin{theorem_}\label{VII:1-9}
Let $f:X\to Y$ be a morphism of $S$-schemes of finite type, and let $\sF$ be
a constructible sheaf of $A$-modules on $X$. There is a dense open subset
$U$ of $S$ such that
\begin{enumerate}[\indent (i)]
  \item Above $U$, the $\eR^i f_\ast \sF$ are constructible, and all but
    finitely many of them vanish.
  \item Formation of the $\eR^i f_\ast \sF$ is compatible with every base
    change $S' \to U\subset S$.
\end{enumerate}
\end{theorem_}

For $S$ the spectrum of a field, one has $U=S$.

If $S$ is the spectrum of a field, a passage-to-the-limit argument gives
compatibility with $S$-base changes for the $\eR^i f_\ast$ for every
quasi-compact quasi-separated morphism of $S$-schemes and every sheaf $\sF$.






\begin{corollary_}\label{VII:1-10}
Let $X$ be a scheme of finite type over a separably closed $k$, and let
$\sF$ be a constructible sheaf of $A$-modules. Then the $\h^i(X,\sF)$ are of
finite type.
\end{corollary_}

This is the special case $S=\spec(k)$, $Y=S$.

\begin{corollary_}\label{VII:1-11}
Let $X$ and $Y$ be two schemes of finite type over a separably closed $k$,
and let $K\in \ob\eD^-(X,A^\circ)$,
$L\in \ob\eD^-(Y,A)$ be complexes of sheaves of right and left modules. The
K\"unneth arrow
\[
  \eR \Gamma(X,K) \lotimes \eR \Gamma(Y,L) \to \eR\Gamma(X\times Y,\pr_1^\ast K \lotimes \pr_2^\ast L)
\]
is an isomorphism.
\end{corollary_}

Proceed as in \cite[XVII 5.4.3]{sga4}:
\[\xymatrix{
  X\times Y \ar[r]^-{\text{pr}_1} \ar[d]^-{\text{pr}_2}
    & X \ar[d]^-a \\
  Y \ar[r]^-b
    & \spec(k)
}\]
base change by $b$ shows that $\eR\pr_{2,\ast}{} \pr_1^\ast K$ is
$b^\ast \eR\Gamma(X,K)$. One then has (cf.
\cite[XVII 5.2.11]{sga4} and the proof)

\begin{align*}
  \eR\Gamma(X\times Y,\pr_1^\ast K\lotimes \pr_2^\ast L)
    &= \eR\Gamma(Y,\eR \pr_{2,\ast}(\pr_1^\ast K\lotimes \pr_2^\ast L)) \\
    &= \eR\Gamma(Y,(\eR\pr_{2,\ast}\pr_1^\ast K)\lotimes L) \\
    &= \eR\Gamma(Y,b^\ast \eR\Gamma(X,K)\lotimes L) \\
    &= \eR\Gamma(X,K)\lotimes \eR\Gamma(Y,L)
\end{align*}






\subsection{}\label{VII:1-12}

Finally, as a local counterpart of \ref{VII:1-9}(ii), we shall prove a
generic local acyclicity theorem (\ref{VII:2-13}).











\section{Generic theorems}\label{VII:2}

In this section, we prove \ref{VII:1-9} and the generic local acyclicity theorem
\ref{VII:2-13}. To prove \ref{VII:1-9}, we immediately reduce to the case where
$S$ is integral. Let $\eta$ be its generic point.





\subsection{}\label{VII:2-1}

We begin by proving \ref{VII:1-9} under the following additional assumptions:
$X$ is smooth over $S$, purely of relative dimension $n$, $A=\dZ/m$, $\sF$ is
locally constant, and $Y=S$. Let $\sF'=\Hom(\sF,\dZ/m)$. We shall show that if
the $\eR^i f_! \sF'$ are locally constant, the conclusions of \ref{VII:1-9}
hold with $U=S$. Recall that if $\sL$ is a constructible locally constant
sheaf, the $\Ext^i(\sL,\sG)$ may be computed fiber by fiber, and are
constructible if $\sG$ is. Recall also that $\dZ/m$ is an injective
$\dZ/m$-module and is the dualizing module, so that
$\sF=\rHom(\sF',\dZ/m)$. Poincaré duality
\[
  \eR f_\ast \rHom(K,\eR f^! L) = \rHom(\eR f_! K,L) \text{,} 
\]
for $K=\sF'$, $L=\dZ/m$, $\eR f^! L = \dZ/m[2n](n)$, therefore gives
\[
  \eR^{2n-i} f_\ast \sF = \Hom(\eR^i f_! \sF',\dZ/m)(-n) \text{.} 
\]

The sheaf on the right is locally constant and its formation commutes with
every base change, which proves the assertion.





\subsection{}\label{VII:2-2}

Let us prove \ref{VII:1-9} under the additional assumptions that $X$ is smooth
over $S$, $\sF$ is locally constant, and $Y=S$.
\begin{enumerate}[a)]
  \item By decomposing $X$ into connected components, we reduce to the case
    where it is purely of some relative dimension $n$ over $S$.
  \item By decomposing $A$ as a product, we reduce to the case where
    $\ell^m A=0$, with $\ell$ a prime invertible on $S$; we then filter $\sF$
    by the $\ell^k \sF$ and use the corresponding spectral sequence to reduce
    to the case $\ell\sF=0$. We may then replace $A$ by $A/\ell A$ and assume
    that $\ell A=0$.
  \item After replacing $S$ by a suitable $U$, we may assume that there exists
    a finite étale Galois covering $X_1/X$, with Galois group $G$, such that
    the inverse image $\sF_1$ of $\sF$ on $X_1$ is a constant sheaf, with
    constant value $F$. Writing $f_1$ for the projection of $X_1$ onto $S$,
    we then have
    \[
      \eR^i {f_1}_\ast \sF_1 = (\eR^i {f_1}_\ast \dZ/\ell)\otimes_{\dZ/\ell} F \text{.} 
    \]
    We also have the Hochschild--Serre spectral sequence (or the Leray
    spectral sequence for the covering $X_1/X$)
    \[
      \underline\h^p(G,\eR^q {f_1}_\ast \sF_1) \Rightarrow \eR^{p+q} f_\ast \sF \text{.} 
    \]
    By \ref{VII:2-1}, after shrinking $U$ if necessary, we may assume that
    the $\eR^i {f_1}_\ast \dZ/\ell$ are locally constant and that their
    formation commutes with every base change. The same property then holds
    for the $\eR^i f_\ast \sF$. Finally, if $\bar\eta$ is a geometric point
    localized at the generic point $\eta$ of $S$, almost all the
    $\eR^i f_\ast\sF$ vanish, because the
    $(\eR^i f_\ast \sF)_{\bar\eta} = \h^i(X_{\bar \eta},\sF)$ do.
\end{enumerate}





\subsection{}\label{VII:2-3}

We prove by induction on $n$ that
\begin{equation*}\tag{$\ast_n$}\label{VII:eq:2-3-1}
  \begin{array}{c}
    \text{The conclusions of \ref{VII:1-9} hold when $\dim{X_\eta}\leqslant n$} \\
    \text{and $f$ is an open immersion with dense image.}
  \end{array}
\end{equation*}

For $n=0$, after shrinking $S$ if necessary, we have $X=Y$: $(\ast_0)$ is
clear. Assume $(\ast_{n-1})$, and prove \eqref{VII:eq:2-3-1}. In
($\ast_{n-1}$), ``open immersion'' may be replaced by ``immersion,'' as is
seen by factoring it as an open immersion followed by a closed immersion.





\begin{lemma_}\label{VII:2-4}
After shrinking $S$ if necessary, the conclusions of \ref{VII:1-9} hold over
$Y'\subset Y$, where the complement $Y_1$ of $Y'$ is finite over $S$.
\end{lemma_}

The assertion is local on $Y$, which we may assume affine:
$Y\subset \dA_S^n$. The induction hypothesis ($\ast_{n-1}$) applies to
\[\xymatrix{
  X \ar[r]^-f \ar[dr] 
    & Y \ar[d]^-{\pr_i} \\
  & \dA_S^1 
}\]
Thus, for each $i$, there exists a dense open subset $U_i$ of $\dA_S^1$ such
that the conclusions of \ref{VII:1-9} hold over $\pr_i^{-1}(U_i)$; they hold
over the union of the $\pr_i^{-1}(U_i)$, and \ref{VII:2-4} follows.





\addtocounter{subsection}{1} % a section is skipped in the original
\subsection{}\label{VII:2-6}

Let us prove \eqref{VII:eq:2-3-1} when $X$ is smooth over $S$ and $\sF$ is
locally constant on $X$. Since the problem is local on $Y$, we may first
assume $Y$ affine and then projective (replacing $Y$ by its closure in a
projective space). Let $i:Y_1\hookrightarrow Y$ and $j:Y' \to Y$ be as in
\ref{VII:2-4}.
\[\xymatrix{
  X \ar@{^{(}->}[r] \ar[dr]_-a 
    & Y \ar[d]^-b 
    & \ar[l]_-i \ar[dl]^-{b_1} Y_1 \\
  & S 
}\] 
After shrinking $S$ if necessary, we know that $j^\ast \eR f_\ast \sF$ is
constructible and that its formation commutes with every base change in $S$;
we also know that $\eR a_\ast \sF = \eR b_\ast \eR f_\ast \sF$ is
constructible and that its formation commutes with every base change.

Apply $\eR b_\ast$ to the triangle defined by the exact sequence
\begin{equation*}\tag{1}\label{VII:eq:2-6-1}
\xymatrix{
  0 \ar[r] 
    & j_! j^\ast \eR f_\ast \sF \ar[r] 
    & \eR f_\ast \sF \ar[r] 
    & i_! i^\ast \eR f_\ast \sF \ar[r] 
    & 0 \text{:} 
}
\end{equation*}
we obtain a triangle
\begin{equation*}\tag{2}\label{VII:eq:2-6-2}
\xymatrix{
  \ar[r] 
    & \eR b_\ast j_! j^\ast \eR f_\ast \sF \ar[r] 
    & \eR a_\ast \sF \ar[r] 
    & {b_1}_\ast i^\ast \eR f_\ast \sF \ar[r] 
    & 
}
\end{equation*}
in which the first two terms are constructible and their formation commutes
with every base change in $S$ (for the first, by the finiteness theorem for
the proper morphism $b$). The same therefore holds for the third. Since $b_1$
is finite, it follows that $i^\ast \eR f_\ast \sF$ is constructible and its
formation commutes with every base change in $S$; the same holds for
$\eR f_\ast \sF$ by \eqref{VII:eq:2-6-1}.





\subsection{}\label{VII:2-7}

Let us prove \eqref{VII:eq:2-3-1} in general. We first reduce to the case where
$X$ contains a dense open subset $V$ smooth over $S$. When $S$ is the spectrum
of a perfect field, it is enough to replace $X$ by $X_\text{red}$ (and $Y$ by
$Y_\text{red}$). In general, we must shrink $S$, make a finite, radicial,
surjective base change $S' \to S$, and replace $X$ and $Y$ by $X_\text{red}$
and $Y_\text{red}$. Since the étale topology is insensitive to finite,
radicial, surjective morphisms, this is harmless. After shrinking $V$, we may
assume that $\sF$ is locally constant on $V$.
\[\xymatrix{
  V \ar@{^{(}->}[r]^j 
    & X \ar@{^{(}->}[r]^-f 
    & Y \text{.} 
}\]
Define $\Delta$ by the triangle
\begin{equation*}\tag{1}\label{VII:eq:2-7-1}
\xymatrix{
  \ar[r] 
    & \sF \ar[r] 
    & \eR j_\ast j^\ast \sF \ar[r] 
    & \Delta \ar[r] 
    & \text{.} 
}
\end{equation*}
The cohomology sheaves of $\Delta$ are supported on $X\setminus V$, and
$\dim(X\setminus V)_\eta < n$. The induction hypothesis therefore allows us
to assume that $\eR f_\ast \Delta$ is constructible. Applying $\eR f_\ast$ to
the triangle \eqref{VII:eq:2-7-1}, we obtain a triangle
\[\xymatrix{
  \ar[r] 
    & \eR f_\ast \sF \ar[r] 
    & \eR(f j)_\ast j^\ast \sF \ar[r] 
    & \eR f_\ast \Delta \ar[r] 
    & 
}\]
in which two of the vertices are constructible and their formation commutes
with every base change in $S$. The same therefore holds for the third.





\subsection{}\label{VII:2-8}

Let us prove \ref{VII:1-9}. The problem is local on $Y$, which we may assume
affine. Taking an affine covering of $X$ and invoking the Leray spectral
sequence, we reduce to the case where $X$ is affine as well. This ensures that
we can factor $f$ as an open immersion followed by a proper morphism:
$f=g j$, hence $\eR f_\ast = \eR g_\ast \eR j_\ast$. The open immersion is
covered by \eqref{VII:eq:2-3-1}, and the proper morphism by the finiteness
theorem.

The following corollaries are proved as in \hyperref[VII:1]{Section 1}.





\begin{corollary_}\label{VII:2-9}
Under the hypotheses of the theorem, for $K$ in $\eD_c^b(X,A)$ or
$\eD_c^b(Y,A)$ respectively, there exists a nonempty open subset $U$ of $S$
over which $\eR f_\ast K$, $\eR f_! K$, $f^\ast K$, and $\eR f^! K$ belong to
$\eD_c^b(Y,A)$ or $\eD_c^b(X,A)$, and their formation commutes with every base
change $S' \to U\subset S$.
\end{corollary_}





\begin{corollary_}\label{VII:2-10}
Let $X$ be as in the theorem, and suppose that $A$ is commutative. Then, if
$\sF$ and $\sG$ are constructible sheaves of $A$-modules on $X$, over a dense
open subset $U$ of $S$, the $\Ext_A^i(\sF,\sG)$ are again constructible and
their formation commutes with every base change $S' \to U\subset S$.
\end{corollary_}





\subsection{}\label{VII:2-11}

If $x$ is a geometric point of a scheme $X$, we denote by $X_x$ the strict
henselization of $X$ at $x$. For $f:X\to S$ and $t$ a geometric point of $S$,
we denote by $X_t$ the geometric fiber of $X$ at $t$. Finally, for $x$ a
geometric point of $X$ and $t$ a geometric point of $S_{f(x)}$, $(X_x)_t$ is
the fiber at $t$ of $X_x \to S_{f(x)}$.





\begin{definition_}\label{VII:2-12}
Let $f:X\to S$ and $K\in \ob\eD^+(X,A)$. We say that $f$ is \emph{locally
acyclic at $x$ relative to $K$} if, for every geometric point $t$ of
$S_{f(x)}$, we have $K_x=\eR\Gamma(X_x,K)\iso \eR\Gamma(X_{x,t},K)$. We say
that $f$ is \emph{locally acyclic relative to $K$} if this holds for every
$x$, and \emph{universally locally acyclic relative to $K$} if it remains true
after every base change $S' \to S$.
\end{definition_}





\begin{theorem_}\label{VII:2-13}
Let $f:X\to S$ be a morphism of finite type and let $\sF$ be constructible on
$X$. There exists a dense open subset $U$ of $S$ over which $f$ is universally
locally acyclic relative to $\sF$.
\end{theorem_}

We shall use the following result, proved in the appendix.





\begin{lemma_}\label{VII:2-14}
Let $X\xrightarrow f S_1 \xrightarrow g S_2$. If $g$ is smooth and $f$ is
universally locally acyclic relative to $K$, then so is $g f$.
\end{lemma_}

We reduce to the case where $S$ is integral, with generic point $\eta$, and
argue by induction on $\dim{X_\eta}$. We first deduce from the induction
hypothesis that
\begin{equation*}\tag{A}\label{VII:eq:2-14-1}
  \begin{array}{c}
    \text{After shrinking $S$, there exists $T\subset X$, finite over $S$, such that $f$} \\
    \text{is universally locally acyclic outside $T$.} 
  \end{array}
\end{equation*}

Since this question is local, we localize on $X$ and factor $f$ as
$X\xrightarrow u \dA_S^1 \to S$. The induction hypothesis applies
to $u$, and we conclude by \ref{VII:2-14} and the usual arguments.

To prove the theorem, we may assume that $X$ is proper, since the problem is
local. We shrink $S$ so that the $\eR^i f_\ast \sF$ are locally constant and
\eqref{VII:eq:2-14-1} applies, and use the following lemma.





\begin{lemma_}\label{VII:2-15}
Let $f:X\to S$ be proper and let $\sF$ be constructible. Suppose that the
$\eR^i f_\ast \sF$ are locally constant and that, outside $T\subset X$ finite
over $S$, $f$ is locally acyclic relative to $\sF$. Then $f$ is locally
acyclic relative to $\sF$.
\end{lemma_}

Let $s$ be a geometric point of $S$, let $S_s$ be the strict localization of
$S$ at $s$, and let $t$ be a geometric point of $S_s$. Let $X_{(s)}$ be the
inverse image of $X$ on $S_s$, let $\bar j:X_t \to X_{(s)}$, and let
$i:X_s \hookrightarrow X_{(s)}$. Again denoting by $\sF$ the inverse image of
$\sF$ on $X_{(s)}$ or $X_t$, we must prove that
$i^\ast \sF \iso i^\ast \eR \bar j_\ast \sF$. Let $\Delta$ be the mapping
cylinder of this map. Its cohomology sheaves are supported on $T_s$. To prove
that $\Delta=0$, it therefore suffices to prove that
$\eR\Gamma(X_s,\Delta)=0$. This is the mapping cylinder of
$(\eR f_\ast \sF)_s = \eR\Gamma(X_s,\sF) \to \eR\Gamma(X_s,i^\ast \eR\bar j_\ast \sF) = \eR\Gamma(X_{(s)},\eR\bar j_\ast \sF) = \eR\Gamma(X_t,\sF) = (\eR f_\ast \sF)_t$.
By hypothesis this specialization morphism is an isomorphism, and
$\Delta=0$.





\begin{corollary_}\label{VII:2-16}
When $S$ is the spectrum of a field, every $S$-scheme $X$ is universally
locally acyclic relative to $K$, for every $K\in\ob\eD^+(X,A)$.
\end{corollary_}

This follows from \ref{VII:2-13} by passage to the limit (which is possible
because the $U$ of \ref{VII:2-13} is always equal to $S$.








\section{Proof of \texorpdfstring{\ref{VII:1-1}}{1.1} and constructibility of the vanishing-cycle sheaves}\label{VII:3}






\subsection{}\label{VII:3-1}

Let $S$ be a strictly local trait (the spectrum of a henselian discrete
valuation ring with separably closed residue field). Denote its closed and
generic points by $s$ and $\eta$, and let $\bar\eta$ be a geometric point
localized at $\eta$ (a separable closure of $k(\eta)$). Let $X$ be over $S$
and let $\sF$ be a sheaf on $X_\eta$. Recall the definition of the
vanishing-cycle sheaves $\eR^i \Psi_\eta(\sF)$ (sheaves on $X_s$, endowed with
an action of $\gal(\bar\eta/\eta)$). Let $i:X_s\hookrightarrow X$,
$j:X_\eta \hookrightarrow X$, and let $\bar j$ be the composite
$X_{\bar\eta}\to X_\eta \hookrightarrow X$; then
\[
  \eR^i \Psi_\eta(\sF) = i^\ast \eR^i \bar j_\ast (\bar j^\ast \sF) \text{.} 
\]

Here is the main theorem of this subsection. It improves
\cite[XIII 2.3.1, 2.4.2]{sga7}.






\begin{theorem_}\label{VII:3-2}
Suppose that $X$ is of finite type over $S$ and that $\sF$ is constructible.
Then the vanishing-cycle sheaves $\eR^i \Psi_\eta(\sF)$ are constructible.
\end{theorem_}

We argue by induction on the dimension $n$ of $X_\eta$. We may moreover
assume---and do assume---that $X_\eta$ is dense in $X$ (replacing $X$ by
$\bar X_\eta$ does not change the $\eR^i \Psi$, which are supported on
$\bar X_\eta$).






\begin{lemma_}\label{VII:3-3}
If \ref{VII:3-2} holds when $X_\eta<n$, and if
$\dim X_\eta=n$, then there exist constructible subsheaves $\sG_i$ of the
$\eR^i \Psi_\eta(\sF)$ such that the supports of the local sections of
$\eR^i\Psi_\eta(\sF)/\sG_i$ are finite.
\end{lemma_}

Let $s'$ be a geometric generic point of the affine line over $s$, and let
$S'$ be the strict localization at $s'$ of the affine line $\dA_S^1$ over
$S$. It is again a strictly local trait, and uniformizers for $S$ are
uniformizers on $S'$.
\[\xymatrix{
  & \spec(k') \ar[rr] \ar[dd] \ar[dr] 
    & & (S',\eta',s') \ar[dr] \ar[dd] \\
  \bar\eta' \ar[ur] \ar[dr] 
    & & \dA_{\bar S}^1 \ar[d]  
    & & \dA_S^1 \ar[dl] \\
  & \bar\eta\ar@{^{(}->}[r] 
    & \bar S \ar[r] 
    & (S,\eta,s) 
}\]
Let $k'=k(\bar\eta)\otimes_{k(\eta)}k(\eta')$. If $\bar S$ is the
normalization of $S$ in $\bar\eta$, then $k'$ is the fraction field of the
strict localization of $\dA_{\bar S}^1$ at $s'$; it is therefore a field, and
$\gal(\bar\eta/\eta)=\gal(k'/\eta')$. Let $\bar\eta'$ be the spectrum of an
algebraic closure of $k'$. The Galois group $P=\gal(\bar\eta'/k')$ is a
pro-$p$ group, where $p$ is the characteristic exponent of $k(s)$.






\begin{lemma_}\label{VII:3-4}
For every scheme $X'$ over $S'$, and every sheaf $\sF$ on
$X_\eta'=X_{\eta'}'$, the vanishing-cycle sheaves for $X'/S'$ and for $X'/S$
are related by
$\eR^i \Psi_\eta(\sF) =\eR^i\Psi_{\eta'}(\sF)^P$.
\end{lemma_}

The following diagram compares the morphisms used in the definition of
$\eR\Psi$ for $X'/S'$ and $X'/S$:
\[\xymatrix{
  X_{\bar\eta'} \ar[rr] \ar[d] 
    & & X_{\eta'}' \ar@{=}[dd] \ar@{^{(}->}[dr] 
    & & S' \\ 
  X_{\eta'}\times_{\eta'}\spec(k') \ar@{=}[d] \ar[urr] 
    & & & X \ar[ur] \ar[dr] \\
  X_\eta' \ar[rr] 
    & & X_\eta' \ar@{^{(}->}[ur] 
    & & S
}\]
and the lemma follows from the Hochschild--Serre spectral sequence, since $P$
is a pro-$p$ group and $p$ is invertible in $A$.

Let us prove \ref{VII:3-3}. The question is local; this allows us to assume
that $X$ is affine, $X\subset \dA_S^n$. Let $f$ be one of the projections
$X\subset \dA_S^n\to \dA_S^1$. Let $X'$ be the ``localization''
$X_{\dA_S^1} S'$ of $X$, and let $\sF'$ be the inverse image of $\sF$ on $X'$.
\[\xymatrix{
  \sF'\text{ on }X' \ar[r] \ar[d]^-\lambda 
    & S' \ar[d]^-\lambda \\
  \sF \text{ on }X \ar[r]^-f 
    & \dA_S^1 \ar[r] 
    & S \text{.} 
}\]
On $X_s'=X_{s'}'$, we have
\[
  \lambda^\ast \eR^i\Psi_\eta(\sF) = \eR^i \Psi_\eta(\sF') = \eR^i\Psi_{\eta'}(\sF')^P \text{:} 
\]
this sheaf is constructible because the induction hypothesis applies to
$X'/S'$, and we use the following lemma, applied to
$X_s\subset \dA_S^n$ and to the vanishing-cycle sheaves.






\begin{lemma_}\label{VII:3-5}
Let $\dA^n$ be affine $n$-space over a field $k$, let $X\subset \dA^n$, let
$\sF$ be a sheaf of $A$-modules on $X$, and let $\bar\eta$ be a geometric
generic point of $\dA^1$. Let $X_{\bar\eta,i}$ be the geometric generic fiber
of $X\subset \dA^n\xrightarrow{\pr_i} \dA^1$, and let
$\sF_{\bar\eta,i}$ be the restriction of $\sF$ to $X_{\bar\eta,i}$. If the
$\sF_{\bar\eta,i}$ are constructible, there exists a constructible
$\sF'\subset \sF$ such that the local sections of $\sF/\sF'$ have finite
support.
\end{lemma_}

By passage to the limit, for each $i$
\begin{enumerate}[a)]
  \item There exists an étale neighborhood $U$ of $\bar\eta$ and a
    constructible sheaf $\sH$ on
    $X_{U,i} = X\times_{\dA^1,\pr_i} U$ whose restriction
    $\sF_{\bar\eta,i}$ to $X_{\bar\eta,i}$ is the given one.
  \item For a suitable $U$, the isomorphism
    $\sH_{\bar\eta,i} \iso \sF_{\bar\eta,i}$ comes from
    $u:\sH\to \sF_{U,i}$, where $\sF_{U,i}$ is the inverse image of $\sF$ on
    $X_{U,i}$. Let $\varphi:X_{U,i} \to X$. The morphism $u$ defines
    $v_i:\varphi_! \sH \to \sF$, with image $\sF_i'$. We have
    $(\sF/\sF_i')_{\bar\eta,i} = 0$, and for $\sF'$ we take the sum of the
    $\sF_i'$.
\end{enumerate}






\subsection{}\label{VII:3-6}

Let us prove \ref{VII:3-2} in dimension $n$. We may assume that $X$ is affine
(because the problem is local on $X_s$), and then projective over $S$ (by
taking the closure of $X$ in a projective embedding). We then have a spectral
sequence \cite[I 2.2.3]{sga7} (the Leray spectral sequence for $\bar j$, taking
into account the base-change theorem for the proper morphism $X\to S$).
\begin{equation*}\tag{1}\label{VII:eq:3-6-1}
  E_2^{p q} = \h^p(X_s,\eR^q\Psi_\eta(\sF)) \Rightarrow \h^{p+q}(X_{\bar\eta},\sF) \text{.} 
\end{equation*}
Let $\sG^q\subset \eR^q \Psi_\eta(\sF)$ be as in \ref{VII:3-3}, and let
$\sH^q$ be the quotient sheaf. We have $\h^i(X_s,\sH^q)=0$ for $i\ne 0$, and
for $\sH^q$, hence $\eR^q\Psi_\eta(\sF)$, to be constructible, it suffices
that $\h^0(X_s,\sH^q)$ be of finite type.

Compute \eqref{VII:eq:3-6-1} modulo modules of finite type (i.e. in the
quotient of the category of $A$-modules by the thick subcategory of
$A$-modules of finite type). The long cohomology exact sequence deduced from
$0 \to \sG^q \to \eR^q\Psi_\eta(\sF) \to \sH^q \to 0$, together with
\ref{VII:1-10} applied to $\sG^q$, gives
$E_2^{p q} \sim \h^p(X_s,\sH^q)$. In particular,
$E_2^{p q}\sim 0$ for $p\ne 0$,
$E_2^{0 q}\sim \h^q(X_{\bar\eta},\sF) \sim 0$, and this completes the proof.

Parallel arguments give the following result, which improves
\cite[XIII 2.1.12, 2.4.2]{sga7}.






\begin{proposition_}\label{VII:3-7}
The formation of the vanishing-cycle sheaves commutes with changes of traits.
\end{proposition_}

Let $g:(S',\eta',s') \to (S,\eta,s)$ be a (surjective) morphism of strictly
local traits, let $X/S$, let $\sF$ be on $X_\eta$, of torsion prime to the
residual characteristic, and let $(X',\sF')$ be their inverse images on $S'$.
Again denoting by $g$ the morphisms parallel to $g$, such as
$X_{s'}' \to X_s$, we must prove that
\begin{equation*}\tag{3.7.1}\label{VII:eq:3-7-1}
\xymatrix{
  g^\ast \eR\Psi_\eta(\sF) \ar[r]^-\sim 
    & \eR\Psi_{\eta'}(\sF') \text{.} 
}
\end{equation*}
This morphism is defined by the commutative diagram
\[\xymatrix{
  X_{s'}' \ar@{^{(}->}[r]^-{i'} \ar[d]^-g 
    & X' \ar[d]^-g 
    & X_{\eta'}' \ar[l]_-{\bar j'} \ar[d]^-g \\
  X_s \ar@{^{(}->}[r]^-i 
    & X 
    & \ar[l]_-{\bar j} X_{\bar\eta} 
}\]

A passage to the limit reduces us to the case where $X$ is of finite type
over $S$, and we argue by induction on $\dim{X_\eta}$. Since the question is
local, we may assume that $X$ is affine, and then projective. Assume
\eqref{VII:eq:3-7-1} when $\dim{X_\eta}<n$. When $\dim{X_\eta}=n$, it then
follows from the following two assertions, where $\Delta$ is the mapping
cylinder of \eqref{VII:eq:3-7-1}.

\begin{enumerate}[(A)]
  \item The support of the local sections of the cohomology sheaves of
    $\Delta$ is finite.
    \phantomsection\hypertarget{VII:3-A}{}
  \item $\eR\Gamma(X_s,\Delta)=0$, hence $\Delta=0$, by
    (\hyperlink{VII:3-A}{A}).
    \phantomsection\hypertarget{VII:3-B}{}
\end{enumerate}





\subsection{}\label{VII:3-8}

The proof of (\hyperlink{VII:3-A}{A}) parallels that of \ref{VII:3-3}. We
localize on $X$ and consider maps $X\to \dA_S^1$; denoting by $S(x)$ the
strict henselization of $\dA_S^1$ at the generic point of the special fiber,
we apply the induction hypothesis to $X_1/S(x)$
($X_1=X\times_{\dA^1} S(x)$) and to $S'(x)$. Taking invariants under a pro-$p$
group, we find that $\Delta$ vanishes on the geometric generic fiber of
$X_{s'} \to \dA_s^1$. Since this holds locally on $X$ and for every
projection, we obtain (\hyperlink{VII:3-A}{A}).






\subsection{}\label{VII:3-9}

For (\hyperlink{VII:3-B}{B}), the properness of $X$ over $S$ ensures that
\[
  \eR\Gamma(X_{s'}',\eR\Psi_{\eta'}(\sF')) = \eR\Gamma(X_{\bar\eta}',\sF') 
\]
and
\[
  \eR\Gamma(X_s,\eR\Psi_\eta(\sF)) = \eR\Gamma(X_{\bar\eta},\sF) \text{.} 
\]
We have
$\eR\Gamma(X_s,\eR\Psi_\eta(\sF))\iso \eR\Gamma(X_{s'}',g^\ast\eR\Psi_\eta(\sF))$
(invariance under change of the separably closed base field); thus
$\eR\Gamma(X_{s'}', \Delta)$ is the mapping cylinder of
$\eR\Gamma(X_{\bar\eta},\sF) \to \eR\Gamma(X_{\bar\eta}',\sF')$, an
isomorphism by the same invariance theorem.






\subsection{}\label{VII:3-10}

Let us prove \ref{VII:1-1}. We reduce to the case where $S$ is connected,
hence integral. If $\dim S=0$ ($S$ is the spectrum of a field), apply
\ref{VII:1-9}, \ref{VII:1-10}. If $S$ has dimension $1$, Theorem
\ref{VII:1-9} ensures that the $\eR^i f_\ast \sF$ are constructible over the
complement of a finite set $T$ of points of $S$. It remains to see that their
restrictions to the $Y_t$ ($t\in T$) are constructible. It suffices to see
this after localization at $t$: we may assume that $S$ is a strictly local
trait. The essential case is the following.






\begin{lemma_}\label{VII:3-11}
Let $X$ be of finite type over $S$, let $j:X_\eta\hookrightarrow X$,
$i:X_s\hookrightarrow X$, and let $\sF$ be constructible on $X_\eta$. Then
$i^\ast\eR j_\ast \sF$ is constructible.
\end{lemma_}

Let $I=\gal(\bar\eta/\eta)$ be the inertia group. It is known to be an
extension of $\widehat\dZ_{p'}(1)$ by a pro-$p$ group $P$ ($p$ being the
characteristic exponent of the residue field). The Hochschild--Serre spectral
sequence gives
\[
  E_2^{p q} = \sH^p(I,\eR^q\Psi_\eta(\sF))\Rightarrow i^\ast \eR^{p+q}j_\ast\sF \text{.} 
\]
Let $R_t^q = \eR^q\Psi_\eta(\sF)^P$ ($t$ for tame). If $\sigma$ is a
generator of $\widehat\dZ_{p'}(1)$, then
$E_2^{0 q} = \ker(\sigma-1,R_t^q)$,
$E_2^{1 q} = \operatorname{coker}(\sigma-1,R_t^q)$, and $E_2^{p q} = 0$ for
$p\ne 0,1$. These sheaves are constructible because the vanishing-cycle
sheaves are, and \ref{VII:3-11} follows.






\subsection{}\label{VII:3-12}

Let us treat the general case, with $S$ still a strictly henselian trait. If
$X=X_\eta$, then $f$ is the composite $X\to Y_\eta \to Y$, and we apply
\ref{VII:1-9} and \ref{VII:3-11}. In the general case, let
$j:X_\eta\hookrightarrow X$ and let $\Delta$ be the mapping cylinder of
$\sF \to \eR j_\ast j^\ast \sF$. Its cohomology sheaves are supported on
$X_s$; by \ref{VII:1-9} applied to $X_s \to Y_s$,
$\eR f_\ast \Delta$ is therefore constructible. The complex
$\eR f_\ast \eR j_\ast j^\ast \sF = \eR(f j)_\ast j^\ast \sF$ is also
constructible, and we conclude from the triangle
\[\xymatrix{
  \ar[r] 
    & \eR f_\ast \sF \ar[r] 
    & \eR f_\ast \eR j_\ast j^\ast \sF \ar[r] 
    & \eR f_\ast \Delta \ar[r] 
    & \text{.} 
}\]










\section{Local biduality}\label{VII:4}






\subsection{}\label{VII:4-1}

Let $S$ be regular of dimension $0$ or $1$, and put $K_S=A$ (the constant
sheaf with value $A$). For $K\in\ob\eD_\text{ctf}^b(S,A)$, put
$D K=\rHom(K,K_S)\in\ob\eD^+(S,A^\circ)$ (a complex of sheaves of right
$A$-modules). We still have $D K\in\ob\eD_\text{ctf}^b(S,A^\circ)$, and an explicit
calculation, possible because $S$ has dimension $1$, shows that
$K\iso D D K$ (for $A=\dZ/n$, \hyperref[V]{Duality} \ref{V:1-4}).





\subsection{}\label{VII:4-2}

We assume that $A$ is commutative, an assumption that is probably
unnecessary. For $a:X\to S$ of finite type over $S$, put
$K_X=\eR a^! K_S$. For $K\in\ob\eD_\text{ctf}^b(X,A)$, put
$D K=\rHom(K,K_X)\in\ob\eD_\text{ctf}^b(X,A)$. Our main result is the
following.





\begin{theorem_}\label{VII:4-3}
We have $K\iso D D K$ ($K\in \ob\eD_\textnormal{ctf}^b(X,A)$).
\end{theorem_}





\begin{lemma_}\label{VII:4-4}
If $a$ is proper, then $\eR a_\ast K \iso \eR a_\ast D D K$.
\end{lemma_}

Poincaré duality
\[
  \eR a_! \rHom(K,\eR a^! K_S) = \rHom(\eR a_! K,K_S) 
\]
gives us an isomorphism $\eR a_\ast D = D \eR a_! = D \eR a_\ast$. The lemma
follows from the commutativity of the diagram (\hyperref[V]{Duality},
\ref{V:1-2})
\[\xymatrix{
  \eR a_\ast K \ar[r] \ar[d]^-\wr 
    & \eR a_\ast D D K \ar[d]^-\wr \\
  D D \eR a_\ast K \ar[r]^-\sim 
    & D \eR a_\ast D K \text{.} 
}\]





\subsection{}\label{VII:4-5}

By localization, it is enough to treat the cases where $S$ is the spectrum
of a field or a trait. When $S$ is the spectrum of a field, we argue by
induction on $\dim X$. Since the problem is local, we may assume that $X$ is
proper. Moreover, the induction hypothesis, by the usual arguments, ensures
that the mapping cylinder $\Delta$ of $K\to D D K$ has skyscraper cohomology
sheaves. The lemma then ensures that $\eR a_\ast \Delta=0$, hence that
$\Delta=0$.





\subsection{}\label{VII:4-6}

When $S$ is a trait, we already know that \ref{VII:4-3} holds on the generic
fiber (\ref{VII:4-5}). We may also assume that $X$ is proper over $S$, and we
argue by induction on $\dim{X_s}$, with $X$ proper over $S$. The mapping
cylinder $\Delta$ again has skyscraper cohomology sheaves on $X_s$, and we
conclude as above.





\subsection{}\label{VII:4-7}

Suppose that $A=\dZ/m$. We then have a variant of the preceding theory in
which $K$ is arbitrary in $\eD_c^b(X,A)$. The complex $\eR a^! \dZ/m$ has
finite injective dimension \cite[XVIII 3.1.7]{sga4}. This ensures that $D K$
still belongs to $\eD_c^b(X,A)$. To prove local biduality over $S$, we use the
fact that $\dZ/m$ is the dualizing $\dZ/m$-module. We then proceed as above.










\section{Appendix, by L.\ Illusie}\label{VII:5}

In this appendix, which takes up certain parts of the late
\cite[II]{sga5}, and a letter from P.\ Deligne to the author, we prove
\ref{VII:2-14}, as well as various generalizations and variants of the
specialization theorems \cite[XVI 2.1]{sga4} and
(\hyperref[I]{Étale cohomology} \ref{I:5-1-7}).

The schemes and morphisms considered will be assumed quasi-compact and
quasi-separated. If $X$ is a scheme and $x$ a geometric point of $X$, we
denote by $X(x)$ the strict localization of $X$ at $x$.

Fix a base scheme $S$ and a sheaf of rings $A$ on $S$ (not necessarily
commutative or noetherian). If $X$ is an $S$-scheme, we write $\eD(X)$ for
$\eD(X,A_X)$.





\subsection{Cohomological properness}\label{VII:5-1}






\begin{proposition}\label{VII:5-1-1}
Let $f:X\to Y$ be an $S$-morphism and let $E\in\ob \eD^+(X)$. The following
conditions are equivalent:
\begin{enumerate}[\indent (i)]
  \item The formation of $\eR f_\ast E$ commutes with every finite base change
    $S'\to S$.
  \item The formation of $\eR f_\ast E$ commutes with every quasi-finite base
    change (or inverse limit of quasi-finite morphisms) $S'\to S$.
  \item For every specialization arrow $i:t\to S(s)$
    (\hyperref[I]{Étale cohomology}, \ref{I:5-1-2}), if $\bar S$ denotes
    the normalization in $k(\bar t)$ of the integral scheme that is the
    closure of $i(t)$ in $S(s)$ (by
    \cite[IV, 6.15.5, 18.6.11, 18.8.16]{ega}, $\bar S$ is therefore local and
    integral, with closed point radicial over $s$), if
    $\bar f:\bar X\to \bar Y$ is the morphism deduced from $f$ by the base
    change $\bar S\to S$, and if $f_s:X_s\to Y_s$ is the fiber of $f$ at $s$,
    then the base-change arrow
    \[
      (\eR \bar f_\ast (E|\bar X))|Y_s \to \eR {f_s}_\ast(E|X_s) 
    \]
    is an isomorphism.
\end{enumerate}
\end{proposition}

The equivalence of (i) and (ii) follows from Zariski's Main Theorem, and it is
clear that (ii) implies (iii). Let us prove that (iii) implies (ii). For every
$S'\to S$ that is an inverse limit of quasi-finite morphisms, and every
geometric point $s$ of $S'$, denote by $C(S',s)$ the cone of the base-change
arrow
\[
  (\eR f'_\ast(E|X'))|Y_s \to \eR{f_s}_\ast (E|X_s) \text{,} 
\]
where $f':X'\to Y'$ is deduced from $f$ by the base change $S'\to S$, and
$f_s:X_s\to Y_s$ is the fiber of $f$ at $s$. We must prove that, for every
$(S',s)$ as above, $C(S',s)$ is acyclic. We shall prove by induction on $N$ that
$\h^i C(S',s) = 0$ for $i\leqslant N$. By the Main Theorem, passage to the
limit, and elimination of radicial morphisms, we may assume that $S$ and $S'$
are strictly local, that $S'$ is integral with closed point $s$, and that
$S'\to S$ is finite. Let $t$ be a geometric generic point of $S'$, and let
$\bar S$ be the normalization of $S'$ in $k(t)$. Let $\bar S_\bullet$ be the
simplicial scheme that is the coskeleton of $\bar S\to S'$
($\bar S_0=\bar S$,\ldots $\bar S_n=(\bar S/S')^{n+1}$,\ldots), and let
$\bar X_\bullet=X'\times_{S'} \bar S_\bullet$. Let
$g_\bullet:\bar X_\bullet \to X'$ be the canonical projection. Put
$E|X'=E'$. The bicomplex ${g_\bullet}_\ast g_\bullet^\ast E'$, whose columns
are the ${g_n}_\ast g_n^\ast E'$, is a resolution of $E'$
\cite[VIII 8]{sga4}. We may therefore represent $C(S',s)$ by a bicomplex whose
columns identify with the $C(\bar S_n,s)$. By (iii),
$C(\bar S_0,s)=C(\bar S,s)$ is acyclic. On the other hand, for $n>0$, we have
$\h^i C(\bar S_n,s)=0$ for $i\leqslant N$ by the induction hypothesis. It
follows that $\h^i C(S',s)=0$ for $i\leqslant N+1$, which completes the proof.





\subsubsection{Remark}\label{VII:5-1-2}

Let $(f,E)$ satisfy condition (i) of \ref{VII:5-1-1}.
\begin{enumerate}[a)]
  \item If $E\in \ob \eD^b(X)$ and $f_\ast$ has finite cohomological
    dimension, then for every $F\in \ob\eD^-(S,A^\circ)$, the canonical arrow
    \begin{equation*}\tag{$*$}\label{VII:eq:5-1-1-1}
      \eR f_\ast(E)\lotimes_A q^\ast F \to \eR f_\ast(E\lotimes_A p^\ast F) \text{,} 
    \end{equation*}
    where $p:X\to S$ and $q:Y\to S$ are the projections, is an isomorphism.
    The same conclusion holds with $E\in\ob\eD^-(X)$.
\end{enumerate} 

Indeed, by dévissage, it is enough to verify the assertion for $F$ of the
form $i_! A_{S'}$, with $i:S'\to S$ quasi-finite. By the Main Theorem, we
reduce to the case where $i$ is finite, and the conclusion is then equivalent
to the commutation of $\eR f_\ast(E)$ with base change by $i$.
\begin{enumerate}[a)]
\setcounter{enumi}{1}
  \item Suppose that $A$ is noetherian and constructible. Then, for every
    constructible $F\in\ob\eD^b(S,A^\circ)$ of finite tor-dimension, the arrow
    \eqref{VII:eq:5-1-1-1} above is an isomorphism.
\end{enumerate}

Indeed, by dévissage, we reduce to the case where $F$ is of the form $i_! M$,
with $i:S'\to S$ locally closed, $A$ locally constant on $S'$, and $M$
constructible, of finite tor-dimension, and with locally constant cohomology
on $S'$. After base change from $S$ to $S'$, we are thus reduced to the case
where $F$ is perfect \cite[I]{sga6}, hence finally, by localization and
dévissage, to the case $F=A$, for which the conclusion is immediate.





\subsubsection{}\label{VII:5-1-3}

Let $f:X\to Y$ be an $S$-morphism and let $E\in\ob\eD^+(X)$. We say that
$(f,E)$ is \emph{cohomologically proper} rel.\ to $S$ if the formation of
$\eR f_\ast E$ commutes with every base change $S'\to S$. Equivalently, after
every base change $S'\to S$, the formation of $\eR f_\ast' E'$ (where
$f'=f\times_S S'$ and $E'=$ the inverse image of $E$ on $X\times_S S'$)
commutes with every finite base change $S'' \to S'$ (cf.
\cite[XII 6.1]{sga4}). Here are some examples of cohomological properness
(for a study of the analogous notion for sheaves of sets or of
noncommutative groups, the reader may consult \cite[XIII]{sga1}).
\begin{enumerate} % originally 1.3.1, 1.3.2, ...
  \item Suppose that $A$ is torsion and $f$ is proper. Then, for every
    $E\in\ob\eD^+(X)$, $(f,E)$ is cohomologically proper rel.\ to $S$
    (proper base-change theorem, \cite[XII 5.1]{sga4}).
  \item Suppose that $S$ is finite, that $X$ and $Y$ are of finite type over
    $S$, and that $A$ is annihilated by an integer invertible on $S$. Then,
    for every $E\in\ob\eD^+(X)$, $(f,E)$ is cohomologically proper rel.\ to
    $S$ (\ref{VII:1-9}).
  \item Suppose that $A$ is constant, noetherian, and annihilated by an
    integer $n$ invertible on $S$, and that
    $f:X=Y\setminus D\to Y$ is the inclusion of the complement of a divisor
    $D$ on $Y$ with normal crossings relative to $S$
    \cite[XIII 2.1]{sga1}. Let $E$ be a sheaf of $A$-modules on $X$ satisfying
    one of the following conditions:
    \begin{enumerate}[(i)]
      \item $E$ is constructible locally constant and tamely ramified along
        $D$;
      \item $E$ is, locally for the étale topology on $Y$ and on $S$, the
        inverse image of a sheaf of $A$-modules on $S$.
    \end{enumerate}
    Then $(f,E)$ is cohomologically proper rel.\ to $S$, and $\eR f_\ast E$
    is constructible (compare \cite[XIII 2.4]{sga1}).
\end{enumerate} 

Let us sketch the proof in case (i). The question is local near a point $y$
of $Y$. We may assume that $D$ is a sum of smooth divisors,
$D=\sum_{i=1}^m D_i$. By the relative Abhyankar lemma
\cite[XIII 5.5]{sga1}, near $y$ we may trivialize $E$ by a finite covering
$\widetilde Y\to Y$ of the form
$\widetilde Y=Y[T_1,\dots,T_r]/(T_1^{n_1}-t_1,\dots,T_r^{n_r}-t_r)$, where the
$t_i$ are local equations of the smooth divisors passing through $y$, and the
$n_i$ are integers prime to the char.\ of $k(y)$. The inverse image
$\widetilde D$ of $D$ in $\widetilde Y$ still has normal crossings rel.\ to $S$, and
if $g:\widetilde X=\widetilde Y\setminus \widetilde D\to X$ is the projection,
$g^\ast E$ is constant. Since $E$ injects into $g_\ast g^\ast E$ and the
quotient is tamely ramified, an easy dévissage reduces us to the case where
$E$ is constant. Denote the projections by $p:X\to S$, $q:Y\to S$, and, for
$1\leqslant i\leqslant m$, let $f_i:Y\setminus D_i \to Y$ be the canonical
inclusion. From the relative purity theorem \cite[XVI 3.7]{sga4}, one easily
deduces by induction on $m$ that, for every locally constant $A$-module $M$ on
$S$, the canonical arrow
\begin{equation*}\tag{5.1.3.1}\label{VII:eq:5-1-3-3-1} % originally 1.3.3.1
  \eR f_\ast (\dZ/n)\lotimes q^\ast M\to \eR f_\ast (p^\ast M) 
\end{equation*}
is an isomorphism, and that on the other hand we have:
\begin{align*}\tag{5.1.3.2}\label{VII:eq:5-1-3-3-2}
  {f_i}_\ast(\dZ/n) &= \dZ/n \text{,} \\
  \eR^1 {f_i}_\ast (\dZ/n) &= (\dZ/n)(-1)_{D_i} \text{,} 
\end{align*}
where a canonical isomorphism is given by the fundamental class $\cl(D_i)$
(\hyperref[IV]{Cycle} \ref{IV:2-1-4}),
\begin{align*}
  \eR^q {f_i}_\ast(\dZ/n) &= 0\text{ for $q>1$,} \\
  \eR f_\ast (\dZ/n) &= \bigotimes_i^\eL \eR {f_i}_\ast (\dZ/n) \text{ , i.e.:} \\
  \eR^1 f_\ast(\dZ/n) &= \bigoplus_{i=1}^m (\dZ/n)(-1)_{D_i} \text{,} \\
  \textstyle\bigwedge^\bullet \eR^1 f_\ast(\dZ/n) &\iso \eR^\bullet f_\ast (\dZ/n) \text{.} 
\end{align*}
The conclusion of 3, in case (i), therefore follows from
\eqref{VII:eq:5-1-3-3-1} and \eqref{VII:eq:5-1-3-3-2}. In case (ii), by
passage to the limit we reduce to the case where $E=p^\ast M$, with $M$
constructible, and then, by dévissage and base change (using (i)), to the
already treated case where $M$ is constant.

It follows from the preceding proof that \eqref{VII:eq:5-1-3-3-1} is an
isomorphism for every $M\in\ob\eD^+(S)$.

One may compare the formulas \eqref{VII:eq:5-1-3-3-2} with the analogous
formulas in integral cohomology when $S=\spec(\dC)$ (see, for example,
\cite[3.1]{de71-2}), which follow from the fact that, locally for the classical
topology, the complement of $D$ has the homotopy type of a torus.

By passage to the limit, one deduces from \eqref{VII:eq:5-1-3-3-2} similar
formulas in $\ell$-adic cohomology, i.e. with $\dZ/n$ replaced by $\dZ_\ell$
($\ell$ a prime invertible on $S$). Suppose that $q:Y\to S$ is proper and
smooth. It then follows from the Weil conjectures \cite{de74,de80} that the
Leray spectral sequence
\[
  E_2^{i j} = \eR^i q_\ast \eR^j f_\ast \dZ_\ell \Rightarrow \eR^\bullet p_\ast \dZ_\ell 
\]
degenerates at $E_3$ modulo torsion \cite[\S 6]{de71-1}. The same phenomenon
occurs in integral cohomology when $S=\spec(\dC)$; see (loc.\ cit.) and
\cite[3.2.13]{de71-2}.





\subsection{Specialization and cospecialization}\label{VII:5-2}

Let $f:X\to S$ and $E\in\ob\eD^+(X)$.






\subsubsection{}\label{VII:5-2-1}

Suppose that the formation of $\eR f_\ast E$ commutes with every finite base
change. For every specialization $u:t\to S(s)$ of geometric points of $S$,
we then define an arrow, called the \emph{specialization arrow},
\begin{equation*}\tag{5.2.1.1}\label{VII:eq:5-2-1-1}
  \spe(u):\eR\Gamma(X_s,E|X_s) \to \eR\Gamma(X_t,E|X_t) \text{,}
\end{equation*}
as follows. Let, as in \ref{VII:5-1-1}(iii), $\bar S$ be the normalization in
$k(t)$ of the closure of $u(t)$ in $S(s)$, and let
$\bar f:\bar X\to \bar S$ be the morphism deduced from $f$ by the base change
$\bar S\to S$. By the immediate implication (i)$\Rightarrow$(iii) of
\ref{VII:5-1-1}, the base-change arrow
\begin{equation*}\tag{5.2.1.2}\label{VII:eq:5-2-1-2}
  \eR\Gamma(\bar X,E|\bar X)=\eR\bar f_\ast (E|\bar X)_s \to \eR\Gamma(X_s,E|X_s) 
\end{equation*}
is an isomorphism. We define \eqref{VII:eq:5-2-1-1} as the composite of the
inverse of \eqref{VII:eq:5-2-1-2} with the restriction arrow
\begin{equation*}\tag{5.2.1.3}\label{VII:eq:5-2-1-3}
  \eR\Gamma(\bar X,E|\bar X) \to \eR\Gamma(X_t,E|X_t) \text{.}
\end{equation*}






\subsubsection{}\label{VII:5-2-2}

Suppose that $f$ is locally acyclic rel.\ to $E$ (\ref{VII:2-12}). For
every specialization $u:t\to S(s)$ of geometric points of $S$, we then define
an arrow, called the \emph{cospecialization arrow},
\begin{equation*}\tag{5.2.2.1}\label{VII:eq:5-2-2-1}
  \cosp(u):\eR\Gamma(X_t,E|X_t) \to \eR\Gamma(X_s,E|X_s) \text{,} 
\end{equation*}
as follows. Let $g:X_t\to X$ and $g':\bar X\to X$ be the canonical arrows.
Viewed vis-\`a-vis these two maps, the canonical arrow
\begin{equation*}\tag{5.2.2.2}\label{VII:eq:5-2-2-2}
  \eR g'_\ast(E|\bar X)\to \eR g_\ast(E|X_t) 
\end{equation*}
is an isomorphism by the local-acyclicity hypothesis. Indeed, we may assume
that $S$ is strictly local and integral, with closed point $s$, and that $t$
is a geometric generic point. Compute the fiber of
\eqref{VII:eq:5-2-2-2} at a geometric point $x$ of $X$, with image $y$ in
$S$. After replacing $S$ by the strict localization of $S$ at $y$, we may
assume that $y=s$. Since $\bar S\to S$ is an inverse limit of finite
morphisms and is radicial at $s$, we find that
$\eR g'_\ast(E|\bar X)_x = E_x$ and that the fiber of
\eqref{VII:eq:5-2-2-2} at $x$ identifies with the restriction
$E_x=\eR\Gamma(X(x),E)\to \eR\Gamma(X(x)_t,E|X(x)_t)$, which is an
isomorphism by hypothesis. Applying $\eR\Gamma(X,-)$ to
\eqref{VII:eq:5-2-2-2}, we deduce that the restriction arrow
\eqref{VII:eq:5-2-1-3} is an isomorphism. We define
\eqref{VII:eq:5-2-2-1} as the composite of the inverse of
\eqref{VII:eq:5-2-1-3} with the base-change arrow
\eqref{VII:eq:5-2-1-2}.






\subsubsection{}\label{VII:5-2-3}

If $s'' \xrightarrow{u'} s' \xrightarrow u s$ are specialization arrows of
geometric points of $S$, one verifies that
\begin{align*}
  \spe(u')\spe(u) &= \spe(u u') \text{,} \\
  \cosp(u) \cosp(u') &= \cosp(u u') \text{,} 
\end{align*}
whenever both sides are defined.






\subsubsection{}\label{VII:5-2-4}

Suppose that the formation of $\eR f_\ast E$ commutes with every finite base
change and that $f$ is locally acyclic rel.\ to $E$. It follows from the
definitions that, for every specialization $u:t\to S(s)$ of geometric points
of $S$, the arrows $\spe(u)$ and $\cosp(u)$ are isomorphisms, inverse to one
another.

The preceding hypotheses hold in particular when $A$ is noetherian and
annihilated by an integer $n$ invertible on $S$, when $f$ is proper and smooth,
and when $E$ has locally constant cohomology (cohomological properness of
proper morphisms (\ref{VII:5-1-3}), and local acyclicity of smooth morphisms
\cite[XV]{sga4}). This recovers the ``specialization theorem''
\cite[XVI 2.1]{sga4}.

On the other hand, \ref{VII:5-1-1} has the following consequence:






\begin{proposition}\label{VII:5-2-5}
Suppose that $f$ is locally acyclic relative to $E$ and that, for every
specialization $u:t\to S(s)$, the cospecialization arrow $\cosp(u)$ is an
isomorphism. Then the formation of $\eR f_\ast E$ commutes with every finite
base change.
\end{proposition}

In view of \ref{VII:5-2-4}, \ref{VII:5-2-5} generalizes
(\hyperref[I]{Étale cohomology} \ref{I:5-1-7}).






\begin{corollary}\label{VII:5-2-6}
Suppose that $f$ is locally acyclic rel.\ to $E$. Let $x$ be a geometric
point of $X$, with image $s$ in $S$, and denote by
$f_{(x)}:X(x)\to S(s)$ the morphism induced by $f$. Then the formation of
$\eR {f_{(x)}}_\ast(E|X(x))$ commutes with every finite base change
$S'\to S(s)$.
\end{corollary}

By \ref{VII:5-2-5}, it is enough to verify that the cospecialization arrows
for $(f_{(x)},E|X(x))$ are isomorphisms. By the transitivity of the
cospecialization arrows (\ref{VII:5-2-3}), it is enough to show that if
$u:t\to S(s)$ is a specialization, then
$\cosp(u):\eR\Gamma(X(x)_t,E|X(x)_t)\to \eR\Gamma(X(x)_s,E|X(x)_s)$ is an
isomorphism. But since $s$ is closed in $S(s)$, $X(x)_s$ is strictly local
with closed point $x$, $\eR\Gamma(X(x)_s,E|X(x)_s)=E_x$, and $\cosp(u)$ is an
isomorphism, inverse to the restriction isomorphism
$E_x\to \eR\Gamma(X(x)_t,E|X(x)_t)$.






\begin{corollary}\label{VII:5-2-7}
Let $f:X\to Y$ be an $S$-morphism, let $p:X\to S$, $q:Y\to S$ be the
projections, and let $E\in\ob\eD^+(X)$. Suppose that $A$ is locally constant,
that $f$ is locally acyclic rel.\ to $E$, and that $q$ is locally acyclic at
every geometric point $y$ of $Y$ rel.\ to every sheaf locally constant
in a neighborhood of $y$. Then $p$ is locally acyclic relative to $E$.
\end{corollary}

Let $x$ be a geometric point of $X$, with images $y$ and $s$ in $Y$ and $S$,
let $t\to S(s)$ be a specialization, and let
$f_{(x)}:X(x)\to Y(y)$ be the morphism induced by $f$. We must show that the
restriction
\begin{equation*}\tag{$*$}\label{VII:eq:5-2-7-1}
  E_x \to \eR\Gamma(X(x)_t,E|X(x)_t) 
\end{equation*}
is an isomorphism. We may assume that $A$ is constant. By \ref{VII:5-2-4} and
\ref{VII:5-2-5}, for every specialization $z\to Y(y)$, the specialization
arrow
\[
  E_x = \eR {f_{(x)}}_\ast (E|X(x))_y \to (\eR {f_{(x)}}_\ast (E|X(x)))_z = \eR\Gamma(X(x)_z,E|X(x)_z) 
\]
is an isomorphism; in other words, $\eR{f_{(x)}}_\ast(E|X(x))$ is constant
with value $E_x$. The acyclicity hypothesis on $q$ therefore plays the
decisive r\^ole and ensures that the restriction
\[
  E_x \to \eR\Gamma(Y(y)_t,(E_x)|Y(y)_t) = \eR\Gamma\left(Y(y)_t,\eR {f_{(x)}}_\ast (E|X(x))|Y(y)_t\right)
\]
is an isomorphism. But this arrow identifies with
\eqref{VII:eq:5-2-7-1}, which proves the assertion.

The hypothesis of \ref{VII:5-2-7} on $q$ holds in particular when $q$ is
smooth and $A$ is annihilated by an integer $n$ invertible on $S$
\cite[XV]{sga4}. On the other hand, if the hypotheses of \ref{VII:5-2-7} hold
after every base change $S'\to S$, then it follows that $p$ is universally
locally acyclic rel.\ to $E$. These remarks imply, in particular, the
statement of \ref{VII:2-14}.






\begin{corollary}\label{VII:5-2-8}
Let $f:X\to S$ and $E\in \ob\eD^+(X)$. Suppose that $f$ is universally locally
acyclic rel.\ to $E$, that the formation of $\eR f_\ast E$ commutes with
every finite base change, and that, for every morphism $s'\to s$, where $s$ is
a geometric point of $S$ and $s'$ is the spectrum of a separably closed
field, the base-change arrow
$\eR\Gamma(X_s,E|X_s)\to \eR\Gamma(X_{s'},E|X_{s'})$ is an isomorphism. Then
$(f,E)$ is cohomologically proper (\ref{VII:5-1-3}).
\end{corollary}

Let $f':X'\to S'$ be the morphism deduced from $f$ by a base change
$S'\to S$; we must show that the formation of $\eR f'_\ast(E|X')$ commutes
with every finite base change. By \ref{VII:5-2-5}, it is enough to show that,
for every specialization $u':t'\to S'(s')$, the cospecialization arrow
$\cosp(u')$:
$\eR\Gamma(X_{t'}',E|X_{t'}') \to \eR\Gamma(X_{s'}',E|X_{s'}')$ is an
isomorphism. Let $s$ (resp.\ $t$) be the geometric point that is the image of
$s'$ (resp.\ $t'$) in $S$, and let $u:t\to S(s)$ be the
specialization defined by $u'$. One readily sees that there is a commutative
square
\[\xymatrix{
  \eR\Gamma(X_t,E|X_t) \ar[r]^-{\cosp(u)} \ar[d] 
    & \eR\Gamma(X_s,E|X_s) \ar[d] \\
  \eR\Gamma(X_{t'}',E|X_{t'}') \ar[r]^-{\cosp(u')} 
    & \eR\Gamma(X_{s'}',E|X_{s'}') 
}\]
in which the vertical arrows are the base-change arrows. By hypothesis, these
are isomorphisms. On the other hand, (\ref{VII:5-2-4}) $\cosp(u)$ is an
isomorphism, hence $\cosp(u')$ is an isomorphism, which completes the proof.

Here is an example of an application of \ref{VII:5-2-8}. Suppose that $A$ is
annihilated by an integer $n$ invertible on $S$, and that $f$ is locally of
finite type and universally locally acyclic. Let $x$ be a geometric point of
$X$, with image $s$ in $S$, and let $f_{(x)}:X(x)\to S(s)$ be the morphism
induced by $f$. It follows from \ref{VII:5-2-6} and \ref{VII:1-9} (by passage
to the limit) that the hypotheses of \ref{VII:5-2-8} are satisfied by the
pair $(f_{(x)},E|X(x))$, which is therefore cohomologically proper.






\subsubsection{}\label{VII:5-2-9}

Suppose that $A$ is constant. Let $f:X\to S$, and let $E\in\ob\eD^+(X)$ be a
complex of finite tor-dimension. We say that $f$ is \emph{strongly locally
acyclic} rel.\ to $E$ if, for every geometric point $x$ of $X$, with image
$s$ in $S$, every specialization $t\to S(s)$, and every $A^\circ$-module $M$,
the restriction arrow
\[
  E_x\lotimes_A M \to \eR\Gamma(X(x)_t,(E|X(x)_t)\lotimes_A M) 
\]
is an isomorphism. I do not know whether ``locally acyclic'' implies
``strongly locally acyclic.''






\begin{proposition}\label{VII:5-2-10}
Under the hypotheses of \ref{VII:5-2-9}, suppose that $A$ is noetherian and
torsion, that $S$ is noetherian, and that $f$ is strongly locally acyclic
rel.\ to $E$. Then, for every diagram with cartesian squares
\[\xymatrix{
  X \ar[d]^-f 
    & \ar[l]_-j X' \ar[d]^-{f'} 
    & \ar[l]_-{j'} X'' \ar[d]^-{f''} \\
  S 
    & \ar[l]_-i S' 
    & \ar[l]_-{i'} S'' \text{,} 
}\]
with $i$ quasi-finite and $i'$ an open immersion, and every
$F\in\ob\eD^+(S'',A^\circ)$, the canonical arrow
\begin{equation*}\tag{$*$}\label{VII:eq:5-2-10-1}
  j^\ast E \lotimes {f'}^\ast \eR i'_\ast F \to \eR j'_\ast ((j j')^\ast E\lotimes {f''}^\ast F) 
\end{equation*}
is an isomorphism.
\end{proposition}

Compare this statement with \cite[XV 1.17]{sga4}.

Let us prove \ref{VII:5-2-10}. We readily reduce to the case where $S=S'$ is
strictly local, and to $F$ constructible and concentrated in degree $0$. We then
resolve $F$ on the right by finite sums of sheaves of the form $u_\ast M$,
where $u:t\to S$ is a geometric point and $M$ is an $A^\circ$-module. The
hypothesis, vis-\`a-vis this resolution, implies that
\eqref{VII:eq:5-2-10-1} is an isomorphism for $F=u_\ast M$, and the result
follows.

% END INCLUDED components/sga4half_7_en.tex

% BEGIN INCLUDED components/sga4half_8_en.tex
% !TEX root = ../SGA4half_English_Reader_1_to_8.tex

\chapter{Derived categories}\label{VIII}
\blfootnote{by Jean-Louis Verdier}

\section{Triangulated categories: definitions and examples}\label{VIII:1}

\subsection{Definition of triangulated categories}\label{VIII:1-1}

\addtocounter{subsubsection}{-1}
\subsubsection{}\label{VIII:1-1-0}

Let $\cA$ be an additive category and let $T$ be an additive automorphism of
$\cA$. For two objects $X$ and $Y$ of $\cA$, we set
\[
  \hom^i(X,Y) = \hom_\cA(X,T^i(Y)) \qquad i\in \dZ
\]
Let $X$, $Y$, and $Z$ be three objects of $\cA$, let
$\alpha\in \hom^i(X,Y)$, and let $\beta\in \hom^j(Y,Z)$. We define the
composite $\beta\circ\alpha$ belonging to $\hom^{i+j}(X,Z)$ by
\[
  \beta\circ \alpha = T^i(\beta)\circ\alpha \text{.}
\]
One checks that this gives a new category whose groups of morphisms between
objects are graded. By an abuse of language, this new category will be called
the category $\cA$ graded by the translation functor $T$.

Let $\cA$ be a category graded by a translation functor $T$. Whenever a
morphism is named or denoted without its degree being specified, it will
always mean a morphism of degree zero.

Let $\cA$ and $\cA'$ be two additive categories graded by the functors $T$
and $T'$, and let $F$ be an additive functor from $\cA$ to $\cA'$. We say
that $F$ \emph{is graded} if an isomorphism of
functors\footnote{Text amended in 1976: the original text required an
equality $F\circ T=T'\circ F$. The case of functors in several variables is
discussed in \cite[XVIII 0.2]{sga4}. It raises a problem of signs.}
has been specified:
\[
  F\circ T = T'\circ F
\]

The definition of morphisms of graded functors is given in
\ref{VIII:4-1-1}.

Let $\cA$ be an additive category graded by the functor $T$. A triangle is
a collection of three objects $X$, $Y$, $Z$ and three morphisms
$u:X\to Y$, $v:Y\to Z$, $w:Z\to T(X)$ ($\deg w=1$). We use either the
notation $(X,Y,Z,u,v,w)$ or the following diagram for a triangle:
\[\xymatrix{
  & Z \ar[dl]_-w \\
  X \ar[rr]^-u
    & & Y \ar[lu]_-v
}\qquad \deg(w)=1
\]
Let $(X,Y,Z,u,v,w)$ and $(X',Y',Z',u',v',w')$ be two triangles of $\cA$.
A morphism from the first triangle to the second is a collection of three
morphisms $f:X\to X'$, $g:Y\to Y'$, $h:Z\to Z'$ such that the following
diagram is commutative:
\[\xymatrix{
  X \ar[r]^-u \ar[d]^-f
    & Y \ar[r]^-v \ar[d]^-g
    & Z \ar[r]^-w \ar[d]^-h
    & T(X) \ar[d]^-{T(f)} \\
  X' \ar[r]^-{u'}
    & Y' \ar[r]^-{v'}
    & Z' \ar[r]^-{w'}
    & T(X')
}\]

\subsubsection{}\label{VIII:1-1-1}

A \emph{triangulated category} is an additive category graded by a
translation functor and equipped with a family of triangles, called the
family of distinguished triangles. This family must also satisfy the
following axioms.

\paragraph{TR1}
\phantomsection\hypertarget{VIII:TR1}{}

Every triangle isomorphic to a distinguished triangle is distinguished.
Every morphism $u:X\to Y$ is contained in a distinguished triangle
$(X,Y,Z,u,v,w)$. The triangle
$(X,X,0,\operatorname{id}_X,0,0)$ is distinguished.

\paragraph{TR2}
\phantomsection\hypertarget{VIII:TR2}{}

A triangle $(X,Y,Z,u,v,w)$ is distinguished if and only if
$(Y,Z,T(X),v,w,-T(u))$ is distinguished.

\paragraph{TR3}
\phantomsection\hypertarget{VIII:TR3}{}

If $(X,Y,Z,u,v,w)$ and $(X',Y',Z',u',v',w')$ are distinguished, then for
every morphism $(f,g):u\to u'$, there exists a morphism $h:Z\to Z'$ such
that $(f,g,h)$ is a morphism of triangles.

\paragraph{TR4}
\phantomsection\hypertarget{VIII:TR4}{}

\[\xymatrix{
  & Z' \ar[dl]|{\deg=1} \\
  X \ar[rr]^-u
    & & Y \ar[ul]_-i
}\quad
\xymatrix{
  & X' \ar[dl]|{\deg(j)=1} \\
  Y \ar[rr]^-v
    & & Z \ar[ul]
}\quad
\xymatrix{
  & Y' \ar[dl]|{\deg=1} \\
  X \ar[rr]^-w
    & & Z \ar[ul]
}\]
If these are three distinguished triangles such that $w=v\circ u$, then
there exist two morphisms $f:Z'\to Y'$, $g:Y' \to X'$ such that
\begin{enumerate}
  \item $(\operatorname{id}_X, v, f)$ is a morphism of triangles;
  \item $(u,\operatorname{id}_Z,g)$ is a morphism of triangles;
  \item $(Z',Y',X',f,g,T(i) j)$ is a distinguished triangle.
\end{enumerate}

If $\cC$ and $\cC'$ are two triangulated categories, an \emph{exact
functor} from $\cC$ to $\cC'$ is an additive graded functor that takes
distinguished triangles to distinguished triangles. Morphisms of exact
functors are morphisms of graded functors.

The following property follows immediately from axioms
\hyperlink{VIII:TR1}{TR1}, \hyperlink{VIII:TR2}{TR2}, and
\hyperlink{VIII:TR3}{TR3}.

\begin{proposition}\label{VIII:1-1-2}
Let $(X,Y,Z,u,v,w)$ be a distinguished triangle and let $M$ be an object
of a triangulated category. Then the following sequences are exact:
\[\xymatrix{
  \cdots \ar[r]
    & \hom^i(M,X) \ar[r]^-{\hom(M,u)}
    & \hom^i(M,Y) \ar[r]^-{\hom(M,v)}
    & \hom^i(M,Z) \ar[r]^-{\hom(M,w)}
    & \hom^{i+1}(M,X) \ar[r]
    & \cdots
}\]
\[\xymatrix{
  \cdots \ar[r]
    & \hom^i(Z,M) \ar[r]^-{\hom(v,M)}
    & \hom^i(Y,M) \ar[r]^-{\hom(u,M)}
    & \hom^i(X,M) \ar[r]^-{\hom(w,M)}
    & \hom^{i+1}(Z,M) \ar[r]
    & \cdots
}\]
\end{proposition}

From this proposition one obtains statements of the following kind:
\begin{itemize}
  \item If $(f,g,h)$ is a morphism of distinguished triangles and $f,g$
    are isomorphisms, then $h$ is an isomorphism.
  \item The distinguished triangles constructed on a morphism
    (\hyperlink{VIII:TR1}{TR1}) are all isomorphic (\emph{but the
    isomorphisms are not, in general, uniquely determined}).
  \item When one of the morphisms in a distinguished triangle has a kernel
    or a cokernel, the latter splits; that is, the triangle is of the form
    \[\xymatrix{
      & Z+T(X) \ar[dl] \\
      X+Y \ar[rr]
        & & Y+Z \ar[ul]
    }\]
  \item The sum of two distinguished triangles is distinguished. (One uses
    axiom \hyperlink{VIII:TR4}{TR4}.)
\end{itemize}

\subsection{Examples}\label{VIII:1-2}

\subsubsection{}\label{VIII:1-2-1}

We first fix some notation. Let $\cA$ be an additive category.
The following category will be denoted by $\eC(\cA)$:
\begin{itemize}
  \item The objects of $\eC(\cA)$ are the complexes of $\cA$, with no
    restriction on degree and with differential of degree $+1$.
  \item The morphisms of $\eC(\cA)$ are the degree-preserving morphisms of
    complexes (those commuting with the differential).
\end{itemize}

Let $T$ be the following functor. For every object $X^\bullet$ of
$\eC(\cA)$,
\begin{align*}
  T(X^\bullet)_i &= (X^\bullet)_{i+1} \\
  d_i(T(X^\bullet)) &= -d_{i+1}(X^\bullet)  \text{.}
\end{align*}
On morphisms the functor $T$ acts as follows: for every morphism
$f:X^\bullet \to Y^\bullet$, the morphism
$T(f):T(X^\bullet) \to T(Y^\bullet)$ is defined by
\[
  T(f)_i = f_{i+1} \text{.}
\]
It is immediate that $T$ is an additive automorphism of $\eC(\cA)$.
The morphisms of degree $n$ defined using this translation functor are
then the morphisms that raise degree by $n$ and commute or anticommute
with the differential according to the parity of $n$. Thus one recovers
the usual definition.

Let $\sC$ denote the following family of triangles. A triangle
$(X^\bullet,Y^\bullet,Z^\bullet,u,v,w)$ belongs to the family
if\footnote{Text amended in 1976 to ensure the validity of
\ref{VIII:1-2-4}.}
\begin{itemize}
  \item $Z^\bullet$ is the simple complex associated with the double
    complex
    \[\xymatrix{
      \cdots \ar[r]
        & 0 \ar[r]
        & X^\bullet \ar[r]^-u
        & Y^\bullet \ar[r]
        & 0 \ar[r]
        & \cdots
    }\]
    in which the objects of $Y^\bullet$ have first degree zero and the
    objects of $X^\bullet$ have first degree $-1$.
  \item $v$ is the image, under the functor: double complex $\mapsto$
    simple complex, of the following morphism of double complexes:
    \[\xymatrix{
      \cdots \ar[r]
        & 0 \ar[r] \ar[d]
        & 0^\bullet \ar[r] \ar[d]
        & Y^\bullet \ar[r] \ar[d]
        & 0 \ar[r] \ar[d]
        & \cdots \\
      \cdots \ar[r]
        & 0 \ar[r]
        & X^\bullet \ar[r]
        & Y^\bullet \ar[r]
        & 0 \ar[r]
        & \cdots
    }\]
  \item $w$ is the negative of the image, under the same functor, of the
    following morphism of double complexes:
    \[\xymatrix{
      \cdots \ar[r]
        & 0 \ar[r]
        & X^\bullet \ar[r]
        & 0 \ar[r]
        & 0 \ar[r]
        & \cdots \\
      \cdots  \ar[r]
        & 0 \ar[r] \ar[u]
        & X^\bullet \ar[r] \ar[u]
        & Y^\bullet \ar[r] \ar[u]
        & 0 \ar[r] \ar[u]
        & \cdots
    }\]
\end{itemize}

\subsubsection{}\label{VIII:1-2-2}

The set of null-homotopic morphisms is a two-sided ideal (if $f$ and $g$
are composable and either $f$ or $g$ is null-homotopic, then the composite
is null-homotopic; if $f$ and $g$ are null-homotopic, then
$f\oplus g$ is null-homotopic). For every pair of objects $X^\bullet$,
$Y^\bullet$, let $\operatorname{Htp}(X^\bullet,Y^\bullet)$ denote the
subgroup of $\hom_{\eC(\cA)}(X^\bullet,Y^\bullet)$ consisting of the
null-homotopic morphisms. The objects of $\eK(\cA)$ are then the objects
of $\eC(\cA)$, and its group of morphisms with source $X^\bullet$ and
target $Y^\bullet$ is
\[
  \hom_{\eC(\cA)}(X^\bullet,Y^\bullet) / \operatorname{Htp}(X^\bullet,Y^\bullet) \text{.}
\]
The composition law on the morphisms of $\eK(\cA)$ is defined by passing
the composition law on the morphisms of $\eC(\cA)$ to the quotient.
The category $\eK(\cA)$ is additive. The translation functor $T$ on
$\eC(\cA)$ passes to the quotient (the translate of a null-homotopic
morphism is null-homotopic). The resulting functor is an additive
automorphism of $\eK(\cA)$, again denoted by $T$. By an abuse of notation,
$\eK(\cA)$ will henceforth denote the category $\eK(\cA)$ graded by the
functor $T$.

A distinguished triangle in $\eK(\cA)$ is any triangle \emph{isomorphic}
to a triangle arising from $\sC$. The family of distinguished triangles
satisfies axioms \hyperlink{VIII:TR1}{TR1},
\hyperlink{VIII:TR2}{TR2}, \hyperlink{VIII:TR3}{TR3}, and
\hyperlink{VIII:TR4}{TR4}. By an abuse of notation, $\eK(\cA)$ will denote
the category $\eK(\cA)$ triangulated by the family of triangles defined
above.

\subsubsection{}\label{VIII:1-2-3}

There is also a whole arsenal of subcategories:
\begin{itemize}
  \item A complex is said to be bounded below if all but finitely many of
    its objects of negative degree are zero. Bounded-above and bounded
    complexes are defined similarly. Let $\eK^+(\cA)$, $\eK^-(\cA)$,
    and $\eK^b(\cA)$ denote the full subcategories of $\eK(\cA)$ generated
    by these families of objects. These subcategories are
    “\emph{triangulated subcategories}” of $\eK(\cA)$: every distinguished
    triangle two of whose objects belong to the subcategory is isomorphic
    to a triangle all three of whose objects belong to the subcategory.
  \item Suppose that $\cA$ is an abelian category. A complex is said to be
    cohomologically bounded below if all but finitely many of its
    cohomology objects of negative degree are zero. Cohomologically
    bounded-above complexes, cohomologically bounded complexes, and
    acyclic complexes are defined similarly. Let
    $\eK^{\infty,+}(\cA)$, $\eK^{\infty,-}(\cA)$,
    $\eK^{\infty,b}(\cA)$, and $\eK^{\infty,\varnothing}(\cA)$ denote the
    full subcategories generated by these families of objects. They are
    triangulated subcategories.

    We shall also use the triangulated subcategories $\eK^{+,b}(\cA)$,
    $\eK^{+,\varnothing}(\cA)$, \ldots etc \ldots (the first superscript
    gives information about the objects of the complexes, and the second
    about the cohomology objects).
  \item Without assuming that $\cA$ is abelian, let $O$ be a set of objects
    of $\cA$ stable under isomorphisms and direct sums. Let $O(\cA)$ denote
    the full subcategory of $\eC(\cA)$ generated by the complexes whose
    objects in every degree belong to $O$. Let $\underline O(\cA)$ denote
    the corresponding triangulated subcategory of $\eK(\cA)$. When $\cA$
    is abelian, one may take for $O$ either the set $I$ of injectives or
    the set $P$ of projectives.
\end{itemize}

\subsubsection{}\label{VIII:1-2-4}

The category $\eK(\cA)$ is uniquely determined by the data of $\eC(\cA)$
and the family $\sC$. More precisely, every graded additive functor from
$\eC(\cA)$ to a triangulated category that takes every triangle of $\sC$
to a distinguished triangle factors uniquely through $\eK(\cA)$, and the
resulting functor is exact.

For every complex $X^\bullet$ of $\eC(\cA)$, let $\bar X^\bullet$ denote
the complex obtained from $X^\bullet$ by setting the differential equal
to zero. A three-term sequence
\[\xymatrix{
  0 \ar[r]
    & X^\bullet \ar[r]
    & Y^\bullet \ar[r]
    & Z^\bullet \ar[r]
    & 0
}\]
is called a semi-split sequence if the sequence
\[\xymatrix{
  0 \ar[r]
    & \bar X^\bullet \ar[r]
    & \bar Y^\bullet \ar[r]
    & \bar Z^\bullet \ar[r]
    & 0
}\]
is split; that is, if it is isomorphic to the sequence
\[\xymatrix{
  0 \ar[r]
    & \bar X^\bullet \ar[r]
    & \bar X^\bullet + \bar Z^\bullet \ar[r]
    & \bar Z^\bullet \ar[r]
    & 0 \text{.}
}\]

Let $0 \to X^\bullet \to Y^\bullet \to Z^\bullet \to 0$ be a semi-split
sequence. Choose a splitting, that is, for every $n$ an isomorphism
\[
  (Y^\bullet)^n \simeq (X^\bullet)^n + (Z^\bullet)^n \text{.}
\]
The matrix of the differential of $Y^\bullet$ in the above direct-sum
decomposition has the form
\[
  \begin{pmatrix}
    d_{X^\bullet} & \varphi_n \\
    0 & d_{Z^\bullet}
  \end{pmatrix}
\]
where $\varphi_n$ is a morphism from $(Z^\bullet)^n$ to
$(X^\bullet)^{n+1}$. The $\varphi_n$ determine a morphism of complexes
\[
  \varphi:Z^\bullet \to T(X^\bullet) \text{.}
\]
When the splitting is changed, the homotopy class of $\varphi$ does not
change. Moreover, if $\underline u$, $\underline v$, and
$\underline\varphi$ denote the images in $\eK(\cA)$ of the morphisms $u$,
$v$, and $\varphi$, then the triangle
$(X^\bullet,Y^\bullet,Z^\bullet,\underline u,\underline v,\underline\varphi)$
is distinguished. Let $\mathsf{SSS}(\cA)$ denote the category of
semi-split sequences of $\eC(\cA)$ and let $\mathsf{TrK}(\cA)$ denote the
category of distinguished triangles of $\eK(\cA)$. The preceding
construction defines a functor
\[
  \rho:\mathsf{SSS}(\cA) \to \mathsf{TrK}(\cA) \text{.}
\]
This functor is essentially surjective.

\subsection{Examples of cohomological and exact functors}\label{VIII:1-3}

\begin{definition}
Let $\cA$ be a triangulated category and let $\cB$ be an abelian category.
An additive functor $F:\cA\to \cB$ is called a \emph{cohomological
functor} if, for every distinguished triangle $(X,Y,Z,u,v,w)$, the
sequence
\[\xymatrix{
  F(X) \ar[r]^-{F(u)}
    & F(Y) \ar[r]^-{F(v)}
    & F(Z)
}\]
is exact.
\end{definition}

The functor $F\circ T^i$ will often be denoted by $F^i$. By axiom
\hyperlink{VIII:TR2}{TR2} for triangulated categories, there is an
unbounded exact sequence
\[\xymatrix{
  \cdots \ar[r]
    & F^i(X) \ar[r]
    & F^i(Y) \ar[r]
    & F^i(Z) \ar[r]
    & F^{i+1}(X) \ar[r]
    & \cdots
}\]

\subsubsection{Examples}\label{VIII:1-3-2}

\begin{enumerate}
  \item Let $\cA$ be a triangulated category. For every object $X$ of
    $\cA$, the functor $\hom_\cA(X,-)$ is cohomological. The functor
    $\hom_\cA(-,X)$ is a cohomological functor
    $\cA^\circ\to\mathsf{Ab}$.
  \item Let $\cA$ be an abelian category and let $\eC(\cA)$ be the category
    of complexes of $\cA$. Define the functor
    $\h^0:\eC(\cA)\to\cA$ as follows. For an object $X^\bullet$ of
    $\eC(\cA)$, set
    \[
      \h^0(X^\bullet) = \ker d^0 / \im{d^{-1}} \text{.}
    \]
    The functor $\h^0$ annihilates null-homotopic morphisms and therefore
    factors uniquely through $\eK(\cA)$. The resulting functor will again
    be denoted by $\h^0:\eK(\cA) \to \cA$.
\end{enumerate}

\subsubsection{Example of an exact bifunctor}\label{VIII:1-3-3}

Let $\cA$, $\cA'$, and $\cA''$ be three additive categories, and let
\[
  P:\cA\times \cA'\to \cA''
\]
be a bilinear functor (that is, additive in each argument separately).
It induces the bilinear functor
\[
  P^\ast:\eC(\cA)\times \eC(\cA') \to \eC(\cA'')
\]
as follows.

Let $X^\bullet$ be an object of $\eC(\cA)$ and let $Y^\bullet$ be an
object of $\eC(\cA')$. Then $P(X^\bullet,Y^\bullet)$ is a double complex
of $\cA''$. Set $P^\ast(X^\bullet,Y^\bullet) = $ the simple complex
associated with $P(X^\bullet,Y^\bullet)$.

Let $f$ be a null-homotopic morphism of $\eC(\cA)$ (respectively
$\eC(\cA')$), and let $Z^\bullet$ be an object of $\eC(\cA')$
(respectively $\eC(\cA)$). The morphism $P^\ast(f,Z^\bullet)$
(respectively $P^\ast(Z^\bullet,f)$) is then null-homotopic. It follows
that $P^\ast$ uniquely defines a functor
\[
  P^\bullet:\eK(\cA) \times \eK(\cA') \to \eK(\cA'') \text{.}
\]
The bifunctor $P^\bullet$ is exact.

In particular, let $\cA$ be an additive category. For $P$ one may take
the functor
\begin{align*}
  \cA^\circ\times \cA &\to \mathsf{Ab} \\
  (X,Y) &\mapsto \hom(X,Y)
\end{align*}

The preceding construction then gives a functor
\[
  \hom^\bullet:\eK(\cA)^\circ \times \eK(\cA) \to \eK(\mathsf{Ab})
\]
which, when composed with the functor
$\h^0:\eK(\mathsf{Ab})\to \mathsf{Ab}$, plainly gives back the functor
$\hom_{\eK(\cA)}$.

\section{Quotient categories}\label{VIII:2}

\subsection{Thick subcategories. Multiplicative systems of morphisms}\label{VIII:2-1}

We begin with two definitions.

\begin{definition}\label{VIII:2-1-1}
A subcategory $\cB$ of a triangulated category $\cA$ is said to be
\emph{thick} if $\cB$ is a full triangulated subcategory of $\cA$ and,
in addition, $\cB$ has the following property:

For every morphism $f:X\to Y$ that factors through an object of $\cB$ and
is contained in a distinguished triangle $(X,Y,Z,f,g,h)$ in which $Z$ is
an object of $\cB$, the source of $f$ and the target of $f$ are objects of
$\cB$.
\end{definition}

The set of thick subcategories of $\cA$ will be denoted by $\sN$.

The intersection of an arbitrary family of thick subcategories is a thick
subcategory.

\begin{definition}
Let $\cA$ be a triangulated category. A set $S$ of morphisms of $\cA$ is
called a \emph{multiplicative system of morphisms} if it has the following
five properties.

\paragraph{FR1}
\phantomsection\hypertarget{VIII:FR1}{}
If $f,g\in S$ and $f$ and $g$ are composable, then $f\circ g\in S$. For
every object $X$ of $\cA$, the identity morphism of $X$ belongs to $S$.

\paragraph{FR2}
\phantomsection\hypertarget{VIII:FR2}{}
In the category $\cA$, every diagram
\[\xymatrix{
  & Y \ar[d]^-{s\in S} \\
  Z \ar[r]^-f
    & X
}\]
can be completed to a commutative diagram
\[\xymatrix{
  P \ar[r]^-g \ar[d]^-{t\in S}
    & Y \ar[d]^-{s\in S} \\
  Z \ar[r]^-f
    & X
}\]
Moreover, the symmetric property holds.

\paragraph{FR3}
\phantomsection\hypertarget{VIII:FR3}{}
If $f$ and $g$ are morphisms, the following properties are equivalent:
\begin{enumerate}[\indent i)]
  \item there exists $s\in S$ such that $s\circ f=s\circ g$;
  \item there exists $t\in S$ such that $f\circ t=g\circ t$.
\end{enumerate}

\paragraph{FR4}
\phantomsection\hypertarget{VIII:FR4}{}
Let $T$ be the translation functor of $\cA$. For every element $s$ of $S$,
one has $T(s)\in S$.

\paragraph{FR5}
\phantomsection\hypertarget{VIII:FR5}{}
Let $(X,Y,Z,u,v,w)$ and $(X',Y',Z',u',v',w')$ be two distinguished
triangles, and let $f$ and $g$ be two elements of $S$ such that the pair
$(f,g)$ is a morphism from $u$ to $u'$. There exists a morphism $h\in S$
such that $(f,g,h)$ is a morphism of triangles.
\end{definition}

A multiplicative system of morphisms is said to be \emph{saturated} if it
has the following property:

A morphism $f$ belongs to $S$ if and only if there exist two morphisms $g$
and $g'$ such that $g\circ f\in S$ and $f\circ g'\in S$.

The intersection of an arbitrary family of saturated multiplicative
systems is a saturated multiplicative system.

The set of saturated multiplicative systems will be denoted by $\sS$.

Let $\cB$ be a thick subcategory. Denote by $\varphi(\cB)$ the set of
morphisms $f$ contained in a distinguished triangle $(X,Y,Z,f,g,h)$ in
which $Z$ is an object of $\cB$. Then $\varphi(\cB)$ is a saturated
multiplicative system.

Let $S$ be a saturated multiplicative system. Denote by $\psi(S)$ the
full subcategory generated by the objects $Z$ contained in a distinguished
triangle $(X,Y,Z,f,g,h)$ in which $f$ is an element of $S$. The
subcategory $\psi(S)$ is thick.

The main result of this subsection is the following: $\varphi$ is an
order-preserving isomorphism from $\sN$ onto $\sS$, where the order
relation is defined by inclusion. Its inverse is precisely $\psi$.

\subsection{Universal problems}\label{VIII:2-2}

Let $\cA$ and $\cA'$ be two triangulated categories and let $F$ be an exact
functor from $\cA$ to $\cA'$. Let $S(F)$ be the set of morphisms of $\cA$
that $F$ takes to isomorphisms, and let $\cB(F)$ be the full subcategory
generated by the objects of $\cA$ that $F$ takes to zero objects of
$\cA'$. Then $S(F)$ is a saturated multiplicative system and $\cB(F)$ is
a thick subcategory. Moreover, $S(F) = \varphi(\cB(F))$. Now let $\cA$ be
a triangulated category, let $\cB$ be a thick subcategory, and let
$S=\varphi(\cB)$ be the corresponding saturated multiplicative system.
The following two universal problems are equivalent.

\paragraph{Problem 1}
\phantomsection\hypertarget{VIII:prob1}{}
Find a triangulated category $\cA/\cB$ and an exact functor
$Q:\cA\to \cA/\cB$ such that every exact functor from $\cA$ to a
triangulated category $\cA'$ that takes the objects of $\cB$ to zero
objects of $\cA'$ admits a unique factorization of the form $G\circ Q$,
where $G$ is exact.

\paragraph{Problem 2}
\phantomsection\hypertarget{VIII:prob2}{}
Find a triangulated category $\cA_S$ and an exact functor
$Q:\cA\to \cA_S$ such that every exact functor from $\cA$ to a
triangulated category $\cA'$ that takes the morphisms of $S$ to
isomorphisms of $\cA'$ admits a unique factorization of the form
$G\circ Q$, where $G$ is exact.

(The author asks the reader to excuse this use of obsolete and reactionary
language.) In the next subsection we shall show that
\hyperlink{VIII:prob2}{Problem 2} has a solution. Consequently, the
corresponding \hyperlink{VIII:prob1}{Problem 1} has the same solution.

\subsection{Calculus of fractions}\label{VIII:2-3}

\subsubsection{}\label{VIII:2-3-1}

Let $\cA$ be a triangulated category and let $S$ be a saturated
multiplicative system. Denote by $\cA[S^{-1}]$ the category of fractions
of $\cA$ with respect to the set of morphisms $S$, and by $Q$ the
canonical functor from $\cA$ to $\cA[S^{-1}]$. The pair
$(\cA[S^{-1}],Q)$ solves the following problem: every functor from $\cA$
to an arbitrary category (not necessarily additive) that takes morphisms
of $S$ to isomorphisms factors uniquely through $(\cA[S^{-1}],Q)$. Such
a pair always exists, without any hypothesis on the set of morphisms $S$
\cite[1.1]{gz67}\footnote{The original text references [C.G.G] here. Most
likely, this is the lost work \emph{Catégories et foncteurs} by Chevalley,
Gabriel and Grothendieck, referenced in\cite{sc72}. }. However, because the
set of morphisms $S$ is a saturated multiplicative system, it has
properties \hyperlink{VIII:FR1}{FR1}, \hyperlink{VIII:FR2}{FR2}, and
\hyperlink{VIII:FR3}{FR3}. Consequently, by \cite{gz67}, the category
$\cA[S^{-1}]$ can be obtained by a right or left calculus of fractions.
We recall how $\cA[S^{-1}]$ is obtained by a right calculus of fractions.
(The left calculus of fractions is obtained by reversing the arrows.)

\subsubsection{}\label{VIII:2-3-2}

The objects of $\cA[S^{-1}]$ are the objects of $\cA$.
\begin{itemize}
  \item For every object $X$ of $\cA$, let $S_X$ denote the category of
    arrows belonging to $S$ and having target $X$. To every object $Y$ of
    $\cA$ one associates, functorially in $Y$, the following set-valued
    functor on $S_X^\circ$ ($S_X^\circ$ denotes the opposite category):
    \[
      H_Y:s\mapsto \hom(\operatorname{source}(s),Y) \text{.}
    \]
    Set
    \[
      \hom_{\cA[S^{-1}]}(X,Y) = \varinjlim_{S_X^\circ} H_Y \text{.}
    \]
    For all matters concerning direct limits, we refer to \cite{ar62}.

    One then observes with pleasure, using \hyperlink{VIII:FR1}{FR1},
    \hyperlink{VIII:FR2}{FR2}, and \hyperlink{VIII:FR3}{FR3}, that the
    category $S_X^\circ$ has properties L1, L2, and L3. Direct limits
    therefore have all the good properties of direct limits over filtered
    ordered sets.
  \item Let $X$, $Y$, and $Z$ be three objects of $\cA$, let $a$ be an
    element of $\hom_{\cA[S^{-1}]}(X,Y)$, and let $b$ be an element of
    $\hom_{\cA[S^{-1}]}(Y,Z)$. Let $s$ be an object of $S_X$, let
    $\underline a$ be an element of
    $\hom(\operatorname{source}(s),Y)$ whose image is $a$, and let an object
    of $S_Y$, $\underline b$ be an element of
    $\hom(\operatorname{source}(t),Z)$ whose image is $b$. There is then a
    diagram
    \begin{equation*}\tag{1}\label{VIII:eq:2-3-2-1}
    \xymatrix{
      & \bullet \ar[dl]_-s \ar[dr]^-{\underline a}
        &
        & \bullet \ar[dl]_-t \ar[dr]^-{\underline b} \\
      X
        &
        & Y
        &
        & Z
    }
    \end{equation*}
    which, by \hyperlink{VIII:FR2}{FR2}, can be completed to a diagram
    \[\xymatrix{
      & & \bullet \ar[dl]_-{t'} \ar[dr]^-{\underline c} \\
      & \bullet \ar[dl]_-s \ar[dr]^-{\underline a}
        & & \bullet \ar[dl]_-t \ar[dr]^-{\underline b} \\
      X
        & & Y
        & & Z
    }\]
    Let $b\circ a$ denote the image in $\hom_{\cA[S^{-1}]}(X,Z)$ of the
    morphism $\underline{b\circ c}$. Using properties
    \hyperlink{VIII:FR1}{FR1}, \hyperlink{VIII:FR2}{FR2}, and
    \hyperlink{VIII:FR3}{FR3}, one checks that $b\circ a$ is independent
    of the chosen representatives $\underline a$ and $\underline b$, and
    also of the way in which diagram \eqref{VIII:eq:2-3-2-1} is completed.
    One further checks that this defines a category $\cA[S^{-1}]$.
\end{itemize}

We denote by $Q$ the evident functor $\cA\to \cA[S^{-1}]$.

Since $\cA$ is additive, $\cA[S^{-1}]$ is an additive category. Moreover,
using \hyperlink{VIII:FR4}{FR4}, one proves that there exists a unique
translation functor on $\cA[S^{-1}]$, denoted by $T$, satisfying
\[
  Q\circ T = T\circ Q \text{.}
\]
Finally, using \hyperlink{VIII:FR5}{FR5}, one proves that there exists a
unique triangulated structure on $\cA[S^{-1}]$ for which $Q$ is exact.
The distinguished triangles of $\cA[S^{-1}]$ are precisely the triangles
isomorphic to the images under $Q$ of the distinguished triangles of
$\cA$. Denoting by $\cA_S$ the category $\cA[S^{-1}]$ equipped with this
triangulated structure, one proves without difficulty that the pair
$(\cA_S,Q)$ is a solution to \hyperlink{VIII:prob2}{Problem 2}.

\begin{definition}\label{VIII:2-3-3}
Let $\cA$ be a triangulated category, let $\cB$ be a thick subcategory of
$\cA$, and let $(Q,\cA/\cB)$ be the solution to
\hyperlink{VIII:prob1}{Problem 1} (\S \ref{VIII:2-2}). The category
$\cA/\cB$ is called the \emph{quotient category} of $\cA$ by $\cB$, and
$Q$ is called the canonical quotient functor. More generally, let $\cN$
be a triangulated subcategory of $\cA$, and let $\underline\cN$ be the
smallest thick subcategory containing $\cN$. By definition, the quotient
category of $\cA$ by $\cN$, denoted $\cA/\cN$, is
$\cA/\underline\cN$.
\end{definition}

\subsection{Properties of quotient categories}\label{VIII:2-4}

\subsubsection{}\label{VIII:2-4-1}

Let $\cA$ be a triangulated category and let $H$ be a cohomological functor
with values in an abelian category $\cG$. Let $\cB_H$ be the full
subcategory generated by the objects all of whose translates are taken by
$H$ to zero objects of $\cG$, and let $S_H$ be the set of morphisms all of
whose translates are taken by $H$ to isomorphisms of $\cG$. Then $\cB_H$
is a thick subcategory of $\cA$, and $S_H$ is a saturated multiplicative
system. Moreover, $S_H=\varphi(\cB_H)$ (\S \ref{VIII:2-1}). The functor
$H$ factors uniquely through $(Q,\cA/\cB_H)$.

\begin{theorem}\label{VIII:2-4-2}
Let $\cA$ be a triangulated category, let $\cB$ be a triangulated
subcategory, let $\cN$ be a thick subcategory of $\cA$, and let
$S=\varphi(\cN)$ be the corresponding saturated multiplicative system.
\begin{enumerate}[(a)]
  \item The category $\cN\cap\cB$ is a thick subcategory of $\cB$. The
    corresponding multiplicative system is $S\cap \cB$.
  \item The following two properties are equivalent:
    \begin{enumerate}[(i)]
      \item For every object $X$ of $\cB$ and every morphism $s:R\to X$,
        where $s$ belongs to $S$, there exists a morphism $s':R'\to R$ such
        that $s\circ s'$ belongs to $S\cap \cB$; and
        $R'\in\operatorname{Ob}\cB$.
      \item Every morphism from an object $X$ of $\cB$ to an object $Y$ of
        $\cN$ factors through an object of $\cN\cap \cB$.
    \end{enumerate}
  \item The properties (i)' and (ii)' obtained from (i) and (ii) by
    reversing the arrows are equivalent.
  \item If properties (i) and (ii) hold, or if properties (i)' and (ii)'
    hold, the canonical functor
    \[\xymatrix{
      \cB/\cN\cap \cB \ar[r] & \cA/\cN
    }\]
    is faithful.

    Since this functor is injective on objects, it realizes
    $\cB/\cN\cap\cB$ as a subcategory of $\cA/\cN$.
  \item If, in addition, $\cB$ is a full subcategory of $\cA$, then
    $\cB/\cN\cap \cB$ is a full subcategory of $\cA/\cN$.
\end{enumerate}
\end{theorem}

\begin{corollary}\label{VIII:2-4-3}
Let $\cA\subset \cB\subset \cC$ be three triangulated categories, with
$\cA$ a thick subcategory of $\cB$ and $\cB$ a thick subcategory of
$\cC$. Then $\cA$ is a thick subcategory of $\cC$, and $\cB/\cA$ is a
thick subcategory of $\cC/\cA$. The canonical functor
$\cC/\cB \to (\cC/\cA)/(\cB/\cA)$ is an isomorphism.
\end{corollary}

\subsection{Properties of the quotient functor}\label{VIII:2-5}

Let $\cA$ be a triangulated category, let $\cB$ be a thick subcategory,
let $S$ be the corresponding multiplicative system, and let
$Q:\cA\to \cA/\cB$ be the canonical quotient functor.

\begin{proposition}\label{VIII:2-5-1}
For objects $X$ and $Y$ of $\cA$, the following assertions are equivalent:
\begin{enumerate}[(a)]
  \item Every morphism $f:X\to Y$ for which there exists a morphism
    $s:Z\to X$ in $S$ satisfying $f\circ s=0$ is zero.
  \item Every morphism $f:X\to Y$ for which there exists a morphism
    $s:Y\to Z$, $s\in S$, satisfying $s\circ f=0$ is zero.
  \item Every morphism $f:X\to Y$ that factors through an object of $\cB$
    is zero.
  \item The canonical map
    \[
      \hom_\cA(X,Y) \to \hom_{\cA/\cB}(Q(X),Q(Y))
    \]
    is injective.
\end{enumerate}
\end{proposition}

\begin{proposition}\label{VIII:2-5-2}
Let $X$ and $Y$ be two objects of $\cA$. The following assertions are
equivalent:
\begin{enumerate}[(a)]
  \item Every diagram ($s\in S$)
    \[\xymatrix{
      & Z \ar[dl]_-s \ar[dr]^-f \\
      X
        & & Y
    }\]
    can be completed to a diagram: $(s,t\in S$):
    \[\xymatrix{
      & Z' \ar[d]_-t \\
      & Z \ar[dl]_-s \ar[dr]^-f \\
      X
        & & Y
    }\]
    in which $f\circ t$ factors through $(X,s\circ t)$.
  \item The assertion obtained by reversing the direction of the arrows in
    (a) and interchanging $X$ and $Y$.
  \item Every diagram
    \[\xymatrix{
      X \ar[r]^-f
        & N \ar[r]^-g
        & Y
    }\]
    in which $g\circ f=0$ and $N$ is an object of $\cB$ can be completed
    to a diagram
    \[\xymatrix{
      & N' \ar[d]^-i \\
      X \ar[ur]^-h \ar[r]^-f
        & N \ar[r]^-g
        & Y
    }\]
    in which $f=i\circ h$ and $g\circ i=0$.
  \item The assertion obtained by reversing the arrows in (c) and
    interchanging $X$ and $Y$.
  \item The canonical map
    \[\xymatrix{
      \hom_\cA(X,Y) \ar[r]^-Q
        & \hom_{\cA/\cB}(Q(X),Q(Y))
    }\]
    is surjective.
\end{enumerate}
\end{proposition}

\begin{proposition}\label{VIII:2-5-3}
Let $X$ be an object of $\cA$. The following assertions are equivalent:
\begin{enumerate}[(a)]
  \item For every object $Y$ of $\cA$, the canonical map
    \[\xymatrix{
      \hom_\cA(X,Y) \ar[r]^-Q
        & \hom_{\cA/\cB}(Q(X),Q(Y))
    }\]
    is an isomorphism.
  \item Every morphism from $X$ to an object of $\cB$ is zero.
\end{enumerate}
Similarly, the following two assertions are equivalent:
\begin{enumerate}[(a)']
  \item For every object $Y$ of $\cA$, the canonical map
    \[\xymatrix{
      \hom_\cA(Y,X) \ar[r]^-Q
        & \hom_{\cA/\cB}(Q(Y),Q(X))
    }\]
    is an isomorphism.
  \item Every morphism from an object of $\cB$ to $X$ is zero.
\end{enumerate}
\end{proposition}

\begin{definition}\label{VIII:2-5-4}
Every object having the equivalent properties (a) and (b) of the preceding
proposition is called an \emph{object free on the left over $\cA/\cB$},
or a \emph{left $Q$-free object}. Similarly, every object having the
equivalent properties (a)' and (b)' is called an \emph{object free on the
right over $\cA/\cB$}, or a \emph{right $Q$-free object}. It is determined
up to isomorphism by the knowledge of
$Q(X)\in\operatorname{Ob}{\cA/\cB}$.
\end{definition}

\subsection{Functors adjoint to the quotient functor}\label{VIII:2-6}

\begin{definition}\label{VIII:2-6-1}
Two triangulated subcategories $\cN$ and $\cN'$ of a triangulated category
$\cA$ are said to be \emph{orthogonal} if, for every object $X$ of $\cN$
and every object $Y$ of $\cN'$, one has
\[
  \hom_\cA(X,Y) = 0
\]
The category $\cN'$ is then said to be right orthogonal to $\cN$, and
$\cN$ left orthogonal to $\cN'$.
\end{definition}

\begin{proposition}\label{VIII:2-6-2}
Let $\cN$ be a triangulated subcategory of a triangulated category $\cA$.
The full category $\cN^\bot$ (respectively $^\bot\cN$) generated by the
objects $X$ of $\cA$ such that, for every object $Y$ of $\cN$,
$\hom_\cA(Y,X)=0$ (respectively $\hom_\cA(X,Y)=0$), is a thick subcategory
of $\cA$. By an abuse of language, the category $\cN^\bot$ is called the
right orthogonal of $\cN$.
\end{proposition}

This proposition allows us to restate Proposition \ref{VIII:2-5-3}.

\begin{proposition}\label{VIII:2-6-3}
Let $\cB\subset\cA$ be a thick subcategory of a triangulated category
$\cA$. The full category generated by the objects free on the right over
$\cA/\cB$ (respectively on the left) is precisely the right
(respectively left) orthogonal of $\cB$.
\end{proposition}

\begin{proposition}\label{VIII:2-6-4}
Let $\cA$ be a triangulated category, let $\cB$ be a thick subcategory of
$\cA$, let $\cB^\bot$ be the right orthogonal of $\cB$, let $S(\cB)$ and
$S(\cB^\bot)$ be the corresponding multiplicative systems, and let $Q$
and $Q^\bot$ be the canonical quotient functors. For an object $X$ of
$\cA$, consider the following properties:
\begin{enumerate}[(i)]
  \item $X$ maps by a morphism of $S(\cB)$ to an object of $\cB^\bot$.
  \item $X$ receives, by a morphism of $S(\cB^\bot)$, an object of
    $\cB$.
  \item The category $S_X(\cB)$ of arrows of $S(\cB)$ with source $X$ has
    a final object.
  \item The category $S^X(\cB^\bot)$ of arrows of $S(\cB^\bot)$ with
    target $X$ has an initial object.
  \item The category $\cB/X$ of objects of $\cB$ over $X$ has a final
    object.
  \item The category $\cB^\bot/X$ of objects of $\cB^\bot$ under $X$ has
    an initial object.
\end{enumerate}
Then
\[\xymatrix@=0.5cm{
  & & & & (iv) \\
  (i) \ar@{<=>}[r]
    & (ii) \ar@{<=>}[r]
    & (iii) \ar@{<=>}[r]
    & (v) \ar@{=>}[ur] \ar@{=>}[dr] \\
  & & & & (vi)
}\]

If, in addition, $^\bot(\cB^\bot)=\cB$, then all the properties are
equivalent.
\end{proposition}

Let $\cB$ be a thick subcategory of a triangulated category $\cA$, let
$i$ be the inclusion functor from $\cB$ into $\cA$, let $Q$ be the
quotient functor from $\cA$ to $\cA/\cB$, let $\cB^\bot$ be the right
orthogonal of $\cB$, and let $i^\bot$ and $Q^\bot$ be the corresponding
functors. Proposition \ref{VIII:2-6-4} immediately gives the following
propositions.

\begin{proposition}\label{VIII:2-6-5}
The following two properties are equivalent:
\begin{enumerate}[(i)]
  \item The functor $i$ has a right adjoint.
  \item The functor $Q$ has a right adjoint.
\end{enumerate}
\end{proposition}

The proposition remains true when “right” is replaced by “left.”

\begin{proposition}\label{VIII:2-6-6}
The following two properties are equivalent:
\begin{enumerate}[(i)]
  \item The functor $i$ has a right adjoint.
  \item The functor $i^\bot$ has a left adjoint. The left orthogonal of the
    category $\cB^\bot$ is equal to $\cB$.
\end{enumerate}
\end{proposition}

\begin{proposition}\label{VIII:2-6-7}
Suppose the properties of Propositions \ref{VIII:2-6-5} and
\ref{VIII:2-6-6} hold. Let $i^\ast$ and $Q^\ast$ be the right adjoints of
$i$ and $Q$. Similarly, let $^\ast i^\bot$ and $^\ast Q^\bot$ be the left
adjoints of $i^\bot$ and $Q^\bot$. All these functors are exact. Moreover:

There is a functorial isomorphism
$Q^\ast\circ Q\iso i^\bot\circ ^\ast i^\bot$ such that the following
diagram is commutative:
\[\xymatrix{
  & \operatorname{id} \ar[dl] \ar[dr] \\
  Q^\ast \circ Q \ar[rr]^\sim
    & & i^\bot\circ ^\ast i^\bot
}\]
(the oblique arrows are defined by the adjunction properties).

Similarly, there is an isomorphism
$^\ast Q^\bot \circ Q^\bot \iso i\circ i^\ast$ such that the following
diagram is commutative:
\[\xymatrix{
  ^\ast Q^\bot\circ Q^\bot \ar[rr]^-\sim \ar[dr]
    & & i\circ i^\ast \ar[dl] \\
  & \operatorname{id}{}
}\]

Finally, there is a functorial morphism of degree $1$,
$\partial:i^\bot\circ ^\ast i^\bot \to i\circ i^\ast$, such that the
following triangle is distinguished $(\deg\partial=1$):
\[\xymatrix{
  i^\bot \circ ^\ast i^\bot \ar[rr]^-\partial
    & & i\circ i^\ast \ar[dl] \\
  & \operatorname{id}{} \ar[ul]
}\]

Such a morphism $\partial$ is unique and satisfies
$\partial(TX)=-T\partial(X)$ ($X\in\operatorname{Ob}\cA$, $T$ the
translation).
\end{proposition}

\section{Derived categories of an abelian category}\label{VIII:3}

\subsection{Derived categories}\label{VIII:3-1}

We use the notation of \ref{VIII:1-2-3}. We shall also use the following
notation.

\subsubsection{Notation}\label{VIII:3-1-1}

Let $\cA$ be an abelian category. Set
\begin{align*}
  \eD(\cA) &= \eK^{\infty,\infty}(\cA)/\eK^{\infty,\varnothing}(\cA) \\
  \eD^b(\cA) &= \eK^{b,b}(\cA)/\eK^{b,\varnothing}(\cA) \\
  \eD^+(\cA) &= \eK^{+,+}(\cA) / \eK^{+,\varnothing}(\cA) \\
  \eD^-(\cA) &= \eK^{-,-}(\cA) / \eK^{-,\varnothing}(\cA) \text{.}
\end{align*}

\begin{proposition}\label{VIII:3-1-2}
The natural functors appearing in the following diagram
\[\xymatrix@=0.5cm{
  & \eD(\cA) \\
  \eD^+(\cA) \ar[ur]
    & & \eD^-(\cA) \ar[ul] \\
  & \eD^b(\cA) \ar[ul] \ar[ur]
}\]
are fully faithful.

The natural functors
\begin{align*}
  \eD^b(\cA) &= \eK^{b,b}(\cA)/\eK^{b,\varnothing}(\cA) \to \eK^{+,b}(\cA)/\eK^{+,\varnothing}(\cA) \\
  \eD^b(\cA) &= \eK^{b,b}(\cA)/\eK^{b,\varnothing}(\cA) \to \eK^{-,b}(\cA) / \eK^{-,\varnothing}(\cA)
\end{align*}
are equivalences of categories.

Since the functors in the diagram
\[\xymatrix@=0.5cm{
  & \eD(\cA) \\
  \eD^+(\cA) \ar[ur]
    & & \eD^-(\cA) \ar[ul] \\
  & \eD^b(\cA) \ar[ul] \ar[ur]
}\]
are injective on objects, they realize the categories $\eD^b(\cA)$,
$\eD^+(\cA)$, and $\eD^-(\cA)$ as full subcategories of $\eD(\cA)$.
\end{proposition}

The proof of this proposition is contained in Theorem \ref{VIII:2-4-2}.

\begin{definition}\label{VIII:3-1-3}
The category $\eD(\cA)$ is called the \emph{derived category} of $\cA$.
The objects of $\eD(\cA)$ isomorphic to objects of $\eD^b(\cA)$ are
called the \emph{bounded} objects of $\eD(\cA)$. The objects of
$\eD(\cA)$ isomorphic to objects of $\eD^+(\cA)$ are called the
\emph{left-bounded} objects, and the objects of $\eD(\cA)$ isomorphic to
objects of $\eD^-(\cA)$ are called the \emph{right-bounded} objects.
\end{definition}

Let $D$ denote the canonical functor
\[
  D : \eC(\cA) \to \eD(\cA) \text{.}
\]
The “restriction” of $D$ to the full subcategory $\cA$ of $\eC(\cA)$
will again be denoted by $D$. The functor $D$ restricted to $\cA$ is fully
faithful. The functor $D$ factors uniquely through $\eK(\cA)$. Let $X$
and $Y$ be two objects of $\cA$.

One readily checks that, for every $m>0$, the groups
\[
  \hom_{\eD(\cA)}(D(X),T^{-m} D(Y)) \qquad \text{($T$ is the translation functor)}
\]
are zero. The cohomological dimension of $\cA$ is the least integer $n$
such that, for every $m>n$,
\[
  \hom_{\eD(\cA)}(D(X),T^m D(Y)) = 0 \qquad\text{for every pair $X,Y$ of objects of $\cA$.}
\]
If there is no such integer, the cohomological dimension of $\cA$ is said
to be infinite.

Let $\eI^+(\cA)$ (respectively $\eI(\cA)$) denote the full triangulated
subcategory of $\eK^+(\cA)$ (respectively of $\eK(\cA)$) defined by the
complexes whose objects are injective in every degree.

\begin{proposition}\label{VIII:3-1-4}
Suppose that the category $\cA$ has enough injectives. The natural functor
\[
  Q^+:\eK^+(\cA) \to \eD^+(\cA)
\]
induces an equivalence of categories
\[
  \eI^+(\cA) \to \eD^+(\cA)  \text{.}
\]
The quasi-inverse functor is a right adjoint of $Q^+$.
\end{proposition}

If, in addition, $\cA$ has finite cohomological dimension, the preceding
assertion remains true after the superscripts $+$ are omitted.

There is an analogous statement concerning projectives.

The proof of Proposition \ref{VIII:3-1-4} relies on
\ref{VIII:2-6-4}.

The triangulated structure of $\eD(\cA)$ is made more explicit by the
following proposition.

Let $\eE(\cA)$ be the category of three-term exact sequences in
$\eC(\cA)$. The category $\mathsf{SSS}(\cA)$ of semi-split exact
sequences (\S\ref{VIII:1-2-4}) is a full subcategory of $\eE(\cA)$.
We defined (loc. cit.) a functor
\[
  \rho:\mathsf{SSS}(\cA) \to \mathsf{TrK}(\cA)
\]
where $\mathsf{TrK}(\cA)$ denotes the category of distinguished triangles
of $\eK(\cA)$. This gives a functor
\[
  \sigma : \mathsf{SSS}(\cA) \to \mathsf{TrD}(\cA)
\]
where $\mathsf{TrD}(\cA)$ denotes the category of distinguished triangles
of $\eD(\cA)$.

\begin{proposition}\label{VIII:3-1-5}
There exists a unique functor
\[
  \Pi:\eE(\cA) \to \mathsf{TrD}(\cA)
\]
of the following form ($\deg \delta(S)=1$):
\[
(S=0 \to X^\bullet \to Y^\bullet \to Z^\bullet \to 0)
\mapsto
\xymatrix{
  & D(Z^\bullet) \ar[dl]_-{\delta(S)} \\
  D(X^\bullet) \ar[rr]^-{D(u)}
    & & D(Y^\bullet) \ar[ul]_-{D(v)}
}
\]
whose restriction to $\mathsf{SSS}(\cA)$ is $\sigma$. \emph{This functor
is essentially surjective.}
\end{proposition}

\begin{proposition}\label{VIII:3-1-6}
The canonical cohomological functor
\[
  \h^0:\eK(\cA) \to \cA
\]
factors uniquely through a functor that we again denote by $\h^0$,
\[
  \h^0:\eD(\cA) \to \cA \text{.}
\]
On $\eD(\cA)$, the functor $\h^0$ and its translates $\h^i$ form a
conservative system; that is, a morphism $f:X^\bullet \to Y^\bullet$ is
an isomorphism if and only if, for every integer $i$,
\[
  \h^i(f):\h^i(X) \to \h^i(Y)
\]
is an isomorphism.
\end{proposition}

\begin{definition}\label{VIII:3-1-7}
A morphism $f:X^\bullet \to Y^\bullet$ of $\eC(\cA)$ (respectively of
$\eK(\cA)$) is called a \emph{quasi-isomorphism} if, for every integer
$i$,
\[
  \h^i(f):\h^i(X) \to \h^i(Y)
\]
is an isomorphism.
\end{definition}

By the preceding proposition, the quasi-isomorphisms are the morphisms
that become isomorphisms in $\eD(\cA)$.

\begin{definition}\label{VIII:3-1-8}
Let $\sigma$ be a set of objects of $\cA$ and let $X^\bullet$ be an
object of $\eC(\cA)$ (respectively $\eK(\cA)$). A right (respectively
left) resolution of type $\sigma$ of $X^\bullet$ is a quasi-isomorphism
$X^\bullet \to V^\bullet$ (respectively $V^\bullet \to X^\bullet$), where
$V^\bullet$ is a complex all of whose objects belong to $\sigma$. We say
\emph{injective resolution} (respectively projective resolution) instead
of right (respectively left) resolution of injective (respectively
projective) type.
\end{definition}

In general, we shall omit the words “right” or “left” when no ambiguity
can result.

\begin{proposition}\label{VIII:3-1-9}
\begin{enumerate}
  \item Suppose that every object of $\cA$ embeds in an object of $\sigma$
    (respectively is a quotient of an object of $\sigma$). Then every
    object of $\eC^+(\cA)$ (respectively $\eC^-(\cA)$) has a right
    (respectively left) resolution of type $\sigma$ in $\eC^+(\cA)$
    (respectively in $\eC^-(\cA)$).
  \item Suppose, in addition, that there exists an integer $n$ such that
    for every exact sequence
    \begin{align*}
      Y^0 &\to Y^1 \to \cdots \to Y^{n-1} \to Y^n \to 0 \\
      \text{(respectively } 0 &\to Y^n \to Y^{n-1} \to \cdots \to Y^1 \to Y^0 \text{)}
    \end{align*}
    of objects of $\cA$ in which, for every $i$, $0\leqslant i\leqslant n-1$,
    $Y^i$ is an object of $\sigma$, the object $Y^n$ is also an object of
    $\sigma$. Then every object of $\eC(\cA)$ has a right (respectively
    left) resolution of type $\sigma$.
\end{enumerate}
\end{proposition}

\subsection{Study of $\ext$}\label{VIII:3-2}

Let $X^\bullet$ and $Y^\bullet$ be two objects of $\eC(\cA)$
(respectively $\eK(\cA)$), and let
$D:\eC(\cA) \to \eD(\cA)$ (respectively
$Q:\eK(\cA) \to \eD(\cA)$) be the canonical functor.

\begin{definition}\label{VIII:3-2-1}
We call the \emph{$i$th hyper-Ext}, and denote by
\[
  (X^\bullet,Y^\bullet)\mapsto \ext^i(X^\bullet,Y^\bullet)
\]
the functor $\hom_{\eD(\cA)}(D(X^\bullet),T^i D(Y^\bullet))$
(respectively
$\hom_{\eD(\cA)}(Q(X^\bullet),T^i Q(Y^\bullet))$).
\end{definition}

\begin{proposition}\label{VIII:3-2-2}
Let $0 \to X^\bullet \to Y^\bullet \to Z^\bullet \to 0$ be an exact
sequence in $\eC(\cA)$, and let $V^\bullet$ be an object of $\eC(\cA)$.
There are unbounded exact sequences
\[\xymatrix@=0.5cm{
  \cdots \ar[r]
    & \ext^i(V^\bullet,X^\bullet) \ar[r]
    & \ext^i(V^\bullet,Y^\bullet) \ar[r]
    & \ext^i(V^\bullet,Z^\bullet) \ar[r]^-\delta
    & \ext^{i+1}(V^\bullet,X^\bullet) \ar[r]
    & \cdots \\
  \cdots \ar[r]
    & \ext^i(Z^\bullet,V^\bullet) \ar[r]
    & \ext^i(Y^\bullet,V^\bullet) \ar[r]
    & \ext^i(X^\bullet,V^\bullet) \ar[r]^-\delta
    & \ext^{i+1}(Z^\bullet,V^\bullet) \ar[r]
    & \cdots
}\]
\end{proposition}

Let $X^\bullet$ and $Y^\bullet$ be two objects of $\eK(\cA)$. Denote by
$\mathsf{Qis}/X^\bullet$ (respectively $\mathsf{Qis}^+/X^\bullet$,
$\mathsf{Qis}^-/X^\bullet$, $\mathsf{Qis}^b/X^\bullet$) the category of
quasi-isomorphisms with target $X^\bullet$ and source in $\eK(\cA)$
(respectively $\eK^+(\cA)$, $\eK^-(\cA)$, $\eK^b(\cA)$). Similarly,
define the categories $Y^\bullet/\mathsf{Qis}$,
$Y^\bullet/\mathsf{Qis}^+$,\ldots (quasi-isomorphisms with source
$Y^\bullet)$.

The following proposition summarizes the various equivalent definitions
of $\ext^i$.

\begin{proposition}\label{VIII:3-2-3}
1. There are isomorphisms of functors
\[\xymatrix{
  \ext^0(X^\bullet,Y^\bullet) = \hom_{\eD(\cA)}(Q(X^\bullet),Q(Y^\bullet)) \ar[r]^-\sim
    & \varinjlim_{(\mathsf{Qis}/X^\bullet)^\circ} \hom_{\eK(\cA)}(-,Y^\bullet) \ar[r]^-\sim
    & \\
  \varinjlim_{Y^\bullet/\mathsf{Qis}} \hom_{\eK(\cA)}(X^\bullet,-) \ar[r]^-\sim
    & \varinjlim_{(\mathsf{Qis}/X^\bullet)^\circ \times Y^\bullet/\mathsf{Qis}} \hom_{\eK(\cA)}(-,-) \text{.}
}\]
2. If $X^\bullet$ is an object of $\eK^-(\cA)$,
\[
  \ext^0(X^\bullet,Y^\bullet) \iso \varinjlim_{(\mathsf{Qis}^-/X^\bullet)^\circ} \hom_{\eK(\cA)}(-,Y^\bullet)
\]
If $Y^\bullet$ is an object of $\eK^+(\cA)$,
\[
  \ext^0(X^\bullet,Y^\bullet) \iso \varinjlim_{Y^\bullet/\mathsf{Qis}} \hom_{\eK(\cA)}(X^\bullet,-)
\]
3. If $X^\bullet$ is an object of $\eK^-(\cA)$ and $Y^\bullet$ is an object
of $\eK^b(\cA)$,
\[
  \ext^0(X^\bullet,Y^\bullet) \iso \varinjlim_{Y^\bullet/\mathsf{Qis}^b} \hom_{\eK(\cA)}(X^\bullet,-)
\]
If $X^\bullet$ is an object of $\eK^b(\cA)$ and $Y^\bullet$ is an object
of $\eK^+(\cA)$,
\[
  \ext^0(X^\bullet,Y^\bullet) \iso \varinjlim_{(\mathsf{Qis}^b/X^\bullet)^\circ} \hom_{\eK(\cA)}(-,Y^\bullet)
\]
4. If the category $\cA$ has enough injectives and $Y^\bullet$ is an object
of $\eK^+(\cA)$, then $Y^\bullet$ has an injective resolution, unique up
to isomorphism in $\eK^+(\cA)$,
\[
  Y^\bullet \to I(Y^\bullet)
\]
and
\[
  \ext^0(X^\bullet,Y^\bullet) \iso \hom_{\eK(\cA)}(X^\bullet,I(Y^\bullet)) \text{.}
\]
There is an analogous statement for projectives.
5. If the category $\cA$
has enough injectives and if $\cA$ has \emph{finite cohomological dimension}, then
every object $Y^\bullet$ of $\eK(\cA)$ has a unique injective resolution
\[
  Y^\bullet \to I(Y^\bullet)
\]
and
\[
  \ext^0(X^\bullet,Y^\bullet) \iso \hom_{\eK(\cA)}(X^\bullet,I(Y^\bullet))
\]
There is an analogous statement for projectives.
\end{proposition}

\paragraph{Remark:}
On making explicit the isomorphism of Proposition
\ref{VIII:3-2-3}(3), one recovers the Yoneda definition.

\section{Derived functors}\label{VIII:4}

\subsection{Definition of derived functors}\label{VIII:4-1}

\begin{definition}\label{VIII:4-1-1}
Let $\cC$ and $\cC'$ be two graded categories (we denote by $T$ the
translation functor of $\cC$ and $\cC'$), and let $F$ and $G$ be two
graded functors from $\cC$ to $\cC'$. A \emph{morphism of graded functors}
is a morphism of functors
\[
  u:F\to G
\]
having the following property:

For every object $X$ of $\cC$, the following diagram is commutative:
\[\xymatrix{
  F(T X) \ar[r]^-{u(T X)}
    & G(T X) \\
  T F(X) \ar[r]^-{T u(X)} \ar[u]_-\wr
    & T G(X) \ar[u]_-\wr
}\]
\end{definition}

Let $\cC$ and $\cC'$ be two triangulated categories. Denote by
$\mathsf{Fex}(\cC,\cC')$ the category of exact functors from $\cC$ to
$\cC'$, the morphisms between two functors being morphisms of graded
functors.

Let $\cA$ and $\cB$ be two abelian categories and let
$\Phi:\eK^\ast(\cA)\to \eK^{\ast'}(\cB)$ be an exact functor ($\ast$ and
$\ast'$ denote one of the signs $+$, $-$, $b$, or $v$, “empty”). By
composition, the canonical functor
$Q:\eK^\ast(\cA)\to \eD^\ast(\cA)$ gives a functor
\[\xymatrix{
  \mathsf{Fex}(\eD^\ast(\cA),\eD^{\ast'}(\cB)) \ar[r]
    & \mathsf{Fex}(\eK^\ast(\cA),\eD^{\ast'}(\cB))
}\]
and hence (also denoting by $Q'$ the canonical functor
$\eK^{\ast'}(\cB) \to \eD^{\ast'}(\cB)$) a functor $\%$
(respectively $\%'$)
$\mathsf{Fex}(\eD^\ast(\cA),\eD^{\ast'}(\cB)) \to \mathsf{Ab}$:
\begin{align*}
  \Psi  &\mapsto \hom(Q'\circ \Phi,\Psi\circ Q) \\
  (\text{respectively }\Psi &\mapsto \hom(\Psi\circ Q,Q'\circ \Phi)\text{ ).}
\end{align*}

\begin{definition}\label{VIII:4-1-2}
We say that $\Phi$ admits a \emph{total right derived functor}
(respectively total left derived functor) if the functor $\%$
(respectively $\%'$) is representable. An object representing $\%$
(respectively $\%'$) is called the total right (respectively left) derived
functor of $\Phi$ and is denoted by $\eR \Phi$ (respectively
$\eL \Phi$).\footnote{A more manageable definition (and one equivalent in
practice) is given in \cite[XVII 1.2]{sga4}.}
\end{definition}

\subsubsection{Notation}\label{VIII:4-1-3}

Let $\cA$ and $\cB$ be two abelian categories and let $f:\cA\to \cB$ be
an additive functor. The \emph{total right derived functor} of
$\eK^+(f):\eK^+(\cA) \to \eK^+(\cB)$ is denoted by $\eR^+ f$. The
functors $\eR^- f$, $\eR f$, $\eL^+ f$, $\eL^- f$, and $\eL f$ are
defined similarly.

Likewise, let $F:\eK^\ast(\cA) \to \cB$ be a cohomological functor.
By composition, the canonical quotient functor
$Q:\eK^\ast(\cA) \to \eD^\ast(\cA)$ defines a functor
\[\xymatrix{
  \mathsf{Foco}(\eD^\ast(\cA),\cB) \ar[r]
    & \mathsf{Foco}(\eK^\ast(\cA),\cB)
}\]
(where $\mathsf{Foco}(-,-)$ denotes the category of cohomological
functors), and hence a functor $\%$ (respectively $\%'$)
$\mathsf{Foco}(\eD^\ast(\cA),\cB) \to \mathsf{Ab}$ given by
$G\mapsto \hom(F,G\circ Q)$ (respectively
$G\mapsto \hom(G\circ Q,F)$).

\begin{definition}\label{VIII:4-1-4}
We say that the functor $F$ admits a \emph{right cohomological derived
functor} (respectively left cohomological derived functor) if the functor
$\%$ (respectively $\%'$) is representable. An object representing $\%$
(respectively $\%'$) is called the right (respectively left)
cohomological derived functor and is denoted by $\eR F$ (respectively
$\eL F$).
\end{definition}

\subsubsection{Notation}\label{VIII:4-1-5}

Let $\cA$ and $\cB$ be two abelian categories and let $f:\cA\to \cB$ be
an additive functor. The right cohomological derived functor of
\[
  \h^0{}\circ \eK^+ f:\eK^+(\cA) \to \cB
\]
is denoted by $\eR^+ f$. When no confusion can result, the functor
$\eR^+ f\circ T^i$ ($T$ is the translation functor) is denoted by
$\eR^i f$. The right cohomological derived functor of
$\h{}\circ \eK f$ is denoted by $\eR f$. When no confusion can result,
the functor $\eR f\circ T^i$ is denoted by $\eR^i f$.

The functors $\eL^- f$, $\eL f$, and $\eL^i f$ are defined similarly.

\subsubsection{Remarks}\label{VIII:4-1-6}

As we shall see in the next subsection, these definitions include the
classical definition of derived functors when the categories have enough
injectives or projectives and the complexes are suitably bounded. They are
nevertheless more general, but without further hypotheses they are not
very manageable. For example, let $f:\cA\to \cB$ be an additive functor
between two abelian categories. Suppose that the functors $\eR^+ f$ and
$\eR f$ exist. I do not know how to prove that $\eR f$ induces on
$\eD^+(\cA)$ the functor $\eR^+ f$. Suppose, in addition, that the functor
$\eR^+ f$ (cohomological derived) exists. I do not know how to prove that
the functor $\h^0{}\circ \eR^+ f$ is isomorphic to $\eR^+ f$. In practice,
however, these properties will hold.

\subsection{Existence of derived functors}\label{VIII:4-2}

\begin{proposition}\label{VIII:4-2-1}
The hypotheses are those of Definition \ref{VIII:4-1-4}. Suppose, in
addition, that $\cA$ is a $\sU$-category, where $\sU$ is a universe such
that the set $\dZ$ of rational integers belongs to $\sU$, and such that
the set of objects of $\cA$ belongs to $\sU$. Suppose also that $\cB$ is
the category of abelian $\sU$-groups, or more generally the category of
sheaves of abelian $\sU$-groups on an arbitrary $\sU$-site. (A $\sU$-site
is a $\sU$-category whose set of objects belongs to $\sU$, equipped with
a topology.) The functor $F$ admits a right cohomological derived functor.
This functor is calculated as follows:

Let $X$ be an object of $\eD^\ast(\cA)$, and let $\underline X$ be the
object of $\eK^\ast(\cA)$ above $X$. There is a functorial isomorphism
\[
  \eR F(X) = \varinjlim_{\underline X/\mathsf{Qis}^\ast} F(-) \text{.}
\]
\end{proposition}

In particular, let $f:\cA\to \cB$ be an additive functor. The functors
$\eR f$ and $\eR^+ f$ exist and are calculated as above. The functor
$\eR f$ induces on $\eD^+(\cA)$ the functor $\eR^+ f$. Since no confusion
can then result, we shall use the notation $\eR^i f$.

\subsubsection{Remarks}\label{VIII:4-2-2}

Proposition \ref{VIII:4-2-1} can be expressed in a more general form in
the setting of triangulated categories. The hypotheses on the category
$\cB$ needed to ensure the validity of Proposition \ref{VIII:4-2-1} are
hypotheses asserting the existence and exactness of direct limits over
pseudo-filtered $\sU$-categories whose sets of objects belong to $\sU$.
Proposition \ref{VIII:4-2-1} introduces an asymmetry between right and
left derived functors, at least in practice.

\begin{theorem}\label{VIII:4-2-3} % labelled 4.2.2 in original
The data are those of Definition \ref{VIII:4-1-2}. Suppose, in addition,
that a full triangulated subcategory of $\eK^\ast(\cA)$, denoted $\cC$, is given,
and denote by $\cN$ the full triangulated subcategory of the acyclic
objects of $\cC$. Suppose that
\begin{enumerate}[(1)]
  \item Every object of $\cN$ is taken by $\Phi$ to an acyclic object of
    $\eK^{\ast'}(\cB)$.
  \item Every object $X$ of $\eK^\ast(\cA)$ is the source (respectively
    target) of a quasi-isomorphism whose target (respectively source) is
    an object of $\cC$.
\end{enumerate}
\begin{enumerate}[a)]
  \item Then $\Phi$ admits a total right (respectively left) derived
    functor.
  \item The functor $\eR \Phi$ (respectively $\eL \Phi$) is obtained as
    follows:

    The restriction of the functor
    $Q'\circ \Phi:\eK^\ast(\cA) \to \eD^{\ast'}(\cB)$ to $\cC$ vanishes on the
    objects of $\cN$ and therefore factors through $\cC/\cN$. This gives
    a functor $\Phi':\cC/\cN\to \eD^{\ast'}(\cB)$. But the natural
    functor $\cC/\cN\to \eD^\ast(\cA)$ is an equivalence of categories
    (Theorem \ref{VIII:2-4-2}). Composing with its quasi-inverse gives the
    functor $\eR\Phi$ (respectively $\eL \Phi$).
  \item Let $X$ be an object of $\eK^\ast(\cA)$, let $Y$ be an object of
    $\cC$, and let $X\to Y$ (respectively $Y\to X$) be a
    quasi-isomorphism. The object $\eR\Phi\circ Q(X)$ (respectively
    $\eL \Phi\circ Q(X)$) is isomorphic to $Q'\circ \Phi(Y)$.
  \item The functor $X\mapsto \eR\Phi\circ Q(X)$ (respectively
    $X\mapsto \eL \Phi\circ Q(X)$) is isomorphic to the functor
    $X\mapsto \varinjlim_{X/\mathsf{Qis}^\ast} Q'\circ\Phi(-)$
    (respectively
    $X\mapsto \varinjlim_{\mathsf{Qis}^\ast/X} Q'\circ \Phi(-)$).
  \item The functor $\h^0{}\circ \Phi:\eK^\ast(\cA)\to \cB$ admits a
    right (respectively left) cohomological derived functor, which is
    precisely $\h^0{}\circ\eR\Phi$ (respectively
    $\h^0{}\circ \eL\Phi$).
\end{enumerate}
\end{theorem}

\begin{corollary}\label{VIII:4-2-3-1} % originally corollary 1
Suppose that $\ast=+$ (respectively $\ast=-$) in the hypotheses of the
preceding theorem and that $\cA$ has enough injectives (respectively
projectives). Then Theorem \ref{VIII:4-2-3} applies with $\cC$ the
category of complexes of injectives (respectively projectives). Suppose
in addition that $\cA$ has finite cohomological dimension. Then the
theorem applies without restriction on $\ast$, with $\cC$ the category
of complexes of injectives (respectively projectives).
\end{corollary}

\begin{definition}\label{VIII:4-2-5} % originally 2.4
Let $\cA$ and $\cB$ be two abelian categories and let $f:\cA\to \cB$ be
an additive functor.

We say that $f$ has \emph{finite cohomological dimension} on the right
(respectively left) if $\eR^+ f$ (respectively $\eL^- f$)
exists and there is an integer $m\geqslant 0$ such that, for every integer
$n>m$, $\eR^n f$ (respectively $\eL^{-n} f$) is zero on the objects
arising from $\cA$ under the functor $D$. The
\emph{cohomological dimension} on the right (respectively left) of $f$ is then the
least integer $m$ having this property. If $\eR^+ f$ (respectively
$\eL^- f$) exists and $f$ does not have finite cohomological dimension,
we say that $f$ has infinite cohomological dimension.

An object $X$ of $\cA$ is said to be \emph{$f$-acyclic} on the right
(respectively on the left) if, for every integer $n>0$,
$\eR^n f(D(X)) = 0$ (respectively $\eL^{-n} f(D(X))=0$).
\end{definition}

\begin{corollary}\label{VIII:4-2-3-2} % originally corollary 2
The data are those of Definition \ref{VIII:4-2-5}. Suppose there exists a
set $\sigma$ of objects of $\cA$, stable under direct sums, having the
following properties:
\begin{enumerate}
  \item Every acyclic complex
    $\in \operatorname{Ob}{\eK^+(\cA)}$ (respectively
    $\in\operatorname{Ob}{\eK^-(\cA)}$) of objects of $\sigma$ is taken by
    $f$ to an acyclic complex.
  \item Every object $X$ of $\cA$ embeds in (respectively is a quotient of)
    an object of $\sigma$. The hypotheses of Theorem \ref{VIII:4-2-3}
    hold with $\ast=+$ (respectively $\ast=-$) and with $\cC$ the category
    of complexes whose objects in every degree belong to $\sigma$. The
    elements of $\sigma$ are right (respectively left) $f$-acyclic
    objects.
\end{enumerate}

If, in addition, $f$ has finite right (respectively left) cohomological
dimension, the set of right (respectively left) $f$-acyclic objects has,
in addition to properties (1) and (2) above, property (2) of Proposition
\ref{VIII:3-1-9}. Theorem \ref{VIII:4-2-3} then applies to the functor
$\eK f$, with $\cC$ the category of complexes whose objects are right
(respectively left) $f$-acyclic in every degree. The functor $\eR f$
(respectively $\eL f$) exists and induces on $\eD^+(\cA)$ the functor
$\eR^+ f$ (respectively $\eL^+ f$), and on $\eD^-(\cA)$ the functor
$\eR^- f$ (respectively $\eL^- f$).
\end{corollary}

\subsection{Product. Composition}\label{VIII:4-3}

Let $\cA$ be an abelian category, let
$F:\eD(\cA) \to \mathsf{Ab}$ be a cohomological functor, and let
$X^\bullet$ and $Y^\bullet$ be two objects of $\eC(\cA)$. The functor $F$
defines a map
\[
  \ext^0(X^\bullet,Y^\bullet) \to \hom_{\mathsf{Ab}}(F(D(X^\bullet)),F(D(Y^\bullet)))
\]
and hence a pairing
$F(D(X^\bullet))\times \ext^0(X^\bullet,Y^\bullet) \to F(D(Y^\bullet))$.

Applying translations to this pairing gives pairings
\[
  F^i(D(X^\bullet))\times \ext^j(X^\bullet,Y^\bullet) \to F^{i+j}(D(Y^\bullet)) \text{.}
\]
By definition, when these pairings are applied to cohomological derived
functors, they are precisely the Yoneda products. The properties of these
products follow immediately from this definition. We shall not develop
them here.

Let $\cA$, $\cA'$, and $\cA''$ be three abelian categories, and let
$f:\cA\to \cA'$ and $g:\cA'\to \cA''$ be two additive functors. Suppose
that the functors $\eR^+ f$ and $\eR^+ g$ exist. I do not know how to
deduce that the functor $\eR^+ g\circ f$ exists, and even if the latter
functor exists, the formula
\[\xymatrix{
  \eR^+ g\circ f\ar[r]^-\sim
    & \eR^+ g \circ \eR^+ f
}\]
is probably not true in general. Nevertheless, one has the following
proposition.

\begin{proposition}\label{VIII:4-3-1}
Suppose that the category $\cA'$ has enough right $g$-acyclic objects and
that $\cA$ has enough right $f$-acyclic objects that $f$ takes to right
$g$-acyclic objects. Then the functor $\eR^+ g\circ f$ exists and there
is an isomorphism
\[\xymatrix{
  \eR^+ g\circ f \ar[r]^-\sim
    & \eR^+ g \circ \eR^+ f \text{.}
}\]
\end{proposition}

There are analogous statements for left derived functors.

\section{Examples}\label{VIII:6} % originally chapter 2, section 3

\subsection{The functor \texorpdfstring{$\mathsf{Rhom}^\bullet$}{Hom*}}\label{VIII:6-1}

Let $\cA$ be an abelian category with enough injectives.

The functor $Y\mapsto \hom^\bullet(X^\bullet,Y^\bullet)$
(\ref{VIII:1-3}), where $X^\bullet$ is an object of $\eK(\cA)$ and
$Y^\bullet$ an object of $\eK^+(\cA)$, has a total right derived functor
(Corollary \ref{VIII:4-2-3-1}), denoted by
$\rhom^\bullet(X^\bullet,-)$. Let $Y$ be an object of $\eD^+(\cA)$.
The functor
\begin{align*}
  \eK(\cA) &\to \eD(\mathsf{Ab}) \\
  X^\bullet &\mapsto \rhom^\bullet(X^\bullet,Y)
\end{align*}
is exact and vanishes on acyclic complexes, hence gives an exact
bifunctor
\[\xymatrix{
  \eD(\cA)^\circ \times \eD^+(\cA) \ar[r]
    & \eD(\mathsf{Ab})
}\]
which we denote by $\mathsf{Rhom}^\bullet$.

When $\cA$ has finite cohomological dimension, the functor
$\mathsf{Rhom}^\bullet$ extends to
$\eD(\cA)^\circ\times \eD(\cA)$ (loc. cit.).

If, in addition, $\cA$ has enough projectives, the functor $\hom$ can be
derived with respect to the first argument. The two definitions coincide
on their common domain of validity.

\subsection{\texorpdfstring{$\Tor$}{Tor} of sheaves}\label{VIII:6-2}

Let $E$ be a ringed site (in particular, a ringed topological space), let
$\sA$ be its sheaf of rings, and let $\mathsf{Mod}(E)$ be the category of
$\sA$-modules. An object $X$ of $\mathsf{Mod}(E)$ is said to be
\emph{$\sA$-flat} if, for every exact sequence
\[\xymatrix{
  0 \ar[r]
    & Y' \ar[r]
    & Y \ar[r]
    & Y'' \ar[r]
    & 0
}\]
the sequence
$0 \to Y'\otimes_\sA X \to Y\otimes_\sA X \to Y''\otimes_\sA X\to 0$
is exact. There are enough $\sA$-flat objects: the direct sums of objects
$\sA_U$, where $U$ is an object of $E$ and $\sA_U$ is the sheaf $\sA$
“extended by zero outside $U$.” Thus Theorem \ref{VIII:4-2-3} applies,
and one can define the functor
\begin{align*}
  \eD^-(E,\sA) \times \eD^-(E,\sA) &\to \eD^-(E,\sA) \\
  (X,Y) &\mapsto X\lotimes Y
\end{align*}
the total tensor product, which is the left derived functor of the tensor
product. Passing to cohomology defines the $\Tor_i^\sA(X,Y)$, the local
hyper-Tor groups.

Since the category $\mathsf{Mod}(E)$ has enough injectives, one can derive
the functor $\hom^\bullet(X,Y)$, the complex of local homomorphisms, and
thus obtain $\mathsf{Rhom}(X,Y)$ ($X\in \eD$, $Y\in \eD^+$).

Let $X$ and $Y$ be two objects of $\eD^-(E,\sA)$ and let $Z$ be an object
of $\eD^+(E,\sA)$. There is a functorial isomorphism
\[\xymatrix{
  \hom_{\eD(E,\sA)}(X\lotimes Y,Z) \ar[r]^-\sim
    &\hom_{\eD(E,\sA)}(X,\mathsf{Rhom}(Y,Z)) \text{.}
}\]

\subsection{Left derived functor of the inverse image}\label{VIII:6-3}

Let $f:E\to E'$ be a map of ringed topological spaces. The inverse-image
functor
\[
  f^\ast:\mathsf{Mod}(E') \to \mathsf{Mod}(E)
\]
is not necessarily exact. But one knows how to define $f^\ast$-flat
objects, and there are enough such objects. Thus one can define the
functor $\eL^- f^\ast$. There is likewise a duality formula:

Let $X$ be an object of $\eD^-(E',\sA')$ and let $Y$ be an object of
$\eD^+(E,\sA)$. There is a functorial isomorphism
\[\xymatrix{
  \hom_{\eD(E,\sA)}(\eL^- f^\ast(X),Y) \ar[r]^-\sim
    & \hom_{\eD(E',\sA')}(X,\eR^+ f_\ast(Y)) \text{.}
}\]

% END INCLUDED components/sga4half_8_en.tex


\bibliographystyle{amsplain}
\bibliography{sga-sources}

\end{document}
